\documentclass[11pt]{article}

\usepackage[margin=1.1in]{geometry}

\usepackage{amssymb,mathtools,natbib}

% Anchored formal environments for causalsmith papers.
% First mandatory argument = obj_id (machine anchor joining the paper to the
% Lean crosswalk; invisible in the rendered PDF). Optional argument = name.
% Each environment also defines \label{obj:<obj_id>} so prose can use
% \cref{obj:T-1}; cleveref derives the printed kind and number from the target.
\usepackage{amsmath}
\usepackage{mathrsfs}
\usepackage{amsthm}
\usepackage{booktabs}
% Long displays may break across pages; let TeX stretch a little harder before an
% overfull \hbox so wide formulas/tables are less likely to run off the page.
\allowdisplaybreaks
\setlength{\emergencystretch}{3em}
\newtheorem{theorem}{Theorem}
\newtheorem{lemma}{Lemma}
\newtheorem{assumption}{Assumption}
\newtheorem{definition}{Definition}
\newtheorem{proposition}{Proposition}
\newtheorem{remark}{Remark}
\newtheorem{citedresult}{Cited result}
% The amsthm counter is `algorithmbox` (not `algorithm`) so a later `\usepackage{algorithm}` cannot clash.
\newcounter{algorithmbox}
\renewcommand{\thealgorithmbox}{\arabic{algorithmbox}}
\NewDocumentEnvironment{theoremv}{m o s}{\begin{theorem}\label{obj:#1}{\normalfont[\texttt{\detokenize{#1}}]\IfValueT{#2}{\ (#2)}\IfBooleanT{#3}{\verificationfootnotemark}\ }}{\end{theorem}}
\NewDocumentEnvironment{lemmav}{m o s}{\begin{lemma}\label{obj:#1}{\normalfont[\texttt{\detokenize{#1}}]\IfValueT{#2}{\ (#2)}\IfBooleanT{#3}{\verificationfootnotemark}\ }}{\end{lemma}}
\NewDocumentEnvironment{assumptionv}{m o s}{\begin{assumption}\label{obj:#1}{\normalfont[\texttt{\detokenize{#1}}]\IfValueT{#2}{\ (#2)}\IfBooleanT{#3}{\verificationfootnotemark}\ }}{\end{assumption}}
\NewDocumentEnvironment{definitionv}{m o s}{\begin{definition}\label{obj:#1}{\normalfont[\texttt{\detokenize{#1}}]\IfValueT{#2}{\ (#2)}\IfBooleanT{#3}{\verificationfootnotemark}\ }}{\end{definition}}
% A `propositionv` is a proved result tiered as a Proposition (theorem-grade, machine-verified).
\NewDocumentEnvironment{propositionv}{m o s}{\begin{proposition}\label{obj:#1}{\normalfont[\texttt{\detokenize{#1}}]\IfValueT{#2}{\ (#2)}\IfBooleanT{#3}{\verificationfootnotemark}\ }}{\end{proposition}}
% A `remarkv` holds NON-load-bearing interpretive / open-question content (e.g. a causalsmith `oeq:`
% open question). It is neither proved nor assumed; no theorem may depend on it.
\NewDocumentEnvironment{remarkv}{m o s}{\begin{remark}\label{obj:#1}{\normalfont[\texttt{\detokenize{#1}}]\IfValueT{#2}{\ (#2)}\IfBooleanT{#3}{\verificationfootnotemark}\ }}{\end{remark}}
% An `algorithmv` is a DEFINITION-kind object presented as a boxed, numbered-step procedure (an
% estimator / selection rule built in stages). Same anchor contract as `definitionv`; its body is
% ordinary LaTeX whose steps are `\begin{enumerate}` items, so tex2html needs no extra machinery.
% Presented in the CONVENTIONAL algorithm style readers expect from `algorithm`+`algorithmic`:
% a rule above, an "Algorithm N (Title)" caption line, a rule under the caption, the numbered
% steps, and a closing rule. It is NOT a float — the anchor contract requires the object to stay
% where the outline places it, and floats drift. The obj-id tag is suppressed here (it is
% bookkeeping, and a caption line carrying it reads as debug output in a finished paper).
\NewDocumentEnvironment{algorithmv}{m o s}%
  {\par\addvspace{\medskipamount}\noindent\rule{\linewidth}{0.8pt}\par\nobreak\vspace{2pt}%
   \refstepcounter{algorithmbox}\label{obj:#1}%
   \noindent\textbf{Algorithm \thealgorithmbox}\IfValueT{#2}{ \textnormal{(#2)}}%
   \IfBooleanT{#3}{\verificationfootnotemark}\par\nobreak\vspace{2pt}%
   \noindent\rule{\linewidth}{0.4pt}\par\nobreak\vspace{2pt}\normalfont}%
  {\par\nobreak\vspace{2pt}\noindent\rule{\linewidth}{0.8pt}\par\addvspace{\medskipamount}}
% A `citedv` is an IMPORTED external result the paper relies on but does NOT prove
% (a causalsmith `gate_class:"cited"` node). It is machine-checked against the cited
% source at F2.5, never proved here; the fixed provenance note makes that explicit
% so the environment can never be read as a contribution of this paper. The \citep
% attribution is supplied by the surrounding prose (P2), keyed to references.bib.
\NewDocumentEnvironment{citedv}{m o s}{\begin{citedresult}\label{obj:#1}{\normalfont[\texttt{\detokenize{#1}}]\IfValueT{#2}{\ (#2)}\IfBooleanT{#3}{\verificationfootnotemark}\ \textit{(imported result; machine-checked against the cited source, not proved here.)}\ }}{\end{citedresult}}
% \leanref{obj_id}{display text}: an inline first-mention link to a from-note object's
% verified Lean code. In the PDF it renders the display text plainly (the Lean link is a
% web-only affordance); tex2html turns it into a drawer-opening span.
\newcommand{\leanref}[2]{#2}
% For a theorem/lemma whose Lean declaration accepts a source-matched published proposition as an
% input, the starred anchored environment puts the mark in its heading; the generated text remains
% outside the frozen mathematical environment. Keep the combined macro for older paper bundles.
\newcommand{\verificationfootnotemark}{\footnotemark}
\newcommand{\verificationfootnotetext}[1]{\footnotetext{#1}}
\newcommand{\verificationfootnote}[1]{\verificationfootnotemark\verificationfootnotetext{#1}}
% xcolor before hyperref (hyperref would otherwise pull in plain `color` for colorlinks, and a
% later xcolor load clashes). The page-1 identity stamp at the end of this file needs a grey.
\usepackage{xcolor}
\usepackage{graphicx}
\usepackage{eso-pic}
% hyperref follows natbib and the environment macros; cleveref must follow hyperref.
\usepackage[colorlinks=true,linkcolor=blue,citecolor=blue,urlcolor=blue]{hyperref}
\usepackage[capitalise,nameinlink,noabbrev]{cleveref}
% Pin the names explicitly: labels are placed inside the underlying amsthm environments, and these
% declarations make the PDF contract independent of cleveref's theorem-name inference. House style:
% reference names always print with a capital first letter ("Theorem 1", "Definition 2"), in any
% sentence position — hence `capitalise` above and capitalized \crefname forms below.
\crefname{theorem}{Theorem}{Theorems}
\Crefname{theorem}{Theorem}{Theorems}
\crefname{lemma}{Lemma}{Lemmas}
\Crefname{lemma}{Lemma}{Lemmas}
\crefname{assumption}{Assumption}{Assumptions}
\Crefname{assumption}{Assumption}{Assumptions}
\crefname{definition}{Definition}{Definitions}
\Crefname{definition}{Definition}{Definitions}
\crefname{proposition}{Proposition}{Propositions}
\Crefname{proposition}{Proposition}{Propositions}
\crefname{remark}{Remark}{Remarks}
\Crefname{remark}{Remark}{Remarks}
\crefname{citedresult}{Cited result}{Cited results}
\Crefname{citedresult}{Cited result}{Cited results}
\crefname{algorithmbox}{Algorithm}{Algorithms}
\Crefname{algorithmbox}{Algorithm}{Algorithms}
% ---------------------------------------------------------------------------
% Page-1 identity stamp (arXiv style) and the pinned title-block date.
%
% Split of responsibility: THIS file owns the FORMAT (placement, font, colour, punctuation),
% the generated `paper_stamp.tex` owns only the VALUES (number+version, area, revised date,
% DOI, link target). P4 refreshes this template on every emit, so a layout fix reaches every
% bundle from a recompile; and because `paper_stamp.tex` is not one of the P2 assembly
% digest's authored inputs, a DOI that arrives after publication needs only that one small
% file rewritten plus a recompile — never a redraft, never a reassembly.
%
% Defaults first, so a bundle WITHOUT a stamp file compiles exactly as it always did:
% `\cspaperdate` falls back to `\today` and no stamp is drawn.
\providecommand*{\cspaperdate}{\today}  % the \date{} target
\providecommand*{\csstampid}{}          % e.g. CSWP-2026-008v3 (empty ⇒ no stamp at all)
\providecommand*{\csstamparea}{}        % e.g. Stat
\providecommand*{\csstampdate}{}        % e.g. 27 Aug 2026
\providecommand*{\csstampdoi}{}         % e.g. 10.5281/zenodo.123 (empty ⇒ DOI part omitted)
\providecommand*{\csstamplink}{}        % href target (DOI url, else the paper's page; may be empty)
\IfFileExists{paper_stamp.tex}{% GENERATED by CausalSmith P4 — paper identity stamp values. Do not hand-edit.
%
% Field order on the stamp (owned by paper_macros.tex):
%   <number>v<version> · doi:<doi> · [<Area>] <date>   — the DOI is never the last token, so a
%   text extractor cannot run it into whatever follows. Without a DOI that part is omitted.
% Format lives in paper_macros.tex; only the values are here, and this file is NOT one of
% the P2 assembly digest's authored inputs. A DOI arriving after publication therefore needs
% this file rewritten plus a `latexmk` recompile — never a redraft or a reassembly.
\renewcommand*{\csstampid}{CSWP-2026-034v5}
\renewcommand*{\csstamparea}{Stat}
\renewcommand*{\csstampdate}{2 Oct 2026}
\renewcommand*{\csstampdoi}{}
\renewcommand*{\csstamplink}{https://causalsmith.org/papers/stat_mar_rareq_logfrontier_v1}
\renewcommand*{\cspaperdate}{2 October 2026}
}{}
\makeatletter
% Field order is deliberate: the DOI is NEVER the last token on the line. A trailing identifier
% runs into whatever a text extractor emits next, which makes it unsafe to match mechanically
% (the Zenodo stamp matcher reads exactly this line); the date is a safe terminator.
%   <number>v<version> · doi:<doi> · [<Area>] <date>      — with a DOI
%   <number>v<version> · [<Area>] <date>                  — without
\newcommand{\cs@stampsep}{\,\textperiodcentered\,}
\newcommand{\cs@stamptext}{%
  \csstampid
  \ifx\csstampdoi\@empty\else\cs@stampsep doi:\csstampdoi\fi
  \ifx\csstamparea\@empty\else\cs@stampsep[\csstamparea]\fi
  \ifx\csstampdate\@empty\else\quad\csstampdate\fi
}
\ifx\csstampid\@empty\else
  \definecolor{csstampgrey}{gray}{0.45}
  % Background shipout material: it occupies no space, so the text block and the \thanks
  % footnote are byte-identical to an unstamped compile. The starred form applies to the
  % NEXT page shipped only — i.e. page 1. Placed in the left margin (the text block starts
  % at 1.1in), reading bottom-to-top, in a Times-like face at 8pt.
  \AddToShipoutPictureBG*{%
    \AtPageLowerLeft{%
      \put(\LenToUnit{0.42in},\LenToUnit{1.1in}){%
        \rotatebox{90}{%
          \fontfamily{ptm}\fontsize{8}{10}\selectfont
          \ifx\csstamplink\@empty
            \textcolor{csstampgrey}{\cs@stamptext}%
          \else
            \expandafter\href\expandafter{\csstamplink}{\textcolor{csstampgrey}{\cs@stamptext}}%
          \fi
        }%
      }%
    }%
  }
\fi
\makeatother


\title{Treatment effect estimation with missing outcomes and many covariate categories}

\author{CausalSmith\thanks{Machine-generated by the CausalSmith research pipeline.}}

\date{\cspaperdate}

\begin{document}

\maketitle

\begin{abstract}
We characterize the precision of population average treatment-effect estimation in balanced randomized trials with selectively observed binary outcomes, a binary surrogate, and a finite baseline alphabet. Outcomes are missing at random conditional on treatment, baseline category, and surrogate, with unknown cell-specific arrival probabilities bounded below by a supplied positive floor \(q\). This floor indexes and tunes the procedure. For \(n\), the sample size, and \(d\), the number of baseline categories, we give an explicit mathematical attainment construction: a deterministic mixed-count estimator with squared-error risk bounded by a universal constant times the comparison rate
\[
r_{n,d,q}
=
\min\left\{1,\frac{1}{nq}
+\left[\frac{d}{nq\log(e+nq)}\right]^2\right\}.
\]
This guarantee holds for every positive arrival floor. Under \(q\log^2(e+nq)\le1/64\), matching lower bounds establish minimax squared-error order \(r_{n,d,q}\) and optimal maximal expected length of order \(\sqrt{r_{n,d,q}}\) for uniformly honest connected intervals at each fixed coverage level strictly between zero and one. At each fixed \(0<q<1/2\), the point-risk comparison also matches, with constants depending on \(q\). Near complete ascertainment, an observed-outcome contrast attains risk at most \(4/n+(1-q)^2\). For \(d\ge1024n^2\), the matching minimax order is \(n^{-1}+(1-q)^2\), making \(1-q=O(n^{-1/2})\) necessary and sufficient along arrival-floor sequences for parametric risk uniformly over all alphabet sizes. A target-preserving synthetic assignment transfers the construction to discrete observational regression contrasts and closes the stated logarithmic minimax gap at each fixed strict-interior treatment-positivity margin.
\end{abstract}

\section{Introduction}

The precision of a randomized trial at a prespecified analysis horizon depends on both the primary outcomes that have arrived and the information recorded for the full sample. Baseline categories, treatment assignments, and short-term measurements can remain available while primary outcomes are being collected. When outcome arrival varies across these observed characteristics, estimating a population treatment effect requires combining population cell weights with outcome information within cells. This paper characterizes the resulting statistical precision in a balanced randomized experiment with binary primary outcomes, a binary surrogate, and a finite baseline alphabet.

The model imposes independent and identically distributed sampling, treatment assignment independent of baseline characteristics and potential outcomes and surrogates, assignment probability one half, and consistency. Outcome arrival is missing at random conditional on treatment, baseline category, and realized surrogate. Each occupied cell has an unknown arrival probability at least a supplied value \(q\), where \(0<q\le1\). This lower bound indexes the model and is used to tune the estimator and interval; individual cell probabilities remain unknown. Baseline probabilities may vary arbitrarily, including zero probabilities. The surrogate supplies observed conditioning information for ascertainment, and the target is the population average treatment effect on the primary outcome. Under these conditions, \cref{obj:lem:unrestricted-cell-identification} expresses the effect as a signed sum of cell membership probabilities multiplied by outcome means among arrivals. Membership information from every sampled individual enters this representation.

Our first result is a constructive squared-error guarantee across the entire positive-arrival domain. Write \(n\) for the sample size, \(d\) for the number of baseline categories, and \(N=nq\) for the effective sample-size scale. The mixed-count estimator in \cref{obj:def:mixed-count-estimator} satisfies a uniform risk bound of order
\[
r_{n,d,q}
=
\min\left\{1,\frac1N+
\left[\frac{d}{N\log(e+N)}\right]^2\right\},
\]
where \(r_{n,d,q}\) is the comparison rate, with a universal comparison constant, as established in \cref{obj:thm:unrestricted-mixed-count-upper}. The construction separates membership, pilot, and outcome estimation into three auxiliary streams. Membership counts estimate cell weights; the pilot selects between empirical arrived-outcome means and a Chebyshev polynomial statistic built from marked factorial counts. Exact conditional averaging over the auxiliary randomization produces a deterministic estimator.

Under the rare-arrival restriction
\[
q\log^2(e+nq)\le\frac1{64},
\]
\cref{obj:thm:matched-minimax-frontier} establishes matching minimax squared-error order \(r_{n,d,q}\), with universal constants. The decision class includes all bounded parameter-independent randomized estimators. Its minimax value agrees with that for deterministic estimators by \cref{obj:lem:randomized-kernel-barycenter}. The lower bound retains the full observed experiment, including memberships whose outcomes have yet to arrive. A one-category family supplies the reciprocal effective-sample-size contribution, while an activated moment-matching construction supplies the category contribution. Along sequences satisfying the rare-arrival restriction and \(N\to\infty\), \cref{obj:prop:phase-boundaries} consequently gives consistency exactly when \(d=o\{N\log(e+N)\}\), and squared-error order \(N^{-1}\) exactly when \(d=O\{\sqrt N\log(e+N)\}\).

The same rare-arrival experiment admits a matching uncertainty-quantification comparison. For each fixed miscoverage level \(\alpha\in(0,1)\), \cref{obj:thm:connected-interval-frontier} establishes optimal maximal expected length of order \(\sqrt{r_{n,d,q}}\) among connected intervals with uniform coverage at least \(1-\alpha\). The comparison constants may depend on \(\alpha\). Coverage and expected length integrate jointly over the sample and any endpoint randomization. A deterministic squared-error envelope supplies an attaining interval, and the converse covers the full randomized endpoint decision class.

Near complete ascertainment, the remaining deficit supplies a complementary precision scale. Throughout the arrival-floor model, \cref{obj:thm:unrestricted-near-complete-arrival-envelope} gives
\[
\frac{1}{256n}
\le \mathfrak R^{+}_{n,d,q}
\le
\min\left\{1,\overline C r_{n,d,q},
\frac4n+(1-q)^2\right\},
\]
where \(\mathfrak R^{+}_{n,d,q}\) is the minimax squared-error risk and \(\overline C\) is a universal constant. The observed-outcome arm contrast attains the last upper bound. For \(d\ge1024n^2\), \cref{obj:thm:unrestricted-large-alphabet-deficit-lower} establishes matching order \(n^{-1}+(1-q)^2\) uniformly in \(q\). Hence, along arrival-floor sequences, an eventual \(n^{-1}\) risk guarantee holding simultaneously over every alphabet size is equivalent to \(1-q_n=O(n^{-1/2})\), by \cref{obj:prop:unrestricted-uniform-near-complete-elbow}. Complete arrival gives \(n^{-1}\) order uniformly over arbitrary alphabet growth.

The construction also connects randomized outcome arrival to observational adjustment. An independent fair assignment retains an observational outcome when the synthetic assignment agrees with the observed treatment. Under a fixed strict-interior treatment-positivity margin, this transformation preserves the baseline-weighted regression contrast and supplies a randomized missing-at-random completion. Averaging the mixed-count estimator over synthetic assignments yields the deterministic observational procedure in \cref{obj:thm:zeng-fixed-positivity-polynomial-frontier}. At each fixed margin strictly between zero and one half, its minimax squared-error order is
\[
\min\left\{1,\frac1n+
\left[\frac{d}{n\log(e+n)}\right]^2\right\},
\]
with margin-dependent constants. Consistency and conditional exchangeability give the observational contrast its population causal interpretation. The same experiment comparison supplies the matched randomized point-risk order at every fixed \(0<q<1/2\) in \cref{obj:thm:unrestricted-fixed-q-minimax-frontier}, with constants depending on that floor.

The statistical mechanism is a weighted-cell-mean problem with many scarce arrived counts. Ordinary cell ratios accumulate a category-dependent bias; on the difficult large-alphabet branch, a Chebyshev correction reduces the aggregate bias to the scale \(d/\{N\log(e+N)\}\). The converse uses moment-matched reciprocal priors whose activated observation laws remain close even though all cell memberships are retained. Each augmented record carries an independent activation flag with probability \(qH\), where \(H\) is the upper endpoint of the latent prior support and has squared-logarithmic order. This explains the rare-arrival condition ensuring a valid same-experiment construction. \Cref{sec:constructive-estimator,obj:thm:full-experiment-lower} present these two sides.

These results connect three strands of existing work. Surrogate-assisted analysis explains how early measurements inform treatment-effect estimation and efficiency under appropriate observation and nuisance conditions \citep[Theorems 2.1--2.2 and 4.1--4.2]{KallusMao2025}. Decaying-labelling analysis establishes asymptotic inference under effective-sample-size and nuisance-rate conditions \citep[Theorems 2.2 and 4.1, equation (2.3)]{ZhangChakraborttyBradic2025}. Our comparisons characterize uniform squared error and honest expected interval length as category frequencies, alphabet size, and ascertainment vary over the specified finite models. For discrete observational adjustment, \citet[Assumption 2, Theorems 1--2, Section 3.2, and Appendix C.5]{ZengBalakrishnanHanKennedy2026} provide the closest rate benchmark. The synthetic-assignment construction supplies an attaining polynomial estimator, while their lower-bound conclusion supports the fixed-positivity comparisons. The polynomial and moment-matching methods build on discrete functional estimation \citep{JiaoVenkatHanWeissman2015,Wu2016,WuYang2019}; the expected-length criterion connects to honest confidence-interval theory \citep{Donoho1994,Cai2004,ArmstrongKolesar2020}.



The exposition proceeds through related work, the experiment and identification conditions, and the point-risk comparisons across ascertainment regimes. It then develops the constructive estimator, honest intervals under rare arrival, and the observational extension through synthetic assignment. Discussion and open questions follow. The appendices collect identification and decision-theory preliminaries, risk certificates, full-experiment lower bounds, interval and observational arguments, the verification note, and proofs of the main results.

\section{Related work}

The availability of short-term measurements before primary outcomes motivates both surrogate analysis and methods for incomplete trial data. \citet{Prentice1989} formulates criteria for surrogate endpoints, while \citet{Fleming1996} emphasizes the substantive evidence required to connect surrogate responses to clinical effects. Short-term outcomes can also inform treatment-effect estimation while primary outcomes are still being collected, as in \citet{VanLancker2020}. Here, the surrogate enters the conditioning information for outcome arrival, and the target remains the population average treatment effect. Identification rests on randomized assignment, consistency, and outcome missingness at random conditional on baseline category, treatment, and surrogate, as specified in \cref{obj:ass:randomized-independence,obj:ass:surrogate-consistency,obj:ass:outcome-consistency,obj:ass:arrival-mar}.

The closest comparison for estimation rates is \citet[Assumption 2, Theorems 1--2, Section 3.2, and Appendix C.5]{ZengBalakrishnanHanKennedy2026}. Their observational experiment comprises independent observations of discrete covariates, treatment, and outcome, with treatment propensities bounded away from zero and one by a fixed positivity margin. Writing \leanref{sym:n}{\ensuremath{n}} for the sample size and \leanref{sym:d}{\ensuremath{d}} for the number of covariate categories, their plug-in analysis gives a squared-error upper bound of order \(n^{-1}+d^2/n^2\). Their minimax lower bound is of order \(n^{-1}+d^2/\{n^2\log^2 n\}\) in the stated regime \(d\lesssim n\log n\), with constants depending on positivity. These conclusions give a logarithmic comparison between the plug-in guarantee and the minimax benchmark. Appendix C.5 develops the moment-matching argument underlying that benchmark. For the corresponding observational regression contrast under fixed strict-interior positivity on occupied categories, \cref{obj:thm:zeng-fixed-positivity-polynomial-frontier} supplies a constructive polynomial-estimation guarantee through synthetic assignment. Its positivity margin, \leanref{sym:\epsilon}{\ensuremath{\epsilon}}, is held fixed in this comparison.

An earlier conference announcement by \citet[abstract 300411]{Kennedy2019DiscreteAbstract} considers an average causal effect with arbitrarily high-dimensional discrete confounders. It reports plug-in risk bounds, a minimax lower bound, and a variation with propensity-score positivity deteriorating with sample size, thereby establishing earlier consideration of the shrinking-positivity observational problem. The theorem-level comparison permitted by the available abstract is recorded in the limitations section.

The randomized analysis addresses a different observation mechanism. In the \leanref{S-1}{randomized fixed-horizon trial with a finite baseline alphabet}, treatment is assigned with known probability one half, and primary outcomes arrive with unknown probabilities conditional on the observed cells. The arrival floor \leanref{sym:q}{\ensuremath{q}} is imposed on occupied cells in \cref{obj:ass:occupied-cell-arrival}. Under the rare-arrival restriction in \cref{obj:ass:rare-arrival-slice}, \cref{obj:thm:matched-minimax-frontier} matches upper and lower squared-error bounds at the rate defined in \cref{obj:def:frontier-rate},
\[
\min\left\{1,\frac{1}{nq}
+\left[\frac{d}{nq\log(e+nq)}\right]^2\right\}.
\]
This comparison separates fixed treatment positivity in the observational experiment from outcome ascertainment in a balanced randomized experiment. The latter analysis quantifies the joint contribution of the effective sample size and the number of categories, while the synthetic-assignment construction connects the two experiments.

Surrogate-assisted efficiency provides a complementary perspective. \citet[Theorems 2.1--2.2 and 4.1--4.2]{KallusMao2025} characterize efficiency bounds with different patterns of surrogate and outcome observation, under treatment unconfoundedness, outcome missingness at random, and the overlap conditions appropriate to each observation pattern. Their labelled-sample asymptotics allow the labelled fraction to vanish while the expected number of labelled observations diverges, with fixed conditional distributions within the labelled and unlabelled populations. Attainment of the efficiency bound requires their nuisance convergence conditions: relevant products of regression and propensity or density-ratio errors must be negligible relative to the inverse square root of the expected labelled sample size. Treatment propensity estimates satisfy the stated boundedness condition, and estimated labelling propensities remain positive. Their moment requirements include uniformly bounded moments of order greater than two for the outcome residual, density ratio, and specified regression terms \citep[Assumption 7 and Appendix D.1, Assumption 9]{KallusMao2025}. These results explain how predictive surrogates improve asymptotic precision under those conditions. The finite-alphabet minimax analysis here characterizes uniform risk as category frequencies and the alphabet size vary over the stated model class.

Decaying outcome observation is also central to \citet[Theorems 2.2 and 4.1, equation (2.3)]{ZhangChakraborttyBradic2025}. Their triangular-array observational framework combines treatment overlap with positive, potentially vanishing labelling probabilities conditional on treatment and covariates. Theorem 2.2 establishes asymptotic normality when both nuisance models are correctly specified, the inverse-propensity effective sample size diverges, and the product-rate condition in equation (2.3) holds. Its conditions also include finite inverse-propensity expectations, bounded outcome variance, conditional residual variance bounded above and away from zero, and the stated weighted nuisance and tail requirements. Their bias-reduced estimator in Theorem 4.1 accommodates nuisance approximation error under its product restriction and alternative sparsity conditions, together with sub-Gaussian design and conditional noise assumptions, covariance lower bounds, propensity moment control, and bounded offset signals \citep[Assumptions 1--3 and 5--8]{ZhangChakraborttyBradic2025}. These conditions support asymptotic inference with estimated nuisance functions. Our randomized rare-arrival comparison instead measures worst-case squared error over a finite categorical model, and the interval analysis imposes uniform coverage over that same experiment.

The approximation methods have established antecedents in the estimation of discrete-distribution functionals. Polynomial estimators and moment-matched lower bounds are developed by \citet{JiaoVenkatHanWeissman2015}, \citet{Wu2016}, and \citet{WuYang2019}. Their pairing reflects the dual relationship between polynomial approximation and separation of distributions with matching moments. Classical approximation theory supplies the reciprocal-function approximation used in the converse argument \citep{Meinardus1967}; explicit reciprocal approximation formulas are also presented in \citet[Section 3]{https://doi.org/10.1007/s11075-026-02364-1}. The application here combines these tools with observed cell memberships and selectively observed outcomes. The roles of reciprocal approximation and moment-matched priors are stated in \cref{obj:lem:reciprocal-best-approximation,obj:lem:moment-matched-reciprocal-priors}, and the resulting comparison concerns the full observed experiment in \cref{obj:thm:full-experiment-lower}.

Bias correction for regression functionals supplies another comparison. Higher-order influence functions use higher-order U-statistics to obtain rate-optimal point and interval procedures for nonlinear functionals under smoothness and nuisance-estimation conditions \citep{RobinsLiTchetgenVanderVaart2008,Robins2009}. Augmented minimax linear estimation corrects a regression plug-in estimator by a minimax linear estimate of its error and obtains semiparametric efficiency under its function-class conditions \citep{Hirshberg2021}. Weak-overlap analysis instead thresholds doubly robust scores and derives asymptotic normality and Wald inference under propensity-tail, outcome-smoothness, and nuisance-rate restrictions \citep{Dorn2025}. The present model replaces smooth nuisance classes by arbitrary masses and heterogeneous arrival probabilities on a finite, growing alphabet. Its guarantees are nonasymptotic worst-case squared-error and honest-length comparisons, and its logarithmic improvement comes from polynomial correction of scarce arrived-cell counts rather than smooth regression approximation.

The interval criterion connects the analysis to uncertainty quantification through maximal expected length subject to uniform coverage. \citet{Donoho1994} relates statistical estimation to optimal recovery, \citet{Cai2004} studies expected-length comparisons and adaptation for confidence intervals, and \citet{ArmstrongKolesar2020} develops honest intervals that account for worst-case bias in nonparametric regression. For the present missing-outcome experiment, \cref{obj:def:interval-length-risk} specifies honesty for connected intervals, and \cref{obj:thm:connected-interval-frontier} gives the corresponding length comparison under rare outcome arrival. This places estimation error and honest interval length within a common finite-alphabet decision problem.

At complete outcome arrival, \citet{Sudijono2026} provides a useful comparison from randomized-experiment theory. They jointly optimize the assignment design and estimator for the sample average treatment effect in a finite population with binary potential outcomes, obtaining a sharp minimax analysis that includes a second-order risk expansion. Our complete-arrival result in \cref{obj:thm:unrestricted-complete-arrival-frontier} concerns a population average treatment effect under independent and identically distributed sampling and a fixed balanced assignment design. The common first-order sampling scale therefore arises in two distinct decision problems: joint design and estimation for a finite-population sample effect, and estimation of a population effect under a specified randomized sampling experiment.


\section{Experiment, identification, and risk}

Throughout, \(n\ge1\) is the sample size, \(d\ge1\) is the number of baseline categories, and \(q\in(0,1]\) is a supplied lower bound on occupied-cell outcome arrival. Generic positive constants \(c\) and \(C\) are independent of these parameters unless dependence is stated explicitly. We write \(u\lesssim v\) for \(u\le Cv\) and \(u\asymp v\) for two such inequalities, uniformly over the stated parameter domain. Scalar errors are measured by absolute value. The arm index \leanref{sym:a}{\ensuremath{a}} and surrogate index \leanref{sym:s}{\ensuremath{s}} range over \(\{0,1\}\), while the baseline index \leanref{sym:x}{\ensuremath{x}} ranges over \(\{0,\ldots,d-1\}\). We use \leanref{sym:N}{\ensuremath{N}}, the effective sample-size scale, for \(N=nq\); it combines the sample size with the guaranteed within-cell arrival probability.

\subsection{Records and cells}

The experiment has a finite baseline category \leanref{sym:X}{\ensuremath{X}}, a binary treatment assignment \leanref{sym:A}{\ensuremath{A}}, binary potential surrogates \leanref{sym:S(a)}{\ensuremath{S(a)}}, binary potential outcomes \leanref{sym:Y(a)}{\ensuremath{Y(a)}}, and an outcome-arrival indicator \leanref{sym:R}{\ensuremath{R}}. The following definitions specify the record space and the coordinates selected by treatment assignment.

\begin{assumptionv}{ass:iid-sampling}[Independent and identically distributed sampling]
The observed records \(O_1,\ldots,O_n\) are independent and identically distributed according to the observed-data law induced by \(P\), the full-data probability law.
\end{assumptionv}

Treatment assignment selects the realized surrogate \leanref{sym:S}{\ensuremath{S}} from the two potential surrogate coordinates.

\begin{definitionv}{synth_2}[Full-data record coordinates]
A full-data record, for alphabet size \(d\in\mathbb N\), consists of a baseline category \(X\), a binary arm indicator \(A\), control and treated surrogate coordinates \(S(0)\) and \(S(1)\), control and treated outcome coordinates \(Y(0)\) and \(Y(1)\), and an arrival flag \(R\), with domains
\[
X\in\{x\in\mathbb Z:0\le x<d\},
\qquad
\bigl(A,S(0),S(1),Y(0),Y(1),R\bigr)\in\{0,1\}^{6}.
\]
\end{definitionv}

The assigned-arm outcome \leanref{sym:Y}{\ensuremath{Y}} is selected in the same way, and the finite record space carries its discrete measurable structure.

\begin{definitionv}{synth_3}[Potential and assigned surrogates]
The potential surrogate \(S(a)\), for \(a\in\{0,1\}\), denotes the record's binary control surrogate when \(a=0\) and binary treated surrogate when \(a=1\). The assigned-arm surrogate \(S\), for any full-data record with binary assignment \(A\in\{0,1\}\), is defined by
\[
S=
\begin{cases}
S(1), & A=1,\\
S(0), & A=0.
\end{cases}
\]
\end{definitionv}

A full-data probability law \leanref{sym:P}{\ensuremath{P}}, defined in \cref{obj:synth_1}, specifies the joint distribution of these coordinates.

\begin{definitionv}{synth_4}[Potential and assigned-arm outcomes]
The binary potential outcomes \(Y(0)\) and \(Y(1)\) are the control and treated outcome coordinates of a full-data record. For a nonnegative integer \(d\), such a record consists of a baseline category \(X\) in an alphabet of size \(d\) and binary coordinates \(A,S(0),S(1),Y(0),Y(1),R\). The assigned-arm outcome \(Y\) is defined by
\[
Y=
\begin{cases}
Y(1), & A=1,\\
Y(0), & A=0.
\end{cases}
\]
The space of full-data records is equipped with the measurable structure in which every subset is measurable.
\end{definitionv}

Outcome ascertainment is described within treatment--baseline--surrogate cells. The cell-index set \leanref{sym:\mathcal J_d}{\ensuremath{\mathcal J_d}} contains a cell \leanref{sym:j}{\ensuremath{j}} for each possible triple of these observed coordinates.

\begin{assumptionv}{ass:randomized-independence}[Independence of treatment assignment]
Under \(P\), the full-data probability law, treatment assignment \(A\) satisfies
\[
A\perp\!\!\!\perp\{X,S(0),S(1),Y(0),Y(1)\}.
\]
Here \(X\) is the baseline covariate, \(S(0),S(1)\) are the potential surrogates, and \(Y(0),Y(1)\) are the potential outcomes.
\end{assumptionv}

Two probabilities distinguish membership from outcome arrival: \leanref{sym:p_{axs}(P)}{\ensuremath{p_{axs}(P)}} is the full-membership cell probability, and \leanref{sym:\rho_{axs}(P)}{\ensuremath{\rho_{axs}(P)}} is the conditional arrival probability in an occupied cell.

\begin{assumptionv}{ass:balanced-randomization}[Balanced treatment assignment]
Under \(P\), the full-data probability law, treatment assignment \(A\) satisfies
\[
P(A=1)=\frac12.
\]
\end{assumptionv}

Thus an occupied cell has positive membership probability, and its arrived-cell probability is the product of its membership and conditional arrival probabilities. For the treatment contrast, the treatment sign \leanref{sym:g(a)}{\ensuremath{g(a)}} assigns opposite signs to the two arms.

\begin{definitionv}{synth_1}[Full-data probability law]
The full-data law \(P\), for a nonnegative integer \(d\), is a probability measure on the finite record space
\[
(X,A,S(0),S(1),Y(0),Y(1),R)
\in \{0,\ldots,d-1\}\times\{0,1\}^{6}.
\]
Here \(X\) is the baseline category, \(A\) is the treatment arm, \(S(0)\) and \(S(1)\) are the potential surrogates, \(Y(0)\) and \(Y(1)\) are the potential outcomes, and \(R\) is the arrival flag. The baseline alphabet is empty when \(d=0\).
\end{definitionv}

The observed record \leanref{sym:O}{\ensuremath{O}} is \((X,A,S,R,RY)\). Its sampled copies \leanref{sym:O_i}{\ensuremath{O_i}} retain cell membership and arrival status for every individual; when \(R=1\), the final coordinate records the outcome. The canonical observed sample law \leanref{sym:\mathbf O_n}{\ensuremath{\mathbf O_n}}, defined in \cref{obj:synth_5}, is the corresponding product distribution.

\begin{assumptionv}{ass:surrogate-consistency}[Consistency of the surrogate]
The realized surrogate \(S\) equals the potential surrogate corresponding to treatment assignment \(A\):
\[
S=S(A)\qquad\text{almost surely under }P,
\]
where \(P\) is the full-data probability law.
\end{assumptionv}

\subsection{Sampling, randomization, and outcome arrival}

Outcome consistency completes the link between observed and potential outcomes.

\begin{assumptionv}{ass:outcome-consistency}[Consistency of the outcome]
The realized outcome \(Y\) equals the potential outcome corresponding to treatment assignment \(A\):
\[
Y=Y(A)\qquad\text{almost surely under }P,
\]
where \(P\) is the full-data probability law.
\end{assumptionv}

Independent and identically distributed sampling is imposed within each triangular-array row, allowing the underlying law and alphabet size to vary across experiments \citep{ZengBalakrishnanHanKennedy2026}. Outcome arrival is governed by the following conditional independence condition.

\begin{assumptionv}{ass:arrival-mar}[Conditionally independent outcome arrival]
Under \(P\), the full-data probability law, the outcome-arrival indicator \(R\) and realized outcome \(Y\) are conditionally independent given baseline \(X\), treatment assignment \(A\), and realized surrogate \(S\):
\[
R\perp\!\!\!\perp Y\mid(X,A,S).
\]
\end{assumptionv}

For later aggregation, we index treatment--baseline--surrogate cells explicitly.

\begin{definitionv}{synth_6}[Cell alphabet and membership]
The cell alphabet \(\mathcal J_d\), for a nonnegative integer \(d\), is
\[
\mathcal J_d
=
\{0,1\}\times\{x\in\mathbb Z:0\le x<d\}\times\{0,1\}.
\]
For a full-data record with baseline category \(X\), binary assignment \(A\), and realized surrogate \(S\), membership in a cell \(j=(a,x,s)\in\mathcal J_d\) is defined by
\[
A=a \quad\land\quad X=x \quad\land\quad S=s.
\]
\end{definitionv}

The next definition records each cell's population mass and its conditional arrival probability.

\begin{definitionv}{synth_16}[Cell probabilities \(p_{axs}(P)\) and \(\rho_{axs}(P)\)]
The full-membership probability of a cell is
\[
p_{axs}(P):=P(A=a,X=x,S=s),
\qquad
(a,x,s)\in\mathcal J_d
=\{0,1\}\times\{0,\ldots,d-1\}\times\{0,1\},
\]
where \(P\) is a full-data law. On every occupied cell, meaning \(p_{axs}(P)>0\), define its conditional outcome-arrival probability by
\[
\rho_{axs}(P)
:=
P(R=1\mid A=a,X=x,S=s)
=
\frac{P(A=a,X=x,S=s,R=1)}{p_{axs}(P)}.
\]
\end{definitionv}

The supplied arrival floor is imposed only on occupied cells.

\begin{assumptionv}{ass:occupied-cell-arrival}[Occupied-cell arrival floor]
For each cell
\[
(a,x,s)\in\mathcal J_d
=\{0,1\}\times\{0,\ldots,d-1\}\times\{0,1\},
\]
let \(p_{axs}(P)\), the full-membership cell probability, and \(\rho_{axs}(P)\), the conditional outcome-arrival probability on an occupied cell, be defined as in \cref{obj:synth_16}. The occupied-cell arrival condition is
\[
p_{axs}(P)>0
\quad\Longrightarrow\quad
\rho_{axs}(P)\ge q.
\]
\end{assumptionv}

The preceding randomization, consistency, and arrival restrictions define the full positive-arrival model.

\begin{definitionv}{def:unrestricted-arrival-model-class}[Arrival-floor model \(\mathcal M^{+}_{n,d,q}\)]
The class \(\mathcal M^{+}_{n,d,q}\) consists of full-data laws \(P\), for integers \(n,d\ge1\) and an arrival floor \(0<q\le1\), satisfying the following conditions:
\begin{itemize}
\item \textbf{(Randomized independence.)} \cref{obj:ass:randomized-independence}.
\item \textbf{(Balanced randomization.)} \cref{obj:ass:balanced-randomization}.
\item \textbf{(Surrogate consistency.)} \cref{obj:ass:surrogate-consistency}.
\item \textbf{(Outcome consistency.)} \cref{obj:ass:outcome-consistency}.
\item \textbf{(Arrival MAR.)} \cref{obj:ass:arrival-mar}.
\item \textbf{(Occupied-cell arrival.)} \cref{obj:ass:occupied-cell-arrival}.
\end{itemize}
The arrival floor ranges over the entire interval \(0<q\le1\). Independent, identically distributed sampling is imposed separately on the experiment through \cref{obj:ass:iid-sampling}.
\end{definitionv}

Outcome arrival missing at random given treatment, baseline, and surrogate equates the outcome distribution among arrivals with the full-cell outcome distribution whenever arrivals have positive probability \citep{Rubin1976,Robins1994,KallusMao2025}. The surrogate therefore serves as an observed conditioning variable for ascertainment: its realized value can capture variation in arrival associated with the assigned treatment and the outcome. Positive ascertainment within occupied cells is imposed next.

\begin{definitionv}{def:ate-functional}[Average treatment effect \(\tau(P)\)]
The average treatment effect is
\[
\tau(P):=E_P[Y(1)-Y(0)],
\]
where \(Y(1)\) and \(Y(0)\) are the potential outcomes under treatment and control.
\end{definitionv}

The occupied-cell arrival floor is specific to this analysis. It guarantees an outcome-arrival probability of at least the supplied value \(q\) in every cell represented in the population, while permitting both cell membership probabilities and arrival probabilities to vary across cells. This condition supplies positive arrived-cell mass wherever the population cell contribution can be nonzero. A prospective analysis must justify this lower bound; the stated tuning and guarantees apply when it is valid.

Together, the randomization, consistency, and arrival restrictions define the arrival-floor model \leanref{sym:\mathcal M^{+}_{n,d,q}}{\ensuremath{\mathcal M^{+}_{n,d,q}}} in \cref{obj:def:unrestricted-arrival-model-class}.

\begin{definitionv}{synth_19}[Treatment sign function]
Define the treatment sign function \(g:\{0,1\}\to\{-1,1\}\) by
\[
g(a):=2a-1,\qquad a\in\{0,1\}.
\]
Thus \(g(0)=-1\) and \(g(1)=1\).
\end{definitionv}

This family covers the full range of positive arrival floors, including complete outcome arrival at \(q=1\). Sampling determines the product experiment induced by each member of the family.

\subsection{Identification and squared-error risk}

The causal target is the population average treatment effect \leanref{sym:\tau(P)}{\ensuremath{\tau(P)}} in \cref{obj:def:ate-functional}.

\begin{definitionv}{def:cell-functional}[Signed cell functional \(\Phi(P)\)]
The total signed observed-cell functional is
\[
\Phi(P):=2\sum_{(a,x,s)\in\mathcal J_d}g(a)c_{axs}(P),
\]
where \(\mathcal J_d\) is the treatment–baseline–surrogate cell-index set and \(g(a)\) is the treatment sign function. On this index set, with \(x\in\{0,\ldots,d-1\}\), the full-membership cell probability \(p_{axs}(P)\) is defined in \cref{obj:synth_16}.

For a cell with positive arrived-cell probability,
\[
P(A=a,X=x,S=s,R=1)>0,
\]
define its arrived-cell outcome mean and weighted contribution by
\[
\mu_{axs}(P):=E_P[Y\mid A=a,X=x,S=s,R=1],
\qquad
c_{axs}(P):=p_{axs}(P)\mu_{axs}(P).
\]
When the arrived-cell probability is zero, including whenever \(p_{axs}(P)=0\), set
\[
\mu_{axs}(P):=0,
\qquad
c_{axs}(P):=0.
\]
Thus \(\Phi\) is defined on every full-data law.
\end{definitionv}

Binary potential outcomes place this target in \([-1,1]\). Its observed-data representation combines full-membership probabilities with outcome means among arrivals. The total signed cell functional \leanref{sym:\Phi(P)}{\ensuremath{\Phi(P)}}, defined in \cref{obj:def:cell-functional}, uses an arrived-cell outcome mean \leanref{sym:\mu_{axs}(P)}{\ensuremath{\mu_{axs}(P)}} and a membership-weighted contribution \leanref{sym:c_{axs}(P)}{\ensuremath{c_{axs}(P)}}, with explicit conventions for cells having zero arrived mass.

\begin{definitionv}{synth_5}[Observed sample and its canonical law]
The random observed sample is \(\mathbf O_n=(O_1,\ldots,O_n)\), where \(O_i=(X_i,A_i,S_i,R_i,R_iY_i)\). For nonnegative integers \(n,d\) and a full-data probability law \(P\) with \(d\) baseline categories, its canonical law is
\[
\mathcal L_P(\mathbf O_n):=\bigl(\mathcal L_P(X,A,S,R,RY)\bigr)^{\otimes n}.
\]
Here \(\mathcal L_P(X,A,S,R,RY)\) is the distribution under \(P\) of the observed record \((X,A,S,R,RY)\), whose last four coordinates are binary and whose final coordinate is the product of arrival and outcome. Thus \(\mathcal L_P(\mathbf O_n)\) is the law of \(n\) independent observed records with this common distribution.
\end{definitionv}

To distinguish the two roles in subsequent formulas, we write \(\mathbf O_n=(O_1,\ldots,O_n)\) for the random observed-record tuple and \(\mathcal L_P(\mathbf O_n)\) for its product law. The canonical-law display above specifies the latter distribution.

For every law in \cref{obj:def:unrestricted-arrival-model-class}, \cref{obj:lem:unrestricted-cell-identification} establishes \(\tau(P)=\Phi(P)\); its proof appears in the appendix. The representation assigns zero contribution to null cells and uses observable arrived-cell means on occupied cells. Full cell membership remains observed even when the outcome is missing, so the representation separates population cell weights from within-cell outcome ascertainment. Balanced assignment supplies the factor of two in the signed contrast.

We evaluate estimation over the entire arrival-floor model through the minimax squared-error risk \leanref{sym:\mathfrak R^{+}_{n,d,q}}{\ensuremath{\mathfrak R^{+}_{n,d,q}}} in \cref{obj:def:unrestricted-minimax-risk}. Its decision class \leanref{sym:\mathcal T_n}{\ensuremath{\mathcal T_n}} contains all parameter-independent bounded statistical decision rules, including a possibly randomized estimator \leanref{sym:\mathsf T}{\ensuremath{\mathsf T}}.

\begin{definitionv}{def:unrestricted-minimax-risk}[Unrestricted minimax squared-error risk]
The minimax squared-error risk \(\mathfrak R^{+}_{n,d,q}\), for \(n,d\in\mathbb N\) and \(q\in\mathbb R\), is defined by
\[
\mathfrak R^{+}_{n,d,q}
:=
\inf_{\mathsf T\in\mathcal T_n}
\sup_{P\in\mathcal M^{+}_{n,d,q}}
\int\!\int_{[-1,1]}
(h-\tau(P))^2\,
\mathsf T(o,dh)\,
\mathcal L_P(\mathbf O_n)(do).
\]
Here \(\mathcal M^{+}_{n,d,q}\) is the full-data model class in \cref{obj:def:unrestricted-arrival-model-class}, and \(\tau(P)\) is the population average treatment effect in \cref{obj:def:ate-functional}. The class \(\mathcal T_n\) consists of all parameter-independent Markov kernels from the observed sample space to \([-1,1]\), including deterministic estimators as Dirac kernels.

The random observed-record tuple \(\mathbf O_n\) and its product law under the sampling in \cref{obj:ass:iid-sampling} are
\[
\mathbf O_n=(O_1,\ldots,O_n),
\qquad
\mathcal L_P(\mathbf O_n)
=
\bigl(\mathcal L_P(O_1)\bigr)^{\otimes n},
\]
where each \(O_i\) is an observed record induced by \(P\). The risk integrates over both the observed sample and the kernel draw.

By \cref{obj:lem:randomized-kernel-barycenter}, the defining infimum equals the infimum over deterministic estimators taking values in \([-1,1]\). At \(q=1\), the same decision class \(\mathcal T_n\) defines the minimax risk over \(\mathcal M^{+}_{n,d,1}\), whose laws satisfy \(R=1\) almost surely.
\end{definitionv}

The inner integral averages over the procedure's conditional output distribution given the sample, and the outer integral averages over sampling. Parameter independence means that the procedure is common to all laws in the specified model. Under squared loss, replacing a randomized output by its conditional mean given the sample weakly reduces risk for each law, as stated in \cref{obj:lem:randomized-kernel-barycenter}. This comparison makes the all-procedure criterion equivalent to its deterministic counterpart.

\subsection{The rare-arrival experiment and comparison rate}

A second model family restricts the arrival floor relative to the effective sample-size scale. Its defining condition uses the logarithmic scale \leanref{sym:\ell}{\ensuremath{\ell}}, which equals \(\log(e+N)\) on the statistical parameter domain fixed above.

\begin{definitionv}{synth_7}[Total real logarithmic scale]
The logarithmic scale \(\ell\), for \(n\in\mathbb N\) and \(q\in\mathbb R\), is defined by
\[
\ell=
\begin{cases}
\log\bigl|\exp(1)+nq\bigr|, & \exp(1)+nq\ne 0,\\
0, & \exp(1)+nq=0,
\end{cases}
\]
where \(\log\) on positive arguments is the natural logarithm.
\end{definitionv}

The rare-arrival regime is selected by the following restriction on \(q\) and this logarithmic scale.

\begin{assumptionv}{ass:rare-arrival-slice}[Rare-arrival restriction]
The arrival floor \(q\) and logarithmic scale \(\ell\) satisfy
\[
q\ell^2\le \frac{1}{64}.
\]
\end{assumptionv}

The rare-arrival restriction is specific to this analysis and selects parameter triples for which the arrival floor is small relative to the squared logarithmic scale. It restricts the guaranteed floor; individual occupied cells may have larger arrival probabilities. The rare-arrival model \leanref{sym:\mathcal M_{n,d,q}}{\ensuremath{\mathcal M_{n,d,q}}} in \cref{obj:def:model-class} combines this parameter restriction with the preceding statistical conditions.

\begin{definitionv}{def:model-class}[Rare-arrival model \(\mathcal M_{n,d,q}\)]
\(\mathcal M_{n,d,q}\) is the class of full-data probability laws \(P\) satisfying the following conditions:
\begin{itemize}
\item \textbf{(Parameters.)} \(n,d\ge1\) and \(0<q\le1\).
\item \textbf{(Sampling.)} \cref{obj:ass:iid-sampling} holds.
\item \textbf{(Randomization.)} \cref{obj:ass:randomized-independence,obj:ass:balanced-randomization} hold.
\item \textbf{(Consistency.)} \cref{obj:ass:surrogate-consistency,obj:ass:outcome-consistency} hold.
\item \textbf{(Arrival.)} \cref{obj:ass:arrival-mar,obj:ass:occupied-cell-arrival} hold.
\item \textbf{(Rare-arrival restriction.)} \cref{obj:ass:rare-arrival-slice} holds.
\end{itemize}
\end{definitionv}

At parameter triples satisfying \cref{obj:ass:rare-arrival-slice}, this model uses the same full-data restrictions as \cref{obj:def:unrestricted-arrival-model-class}, together with the stated sampling condition. The observed experiment \leanref{sym:\mathcal E_{n,d,q}}{\ensuremath{\mathcal E_{n,d,q}}} in \cref{obj:def:observed-experiment} is explicitly indexed by this rare-arrival model.

\begin{definitionv}{def:observed-experiment}[Observed sample experiment]
The observed experiment \(\mathcal E_{n,d,q}\), for \(n,d\in\mathbb N\) and \(q\in\mathbb R\), is the family
\[
\mathcal E_{n,d,q}
:=
\bigl\{\mathcal L_P(\mathbf O_n):P\in\mathcal M_{n,d,q}\bigr\},
\]
where \(\mathcal M_{n,d,q}\) is the model class defined in \cref{obj:def:model-class}. Here \(\mathbf O_n\), the random observed-record tuple, and its law are specified by
\[
\mathbf O_n=(O_i)_{i=1}^{n},
\qquad
O_i=(X_i,A_i,S_i,R_i,R_iY_i),
\qquad
\mathcal L_P(\mathbf O_n)
=
\bigl(\mathcal L_P(X,A,S,R,RY)\bigr)^{\otimes n}.
\]
Thus the records are independent and identically distributed with the observed marginal law induced by \(P\). Each record retains the baseline category \(X\in\{0,\ldots,d-1\}\), assignment \(A\), realized surrogate \(S\), arrival indicator \(R\), and arrival-multiplied outcome \(RY\); the last four coordinates are binary. The record space carries the discrete sigma-algebra.
\end{definitionv}

The decision problem retains every sampled cell membership, including those whose outcomes have yet to arrive. Its minimax squared-error risk \leanref{sym:\mathfrak R_{n,d,q}}{\ensuremath{\mathfrak R_{n,d,q}}}, defined in \cref{obj:def:minimax-risk}, uses the same decision class and joint sample-and-procedure expectation as the unrestricted criterion.

\begin{definitionv}{def:minimax-risk}[Minimax squared-error risk]
The minimax squared-error risk \(\mathfrak R_{n,d,q}\), for \(n,d\in\mathbb N\) and \(q\in\mathbb R\), is
\[
\mathfrak R_{n,d,q}
:=
\inf_{\mathsf T\in\mathcal T_n}
\sup_{P\in\mathcal M_{n,d,q}}
\int\!\int_{[-1,1]}
(h-\tau(P))^2\,
\mathsf T(o,dh)\,
\mathcal L_P(\mathbf O_n)(do).
\]
Here \(\mathcal M_{n,d,q}\) is the full-data model class of \cref{obj:def:model-class}, and \(\tau(P)\) is the population average treatment effect of \cref{obj:def:ate-functional}. The observed sample \(\mathbf O_n=(O_i)_{i=1}^n\) denotes the random tuple of observed records. Its law, under \cref{obj:ass:iid-sampling}, is
\[
\mathcal L_P(\mathbf O_n)
=
\bigotimes_{i=1}^{n}\mathcal L_P(O_i).
\]
The decision class \(\mathcal T_n\) consists of all parameter-independent Markov kernels from the observed-sample space to \([-1,1]\), including deterministic estimators represented by Dirac kernels.

The notation \(E_P[(\mathsf T-\tau(P))^2]\) denotes the displayed joint sample-and-kernel risk. By \cref{obj:lem:randomized-kernel-barycenter}, the minimax value also equals the infimum of the maximal squared-error risk over deterministic estimators taking values in \([-1,1]\).
\end{definitionv}

The two risk definitions distinguish the parameter domains over which the subsequent comparisons are stated. For the rare-arrival analysis, the comparison-rate expression \leanref{sym:r_{n,d,q}}{\ensuremath{r_{n,d,q}}} in \cref{obj:def:frontier-rate} combines the effective sample-size scale with the number of baseline categories.

\begin{definitionv}{def:frontier-rate}[Frontier rate \(r_{n,d,q}\)]
The frontier rate is
\[
r_{n,d,q}
:=
\min\left\{
1,\,
\frac{1}{nq}
+
\left[\frac{d}{nq\log(e+nq)}\right]^2
\right\}.
\]
\end{definitionv}

In terms of the setup notation, the expression is \(\min\{1,N^{-1}+[d/(N\ell)]^2\}\). Its first term records the reciprocal effective sample-size scale, while its second term expresses the joint dependence on alphabet size and outcome ascertainment. The truncation at one matches the uniform squared-error bound of the estimator that always reports zero for a target in \([-1,1]\). The estimation results that follow compare the defined risks with these scales under their stated ascertainment regimes.


\section{Point estimation across ascertainment regimes}

The squared-error comparisons depend jointly on outcome ascertainment and the number of baseline categories. For every arrival floor \(0<q\le1\), the unrestricted minimax risk \leanref{sym:\mathfrak R^{+}_{n,d,q}}{\ensuremath{\mathfrak R^{+}_{n,d,q}}} of \cref{obj:def:unrestricted-minimax-risk} admits an upper envelope combining the mixed-count rate in \cref{obj:def:frontier-rate} with a dimension-free observed-outcome bound. The latter separates a sampling contribution of order \(n^{-1}\) from the squared ascertainment deficit \((1-q)^2\). We first specify the estimators supporting this envelope, then give matched comparisons for large alphabets, rare arrival, and a fixed arrival floor \(0<q<1/2\).

\subsection{An all-arrival envelope and its estimators}

The clipped observed-outcome arm contrast \leanref{sym:\widehat\tau^{\mathrm{HT}}_n}{\ensuremath{\widehat\tau^{\mathrm{HT}}_n}} uses balanced assignment to form a signed average of arrived outcomes. Its projection \leanref{sym:\Pi_{[-1,1]}}{\ensuremath{\Pi_{[-1,1]}}} keeps the estimate in the range of the population treatment effect. The following definition also fixes the clipping convention used by the mixed-count estimator.

\begin{definitionv}{def:complete-arrival-ht-estimator}[Clipped observed-outcome arm contrast]
The clipped observed-outcome arm contrast is the map
\[
\widehat\tau^{\mathrm{HT}}_n(\mathbf O_n)
:=
\Pi_{[-1,1]}\!\left\{
\frac{2}{n}\sum_{i=1}^n
(2A_i-1)\mathbf{1}\{R_i=1,\ R_iY_i=1\}
\right\},
\qquad n\ge1.
\]
Here \(n,d\) are nonnegative integers, and \(\mathbf O_n\), the observed sample, consists of records
\(O_i=(X_i,A_i,S_i,R_i,R_iY_i)\), with \(X_i\in\{0,\ldots,d-1\}\) and all four remaining coordinates binary; the baseline support is empty when \(d=0\). The notation \(R_iY_i\) denotes the stored arrived-outcome coordinate; the indicator equals one exactly when both this coordinate and the arrival coordinate \(R_i\) equal one. Projection and clipping coincide:
\[
\Pi_{[-1,1]}(u)=\operatorname{clip}_{[-1,1]}(u)
=\max\{-1,\min\{1,u\}\},
\qquad u\in\mathbb R.
\]
For \(n=0\), the map is defined to equal zero.
\end{definitionv}

The \leanref{S-2}{mixed-count construction} combines cell membership information with cell-specific arrived-outcome statistics. It uses the effective sample size \(N=nq\) and the logarithmic scale \(\ell=\log(e+N)\) in \cref{obj:synth_7}. Its polynomial degree \leanref{sym:k}{\ensuremath{k}} grows with this logarithmic scale, according to the following total tuning rule.

\begin{definitionv}{synth_8}[Polynomial degree tuning]
The degree \(k\), for \(n\in\mathbb N\) and \(q\in\mathbb R\), is defined by
\[
k=
\begin{cases}
\displaystyle \max\left\{0,\left\lfloor\frac{\log|e+nq|}{64}\right\rfloor\right\},
& e+nq\ne 0,\\
0, & e+nq=0.
\end{cases}
\]
\end{definitionv}

On the polynomial branch, the first-kind Chebyshev polynomial \leanref{sym:T_k}{\ensuremath{T_k}} determines the needle polynomial \leanref{sym:Q_k(t)}{\ensuremath{Q_k(t)}}. Here \leanref{sym:t}{\ensuremath{t}} denotes an arrived-count intensity, and the threshold \leanref{sym:B}{\ensuremath{B}} sets the scale of the construction. The next two definitions specify the polynomial and its coefficient sequence \leanref{sym:a_v}{\ensuremath{a_v}}, where \leanref{sym:v}{\ensuremath{v}} indexes the degree.

\begin{definitionv}{synth_18}[Chebyshev needle polynomial]
Put \(N=nq\), \(\ell=\log(e+N)\), \(k=\lfloor\ell/64\rfloor\), and \(B=256\ell\). Let \(T_k\) be the first-kind Chebyshev polynomial, characterized by \(T_k(\cos u)=\cos(ku)\). On the branch \(\ell\ge128\) and \(d\ge\sqrt N\), where \(k\ge2\), define
\[
Q_k(t):=
\begin{cases}
\displaystyle\frac{1-T_k(1-2t/B)}{2k^2t/B},&t>0,\\[6pt]
1,&t=0.
\end{cases}
\]
Here \(t\ge0\) is an arrived-count intensity, and the value at zero is the continuous polynomial extension of the displayed quotient. On every other branch, set \(Q_k(t):=1\) for every \(t\ge0\).
\end{definitionv}

\begin{definitionv}{synth_13}[Polynomial coefficient rule]
The coefficients \(a_v\), for \(n,d,v\in\mathbb N\) and \(q\in\mathbb R\), are defined by
\[
a_v=
\begin{cases}
[t^v]\bigl(1-Q_k(t)\bigr),
& 128\le \ell \ \text{and}\ \sqrt{N}\le d,\\
0,
& \text{otherwise}.
\end{cases}
\]
Here \([t^v]\) denotes the coefficient of \(t^v\), \(N=nq\) is the effective sample size, \(\ell\) is the logarithmic scale, and \(Q_k\) is the Chebyshev needle polynomial. On the indicated branch, this polynomial is the needle radius divided by \(2k^2\), multiplied by the polynomial quotient of
\[
1-T_k\!\left(1-\frac{2t}{\text{needle radius}}\right)
\]
upon division by \(t\), with the remainder discarded. The degree \(k\) and needle radius are the construction's values at \((n,q)\), and \(T_k\) is the first-kind Chebyshev polynomial.
\end{definitionv}

The estimator assigns observations to membership, pilot, and estimation streams. The pilot arrived count \leanref{sym:C'_j}{\ensuremath{C'_j}} selects the cell statistic, while the estimation arrived count \leanref{sym:C_j}{\ensuremath{C_j}} and arrived-one count \leanref{sym:U_j}{\ensuremath{U_j}} supply its outcome information. Their definitions retain the observed cell memberships and distinguish the two uses of arrived outcomes.

\begin{definitionv}{synth_17}[Arrived-outcome counts \(C'_j\) and \(C_j\)]
The pilot and estimation arrived-outcome counts are
\[
C'_j:=\sum_{i\in\mathcal I_{\mathrm{pil}}}
\mathbf1\{A_i=a,\ X_i=x,\ S_i=s,\ R_i=1\},
\qquad
C_j:=\sum_{i\in\mathcal I_{\mathrm{est}}}
\mathbf1\{A_i=a,\ X_i=x,\ S_i=s,\ R_i=1\},
\quad j=(a,x,s)\in\mathcal J_d.
\]
Here \(\mathbf O_n\), the random observed sample, consists of the records
\[
O_i=(X_i,A_i,S_i,R_i,R_iY_i),
\qquad
\mathbf O_n=(O_i)_{i=1}^n,
\qquad
\mathcal L_P(\mathbf O_n)=\bigl(\mathcal L_P(O_1)\bigr)^{\otimes n}.
\]
Given a prefix length \(K_0\le n\) and an assignment of its indices to membership, pilot, and estimation streams, \(\mathcal I_{\mathrm{pil}}\) and \(\mathcal I_{\mathrm{est}}\) denote the sets of prefix indices assigned to the pilot and estimation streams, respectively. The cell alphabet and observed-sample law are defined in \cref{obj:synth_6,obj:synth_5}; the prefix and stream assignments are those used in \cref{obj:def:mixed-count-estimator}.
\end{definitionv}

\begin{definitionv}{synth_10}[Estimation arrived-one count]
The estimation arrived-one count \(U_j\) is defined for any \(k,d\in\mathbb N\), collection of \(k\) observed records, assignment of each record to a stream in \(\{0,1,2\}\), and cell \(j=(a,x,s)\in\{0,1\}\times\{0,\ldots,d-1\}\times\{0,1\}\) by
\[
U_j
:=
\#\bigl\{\,0\le i<k:
i\text{ is assigned to stream }2,\ 
A_i=a,\ X_i=x,\ S_i=s,\ R_i=1,\ (RY)_i=1
\,\bigr\}.
\]
Here \(O_i=(X_i,A_i,S_i,R_i,(RY)_i)\) denotes the \(i\)th observed record, with baseline coordinate \(X_i\in\{0,\ldots,d-1\}\) and binary coordinates \(A_i,S_i,R_i,(RY)_i\).
\end{definitionv}

The selected cell statistic \leanref{sym:G_j}{\ensuremath{G_j}} uses a polynomial expression when the pilot count is small and an empirical arrived-outcome mean on the remaining cells. The associated factorial statistic \leanref{sym:H_j}{\ensuremath{H_j}} records the polynomial expression separately, and the elementary ratio \leanref{sym:D_j}{\ensuremath{D_j}} specifies the empirical-mean branch, including its value when the arrived count is zero.

\begin{definitionv}{synth_15}[Selected polynomial cell statistic]
The selected cell statistic \(G_j\), for natural-number parameters \(n,d\), a real parameter \(q\), a finite collection of observed records with baseline coordinate in \(\{0,\ldots,d-1\}\) and binary assignment, surrogate, arrival, and arrived-outcome coordinates, an assignment of each record to one of three streams, and a cell \(j\in\{0,1\}\times\{0,\ldots,d-1\}\times\{0,1\}\), is defined by
\[
G_j=
\begin{cases}
\displaystyle
\sum_{v=1}^{\max\{k-1,0\}}
a_v U_j
\prod_{t=0}^{v-2}\max\{C_j-1-t,0\},
&
128\le \ell,\quad \sqrt{N}\le d,\quad C'_j\le B/4,
\\[6pt]
\displaystyle U_j/C_j,
&
C_j>0
\ \text{and the preceding branch condition fails},
\\[4pt]
0,
&
C_j=0
\ \text{and the preceding branch condition fails}.
\end{cases}
\]
Here \(N=nq\) is the effective sample size, \(\ell\) is the construction's logarithmic scale, \(B=256\ell\) is its threshold, \(k\) is its polynomial degree, and \(a_v\) are its polynomial coefficients. The quantities \(C'_j\), \(C_j\), and \(U_j\) are, respectively, the pilot-stream arrived count, the estimation-stream arrived count, and the estimation-stream arrived count with outcome equal to one in cell \(j\). An empty sum equals zero and an empty product equals one.
\end{definitionv}

\begin{definitionv}{synth_14}[Factorial cell statistic]
The cell statistic \(H_j\), for a finite collection of observed records, assignments to streams \(0,1,2\), a cell \(j\in\{0,1\}\times\{0,\ldots,d-1\}\times\{0,1\}\), \(n,d\in\mathbb N\), and \(q\in\mathbb R\), is
\[
H_j=
\begin{cases}
\displaystyle
\sum_{v=1}^{\max\{k-1,0\}}
a_v U_j
\prod_{r=1}^{v-1}\max\{C_j-r,0\},
& \ell\ge128\ \text{and}\ \sqrt N\le d,\\[6pt]
0,&\text{otherwise}.
\end{cases}
\]
Here \(N=nq\) is the effective sample size, \(\ell\) is the logarithmic scale, and the degree parameter is
\[
k=\max\{\lfloor\ell/64\rfloor,0\}.
\]
The estimation-stream arrival count \(C_j\) and arrived-one count \(U_j\) are
\[
C_j=\sum_i
\mathbf 1\{\text{record }i\text{ is assigned to stream }2,\ 
(A_i,X_i,S_i)=j,\ R_i=1\},
\qquad
U_j=\sum_i
\mathbf 1\{\text{record }i\text{ is assigned to stream }2,\ 
(A_i,X_i,S_i)=j,\ R_i=1,\ (RY)_i=1\}.
\]
The coefficients associated with the Chebyshev polynomial \(Q_k\) are
\[
a_v=
\begin{cases}
[t^v]\bigl(1-Q_k(t)\bigr),
&\ell\ge128\ \text{and}\ \sqrt N\le d,\\
0,&\text{otherwise},
\end{cases}
\]
where \([t^v]\) denotes the coefficient of \(t^v\). Empty sums equal zero and empty products equal one.
\end{definitionv}

\begin{definitionv}{synth_12}[Elementary cell ratio statistic]
The elementary cell statistic \(D_j\), for any finite collection of observed records, any assignment of those records to three streams indexed by \(0,1,2\), and any cell \(j\in\{0,1\}\times\{0,\ldots,d-1\}\times\{0,1\}\), is
\[
D_j=
\begin{cases}
U_j/C_j, & C_j>0,\\
0, & C_j=0.
\end{cases}
\]
Here the records have baseline coordinate \(X\in\{0,\ldots,d-1\}\) and binary coordinates \(A,S,R,RY\), and the counts are
\[
\begin{aligned}
C_j&=\#\{\text{records assigned to stream }2:
                 (A,X,S)=j,\ R=1\},\\
U_j&=\#\{\text{records assigned to stream }2:
                 (A,X,S)=j,\ R=1,\ RY=1\}.
\end{aligned}
\]
\end{definitionv}

Cell weights use the membership count \leanref{sym:M_j}{\ensuremath{M_j}}, which includes observations with either arrival status. The weight \leanref{sym:W_j}{\ensuremath{W_j}} divides this count by the stream intensity \leanref{sym:m}{\ensuremath{m}}. The following definitions therefore separate estimation of cell mass from estimation of its outcome mean.

\begin{definitionv}{synth_9}[Membership stream cell count]
The membership count \(M_j\), for any nonnegative integers \(k,d\), a collection of \(k\) observed records \(O_i\) with baseline coordinates in \(\{0,\ldots,d-1\}\) and binary assignment, surrogate, arrival, and recorded-outcome coordinates, an assignment of each record to one of the streams \(0,1,2\), and a cell \(j=(a,x,s)\in\{0,1\}\times\{0,\ldots,d-1\}\times\{0,1\}\), is
\[
M_j
:=
\#\bigl\{i:0\le i<k,\ 
i\text{ is assigned to stream }0,\ 
A_i=a,\ X_i=x,\ S_i=s\bigr\}.
\]
\end{definitionv}

\begin{definitionv}{synth_11}[Cell membership weight]
The cell-weight statistic \(W_j\) is defined for nonnegative integers \(k,d,n\), a collection of \(k\) observed records \(O_i\) with baseline coordinate \(X_i\in\{0,\ldots,d-1\}\) and binary coordinates \(A_i,S_i,R_i,RY_i\), an assignment of each record to a stream labelled \(0,1,2\), and a cell
\[
j\in\{0,1\}\times\{0,\ldots,d-1\}\times\{0,1\}.
\]
Using the membership count and stream intensity
\[
M_j=\#\{i\in\{0,\ldots,k-1\}: O_i\text{ is assigned to stream }0,\ (A_i,X_i,S_i)=j\},
\qquad
m=\frac{n}{6},
\]
define
\[
W_j=\frac{M_j}{m}.
\]
The count \(M_j\) is interpreted as a real number in this quotient, and division by zero is assigned the value zero, so \(W_j=0\) when \(n=0\). The index sets \(\{0,\ldots,k-1\}\) and \(\{0,\ldots,d-1\}\) are empty when their respective sizes are zero.
\end{definitionv}

These ingredients define the mixed-count estimator \leanref{sym:\widehat\tau^{\mathrm{mix}}_{n,d,q}}{\ensuremath{\widehat\tau^{\mathrm{mix}}_{n,d,q}}}. A Poisson prefix length \leanref{sym:K_0}{\ensuremath{K_0}} and independent stream assignments produce an auxiliary estimate; averaging over this randomization yields a deterministic function of the observed sample.

\begin{algorithmv}{def:mixed-count-estimator}[Mixed-count estimator $\widehat\tau^{\mathrm{mix}}_{n,d,q}$]
Given the observed sample \(\mathbf O_n\), parameters \(n,d,q\), effective sample size \(N=nq\), logarithmic scale \(\ell\), and the declared cell-statistic and polynomial formulas, define the estimator by the following steps.
\begin{enumerate}
\item Draw an auxiliary prefix length independently of the data:
\[
K_0\sim\operatorname{Poisson}(n/2).
\]
When \(K_0\le n\), independently assign each of the first \(K_0\) observations to the membership, pilot, or estimation stream, each with probability \(1/3\). Denote the prefix and stream-assignment randomization by \(\omega\).

\item On the event \(K_0>n\), set
\[
\widetilde\tau(\mathbf O_n;\omega)=0.
\]
When \(K_0\le n\), form the declared statistics \(M_j,C'_j,C_j,U_j,W_j,D_j\). If \(N\le1\), set
\[
\widetilde\tau(\mathbf O_n;\omega)=0.
\]

\item When \(K_0\le n\), \(N>1\), and either \(\ell<128\) or \(d<\sqrt N\), set
\[
\widetilde\tau(\mathbf O_n;\omega)
=
\operatorname{clip}_{[-1,1]}
\left\{
2\sum_{j=(a,x,s)}g(a)W_jD_j
\right\}.
\]
On this branch, the Chebyshev quotient and its coefficients are left unevaluated.

\item When \(K_0\le n\), \(N>1\), \(\ell\ge128\), and \(d\ge\sqrt N\), form the declared total quantities \(Q_k,a_v,H_j,G_j\) and set
\[
\widetilde\tau(\mathbf O_n;\omega)
=
\operatorname{clip}_{[-1,1]}
\left\{
2\sum_{j=(a,x,s)}g(a)W_jG_j
\right\}.
\]

\item Output the conditional average
\[
\widehat\tau^{\mathrm{mix}}_{n,d,q}(\mathbf O_n)
:=
E_\omega\!\left[
\widetilde\tau(\mathbf O_n;\omega)\mid\mathbf O_n
\right].
\]
This average is a finite sum over prefix lengths \(K_0\le n\) and their stream assignments, together with the zero contribution from the Poisson tail. It defines a total, deterministic, sample-measurable estimator whose squared-error risk is no larger than that of the auxiliary randomized output.
\end{enumerate}
\end{algorithmv}

The resulting deterministic risk envelope \leanref{sym:\mathfrak U_{n,d,q}}{\ensuremath{\mathfrak U_{n,d,q}}} is specified branch by branch in the next result. Its certificate uses the effective stream size \leanref{sym:N_0}{\ensuremath{N_0}}, equal to \(mq\). Together with the observed-outcome contrast, it provides two upper bounds applicable throughout the unrestricted arrival domain.

\begin{theoremv}{thm:unrestricted-near-complete-arrival-envelope}[Near-complete arrival risk envelope]
There exists a universal constant \(\overline C>0\) with the following properties. For every \(n,d\ge1\) and \(0<q\le1\), let \(r_{n,d,q}\) be the rate in \cref{obj:def:frontier-rate}, let \(\ell\) be the logarithmic scale in \cref{obj:synth_7}, let \(k\) be the polynomial degree specified in \cref{obj:synth_8}, and let \(B\) be the threshold specified in \cref{obj:synth_18}. Define the effective sample size \(N\), stream intensity \(m\), effective stream size \(N_0\), and prefix-tail constant \(\gamma_0\) by
\[
N=nq,\qquad m=\frac n6,\qquad N_0=mq,\qquad
\gamma_0=\log 2-\frac12.
\]
Define the deterministic squared-error envelope \(\mathfrak U_{n,d,q}\) as follows. When \(N\le1\), set
\[
\mathfrak U_{n,d,q}=1.
\]
When \(N>1\) and the estimator in \cref{obj:def:mixed-count-estimator} uses its elementary branch, set
\[
\mathfrak U_{n,d,q}
=
\min\left\{
4,\,
\left(\frac{8d}{eN_0}\right)^2
+4\left(\frac4{N_0}+\frac1m\right)
+4e^{-n\gamma_0}
\right\}.
\]
When \(N>1\) and that estimator uses its polynomial branch, define the certificate quantities
\[
b_{n,d,q}
=
2\left(\frac{4dB}{N_0k^2}+3e^{-16\ell}\right),
\]
\[
v_{n,d,q}
=
8\left(\frac4{N_0}+\frac1m\right)
+64d\left(e^{\ell/8}+1\right)
\left(\frac{B^2}{N_0^2}+\frac{B}{mN_0}\right)
+32e^{-32\ell}\left(1+\frac1m\right),
\]
and set
\[
\mathfrak U_{n,d,q}
=
\min\left\{
4,\,
b_{n,d,q}^{\,2}+v_{n,d,q}+4e^{-n\gamma_0}
\right\}.
\]

For every full-data law \(P\in\mathcal M^{+}_{n,d,q}\), the unrestricted-arrival model class in \cref{obj:def:unrestricted-arrival-model-class}, suppose that \(\mathbf O_n\), the observed sample, satisfies \cref{obj:ass:iid-sampling} under the canonical product law of \(n\) observed records induced by \(P\). The mixed-count estimator of \cref{obj:def:mixed-count-estimator} satisfies
\[
E_P\!\left[
\left\{
\widehat\tau^{\mathrm{mix}}_{n,d,q}(\mathbf O_n)-\tau(P)
\right\}^{2}
\right]
\le \mathfrak U_{n,d,q}
\le \overline C\,r_{n,d,q},
\]
where \(\tau(P)\) is the average treatment effect in \cref{obj:def:ate-functional}. It also satisfies
\[
E_P\!\left[
\left\{
\widehat\tau^{\mathrm{mix}}_{n,d,q}(\mathbf O_n)-\tau(P)
\right\}^{2}
\right]
\le
E_P\!\left[
E_\omega\!\left[
\left\{
\widetilde\tau(\mathbf O_n;\omega)-\tau(P)
\right\}^{2}
\right]
\right],
\]
where \(\widetilde\tau(\mathbf O_n;\omega)\) is the auxiliary randomized output in \cref{obj:def:mixed-count-estimator}, and the inner expectation uses its specified auxiliary randomization, including the zero output when the Poisson prefix length exceeds \(n\).

For every \(n,d\ge1\) and \(0<q\le1\), the minimax squared-error risk \(\mathfrak R^{+}_{n,d,q}\) over the parameter-independent bounded estimator kernels in \cref{obj:def:unrestricted-minimax-risk} obeys the combined envelope
\[
\frac{1}{256n}
\le
\mathfrak R^{+}_{n,d,q}
\le
\min\left\{
1,\,
\overline C\,r_{n,d,q},\,
\frac4n+(1-q)^2
\right\}.
\]
The clipped observed-outcome contrast \(\widehat\tau^{\mathrm{HT}}_n\) of \cref{obj:def:complete-arrival-ht-estimator} satisfies
\[
\sup_{P\in\mathcal M^{+}_{n,d,q}}
E_P\!\left[
\left\{
\widehat\tau^{\mathrm{HT}}_n(\mathbf O_n)-\tau(P)
\right\}^{2}
\right]
\le
\frac4n+(1-q)^2,
\]
where each expectation is taken under the canonical product law of \(n\) observed records induced by \(P\).
\end{theoremv}

The universal comparison constant \leanref{sym:\overline C}{\ensuremath{\overline C}} is independent of sample size, alphabet size, and arrival floor. The two upper bounds express different ways to control squared error: the mixed-count bound accounts for category complexity through \(r_{n,d,q}\), whereas the observed-outcome bound accounts for missing outcomes through the ascertainment deficit. As \(q\) approaches one, the latter supplies a dimension-free guarantee. Conditional averaging in \cref{obj:def:mixed-count-estimator} preserves the mixed-count guarantee while eliminating the auxiliary randomization from the reported estimate.

\subsection{Large alphabets and the uniform near-complete criterion}

For alphabets at least quadratic in sample size, the sampling and ascertainment-deficit terms also give a matching lower bound. The next result establishes their joint order with universal constants throughout the stated large-alphabet region.

\begin{theoremv}{thm:unrestricted-large-alphabet-deficit-lower}[Unrestricted risk with large alphabets]
Let \(n,d\) be nonnegative integers and \(q\) a real number satisfying:
\begin{itemize}
\item \textbf{(Sample size.)} \(n\ge 1\).
\item \textbf{(Alphabet size.)} \(d\ge 1024n^2\).
\item \textbf{(Arrival floor.)} \(0<q\le 1\).
\end{itemize}
Then \(\mathfrak R^{+}_{n,d,q}\), the unrestricted minimax squared-error risk defined in \cref{obj:def:unrestricted-minimax-risk}, satisfies
\[
\max\left\{\frac{1}{256n},\frac{(1-q)^2}{1024}\right\}
\le \mathfrak R^{+}_{n,d,q}
\le \frac{4}{n}+(1-q)^2.
\]
These bounds also give the envelope
\(\frac{1}{2048}\left\{\frac{1}{n}+(1-q)^2\right\}
\le \mathfrak R^{+}_{n,d,q}
\le 4\left\{\frac{1}{n}+(1-q)^2\right\}\).
\end{theoremv}

The comparison identifies two sources of difficulty on the same scale: ordinary sampling variation and the remaining ascertainment deficit. When \((1-q)^2\) exceeds \(n^{-1}\), it determines the risk order in this region; when it is at most of order \(n^{-1}\), sampling variation determines that order. Thus the dimension-free observed-outcome guarantee in \cref{obj:thm:unrestricted-near-complete-arrival-envelope} is sharp up to universal constants for \(d\ge1024n^2\).

A guarantee uniform over every alphabet size must cover this large-alphabet region. Combining its lower bound with the all-arrival upper envelope yields an exact sequence-level criterion for parametric squared-error risk.

\begin{propositionv}{prop:unrestricted-uniform-near-complete-elbow}[Uniform near-complete arrival elbow]
Let \((q_n)_{n\ge0}\) satisfy \(0<q_n\le1\) for every \(n\), and let \(\mathfrak R^{+}_{n,d,q}\) be the unrestricted-arrival minimax squared-error risk defined in \cref{obj:def:unrestricted-minimax-risk}. For every integer \(n\ge1\),
\[
n\mathfrak R^{+}_{n,d,q_n}\le 4+n(1-q_n)^2
\qquad\text{for every integer }d\ge1.
\]
Moreover, for every real number
\[
b<\frac{1}{2048}\bigl\{1+n(1-q_n)^2\bigr\},
\]
there exists an integer \(d\ge1\) such that
\[
b<n\mathfrak R^{+}_{n,d,q_n}.
\]

The following two conditions are equivalent:
\begin{itemize}
\item \textbf{(Arrival rate.)} There exists a real constant \(M\ge0\) such that, for all sufficiently large integers \(n\),
\[
1-q_n\le \frac{M}{\sqrt n}.
\]
\item \textbf{(Uniform risk bound.)} There exists a real constant \(D\) such that, for all sufficiently large integers \(n\) and every integer \(d\ge1\),
\[
\mathfrak R^{+}_{n,d,q_n}\le \frac{D}{n}.
\]
\end{itemize}

For every pair of integers \(n,d\ge1\),
\[
\frac{1}{256n}\le \mathfrak R^{+}_{n,d,q_n}.
\]
If the sequence \(\bigl(\sqrt n(1-q_n)\bigr)_{n\ge0}\) is unbounded above, then for every real constant \(C\), there are infinitely many integers \(n\) such that
\[
C<n\mathfrak R^{+}_{n,\,1024n^2,\,q_n}.
\]
\end{propositionv}

The quantifier over all \(d\) is central to this criterion. A deficit of order \(n^{-1/2}\) is necessary and sufficient for an eventual \(n^{-1}\) upper bound holding simultaneously over every alphabet size. The explicit sequence \(d_n=1024n^2\) witnesses the converse when the rescaled deficit is unbounded. For any fixed deficit multiplier, the sufficient window also gives the following numerical comparison along arbitrary alphabet-size sequences.

\begin{propositionv}{prop:unrestricted-near-complete-phase-boundary}[Near-complete arrival phase boundary]
Let \(\mathfrak R^{+}_{n,d,q}\) be the unrestricted minimax squared-error risk defined in \cref{obj:def:unrestricted-minimax-risk}. Fix a real constant \(M\) and sequences \((d_n)_{n\ge0}\) and \((q_n)_{n\ge0}\) satisfying:
\begin{itemize}
\item \textbf{(Constant.)} \(M\ge0\).
\item \textbf{(Alphabet size.)} \(d_n\) is a natural number with \(d_n\ge1\) for every \(n\).
\item \textbf{(Arrival floor.)} \(q_n\in(0,1]\) for every \(n\).
\item \textbf{(Arrival deficit.)} For all sufficiently large \(n\),
\[
1-q_n\le \frac{M}{\sqrt n}.
\]
\end{itemize}
Then, for all sufficiently large \(n\),
\[
\frac{1}{256n}
\le \mathfrak R^{+}_{n,d_n,q_n}
\le \frac{4+M^2}{n}.
\]
\end{propositionv}

At complete ascertainment, the deficit is zero. The following endpoint specialization records the observed-outcome contrast as a direct witness to the dimension-free parametric benchmark.

\begin{theoremv}{thm:unrestricted-complete-arrival-frontier}[Complete-arrival minimax risk frontier]
For every sample size \(n\ge 1\) and alphabet size \(d\ge 1\), let \(\mathcal M^{+}_{n,d,1}\) be the complete-arrival model class of \cref{obj:def:unrestricted-arrival-model-class}, and let \(\mathbf O_n\) consist of \(n\) independent, identically distributed observed records induced by the full-data law \(P\). Let \(\mathfrak R^{+}_{n,d,1}\) be the minimax squared-error risk of \cref{obj:def:unrestricted-minimax-risk}, and let \(\widehat\tau^{\mathrm{HT}}_n\) be the clipped observed-outcome arm contrast of \cref{obj:def:complete-arrival-ht-estimator}. Writing \(\tau(P)\) for the population average treatment effect, we have
\[
\frac{1}{256n}
\le \mathfrak R^{+}_{n,d,1}
\le
\sup_{P\in\mathcal M^{+}_{n,d,1}}
E_P\!\left[
\left\{\widehat\tau^{\mathrm{HT}}_n(\mathbf O_n)-\tau(P)\right\}^{2}
\right]
\le \frac{4}{n}.
\]
Consequently, at the complete-arrival endpoint \(q=1\), the minimax squared-error risk has order \(n^{-1}\) with universal constants, uniformly over \(d\ge 1\).
\end{theoremv}

Complete arrival makes the observed-outcome arm contrast sufficient to attain the minimax order for every alphabet size. The sampling lower bound remains of the same order, so the endpoint comparison is uniform even when the alphabet grows arbitrarily quickly. Its sequence formulation follows directly.

\begin{propositionv}{prop:unrestricted-complete-arrival-phase-boundary}[Complete-arrival point-risk boundary]
There exist real constants \(0<c\le C\) such that, for every sequence of alphabet sizes \(d_n\ge1\), the unrestricted minimax squared-error risk \(\mathfrak R^{+}_{n,d_n,1}\) defined in \cref{obj:def:unrestricted-minimax-risk} at complete arrival \(q=1\) satisfies, for all sufficiently large \(n\),
\[
\frac{c}{n}\le \mathfrak R^{+}_{n,d_n,1}\le \frac{C}{n}.
\]
The comparison constants are independent of the chosen alphabet-size sequence.
\end{propositionv}

\subsection{Rare arrival and effective sample size}

Under \cref{obj:ass:rare-arrival-slice}, the rate in \cref{obj:def:frontier-rate} gives a matched comparison for the rare-arrival risk of \cref{obj:def:minimax-risk}. This comparison retains the full observed experiment, including memberships whose outcomes have yet to arrive, and applies to the entire decision class in that risk definition.

\begin{theoremv}{thm:matched-minimax-frontier}[Matched minimax frontier]
There exist universal constants \(0<\underline c\le\overline C<\infty\) such that, for every \(n,d\in\mathbb N\) and \(q\in\mathbb R\) satisfying
\begin{itemize}
\item \textbf{(Sample size and dimension.)} \(n\ge1\) and \(d\ge1\).
\item \textbf{(Arrival parameter.)} \(0<q\le1\).
\item \textbf{(Rare-arrival restriction.)} \cref{obj:ass:rare-arrival-slice} holds.
\end{itemize}
the minimax squared-error risk \(\mathfrak R_{n,d,q}\) defined in \cref{obj:def:minimax-risk} and the frontier rate \(r_{n,d,q}\) defined in \cref{obj:def:frontier-rate} satisfy
\[
\underline c\,r_{n,d,q}
\le \mathfrak R_{n,d,q}
\le \overline C\,r_{n,d,q}.
\]
\end{theoremv}

The universal lower constant \leanref{sym:\underline c}{\ensuremath{\underline c}} makes the two components of the rate jointly informative about minimax difficulty. The reciprocal effective sample size captures the sampling scale, while the squared category term captures the cost of estimating the treatment contrast across many cells:
\[
r_{n,d,q}
=
\min\left\{1,\frac1N+\left(\frac{d}{N\ell}\right)^2\right\}.
\]
At fixed alphabet size, the category term becomes asymptotically negligible as \(N\) grows. The next result also establishes this effective-sample-size order within the submodel having a constant surrogate in both treatment arms.

\begin{propositionv}{prop:fixed-d-effective-sample-reduction}[Fixed-alphabet effective sample reduction]
Fix an integer \(d\ge 1\). Let \(n(t)\) and \(q(t)\) be sequences indexed by \(t\in\mathbb N\) satisfying:
\begin{itemize}
\item \textbf{(Sample size and arrival.)} For every \(t\), \(n(t)\in\mathbb N\), \(n(t)\ge 1\), and \(0<q(t)\le 1\).
\item \textbf{(Rare-arrival slice.)} For every \(t\), the pair \((n(t),q(t))\) satisfies \cref{obj:ass:rare-arrival-slice}.
\item \textbf{(Effective sample size.)} The effective sample size
\[
N(t)=n(t)q(t)
\]
satisfies \(N(t)\to\infty\) as \(t\to\infty\).
\end{itemize}
Then the frontier rate and minimax squared-error risk of \cref{obj:def:frontier-rate,obj:def:minimax-risk} satisfy
\[
r_{n(t),d,q(t)}\asymp N(t)^{-1},
\qquad
\mathfrak R_{n(t),d,q(t)}\asymp N(t)^{-1}
\qquad (t\to\infty).
\]
The minimax squared-error risk with the same estimator class and loss as in \cref{obj:def:minimax-risk}, taken over the restricted full-data laws
\[
\bigl\{P\in\mathcal M_{n(t),d,q(t)}:
P\bigl(S(0)=0,\ S(1)=0\bigr)=1\bigr\},
\]
is also of order \(N(t)^{-1}\) as \(t\to\infty\). Each order comparison means eventual upper and lower bounds by positive constant multiples of \(N(t)^{-1}\).
\end{propositionv}

Allowing the alphabet to grow separates consistency from effective-sample-size precision. The category term vanishes when \(d\) is small relative to \(N\log(e+N)\), and it remains at most of order \(N^{-1}\) when \(d\) is at most of order \(\sqrt N\log(e+N)\). The matched comparison turns these rate calculations into exact equivalences for the minimax risk.

\begin{propositionv}{prop:phase-boundaries}[Rare-arrival minimax phase boundaries]
Let \((n_t)_{t\in\mathbb N}\) and \((d_t)_{t\in\mathbb N}\) be sequences of natural numbers, and let \((q_t)_{t\in\mathbb N}\) be a sequence of real numbers. Define the effective sample size by
\[
N_t=n_tq_t.
\]
Suppose that:
\begin{itemize}
\item \textbf{(Positive sizes.)} \(n_t\ge 1\) and \(d_t\ge 1\) for every \(t\in\mathbb N\).
\item \textbf{(Arrival probabilities.)} \(0<q_t\le 1\) for every \(t\in\mathbb N\).
\item \textbf{(Rare-arrival slice.)} \((n_t,q_t)\) satisfies \cref{obj:ass:rare-arrival-slice} for every \(t\in\mathbb N\).
\item \textbf{(Effective sample growth.)} \(N_t\to\infty\) as \(t\to\infty\).
\end{itemize}
Then the minimax squared-error risk \(\mathfrak R_{n_t,d_t,q_t}\) of \cref{obj:def:minimax-risk} satisfies both equivalences:
\[
\mathfrak R_{n_t,d_t,q_t}\longrightarrow 0
\quad\Longleftrightarrow\quad
d_t=o\!\left(N_t\log(e+N_t)\right),
\]
\[
\mathfrak R_{n_t,d_t,q_t}=\Theta\!\left(N_t^{-1}\right)
\quad\Longleftrightarrow\quad
d_t=O\!\left(\sqrt{N_t}\log(e+N_t)\right),
\]
where all limits and asymptotic comparisons are taken as \(t\to\infty\).
\end{propositionv}

These thresholds quantify how many baseline categories can be accommodated at a given outcome-maturity floor. They apply along sequences satisfying the rare-arrival restriction and growing effective sample size. In particular, they distinguish the broader category-growth range permitting consistency from the smaller range preserving squared-error risk of order \(N_t^{-1}\).

\subsection{A fixed strict-interior arrival floor}

A separate matched comparison applies to the unrestricted model when the arrival floor is fixed at \(0<q<1/2\). Its lower comparison uses the discrete observational lower bound of \citet[Theorem~2, Section~3.2, p.~9, and its proof in Section~C.5, pp.~39--44]{ZengBalakrishnanHanKennedy2026}, including the restriction setting the control regression to zero. The resulting comparison holds for every \(n,d\ge1\), with constants depending on the fixed arrival floor.

\begin{theoremv}{thm:unrestricted-fixed-q-minimax-frontier}[Unrestricted fixed-arrival-floor minimax frontier]*
Fix an arrival floor \(q\) with \(0<q<1/2\). There exist real constants \(0<\underline c(q)\le \overline C(q)<\infty\), depending only on \(q\), such that for every pair of integers \(n,d\ge1\), the unrestricted minimax squared-error risk \(\mathfrak R^{+}_{n,d,q}\) defined in \cref{obj:def:unrestricted-minimax-risk} satisfies
\[
\underline c(q)
\min\left\{1,\frac1n+\left(\frac{d}{n\log(e+n)}\right)^2\right\}
\le \mathfrak R^{+}_{n,d,q}
\le \overline C(q)
\min\left\{1,\frac1n+\left(\frac{d}{n\log(e+n)}\right)^2\right\}.
\]
With the same constants, the frontier rate \(r_{n,d,q}\) defined in \cref{obj:def:frontier-rate} satisfies
\[
\underline c(q)\,r_{n,d,q}
\le \mathfrak R^{+}_{n,d,q}
\le \overline C(q)\,r_{n,d,q}.
\]
\end{theoremv}
% CAUSALSMITH-CITED-SCOPE-BEGIN thm:unrestricted-fixed-q-minimax-frontier
\verificationfootnotetext{\textbf{Formalization scope.} This result uses the cited conclusion \citep{ZengBalakrishnanHanKennedy2026} (Theorem 2, Section 3.2, p. 9, together with its proof in Section C.5, pp. 39--44, which sets the control regression to zero). The cited source proof is not formalized here: Lean verifies its use and all remaining steps. If that cited conclusion is false, the portions of this result that depend on it are not certified.}
% CAUSALSMITH-CITED-SCOPE-END thm:unrestricted-fixed-q-minimax-frontier

For a fixed arrival floor, effective sample size is proportional to sample size, so the comparison can be written using \(n\) and \(\log(e+n)\). The lower constant \leanref{sym:\underline c(q)}{\ensuremath{\underline c(q)}} and upper constant \leanref{sym:\overline C(q)}{\ensuremath{\overline C(q)}} may depend on that floor. This is a second logarithmic matched regime: \cref{obj:thm:matched-minimax-frontier} gives universal constants under the rare-arrival restriction, while the present result gives a comparison over the unrestricted model at each fixed \(q\in(0,1/2)\).

The fixed-floor comparison yields corresponding sequence-level thresholds. The next statement records both the exact consistency condition and the necessary and sufficient category-growth order for parametric squared-error risk.

\begin{propositionv}{prop:unrestricted-fixed-q-phase-boundaries}[Unrestricted fixed-arrival-floor phase boundaries]*
Fix \(0<q<1/2\), where \(q\) is the arrival floor, and let \(\mathfrak R^{+}_{n,d,q}\) be the unrestricted-arrival minimax squared-error risk defined in \cref{obj:def:unrestricted-minimax-risk}. For every sequence of integers \(d_n\ge1\),
\[
\mathfrak R^{+}_{n,d_n,q}\longrightarrow0
\quad\Longleftrightarrow\quad
d_n=o\!\left(n\log(e+n)\right),
\qquad n\to\infty.
\]

There exists a real constant \(c>0\) such that, for every real \(M>0\), there exists a real constant \(C\ge c\) with the following property: for every sequence of integers \(d_n\ge1\), if
\[
d_n\le M\sqrt n\,\log(e+n)
\qquad\text{for all sufficiently large }n,
\]
then
\[
\frac{c}{n}\le\mathfrak R^{+}_{n,d_n,q}\le\frac{C}{n}
\qquad\text{for all sufficiently large }n.
\]
Here \(c\) may depend on the fixed \(q\), and \(C\) may depend on \(q\) and \(M\); the thresholds for sufficiently large \(n\) may depend on the sequence.

Conversely, for every sequence of integers \(d_n\ge1\), if there exist real constants \(0<c_*\le C_*\) such that
\[
\frac{c_*}{n}\le\mathfrak R^{+}_{n,d_n,q}\le\frac{C_*}{n}
\qquad\text{for all sufficiently large }n,
\]
then
\[
d_n=O\!\left(\sqrt n\,\log(e+n)\right),
\qquad n\to\infty.
\]
\end{propositionv}
% CAUSALSMITH-CITED-SCOPE-BEGIN prop:unrestricted-fixed-q-phase-boundaries
\verificationfootnotetext{\textbf{Formalization scope.} This result uses the cited conclusion \citep{ZengBalakrishnanHanKennedy2026} (Theorem 2, Section 3.2, p. 9, together with its proof in Section C.5, pp. 39--44, which sets the control regression to zero). The cited source proof is not formalized here: Lean verifies its use and all remaining steps. If that cited conclusion is false, the portions of this result that depend on it are not certified.}
% CAUSALSMITH-CITED-SCOPE-END prop:unrestricted-fixed-q-phase-boundaries

The fixed-floor thresholds express consistency on the \(n\log(e+n)\) category scale and parametric precision on the \(\sqrt n\log(e+n)\) scale. For a prescribed category-growth multiplier, the eventual upper constant may depend on that multiplier and on \(q\). This dependence differs from the universal constants in the rare-arrival comparison and from the explicit deficit-dependent constant in the near-complete window.

The following table collects the comparisons while keeping the universal upper envelope, the matched large-alphabet region, and the two logarithmic matched regimes separate. Here \(N=nq\), \(\ell=\log(e+N)\), and \(r_{n,d,q}\) is the rate in \cref{obj:def:frontier-rate}.

\begin{table}[htbp]
\centering
\small
\caption{Squared-error comparisons across ascertainment regimes}
\begin{tabular}{p{0.22\linewidth}p{0.31\linewidth}p{0.18\linewidth}p{0.21\linewidth}}
\hline
Domain & Risk comparison & Constants & Supporting results \\
\hline
\(n,d\ge1,\ 0<q\le1\)
&
\(\displaystyle \frac1{256n}\le\mathfrak R^+_{n,d,q}
\le\min\{1,\overline C r_{n,d,q},4/n+(1-q)^2\}\)
&
Universal
&
\cref{obj:thm:unrestricted-near-complete-arrival-envelope}
\\[4pt]
\(d\ge1024n^2,\ 0<q\le1\)
&
\(\mathfrak R^+_{n,d,q}\asymp n^{-1}+(1-q)^2\)
&
Universal
&
\cref{obj:thm:unrestricted-large-alphabet-deficit-lower}
\\[4pt]
\(1-q_n\le M/\sqrt n\) eventually, all \(d\ge1\)
&
\(\displaystyle \frac1{256n}\le\mathfrak R^+_{n,d,q_n}\le\frac{4+M^2}{n}\); uniform \(n^{-1}\) order iff \(1-q_n=O(n^{-1/2})\)
&
Explicit dependence on \(M\)
&
\cref{obj:prop:unrestricted-uniform-near-complete-elbow,obj:prop:unrestricted-near-complete-phase-boundary}
\\[4pt]
\(q=1,\ d\ge1\)
&
\(\mathfrak R^+_{n,d,1}\asymp n^{-1}\), uniformly in \(d\)
&
Universal
&
\cref{obj:thm:unrestricted-complete-arrival-frontier,obj:prop:unrestricted-complete-arrival-phase-boundary}
\\[4pt]
Rare-arrival restriction
&
\(\displaystyle \mathfrak R_{n,d,q}\asymp
\min\{1,N^{-1}+[d/(N\ell)]^2\}\)
&
Universal
&
\cref{obj:thm:matched-minimax-frontier,obj:prop:phase-boundaries}
\\[4pt]
Fixed \(0<q<1/2\)
&
\(\displaystyle \mathfrak R^+_{n,d,q}\asymp
\min\{1,n^{-1}+[d/\{n\log(e+n)\}]^2\}\)
&
Depend on fixed \(q\)
&
\cref{obj:thm:unrestricted-fixed-q-minimax-frontier,obj:prop:unrestricted-fixed-q-phase-boundaries}
\\
\hline
\end{tabular}
\end{table}


\section{A constructive estimator}\label{sec:constructive-estimator}

The risk comparisons admit an explicit statistical construction. The mixed-count estimator in \cref{obj:def:mixed-count-estimator} combines information about cell membership with information about arrived outcomes, assigning these two tasks to separate parts of each cell contribution. Its elementary branch uses empirical cell means; its polynomial branch modifies the contribution of cells with few arrivals. We collect the ingredients by their statistical role before describing the resulting estimate and its risk guarantee.

\subsection{Mechanism of the mixed-count correction}

The canonical formulas are collected once in \cref{obj:def:mixed-count-estimator} and its immediate ingredients \cref{obj:synth_7,obj:synth_8,obj:synth_9,obj:synth_10,obj:synth_11,obj:synth_12,obj:synth_13,obj:synth_14,obj:synth_15,obj:synth_17,obj:synth_18}. Their statistical roles are simple. Membership-labelled records estimate the population cell weights, while pilot and estimation labels separate branch selection from estimation of each arrived-cell mean. An ordinary ratio assigns zero to a cell with no arrived outcome. Across many rare cells, those zero-count events create the category-dependent bias that appears in the elementary certificate.

On the large-alphabet branch, the polynomial extension of \(\{1-Q_k(t)\}/t\) supplies the reciprocal correction implicit in a cell mean. Marked factorial counts estimate its monomials without discarding the association between an outcome and its cell: for cell intensity \(t_j\) and arrived-cell mean \(\mu_j\), the factorial statistic satisfies \(E(H_j)=\mu_j\{1-Q_k(t_j)\}\). Thus \(Q_k\) controls the residual bias, and the bound on \(t_j|Q_k(t_j)|\) controls its membership-weighted aggregate. With degree \(k\) proportional to \(\log(e+N)\), this changes the aggregate bias scale from the elementary scarce-cell scale to \(d/\{N\log(e+N)\}\), while a pilot count sends well-observed cells back to the stable empirical ratio. Independent membership, pilot, and estimation streams justify the comparison calculation; the capped-prefix transfer and conditional averaging turn that comparison into the deterministic fixed-sample estimator.

The construction is explicit on every parameter triple. Its polynomial branch activates at \(\log(e+nq)\ge128\), equivalently \(nq\ge e^{128}-e\). Exact evaluation averages \(3^{k_0}\) assignments at each retained prefix length \(k_0\); the limitations section discusses this computational scale separately from the attainment guarantee.

\subsection{Risk control and the deterministic envelope}

The construction supplies a squared-error guarantee throughout the unrestricted arrival class in \cref{obj:def:unrestricted-arrival-model-class}. Its deterministic envelope \(\mathfrak U_{n,d,q}\) depends on the sample size, alphabet size, and arrival floor. The following result gives the envelope explicitly, compares it with the rate in \cref{obj:def:frontier-rate}, and records the risk improvement from conditional averaging.

\begin{theoremv}{thm:unrestricted-mixed-count-upper}[Unrestricted mixed-count risk bound]
There exists a universal constant \(0<\overline C<\infty\) such that, for every sample size \(n\), alphabet size \(d\), arrival floor \(q\), and full-data law \(P\) satisfying
\begin{itemize}
\item \textbf{(Domain.)} \(n,d\ge1\) and \(0<q\le1\).
\item \textbf{(Model.)} \(P\in\mathcal M^{+}_{n,d,q}\), the model class defined in \cref{obj:def:unrestricted-arrival-model-class}.
\item \textbf{(Sampling.)} The observed sample \(\mathbf O_n\) has the canonical product law induced by \(P\), with its coordinate observations satisfying \cref{obj:ass:iid-sampling}.
\end{itemize}
the deterministic estimator \(\widehat\tau^{\mathrm{mix}}_{n,d,q}\) of \cref{obj:def:mixed-count-estimator} satisfies
\[
E_P\!\left[
 \bigl(\widehat\tau^{\mathrm{mix}}_{n,d,q}(\mathbf O_n)-\tau(P)\bigr)^2
\right]
\le \mathfrak U_{n,d,q}
\le \overline C\,r_{n,d,q},
\]
where \(\tau(P)\) is the average treatment effect defined in \cref{obj:def:ate-functional}, \(r_{n,d,q}\) is the rate defined in \cref{obj:def:frontier-rate}, and \(\mathfrak U_{n,d,q}\) is the deterministic risk envelope specified below.

Set
\[
N=nq,\qquad m=\frac n6,\qquad N_0=mq=\frac N6,
\qquad \gamma_0=\log 2-\frac12.
\]
Use the logarithmic scale \(\ell\) defined in \cref{obj:synth_7}, the polynomial degree \(k\) defined in \cref{obj:synth_8}, and the threshold \(B\) defined in \cref{obj:synth_18}. Define the certificate quantities
\[
b_{n,d,q}
=2\left(\frac{4dB}{N_0k^2}+3e^{-16\ell}\right),
\]
\[
v_{n,d,q}
=8\left(\frac4{N_0}+\frac1m\right)
+64d\left(e^{\ell/8}+1\right)
 \left(\frac{B^2}{N_0^2}+\frac{B}{mN_0}\right)
+32e^{-32\ell}\left(1+\frac1m\right).
\]
For \(N\le1\), set \(\mathfrak U_{n,d,q}=1\). For \(N>1\) on the elementary branch of \cref{obj:def:mixed-count-estimator}, set
\[
\mathfrak U_{n,d,q}
=\min\left\{
4,\,
\left(\frac{8d}{eN_0}\right)^2
+4\left(\frac4{N_0}+\frac1m\right)
+4e^{-n\gamma_0}
\right\}.
\]
For \(N>1\) on its polynomial branch, set
\[
\mathfrak U_{n,d,q}
=\min\left\{
4,\,
b_{n,d,q}^{\,2}+v_{n,d,q}+4e^{-n\gamma_0}
\right\}.
\]

Moreover, the estimator's conditional average over the auxiliary randomization satisfies
\[
E_P\!\left[
 \bigl(\widehat\tau^{\mathrm{mix}}_{n,d,q}(\mathbf O_n)-\tau(P)\bigr)^2
\right]
\le
E_P\!\left[
 E_{\omega}\!\left[
  \bigl(\widetilde\tau(\mathbf O_n;\omega)-\tau(P)\bigr)^2
  \,\middle|\,\mathbf O_n
 \right]
\right],
\]
where \(\omega\) is the auxiliary prefix and stream-assignment randomization from \cref{obj:def:mixed-count-estimator}, and \(\widetilde\tau(\mathbf O_n;\omega)\) is its output, including the zero output when the Poisson prefix length exceeds \(n\).
\end{theoremv}

The bound separates effective-sample-size error from category complexity through
\[
r_{n,d,q}
=\min\left\{1,\frac1N+
\left(\frac{d}{N\ell}\right)^2\right\},
\]
as defined in \cref{obj:def:frontier-rate}. The polynomial certificate reflects the mechanism of the construction: its leading squared contribution involves \(dB/(N_0k^2)\), which scales as \(d/(N\ell)\) on the polynomial branch. Membership weights and arrived-outcome statistics thus enter a common aggregate guarantee even when many cells have few arrived observations. The final inequality in \cref{obj:thm:unrestricted-mixed-count-upper} also shows why the conditional average is the reported estimator: averaging the bounded auxiliary outputs preserves their risk guarantee and can reduce squared-error risk.

The detailed cellwise moment bounds and their aggregation are collected in \cref{obj:thm:unrestricted-uniform-mixed-count-certificate}. For the subsequent interval construction, the essential feature is that \(\mathfrak U_{n,d,q}\) is a known, deterministic upper bound on squared-error risk, uniform over the stated model class. It supplies the radius \(\sqrt{\mathfrak U_{n,d,q}/\alpha}\), where \(\alpha\in(0,1)\) is the miscoverage level, used by the clipped interval in \cref{obj:def:risk-envelope-interval}.


\section{Honest intervals under rare outcome arrival}

The deterministic squared-error envelope in \cref{obj:thm:unrestricted-mixed-count-upper} also supports uncertainty quantification for the average treatment effect. The relevant criterion combines uniform coverage over the rare-arrival model with maximal expected interval length. Fix a miscoverage level \leanref{sym:\alpha}{\ensuremath{\alpha\in(0,1)}}, so that the required coverage is \(1-\alpha\). To include procedures with auxiliary randomization, both coverage and expected length are evaluated under the joint distribution of the observed sample and the reported endpoints. The following definition specifies this decision problem.

\begin{definitionv}{def:interval-length-risk}[Minimax expected interval length]
The minimax maximal expected interval length, \(\mathfrak L_{n,d,q,\alpha}\), is defined for \(n,d\in\mathbb N\) and \(q,\alpha\in\mathbb R\) as follows. Let \(\mathcal M_{n,d,q}\) be the model class in \cref{obj:def:model-class}, let \(\mathbf O_n\) be the observed sample with law \(\mathcal L_P(\mathbf O_n)\) from \cref{obj:def:observed-experiment}, and let \(\tau(P)\) be the average treatment effect in \cref{obj:def:ate-functional}.

The class \(\mathcal C_{n,d,q,\alpha}\) consists of parameter-independent Markov kernels \(\mathsf I\) from the observed-sample space to endpoint pairs \((l,u)\) satisfying \(-1\le l\le u\le1\), with uniform joint sample-and-randomization coverage: for every \(P\in\mathcal M_{n,d,q}\),
\[
\int\!\int
\mathbf1\{l\le\tau(P)\le u\}\,
\mathsf I(o,d(l,u))\,
\mathcal L_P(\mathbf O_n)(do)
\ge 1-\alpha.
\]
The realized connected interval is \([l,u]\), with length \(u-l\). Define
\[
\mathfrak L_{n,d,q,\alpha}
:=
\inf_{\mathsf I\in\mathcal C_{n,d,q,\alpha}}
\sup_{P\in\mathcal M_{n,d,q}}
\int\!\int (u-l)\,
\mathsf I(o,d(l,u))\,
\mathcal L_P(\mathbf O_n)(do),
\]
where the infimum and supremum are real-valued, with the convention that an infimum or supremum over an empty set equals zero. In particular, \(\mathfrak L_{n,d,q,\alpha}=0\) when \(\mathcal C_{n,d,q,\alpha}\) is empty. Both coverage and expected length integrate over the sample and the kernel draw. Deterministic connected-interval procedures are included through Dirac kernels.
\end{definitionv}

Uniform honesty requires the coverage inequality separately for every law in \cref{obj:def:model-class}. For a randomized endpoint procedure \leanref{sym:\mathsf I}{\ensuremath{\mathsf I}}, this requirement averages over its auxiliary draw as well as sampling variation. The same averaging governs length. Thus the admissible class \leanref{sym:\mathcal C_{n,d,q,\alpha}}{\ensuremath{\mathcal C_{n,d,q,\alpha}}} treats deterministic and randomized procedures within one criterion, and the minimax length \leanref{sym:\mathfrak L_{n,d,q,\alpha}}{\ensuremath{\mathfrak L_{n,d,q,\alpha}}} measures the smallest attainable worst-case expected width at the specified coverage.

A concrete honest interval uses the mixed-count estimator \leanref{sym:\widehat\tau^{\mathrm{mix}}_{n,d,q}}{\ensuremath{\widehat\tau^{\mathrm{mix}}_{n,d,q}}} from \cref{obj:def:mixed-count-estimator} and its deterministic risk envelope \leanref{sym:\mathfrak U_{n,d,q}}{\ensuremath{\mathfrak U_{n,d,q}}} in \cref{obj:thm:unrestricted-mixed-count-upper}. Both are calibrated to the supplied arrival floor \(q\). The envelope determines a radius before observing the sample; intersecting the resulting interval with the target range gives the following construction. Its coverage guarantee therefore requires the supplied floor to be no larger than every occupied-cell arrival probability.

\begin{definitionv}{def:risk-envelope-interval}[Risk-envelope interval \(I^{\mathrm{mix}}_\alpha\)]
The risk-envelope interval is
\[
I^{\mathrm{mix}}_\alpha
:=
\left[
\widehat\tau^{\mathrm{mix}}_{n,d,q}
-
\sqrt{\frac{\mathfrak U_{n,d,q}}{\alpha}},
\,
\widehat\tau^{\mathrm{mix}}_{n,d,q}
+
\sqrt{\frac{\mathfrak U_{n,d,q}}{\alpha}}
\right]\cap[-1,1],
\]
where \(\widehat\tau^{\mathrm{mix}}_{n,d,q}\) is the mixed-count estimator and \(\mathfrak U_{n,d,q}\) is its explicit deterministic squared-error envelope.
\end{definitionv}

The interval \leanref{sym:I^{\mathrm{mix}}_\alpha}{\ensuremath{I^{\mathrm{mix}}_\alpha}} is deterministic conditional on the observed sample and therefore enters \cref{obj:def:interval-length-risk} through a Dirac endpoint rule. Its radius converts uniform squared-error control into coverage at least \(1-\alpha\). The next result compares its maximal expected length with the best uniformly honest connected interval, using the comparison rate \leanref{sym:r_{n,d,q}}{\ensuremath{r_{n,d,q}}} defined in \cref{obj:def:frontier-rate}.

\begin{theoremv}{thm:connected-interval-frontier}[Connected interval length frontier]
For every miscoverage level \(\alpha\in(0,1)\), there exist constants \(0<\underline c_\alpha\le\overline C_\alpha<\infty\) such that the following holds for every \(n,d\in\mathbb N\) and \(q\in\mathbb R\) satisfying:
\begin{itemize}
\item \textbf{(Sample size and dimension.)} \(n\ge1\) and \(d\ge1\).
\item \textbf{(Arrival floor.)} \(0<q\le1\).
\item \textbf{(Rare-arrival slice.)} The restriction in \cref{obj:ass:rare-arrival-slice} holds.
\end{itemize}
Let \(r_{n,d,q}\) be the frontier rate in \cref{obj:def:frontier-rate}, let \(\mathfrak L_{n,d,q,\alpha}\) be the minimax expected interval length in \cref{obj:def:interval-length-risk}, and let \(I^{\mathrm{mix}}_\alpha\) be the interval in \cref{obj:def:risk-envelope-interval}. Over the model class \(\mathcal M_{n,d,q}\) in \cref{obj:def:model-class},
\[
\underline c_\alpha\sqrt{r_{n,d,q}}
\le \mathfrak L_{n,d,q,\alpha}
\le \sup_{P\in\mathcal M_{n,d,q}}
E_P\!\left[\operatorname{length}(I^{\mathrm{mix}}_\alpha)\right]
\le \overline C_\alpha\sqrt{r_{n,d,q}}.
\]
Here expectations and probabilities are taken under the law of the independent observed sample induced by \(P\), and interval length is defined by
\[
\operatorname{length}([l,u])=u-l.
\]
Moreover, for every \(P\in\mathcal M_{n,d,q}\), the average treatment effect \(\tau(P)\) from \cref{obj:def:ate-functional} satisfies
\[
\Pr_P\!\left\{\tau(P)\in I^{\mathrm{mix}}_\alpha\right\}\ge1-\alpha.
\]
\end{theoremv}

The upper bound reflects the square-root conversion from squared-error risk to an interval radius: the deterministic envelope controls coverage, while clipping bounds the realized length by the width of the untruncated interval. The converse connects statistical indistinguishability with the geometry of a connected interval. Coverage transfers between competing laws in the joint sample-and-endpoint distribution, and a realized interval that meets several separated target neighborhoods must span their separation. Working with this joint distribution retains the auxiliary endpoint draw throughout both the coverage comparison and the length calculation. The parameter-independent decision-rule comparison in \cref{obj:lem:randomized-kernel-tv-comparison} supplies the corresponding control when an endpoint procedure is appended to the observed experiment.

For each fixed coverage level \(1-\alpha\in(0,1)\), \cref{obj:thm:connected-interval-frontier} establishes expected-length order \(\sqrt{r_{n,d,q}}\), with constants that may depend on \(\alpha\), uniformly over the sample sizes, dimensions, and arrival floors satisfying its conditions. Alongside the squared-error order in \cref{obj:thm:matched-minimax-frontier}, this gives an uncertainty-quantification interpretation of the same comparison rate: its square root is the attainable and necessary order of maximal expected width for uniformly honest connected intervals. In particular, along sequences satisfying \cref{obj:ass:rare-arrival-slice}, the optimal expected length tends to zero exactly when the rate in \cref{obj:def:frontier-rate} tends to zero.


\section{An observational extension through synthetic assignment}

The \leanref{S-4}{discrete observational experiment} connects adjustment for many covariate categories to the randomized experiment with missing outcomes. Its target is a baseline-weighted regression contrast under a fixed strict positivity margin. An independent synthetic assignment converts observational records into randomized records whose outcomes arrive when the synthetic and observational treatments agree. This construction preserves the contrast and transfers the mixed-count estimator to observational data.

We begin by specifying the observational law and its occupied-category positivity condition.
\begin{definitionv}{def:zeng-discrete-ate-class}[Observational law class \(\mathcal D_{d,\epsilon}\)]
The class \(\mathcal D_{d,\epsilon}\), for \(0<\epsilon<1/2\), consists of one-record laws \(P_Z\) of the baseline category, observational treatment, and outcome
\[
(X,B^{\mathrm{obs}},Y)\in\{1,\ldots,d\}\times\{0,1\}^2
\]
with baseline probabilities
\[
p_x=P_Z(X=x)
\]
and, on each occupied category \(p_x>0\), conditional treatment probabilities and outcome means
\[
\pi_x=P_Z(B^{\mathrm{obs}}=1\mid X=x),
\qquad
\mu_{bx}=E_{P_Z}[Y\mid X=x,B^{\mathrm{obs}}=b],
\]
satisfying
\[
\epsilon\le\pi_x\le1-\epsilon,
\qquad
0\le\mu_{bx}\le1
\quad (b\in\{0,1\}).
\]
The conditional quantities \(\pi_x\) and \(\mu_{bx}\) are specified on occupied categories.
\end{definitionv}

For a law \leanref{sym:P_Z}{\ensuremath{P_Z}} in the class \leanref{sym:\mathcal D_{d,\epsilon}}{\ensuremath{\mathcal D_{d,\epsilon}}}, both values of the observational treatment \leanref{sym:B^{\mathrm{obs}}}{\ensuremath{B^{\mathrm{obs}}}} have positive conditional probability in every occupied category. The margin bounds these probabilities away from zero while allowing arbitrary baseline probabilities, including null categories. Consequently, both conditional outcome means are defined wherever the baseline distribution assigns positive mass.

The estimand aggregates the within-category regression contrasts using that baseline distribution.
\begin{definitionv}{def:zeng-ate-functional}[Observed-law contrast $\theta_Z$]
The total observed-law functional \(\theta_Z\) is defined, for every \(P_Z\in\mathcal D_{d,\epsilon}\), by
\[
\theta_Z(P_Z)
:=
\sum_{x:p_x>0}p_x(\mu_{1x}-\mu_{0x}).
\]
Null baseline categories contribute zero.
\end{definitionv}

The contrast \leanref{sym:\theta_Z(P_Z)}{\ensuremath{\theta_Z(P_Z)}} is a functional of the observational law. Under outcome consistency, \(Y=Y(B^{\mathrm{obs}})\), and conditional exchangeability,
\[
\{Y(0),Y(1)\}\perp\!\!\!\perp B^{\mathrm{obs}}\mid X,
\]
positivity identifies each occupied-category regression mean with the corresponding conditional potential-outcome mean. The contrast then equals the population average treatment effect \(E[Y(1)-Y(0)]\), giving the usual causal interpretation of covariate adjustment \citep{Rosenbaum1983}.

Let \leanref{sym:\mathbf Z_n}{\ensuremath{\mathbf Z_n}} denote the sample of \(n\) independent observational records. The decision problem allows estimation procedures to use independent auxiliary randomness as well as the sample.
\begin{definitionv}{def:zeng-minimax-risk}[Observational minimax risk \(\mathfrak R^Z_{n,d,\epsilon}\)]
The minimax squared-error risk for the observational contrast is
\[
\mathfrak R^Z_{n,d,\epsilon}
:=\inf_{\widehat\theta}
\sup_{P_Z\in\mathcal D_{d,\epsilon}}
E_{P_Z}\!\left[
(\widehat\theta(\mathbf Z_n)-\theta_Z(P_Z))^2
\right],
\]
where \(\mathbf Z_n\) denotes the independent, identically distributed sample of \(n\) observational records with law \(P_Z\). The infimum ranges over all possibly randomized estimators measurable with respect to \(\mathbf Z_n\) and auxiliary randomness independent of \(P_Z\); the expectation integrates both sources of randomness.
\end{definitionv}

The observational minimax risk \leanref{sym:\mathfrak R^Z_{n,d,\epsilon}}{\ensuremath{\mathfrak R^Z_{n,d,\epsilon}}} therefore measures the smallest worst-case squared error over the full decision class in \cref{obj:def:zeng-minimax-risk}. To connect this problem to the randomized construction, we present together the transformation in \cref{obj:lem:synthetic-randomization-kernel} and the estimation procedure in \cref{obj:thm:zeng-fixed-positivity-polynomial-frontier}.

\paragraph{Synthetic assignment and observational estimation.}
For an observational record \((x,b,y)\), draw a fair coin \(U\), independently of the record, and define
\[
o(x,b,y;u)
:=
\bigl(x-1,u,0,1\{b=u\},1\{b=u\}y\bigr),
\qquad u\in\{0,1\}.
\]
The five coordinates are the randomized baseline, assignment, surrogate, arrival indicator, and arrived outcome. For a sample realization \(\mathbf Z_n=((x_i,b_i,y_i))_{i=1}^n\), use independent fair coins \(U_1,\ldots,U_n\) and set
\[
o_i(u_i):=o(x_i,b_i,y_i;u_i).
\]
The transformation uses the observed record and a known coin distribution, so its rule is independent of the unknown observational law.

The induced randomized full-data law \(P_Z^*\) in \cref{obj:lem:synthetic-randomization-kernel} can be specified through a latent representation. Draw the baseline with probabilities \(p_x\). Conditional on an occupied category \(x\), draw \(B^{\mathrm{obs}}\) with success probability \(\pi_x\), and draw independent binary potential outcomes with means \(\mu_{0x}\) and \(\mu_{1x}\), independently of \(B^{\mathrm{obs}}\). Draw an independent fair assignment \(A\), retain \(X-1\) as the randomized baseline, and set
\[
S(0)=S(1)=0,\qquad
Y=Y(A),\qquad
R=1\{B^{\mathrm{obs}}=A\}.
\]
This representation reproduces the synthetic observed-record distribution. Its arm-specific arrival probabilities are
\[
P_Z^*(R=1\mid X-1=x-1,A=1,S=0)=\pi_x,
\qquad
P_Z^*(R=1\mid X-1=x-1,A=0,S=0)=1-\pi_x,
\]
and its target satisfies
\[
\tau(P_Z^*)=\theta_Z(P_Z).
\]
Both arrival probabilities are at least \(\epsilon\). Independence of the latent observational treatment and potential outcomes conditional on the baseline supplies the missing-at-random condition in this constructed law. The causal interpretation of the original observational contrast continues to rest on consistency and conditional exchangeability.

Apply the mixed-count estimator
\leanref{sym:\widehat\tau^{\mathrm{mix}}_{n,d,q}}{\ensuremath{\widehat\tau^{\mathrm{mix}}_{n,d,\epsilon}}}
from \cref{obj:def:mixed-count-estimator} to the synthetic sample. Averaging over all fair-coin assignments defines the deterministic observational estimator
\leanref{sym:\widehat\theta^{\mathrm{poly}}_{n,d,\epsilon}}{\ensuremath{\widehat\theta^{\mathrm{poly}}_{n,d,\epsilon}}}:
\[
\widehat\theta^{\mathrm{poly}}_{n,d,\epsilon}(\mathbf Z_n)
:=
\frac{1}{2^n}
\sum_{(u_1,\ldots,u_n)\in\{0,1\}^n}
\widehat\tau^{\mathrm{mix}}_{n,d,\epsilon}
\bigl(o_1(u_1),\ldots,o_n(u_n)\bigr).
\]
This finite average defines the procedure for every observational sample. Conditional averaging weakly reduces squared-error risk relative to retaining the auxiliary coin draws.

The resulting risk bound is uniform over sample size and alphabet size. At each fixed strict positivity margin, it also yields a matching all-estimator comparison.
\begin{theoremv}{thm:zeng-fixed-positivity-polynomial-frontier}[Fixed positivity polynomial frontier]*
There exists a universal constant \(\overline C>0\) such that the following holds for every pair of integers \(n,d\ge1\) and every strict positivity margin \(0<\epsilon<1/2\). Let \(\mathcal D_{d,\epsilon}\) be the observational class in \cref{obj:def:zeng-discrete-ate-class}, and let \(\mathbf Z_n=((X_i,B_i^{\mathrm{obs}},Y_i))_{i=1}^n\) be an independent sample from \(P_Z\in\mathcal D_{d,\epsilon}\). Write the observational baseline categories as \(1,\ldots,d\). For a realization \(\mathbf Z_n=((x_i,b_i,y_i))_{i=1}^n\), define the synthetic records, with baseline coordinates indexed from zero, by
\[
o_i(u_i)
=
\bigl(x_i-1,u_i,0,1\{b_i=u_i\},1\{b_i=u_i\}y_i\bigr),
\qquad u_i\in\{0,1\},
\]
and define the deterministic estimator
\[
\widehat\theta^{\mathrm{poly}}_{n,d,\epsilon}(\mathbf Z_n)
=
\frac{1}{2^n}
\sum_{(u_1,\ldots,u_n)\in\{0,1\}^n}
\widehat\tau^{\mathrm{mix}}_{n,d,\epsilon}
\bigl(o_1(u_1),\ldots,o_n(u_n)\bigr),
\]
where \(\widehat\tau^{\mathrm{mix}}_{n,d,\epsilon}\) is the mixed-count estimator in \cref{obj:def:mixed-count-estimator}. Then
\[
\sup_{P_Z\in\mathcal D_{d,\epsilon}}
E_{P_Z}\!\left[
\left(
\widehat\theta^{\mathrm{poly}}_{n,d,\epsilon}(\mathbf Z_n)
-\theta_Z(P_Z)
\right)^2
\right]
\le
\overline C\,r_{n,d,\epsilon},
\]
where \(\theta_Z(P_Z)\) is the observational contrast in \cref{obj:def:zeng-ate-functional} and \(r_{n,d,\epsilon}\) is the rate in \cref{obj:def:frontier-rate}.

Moreover, for each fixed \(0<\epsilon<1/2\), there exist constants \(0<c_\epsilon\le C_\epsilon<\infty\) such that, for every pair of integers \(n,d\ge1\), the minimax risk in \cref{obj:def:zeng-minimax-risk} satisfies
\[
c_\epsilon
\min\left\{
1,\frac1n+
\left[\frac{d}{n\log(e+n)}\right]^2
\right\}
\le
\mathfrak R^Z_{n,d,\epsilon}
\le
C_\epsilon
\min\left\{
1,\frac1n+
\left[\frac{d}{n\log(e+n)}\right]^2
\right\}.
\]
\end{theoremv}
% CAUSALSMITH-CITED-SCOPE-BEGIN thm:zeng-fixed-positivity-polynomial-frontier
\verificationfootnotetext{\textbf{Formalization scope.} This result uses the cited conclusion \citep{ZengBalakrishnanHanKennedy2026} (Theorem 2, Section 3.2, p. 9, together with its proof in Section C.5, pp. 39--44, which sets the control regression to zero). The cited source proof is not formalized here: Lean verifies its use and all remaining steps. If that cited conclusion is false, the portions of this result that depend on it are not certified.}
% CAUSALSMITH-CITED-SCOPE-END thm:zeng-fixed-positivity-polynomial-frontier

The statistical mechanism is the conversion of observational treatment imbalance into selective outcome arrival under independent assignment. Positivity supplies a common arrival floor, and the target-preserving transformation makes the randomized MAR risk certificate applicable. Averaging over the synthetic assignments then removes their auxiliary variation. For fixed \(\epsilon\), the two terms in the comparison describe sampling error and the cost of adjusting for many categories; the truncation captures the bounded scale of the contrast.

Under consistency and conditional exchangeability, \cref{obj:thm:zeng-fixed-positivity-polynomial-frontier} therefore characterizes minimax squared-error risk for the observational average treatment effect. As \(n\) tends to infinity at a fixed strict positivity margin, uniform mean-square consistency is attainable exactly when
\[
d=o\{n\log(e+n)\}.
\]
The constants in the two-sided comparison may depend on that margin, while the constructive bound retains the dependence on \(\epsilon\) through the rate in \cref{obj:def:frontier-rate}.

The observational lower-bound input comes from \citet[Theorem~2, Section~3.2, p.~9, together with its proof in Section~C.5, pp.~39--44]{ZengBalakrishnanHanKennedy2026}. Their argument sets the control regression to zero on occupied categories and supplies the lower order
\[
\frac1n+\left(\frac{d}{n\log n}\right)^2
\]
for fixed strict positivity, sufficiently large \(n\), and \(1\le d\le c_0n\log n\), where \(c_0>0\) is a constant from that result. The comparison in \cref{obj:thm:zeng-fixed-positivity-polynomial-frontier} extends the risk characterization to every pair \(n,d\ge1\), with a total estimator and the bounded-risk truncation.

Finally, the same construction fixes the direction of the experiment comparison supporting the fixed-arrival analysis. At the fixed arrival floor \(q=\epsilon\in(0,1/2)\), the induced laws belong to the randomized class in \cref{obj:def:unrestricted-arrival-model-class}. Any estimator for their randomized observed samples can be applied to observational records by appending independent synthetic assignments, with the same target and risk. The observational sample is therefore at least as informative as its synthetic image. Taking worst-case risk over the containing randomized class gives a risk at least as large as the observational minimax risk, so the published observational lower bound transfers in this direction to the fixed-arrival lower comparison in \cref{obj:thm:unrestricted-fixed-q-minimax-frontier}. The formal target-preserving map is given in \cref{obj:lem:synthetic-randomization-kernel}.


\section{Discussion and extensions}

The precision comparisons distinguish three features of the observed experiment: the availability of primary outcomes, the number of covariate categories across which ascertainment must be accommodated, and the remaining ascertainment deficit. Their relative importance changes with the arrival regime. The resulting interpretation is a comparison of attainable precision over the stated model classes, with each regime retaining its own conditions and quantifiers.

\subsection{Informative outcomes and adjustment complexity}

The effective sample size \(N=nq\) measures outcome availability at the guaranteed occupied-cell arrival floor. Because individual cells may have arrival probabilities exceeding that floor, it is a conservative scale for outcome information. Under \cref{obj:ass:rare-arrival-slice}, the matched comparison in \cref{obj:thm:matched-minimax-frontier} separates the sampling contribution \(N^{-1}\) from the category contribution
\[
\left\{\frac{d}{N\log(e+N)}\right\}^{2},
\]
with their sum truncated at a constant. Thus increasing outcome availability improves both components, whereas increasing the alphabet size raises the cost of estimating the aggregate contrast across sparsely observed cells. The comparison incorporates the memberships observed for individuals whose primary outcomes have yet to arrive.

Along sequences satisfying the rare-arrival restriction and \(N\to\infty\), \cref{obj:prop:phase-boundaries} gives two distinct measures of manageable adjustment complexity. Consistency corresponds to \(d=o\{N\log(e+N)\}\), while squared-error risk of order \(N^{-1}\) corresponds to \(d=O\{\sqrt N\log(e+N)\}\). These thresholds distinguish category growth compatible with learning the treatment effect from category growth compatible with preserving the effective-sample-size precision benchmark. At each fixed coverage level, \cref{obj:thm:connected-interval-frontier} gives the corresponding interpretation for uniformly honest connected intervals: optimal maximal expected length has the order of the square root of the matched squared-error rate.

Near complete ascertainment, the deficit \(1-q\) provides a complementary measure of difficulty. The observed-outcome contrast attains the upper bound \(4/n+(1-q)^2\) throughout the unrestricted model in \cref{obj:thm:unrestricted-near-complete-arrival-envelope}. For \(d\ge1024n^2\), \cref{obj:thm:unrestricted-large-alphabet-deficit-lower} establishes the matching order \(n^{-1}+(1-q)^2\). At complete arrival, \cref{obj:thm:unrestricted-complete-arrival-frontier} gives squared-error order \(n^{-1}\) uniformly over every alphabet size. Together, these results explain how sufficiently complete ascertainment supports precision governed by sampling variation even under arbitrary category growth.

The quantifiers distinguish a uniform guarantee from a guarantee along a specified dimension sequence. In \cref{obj:prop:unrestricted-uniform-near-complete-elbow}, a single risk constant and an eventual sample-size threshold apply simultaneously to every \(d\ge1\). Such an \(n^{-1}\) guarantee is equivalent to \(1-q_n=O(n^{-1/2})\). A specified sequence can satisfy a different sufficient condition through its category growth: at any fixed \(q\in(0,1/2)\), \cref{obj:prop:unrestricted-fixed-q-phase-boundaries} characterizes \(n^{-1}\) risk by \(d_n=O\{\sqrt n\log(e+n)\}\). The comparison constants may depend on the fixed floor and the category-growth multiplier. The near-complete criterion therefore describes precision uniformly over arbitrary alphabets, while the fixed-floor criterion describes precision along sequences with controlled adjustment complexity.

\subsection{Surrogates and prospective trial interpretations}

The surrogate serves as an observed variable for conditioning outcome arrival. Under \cref{obj:ass:arrival-mar}, arrivals represent the outcome distribution within treatment--baseline--surrogate cells; \cref{obj:ass:occupied-cell-arrival} supplies positive ascertainment in every occupied cell. Its realized value may depend on assignment through \cref{obj:ass:surrogate-consistency}. The target remains the population average treatment effect on the primary outcome in \cref{obj:def:ate-functional}, with the cell aggregation specified in \cref{obj:def:cell-functional}.

The constant-surrogate specialization provides a useful benchmark for this interpretation. When both potential surrogates equal zero, the arrival condition reduces to conditioning on treatment and baseline category. For fixed alphabet size, sequences in the rare-arrival slice with diverging effective sample size retain minimax squared-error order \(N^{-1}\) within this specialization, as established by \cref{obj:prop:fixed-d-effective-sample-reduction}. That benchmark isolates the precision scale associated with outcome availability in a model where the surrogate supplies a single category.

Clinical-trial applications can be interpreted prospectively through these conditions. At a prespecified analysis horizon, a binary primary endpoint would supply the outcome, an available short-term measurement would supply the binary surrogate, and recorded baseline categories would determine the adjustment alphabet. Applying the comparisons would require independent sampling, balanced randomized assignment, consistency, and the stated conditional arrival and occupied-cell positivity conditions. The delayed binary endpoints considered by \citet[Abstract]{MarschnerBecker2001} provide motivation for this interpretation. The use of short-term outcomes and baseline covariates for interim decisions in \citet[Abstract]{VanLancker2020} supplies related motivation. Here, these examples connect the statistical observation scheme to prospective trial settings, while the precision conclusions retain the conditions of the cited internal results.

Further questions concerning scope and research priorities are collected in \cref{sec:discussion-limitations-and-open-questions}.


\section{Limitations and open questions}
\label{sec:discussion-limitations-and-open-questions}

The risk comparisons leave several questions about joint changes in outcome arrival and alphabet size. The unrestricted upper bound in \cref{obj:thm:unrestricted-mixed-count-upper} provides an envelope across arrival regimes, while \cref{obj:thm:matched-minimax-frontier,obj:thm:unrestricted-fixed-q-minimax-frontier} give matching comparisons under their respective conditions. Matching the unrestricted envelope throughout intermediate moving-arrival and alphabet regimes remains open. This requires lower bounds that track simultaneous changes in the arrival floor and the number of baseline categories.

A related question concerns the distinction between a uniform near-complete requirement and a requirement for a particular alphabet sequence. \Cref{obj:prop:unrestricted-uniform-near-complete-elbow} characterizes the uniform comparison, and \cref{obj:thm:unrestricted-large-alphabet-deficit-lower} supplies the large-alphabet obstruction. These results leave open the precise near-complete requirement along individual smaller-alphabet sequences. A sequence-specific characterization would determine how the necessary degree of ascertainment depends on the growth of the alphabet.

For uncertainty quantification, \cref{obj:thm:connected-interval-frontier} compares uniformly honest connected intervals under the rare-arrival restriction in \cref{obj:ass:rare-arrival-slice}. Extending matching expected-length comparisons beyond that restriction is another open problem. Such an extension would require interval lower bounds appropriate to the additional arrival regimes, together with procedures attaining the corresponding lengths under uniform coverage.

The observational comparison in \cref{obj:thm:zeng-fixed-positivity-polynomial-frontier} concerns a fixed positivity margin strictly between zero and one half, as specified in \cref{obj:def:zeng-discrete-ate-class}. The endpoint at which every occupied category has treatment propensity one half calls for a separate risk characterization. Determining how this endpoint relates to the strict-interior comparisons would clarify the role of treatment-propensity variation in the observational problem.

The conditional arrival assumption in \cref{obj:ass:arrival-mar} also defines an important boundary of the model. Extending the analysis beyond missing at random requires specifying how arrival may depend on the outcome after conditioning on the recorded covariates, assignment, and surrogate. Alternative identifying restrictions or sensitivity parameters would then determine the target and the precision comparison. The present risk results do not supply those extensions.

Statistical attainment and practical evaluation are separate questions. The procedure in \cref{obj:def:mixed-count-estimator} averages \(3^{k_0}\) stream assignments at each retained prefix length \(k_0\), and the observational procedure in \cref{obj:thm:zeng-fixed-positivity-polynomial-frontier} additionally averages over \(2^n\) synthetic assignments. Their risk guarantees concern those exact procedures. Moreover, the polynomial branch begins only when \(\log(e+nq)\ge128\), or \(nq\ge e^{128}-e\). These scales make the displayed formulas mathematical attainment constructions rather than evidence of numerical usefulness at ordinary sample sizes. Computational evaluation should seek more efficient representations and quantify additional error from approximations. Prospective clinical evaluation must also justify the supplied arrival floor used for tuning and coverage, examine the choice of baseline categories and the conditional arrival assumption, and compare with established trial analyses.

\citet[abstract 300411]{Kennedy2019DiscreteAbstract} earlier announced average-causal-effect estimation with arbitrarily high-dimensional discrete confounders, including a variation with sample-size-dependent weakening of propensity-score positivity. The available abstract does not state its precise model class, rate, loss, or attaining construction. It therefore supports priority for studying that shrinking-positivity observational setting but cannot settle theorem-level overlap with the fixed-margin observational result here. A comparison with a complete theorem statement or manuscript remains open.


\appendix

\section*{Appendices}

\section{Identification and decision-theory preliminaries}

This appendix collects the decision-theoretic comparisons, cell-identification identities, and synthetic-assignment map used in the risk analysis. The generic experiment and testing notation introduced below is local to this appendix.

\subsection{Finite decision experiments}

For the decision experiment in \cref{obj:lem:randomized-kernel-barycenter}, let \(\Lambda\) denote a parameter set, \(\mathcal Z\) a finite observation space with measurable singletons, \(Q_\lambda\) the observation law at parameter \(\lambda\), and \(\vartheta(\lambda)\in[-1,1]\) the scalar target. A randomized procedure is a parameter-independent Markov kernel \(\kappa\): conditional on observation \(o\), it draws an action from \(\kappa(o,\cdot)\). Its squared risk averages the squared estimation error over both the observation and the action. A deterministic procedure is represented by a kernel concentrated at one action for each observation. For actions in \([-1,1]\), the barycenter is the conditional mean of the action. The following comparison establishes the risk effect of replacing a randomized action by this mean, together with the corresponding clipping and uniform-seed representations.

\begin{lemmav}{lem:randomized-kernel-barycenter}[Barycenters and uniform randomization]
Let \(\Lambda\) be an arbitrary parameter set, and suppose:
\begin{itemize}
\item \textbf{(Sample space.)} \(\mathcal Z\) is a finite measurable space with measurable singletons.
\item \textbf{(Experiment.)} For each \(\lambda\in\Lambda\), \(Q_\lambda\) is a probability law on \(\mathcal Z\).
\item \textbf{(Target.)} For each \(\lambda\in\Lambda\), the target satisfies \(\vartheta(\lambda)\in[-1,1]\).
\end{itemize}
For a parameter-independent Markov kernel \(\kappa\) from \(\mathcal Z\) to \(\mathbb R\), supported on \([-1,1]\) at every observation, define its barycenter by
\[
b_\kappa(o)=\int_{\mathbb R}h\,\kappa(o,dh).
\]
Then \(b_\kappa(o)\in[-1,1]\) for every \(o\in\mathcal Z\), and, for every \(\lambda\in\Lambda\),
\[
\int_{\mathcal Z}\bigl(b_\kappa(o)-\vartheta(\lambda)\bigr)^2\,Q_\lambda(do)
\le
\int_{\mathcal Z}\int_{\mathbb R}
\bigl(h-\vartheta(\lambda)\bigr)^2\,\kappa(o,dh)\,Q_\lambda(do).
\]
The minimax squared risk over these bounded kernels equals that over deterministic maps into \([-1,1]\):
\[
\inf_{\kappa}\sup_{\lambda\in\Lambda}
\int_{\mathcal Z}\int_{\mathbb R}
\bigl(h-\vartheta(\lambda)\bigr)^2\,\kappa(o,dh)\,Q_\lambda(do)
=
\inf_{b:\mathcal Z\to[-1,1]}\sup_{\lambda\in\Lambda}
\int_{\mathcal Z}\bigl(b(o)-\vartheta(\lambda)\bigr)^2\,Q_\lambda(do),
\]
where the first infimum ranges over the bounded kernels just specified.

Define clipping by
\[
\operatorname{clip}_{[-1,1]}(h)=\max\{-1,\min\{1,h\}\}.
\]
For every parameter-independent Markov kernel \(\kappa\) from \(\mathcal Z\) to \(\mathbb R\), and every \(\lambda\in\Lambda\), its image under clipping satisfies
\[
\int_{\mathcal Z}\int_{\mathbb R}
\bigl(\operatorname{clip}_{[-1,1]}(h)-\vartheta(\lambda)\bigr)^2
\,\kappa(o,dh)\,Q_\lambda(do)
\le
\int_{\mathcal Z}\int_{\mathbb R}
\bigl(h-\vartheta(\lambda)\bigr)^2\,\kappa(o,dh)\,Q_\lambda(do),
\]
with both risks interpreted as nonnegative extended integrals. The extended minimax squared risk over all such real-valued Markov kernels, converted to a real value, equals the bounded-kernel minimax squared risk above.

The uniform law on \([0,1]\), defined by restricting Lebesgue measure to this interval, is a probability law. Every bounded kernel \(\kappa\) admits a measurable map
\(f:\mathcal Z\times[0,1]\to[-1,1]\) such that, for every \(o\in\mathcal Z\), \(\kappa(o,\cdot)\) is the distribution of \(f(o,U)\), where \(U\) has this uniform law. Conversely, every such measurable map induces a bounded Markov kernel with precisely these conditional distributions. The same uniform seed law serves every parameter.
\end{lemmav}

\begin{proof}[Proof of \cref{obj:lem:randomized-kernel-barycenter}]
\begin{enumerate}
\item Fix a bounded kernel \(\kappa\) and an observation \(o\in\mathcal Z\). Its support condition gives
\[
\kappa(o,[-1,1])=1,
\qquad
-1\le h\le 1
\quad\text{for }\kappa(o,dh)\text{-almost every }h.
\]
The coordinate function is therefore integrable. Integrating these bounds against the probability measure \(\kappa(o,\cdot)\) yields
\[
-1=\int_{\mathbb R}(-1)\,\kappa(o,dh)
\le b_\kappa(o)
\le \int_{\mathbb R}1\,\kappa(o,dh)=1.
\]
Every map on the finite space \(\mathcal Z\) is measurable because its singletons are measurable. Thus \(b_\kappa\) is a measurable map into \([-1,1]\), determined by \(\kappa\). % lean: CausalSmith.Stat.MarRareqLogfrontier.kernel_barycenter_mem

\item Fix \(\lambda\in\Lambda\). Bounded support gives square integrability of the coordinate and its translate by \(\vartheta(\lambda)\), and
\[
\int_{\mathbb R}\bigl(h-\vartheta(\lambda)\bigr)\,\kappa(o,dh)
=b_\kappa(o)-\vartheta(\lambda).
\]
The second-moment formula for the variance of the translated coordinate, together with invariance of variance under subtraction of a constant, gives
\[
\begin{aligned}
\int_{\mathbb R}\bigl(h-b_\kappa(o)\bigr)^2\,\kappa(o,dh)
&=
\int_{\mathbb R}\bigl(h-\vartheta(\lambda)\bigr)^2\,\kappa(o,dh)\\
&\quad-
\left(
\int_{\mathbb R}\bigl(h-\vartheta(\lambda)\bigr)\,\kappa(o,dh)
\right)^2.
\end{aligned}
\]
Substituting the preceding mean identity and rearranging yields
\[
\int_{\mathbb R}\bigl(h-\vartheta(\lambda)\bigr)^2\,\kappa(o,dh)
=
\bigl(b_\kappa(o)-\vartheta(\lambda)\bigr)^2
+
\int_{\mathbb R}\bigl(h-b_\kappa(o)\bigr)^2\,\kappa(o,dh).
\]
% lean: CausalSmith.Stat.MarRareqLogfrontier.kernel_barycenter_variance_identity
The last integral is nonnegative, so
\[
\bigl(b_\kappa(o)-\vartheta(\lambda)\bigr)^2
\le
\int_{\mathbb R}\bigl(h-\vartheta(\lambda)\bigr)^2\,\kappa(o,dh).
\]
% lean: CausalSmith.Stat.MarRareqLogfrontier.kernel_barycenter_sq_le
Both sides are integrable as functions of \(o\), since \(\mathcal Z\) is finite and \(Q_\lambda\) is a probability measure. Integrating gives
\[
\int_{\mathcal Z}\bigl(b_\kappa(o)-\vartheta(\lambda)\bigr)^2\,Q_\lambda(do)
\le
\int_{\mathcal Z}\int_{\mathbb R}
\bigl(h-\vartheta(\lambda)\bigr)^2\,\kappa(o,dh)\,Q_\lambda(do).
\]
% lean: CausalSmith.Stat.MarRareqLogfrontier.kernel_barycenter_risk_le

\item Define the estimator classes and their real-valued risks by
\[
\begin{aligned}
\mathcal K_{\mathrm{bd}}
&:=
\left\{\kappa:\mathcal Z\leadsto\mathbb R:
\kappa\text{ is a parameter-independent Markov kernel and }
\kappa(o,[-1,1])=1\text{ for every }o\right\},\\
\mathcal B_{\mathrm{det}}
&:=\{b:\mathcal Z\to[-1,1]\},\\
\mathcal R_{\mathrm{bd}}(\kappa,\lambda)
&:=
\int_{\mathcal Z}\int_{\mathbb R}
\bigl(h-\vartheta(\lambda)\bigr)^2\,\kappa(o,dh)\,Q_\lambda(do),
\qquad (\kappa,\lambda)\in\mathcal K_{\mathrm{bd}}\times\Lambda,\\
\mathcal R_{\mathrm{det}}(b,\lambda)
&:=
\int_{\mathcal Z}\bigl(b(o)-\vartheta(\lambda)\bigr)^2\,Q_\lambda(do),
\qquad (b,\lambda)\in\mathcal B_{\mathrm{det}}\times\Lambda.
\end{aligned}
\]
All real suprema over an empty parameter set are taken to be zero.

For every bounded kernel, its support condition and the target bound imply
\[
-2\le h-\vartheta(\lambda)\le2,
\qquad
0\le\bigl(h-\vartheta(\lambda)\bigr)^2\le4
\quad\text{for }\kappa(o,dh)\text{-almost every }h.
\]
Integrating first against \(\kappa(o,\cdot)\) and then against \(Q_\lambda\) gives
\[
0\le
\int_{\mathbb R}\bigl(h-\vartheta(\lambda)\bigr)^2\,\kappa(o,dh)
\le4,
\qquad
0\le\mathcal R_{\mathrm{bd}}(\kappa,\lambda)\le4.
\]
% lean: CausalSmith.Stat.MarRareqLogfrontier.bounded_kernelRisk_bounds

For every \(b\in\mathcal B_{\mathrm{det}}\), define its deterministic kernel by
\[
\kappa_b(o,\cdot):=\delta_{b(o)},
\qquad o\in\mathcal Z.
\]
Measurability of \(b\) makes this a Markov kernel, and its interval-valued range gives
\[
\kappa_b(o,[-1,1])=1,
\qquad
\mathcal R_{\mathrm{bd}}(\kappa_b,\lambda)
=\mathcal R_{\mathrm{det}}(b,\lambda).
\]
% lean: CausalSmith.Stat.MarRareqLogfrontier.boundedMapKernel_risk
In particular, both estimator classes are nonempty: take
\[
b_{\mathrm{zero}}:\mathcal Z\to[-1,1],
\qquad b_{\mathrm{zero}}(o):=0,
\]
and its kernel \(\kappa_{b_{\mathrm{zero}}}\). The risk identity also gives
\(0\le\mathcal R_{\mathrm{det}}(b,\lambda)\le4\).
Consequently, all worst-case risks lie in \([0,4]\), including when
\(\Lambda\) is empty, and their infima are well defined.

The deterministic-kernel construction gives the first comparison:
\[
\inf_{\kappa\in\mathcal K_{\mathrm{bd}}}
\sup_{\lambda\in\Lambda}\mathcal R_{\mathrm{bd}}(\kappa,\lambda)
\le
\inf_{b\in\mathcal B_{\mathrm{det}}}
\sup_{\lambda\in\Lambda}\mathcal R_{\mathrm{det}}(b,\lambda).
\]
Conversely, each
\(\kappa\in\mathcal K_{\mathrm{bd}}\) has the barycenter map
\(b_\kappa\in\mathcal B_{\mathrm{det}}\), whose pointwise risk inequality gives
\[
\sup_{\lambda\in\Lambda}\mathcal R_{\mathrm{det}}(b_\kappa,\lambda)
\le
\sup_{\lambda\in\Lambda}\mathcal R_{\mathrm{bd}}(\kappa,\lambda).
\]
For a nonempty parameter set this follows from monotonicity of bounded
suprema; for an empty parameter set both sides equal zero. Taking infima gives
\[
\inf_{b\in\mathcal B_{\mathrm{det}}}
\sup_{\lambda\in\Lambda}\mathcal R_{\mathrm{det}}(b,\lambda)
\le
\inf_{\kappa\in\mathcal K_{\mathrm{bd}}}
\sup_{\lambda\in\Lambda}\mathcal R_{\mathrm{bd}}(\kappa,\lambda).
\]
The two comparisons establish the asserted equality. % lean: CausalSmith.Stat.MarRareqLogfrontier.bounded_kernel_minimax_eq_deterministic

\item Fix an arbitrary real-valued kernel \(\kappa\) and
\(\lambda\in\Lambda\). First suppose \(h\le1\). If \(-1\le h\), clipping fixes
\(h\), so the squared losses agree. If \(h\le-1\), clipping gives \(-1\), and
\[
\begin{aligned}
\bigl(h-\vartheta(\lambda)\bigr)^2
-\bigl(-1-\vartheta(\lambda)\bigr)^2
&=(-1-h)\bigl(1-h+2\vartheta(\lambda)\bigr)\\
&\ge0,
\end{aligned}
\]
because both factors are nonnegative. In the remaining case \(1\le h\),
clipping gives \(1\), and
\[
\begin{aligned}
\bigl(h-\vartheta(\lambda)\bigr)^2
-\bigl(1-\vartheta(\lambda)\bigr)^2
&=(h-1)\bigl(h+1-2\vartheta(\lambda)\bigr)\\
&\ge0.
\end{aligned}
\]
Thus, for every \(h\in\mathbb R\),
\[
\bigl(\operatorname{clip}_{[-1,1]}(h)-\vartheta(\lambda)\bigr)^2
\le
\bigl(h-\vartheta(\lambda)\bigr)^2.
\]
% lean: CausalSmith.Stat.MarRareqLogfrontier.realKernelClip_sq_le

Define the clipped kernel by the pushforward
\[
\kappa^{\mathrm{cl}}(o,\cdot)
:=
\kappa(o,\cdot)\circ\operatorname{clip}_{[-1,1]}^{-1},
\qquad o\in\mathcal Z,
\]
and define the extended risk for any real-valued Markov kernel by
\[
\mathcal R_{\mathrm{ext}}(\kappa,\lambda)
:=
\int_{\mathcal Z}\int_{\mathbb R}
\bigl(h-\vartheta(\lambda)\bigr)^2\,\kappa(o,dh)\,Q_\lambda(do)
\in[0,\infty],
\]
where both integrals are nonnegative extended integrals. Clipping is measurable.
The pushforward integration formula and monotonicity of these integrals give
\[
\begin{aligned}
\mathcal R_{\mathrm{ext}}(\kappa^{\mathrm{cl}},\lambda)
&=
\int_{\mathcal Z}\int_{\mathbb R}
\bigl(\operatorname{clip}_{[-1,1]}(h)-\vartheta(\lambda)\bigr)^2
\,\kappa(o,dh)\,Q_\lambda(do)\\
&\le \mathcal R_{\mathrm{ext}}(\kappa,\lambda).
\end{aligned}
\]
This comparison includes infinite original risks. % lean: CausalSmith.Stat.MarRareqLogfrontier.realKernelClip_risk_le

\item Define the full kernel class by
\[
\mathcal K_{\mathrm{real}}
:=
\{\kappa:\mathcal Z\leadsto\mathbb R:
\kappa\text{ is a parameter-independent Markov kernel}\}.
\]
For a bounded kernel, squared loss is nonnegative and integrable at every
observation. Its inner integral is integrable over the finite observation
space. Applying the identification of a nonnegative integrable function's real
and extended integrals at both levels gives
\[
\mathcal R_{\mathrm{ext}}(\kappa,\lambda)
=\mathcal R_{\mathrm{bd}}(\kappa,\lambda)
\quad\text{in }[0,\infty],
\qquad \kappa\in\mathcal K_{\mathrm{bd}}.
\]
% lean: CausalSmith.Stat.MarRareqLogfrontier.bounded_realKernelRisk_eq_ofReal

Inclusion of bounded kernels into the full class therefore gives
\[
\inf_{\kappa\in\mathcal K_{\mathrm{real}}}
\sup_{\lambda\in\Lambda}\mathcal R_{\mathrm{ext}}(\kappa,\lambda)
\le
\inf_{\kappa\in\mathcal K_{\mathrm{bd}}}
\sup_{\lambda\in\Lambda}\mathcal R_{\mathrm{bd}}(\kappa,\lambda),
\]
with the right side here evaluated in \([0,\infty]\).

For the reverse comparison, clipping takes values in \([-1,1]\), since
\[
-1\le \max\{-1,\min\{1,h\}\}\le1
\qquad(h\in\mathbb R).
\]
A measurable pushforward preserves probability, and
\[
\operatorname{clip}_{[-1,1]}^{-1}([-1,1])=\mathbb R,
\qquad
\kappa^{\mathrm{cl}}(o,[-1,1])=\kappa(o,\mathbb R)=1.
\]
Thus \(\kappa^{\mathrm{cl}}\in\mathcal K_{\mathrm{bd}}\) for every
\(\kappa\in\mathcal K_{\mathrm{real}}\). The preceding risk comparison yields
\[
\inf_{\kappa\in\mathcal K_{\mathrm{bd}}}
\sup_{\lambda\in\Lambda}\mathcal R_{\mathrm{bd}}(\kappa,\lambda)
\le
\inf_{\kappa\in\mathcal K_{\mathrm{real}}}
\sup_{\lambda\in\Lambda}\mathcal R_{\mathrm{ext}}(\kappa,\lambda).
\]
Hence the two extended minimax values agree.

For the conversion to a real value, use
\[
\operatorname{rv}:[0,\infty]\to\mathbb R,
\qquad
\operatorname{rv}|_{[0,\infty)}=\operatorname{id},
\qquad
\operatorname{rv}(\infty)=0.
\]
The bounds \(0\le\mathcal R_{\mathrm{bd}}\le4\) place every bounded-kernel
supremum, and their infimum, in \([0,4]\). On this interval the real and
extended orders agree, so conversion commutes with these suprema and the
infimum. For an empty parameter set both kinds of supremum equal zero.
It follows that
\[
\operatorname{rv}\!\left(
\inf_{\kappa\in\mathcal K_{\mathrm{real}}}
\sup_{\lambda\in\Lambda}\mathcal R_{\mathrm{ext}}(\kappa,\lambda)
\right)
=
\inf_{\kappa\in\mathcal K_{\mathrm{bd}}}
\sup_{\lambda\in\Lambda}\mathcal R_{\mathrm{bd}}(\kappa,\lambda),
\]
where the right side is now evaluated in \(\mathbb R\). % lean: CausalSmith.Stat.MarRareqLogfrontier.real_kernel_minimax_eq_bounded

\item Define the seed law, viewed as a measure on \([0,1]\), by
\[
\nu_{\mathrm{seed}}
:=\left.\operatorname{Leb}_{\mathbb R}\right|_{[0,1]}.
\]
Its total mass is
\[
\nu_{\mathrm{seed}}([0,1])
=\operatorname{Leb}_{\mathbb R}([0,1])
=1-0=1,
\]
so it is a probability law. % lean: CausalSmith.Stat.MarRareqLogfrontier.uniformSeedLaw_probability

\item Fix \(\kappa\in\mathcal K_{\mathrm{bd}}\). The measurable randomization
theorem for Markov kernels with a standard Borel output space, applied to the
real output space, supplies a jointly measurable map satisfying
\[
\widetilde f:\mathcal Z\times[0,1]\to\mathbb R,
\qquad
\kappa(o,\cdot)
=
\nu_{\mathrm{seed}}\circ\widetilde f(o,\cdot)^{-1}
\quad(o\in\mathcal Z).
\]
Define the interval-valued realization by
\[
f:\mathcal Z\times[0,1]\to[-1,1],
\qquad
f(o,v):=\operatorname{clip}_{[-1,1]}\bigl(\widetilde f(o,v)\bigr).
\]
It is jointly measurable by measurability of clipping. For
\(h\in[-1,1]\), clipping fixes \(h\); the bounded support condition consequently
gives
\[
\operatorname{clip}_{[-1,1]}(h)=h
\quad\text{for }\kappa(o,dh)\text{-almost every }h,
\qquad
\kappa(o,\cdot)\circ\operatorname{clip}_{[-1,1]}^{-1}
=\kappa(o,\cdot).
\]
% lean: CausalSmith.Stat.MarRareqLogfrontier.boundedKernel_map_clip
Composition of measurable pushforwards now yields
\[
\begin{aligned}
\nu_{\mathrm{seed}}\circ f(o,\cdot)^{-1}
&=
\bigl(\nu_{\mathrm{seed}}\circ\widetilde f(o,\cdot)^{-1}\bigr)
\circ\operatorname{clip}_{[-1,1]}^{-1}\\
&=\kappa(o,\cdot)\circ\operatorname{clip}_{[-1,1]}^{-1}
=\kappa(o,\cdot).
\end{aligned}
\]
Thus \(f(o,U)\), for \(U\sim\nu_{\mathrm{seed}}\), has the required distribution
at every observation. % lean: CausalSmith.Stat.MarRareqLogfrontier.boundedKernel_uniformSeed_representation

\item Conversely, let \(f:\mathcal Z\times[0,1]\to[-1,1]\) be jointly
measurable, and define
\[
\kappa_f(o,\cdot)
:=
\nu_{\mathrm{seed}}\circ f(o,\cdot)^{-1},
\qquad o\in\mathcal Z.
\]
Every section \(f(o,\cdot)\) is measurable, so its pushforward is a probability
measure. Finiteness of \(\mathcal Z\) and measurability of its singletons give
the kernel measurability condition. Moreover,
\[
f(o,\cdot)^{-1}([-1,1])=[0,1],
\qquad
\kappa_f(o,[-1,1])
=\nu_{\mathrm{seed}}([0,1])=1.
\]
Hence \(\kappa_f\) is a bounded Markov kernel with exactly the displayed
conditional distributions. Both constructions use the same law
\(\nu_{\mathrm{seed}}\), independently of the parameter. For each
\(\lambda\), the auxiliary seed can be sampled independently of the
observation under the product law
\[
Q_\lambda\otimes\nu_{\mathrm{seed}}.
\]
% lean: CausalSmith.Stat.MarRareqLogfrontier.boundedKernel_of_uniformSeed
\end{enumerate}
\end{proof}

The comparison holds at every parameter, so conditional averaging also preserves a uniform squared-risk bound. It supplies the decision-theoretic justification for the auxiliary average in \cref{obj:def:mixed-count-estimator} and the synthetic-assignment average in \cref{obj:thm:zeng-fixed-positivity-polynomial-frontier}. Clipping weakly reduces squared loss for targets in the stated interval; on the real line, this clipping map coincides with Euclidean projection onto that interval. The uniform-seed representation places bounded randomized procedures and procedures implemented with independent auxiliary randomization in the same decision class.

For the testing comparison in \cref{obj:lem:randomized-kernel-tv-comparison}, let \(Q_0,Q_1\) denote two laws on \(\mathcal Z\), and let \(\mathcal H\) be a measurable action space. The notation \(Q_i\ltimes\kappa\) denotes the joint observation--action law, whereas \(Q_i\kappa\) denotes its action marginal. Total variation, written \(\operatorname{TV}\), is the supremum of absolute probability differences over measurable events. For binary actions, a randomized test has conditional probability \(\psi(o)=\kappa(o,\{1\})\) of selecting alternative \(1\). The next comparison applies both to events that retain the observation and to events determined by the action alone.

\begin{lemmav}{lem:randomized-kernel-tv-comparison}[Total variation under randomization]
Let \(Q_0,Q_1\) be probability laws on a finite measurable space \(\mathcal Z\) with measurable singletons, and let \(\kappa\) be a common Markov kernel from \(\mathcal Z\) to a measurable action space \(\mathcal H\). For \(i\in\{0,1\}\), write \(Q_i\ltimes\kappa\) for the joint law obtained by drawing an input from \(Q_i\) and then an action from \(\kappa\) conditional on that input, and write \(Q_i\kappa\) for its action marginal. For two probability laws on the same measurable space, define their total variation distance as the supremum of their absolute probability differences over measurable events. Then
\[
\operatorname{TV}(Q_0\kappa,Q_1\kappa)
\le
\operatorname{TV}(Q_0\ltimes\kappa,Q_1\ltimes\kappa)
=
\operatorname{TV}(Q_0,Q_1).
\]
For every measurable event \(T\subseteq\mathcal Z\times\mathcal H\),
\[
(Q_0\ltimes\kappa)(T)+(Q_1\ltimes\kappa)(T^c)
\ge 1-\operatorname{TV}(Q_0,Q_1),
\]
and for every measurable event \(T\subseteq\mathcal H\),
\[
(Q_0\kappa)(T)+(Q_1\kappa)(T^c)
\ge 1-\operatorname{TV}(Q_0,Q_1).
\]
\end{lemmav}

\begin{proof}[Proof of \cref{obj:lem:randomized-kernel-tv-comparison}]
\begin{enumerate}
\item We first establish the bounded-function estimate used for the kernel probabilities. For any measurable function
\[
\varphi:\mathcal Z\longrightarrow[0,1],
\qquad
L_{\varphi,t}:=\{z\in\mathcal Z:t\le\varphi(z)\}
\quad(t\in\mathbb R),
\]
the level sets are measurable, and the definition of total variation gives
\[
\bigl|Q_0(L_{\varphi,t})-Q_1(L_{\varphi,t})\bigr|
\le \operatorname{TV}(Q_0,Q_1).
\]
The layer-cake identity gives, for \(i\in\{0,1\}\),
\[
\int_{\mathcal Z}\varphi(z)\,Q_i(dz)
=
\int_{(0,1]}Q_i(L_{\varphi,t})\,dt.
\]
All integrands are measurable and bounded. Subtracting these identities and integrating the preceding event bound yields
\[
\begin{aligned}
\left|
\int_{\mathcal Z}\varphi\,dQ_0-
\int_{\mathcal Z}\varphi\,dQ_1
\right|
&=
\left|\int_{(0,1]}
\bigl[Q_0(L_{\varphi,t})-Q_1(L_{\varphi,t})\bigr]\,dt\right|\\
&\le
\int_{(0,1]}
\bigl|Q_0(L_{\varphi,t})-Q_1(L_{\varphi,t})\bigr|\,dt\\
&\le \operatorname{TV}(Q_0,Q_1).
\end{aligned}
\]
% lean: Causalean.Stat.tvDist_integral_range
For any measurable event \(T_{\mathrm{act}}\subseteq\mathcal H\), define
\[
\varphi_{\mathrm{act}}(z):=\kappa(z,T_{\mathrm{act}})
\quad(z\in\mathcal Z).
\]
This function is measurable and takes values in \([0,1]\), and
\[
(Q_i\kappa)(T_{\mathrm{act}})
=
\int_{\mathcal Z}\varphi_{\mathrm{act}}\,dQ_i
\quad(i\in\{0,1\}).
\]
Applying the bounded-function estimate and taking the supremum over measurable action events proves
\[
\operatorname{TV}(Q_0\kappa,Q_1\kappa)
\le \operatorname{TV}(Q_0,Q_1).
\]
% lean: Causalean.Stat.tvDist_bind_le

\item For any measurable joint event \(T_{\mathrm{joint}}\subseteq\mathcal Z\times\mathcal H\), define its sections and their kernel probabilities by
\[
T_{\mathrm{joint},z}
:=
\{h\in\mathcal H:(z,h)\in T_{\mathrm{joint}}\},
\qquad
\varphi_{\mathrm{joint}}(z)
:=
\kappa(z,T_{\mathrm{joint},z})
\quad(z\in\mathcal Z).
\]
The sections are measurable, and \(\varphi_{\mathrm{joint}}\) is measurable with values in \([0,1]\). The joint-law identity is
\[
(Q_i\ltimes\kappa)(T_{\mathrm{joint}})
=
\int_{\mathcal Z}\varphi_{\mathrm{joint}}\,dQ_i
\quad(i\in\{0,1\}).
\]
Thus the bounded-function estimate, followed by the supremum over joint events, gives
\[
\operatorname{TV}(Q_0\ltimes\kappa,Q_1\ltimes\kappa)
\le \operatorname{TV}(Q_0,Q_1).
\]
For the reverse inequality, every measurable input event \(T_{\mathrm{in}}\subseteq\mathcal Z\) satisfies
\[
(Q_i\ltimes\kappa)(T_{\mathrm{in}}\times\mathcal H)
=
Q_i(T_{\mathrm{in}})
\quad(i\in\{0,1\}),
\]
because \(\kappa(z,\mathcal H)=1\). Consequently,
\[
\begin{aligned}
|Q_0(T_{\mathrm{in}})-Q_1(T_{\mathrm{in}})|
&=
\left|
(Q_0\ltimes\kappa)(T_{\mathrm{in}}\times\mathcal H)
-
(Q_1\ltimes\kappa)(T_{\mathrm{in}}\times\mathcal H)
\right|\\
&\le
\operatorname{TV}(Q_0\ltimes\kappa,Q_1\ltimes\kappa).
\end{aligned}
\]
Taking the supremum over input events proves
\[
\operatorname{TV}(Q_0\ltimes\kappa,Q_1\ltimes\kappa)
=
\operatorname{TV}(Q_0,Q_1).
\]
% lean: Causalean.Stat.tvDist_compProd_eq
Combining this equality with the action-marginal bound gives
\[
\operatorname{TV}(Q_0\kappa,Q_1\kappa)
\le
\operatorname{TV}(Q_0\ltimes\kappa,Q_1\ltimes\kappa)
=
\operatorname{TV}(Q_0,Q_1).
\]
% lean: Causalean.Stat.tvDist_bind_le_compProd

\item The joint laws are probability laws. Hence, for any measurable joint event \(T_{\mathrm{joint}}\),
\[
(Q_1\ltimes\kappa)(T_{\mathrm{joint}}^c)
=
1-(Q_1\ltimes\kappa)(T_{\mathrm{joint}}),
\]
and the event bound defining their total variation gives
\[
\begin{aligned}
&(Q_0\ltimes\kappa)(T_{\mathrm{joint}})
 +(Q_1\ltimes\kappa)(T_{\mathrm{joint}}^c)\\
&\qquad=
1+(Q_0\ltimes\kappa)(T_{\mathrm{joint}})
 -(Q_1\ltimes\kappa)(T_{\mathrm{joint}})\\
&\qquad\ge
1-\operatorname{TV}(Q_0\ltimes\kappa,Q_1\ltimes\kappa)
=
1-\operatorname{TV}(Q_0,Q_1).
\end{aligned}
\]
% lean: Causalean.Stat.one_sub_tvDist_le_compProd_test

\item Likewise, the action marginals are probability laws, so every measurable action event \(T_{\mathrm{act}}\) satisfies
\[
(Q_1\kappa)(T_{\mathrm{act}}^c)
=
1-(Q_1\kappa)(T_{\mathrm{act}}).
\]
Using their total-variation event bound and then the action-marginal contraction established above,
\[
\begin{aligned}
(Q_0\kappa)(T_{\mathrm{act}})
 +(Q_1\kappa)(T_{\mathrm{act}}^c)
&=
1+(Q_0\kappa)(T_{\mathrm{act}})
 -(Q_1\kappa)(T_{\mathrm{act}})\\
&\ge
1-\operatorname{TV}(Q_0\kappa,Q_1\kappa)\\
&\ge
1-\operatorname{TV}(Q_0,Q_1).
\end{aligned}
\]
% lean: Causalean.Stat.one_sub_tvDist_le_bind_test
\end{enumerate}
\end{proof}

Retaining the observation alongside the randomized action preserves total variation, while taking the action marginal weakly reduces it. When the event represents selection of alternative \(1\), the displayed probability sum is the sum of the two testing errors. These comparisons allow the lower-bound analysis in \cref{obj:thm:full-experiment-lower} and the interval analysis in \cref{obj:thm:connected-interval-frontier} to cover parameter-independent randomized procedures using the same observation-law comparison.

\subsection{Identification through observed cells}

The cell functional in \cref{obj:def:cell-functional} expresses the target through membership probabilities and arrived-outcome means, with zero contributions assigned to null cells. Its connection to the average treatment effect uses randomized independence, balanced assignment, consistency, conditional missing at random, and positive arrival in occupied cells, as specified in \cref{obj:ass:randomized-independence,obj:ass:balanced-randomization,obj:ass:surrogate-consistency,obj:ass:outcome-consistency,obj:ass:arrival-mar,obj:ass:occupied-cell-arrival}. The first identity applies to the unrestricted-arrival class.

\begin{lemmav}{lem:unrestricted-cell-identification}[Identification by cell contributions]
Let \(n,d\in\mathbb N\), \(q\in\mathbb R\), and let \(P\in\mathcal M^{+}_{n,d,q}\) satisfy \cref{obj:def:unrestricted-arrival-model-class}. The average treatment effect \(\tau(P)\) and the total observed-cell functional \(\Phi(P)\), defined in \cref{obj:def:ate-functional,obj:def:cell-functional}, satisfy
\[
\tau(P)=\Phi(P)
=2\sum_{(a,x,s)\in\mathcal J_d}g(a)c_{axs}(P).
\]
\end{lemmav}

\begin{proof}[Proof of \cref{obj:lem:unrestricted-cell-identification}]
Fix \(P\in\mathcal M^{+}_{n,d,q}\). We use the conditions in \cref{obj:def:unrestricted-arrival-model-class}.

\begin{enumerate}
\item By \cref{obj:synth_4}, the potential outcomes are binary. Linearity of integration for these bounded variables therefore gives
\[
\tau(P)
=E_P[Y(1)]-E_P[Y(0)]
=P(Y(1)=1)-P(Y(0)=1),
\]
using \cref{obj:def:ate-functional}. % lean: CausalSmith.Stat.MarRareqLogfrontier.ate_eq_potentialOutcomeMass

\item We establish, for every \((a,x,s)\in\mathcal J_d\), the identity
\[
c_{axs}(P)=P(A=a,X=x,S=s,Y=1).
\]
First suppose \(p_{axs}(P)>0\). By \cref{obj:synth_16,obj:ass:occupied-cell-arrival} and \(q>0\),
\[
P(A=a,X=x,S=s,R=1)
=p_{axs}(P)\rho_{axs}(P)
\ge q\,p_{axs}(P)>0.
\]
Conditional independence in \cref{obj:ass:arrival-mar}, applied within this occupied cell, yields the cross-multiplied identity
\[
\begin{aligned}
&P(A=a,X=x,S=s,R=1,Y=1)\,p_{axs}(P)\\
&\qquad =
P(A=a,X=x,S=s,R=1)\,
P(A=a,X=x,S=s,Y=1).
\end{aligned}
\]
Since \(Y\) is binary and the arrived-cell probability is positive, \cref{obj:def:cell-functional} gives
\[
\begin{aligned}
c_{axs}(P)
&=p_{axs}(P)
\frac{P(A=a,X=x,S=s,R=1,Y=1)}
     {P(A=a,X=x,S=s,R=1)}\\
&=\frac{P(A=a,X=x,S=s,R=1,Y=1)\,p_{axs}(P)}
        {P(A=a,X=x,S=s,R=1)}\\
&=P(A=a,X=x,S=s,Y=1).
\end{aligned}
\]

In the remaining case, nonnegativity implies \(p_{axs}(P)=0\). Monotonicity of probability gives
\[
\begin{aligned}
0&\le P(A=a,X=x,S=s,Y=1)\le p_{axs}(P)=0,\\
0&\le P(A=a,X=x,S=s,R=1)\le p_{axs}(P)=0.
\end{aligned}
\]
Thus both probabilities vanish, and the zero-arrival convention in
\cref{obj:def:cell-functional} gives \(c_{axs}(P)=0\), proving the identity also in this case. % lean: CausalSmith.Stat.MarRareqLogfrontier.cellContribution_eq_cellOutcomeMass

\item For each \(a\in\{0,1\}\), define the finite measure obtained by restricting \(P\) to the arm-\(a\), outcome-one event:
\[
\nu_a^{\mathrm{out}}(\,\cdot\,)
:=P\bigl(\,\cdot\,\cap\{A=a,Y=1\}\bigr).
\]
The events \(\{X=x,S=s\}\), with \(x\in\{0,\ldots,d-1\}\) and \(s\in\{0,1\}\), are disjoint and cover the full-data record space. Finite additivity of this restricted measure therefore yields
\[
\begin{aligned}
&\sum_{x=0}^{d-1}\sum_{s\in\{0,1\}}
P(A=a,X=x,S=s,Y=1)\\
&\qquad =
\sum_{x=0}^{d-1}\sum_{s\in\{0,1\}}
\nu_a^{\mathrm{out}}(X=x,S=s)
=P(A=a,Y=1).
\end{aligned}
\] % lean: CausalSmith.Stat.MarRareqLogfrontier.sum_cellOutcomeMass_eq_armOutcomeMass

\item Expanding \cref{obj:def:cell-functional} over the two treatment arms, using \(g(1)=1\) and \(g(0)=-1\) from \cref{obj:synth_19}, and substituting the preceding identities gives
\[
\begin{aligned}
\Phi(P)
&=2\sum_{x=0}^{d-1}\sum_{s\in\{0,1\}}
\bigl(c_{1xs}(P)-c_{0xs}(P)\bigr)\\
&=2\bigl(P(A=1,Y=1)-P(A=0,Y=1)\bigr).
\end{aligned}
\] % lean: CausalSmith.Stat.MarRareqLogfrontier.cellFunctional_eq_armOutcomeMass

\item Fix \(a\in\{0,1\}\). Projecting the independence asserted in
\cref{obj:ass:randomized-independence} onto the selected potential outcome gives
\[
A\perp\!\!\!\perp Y(a),
\qquad
P(A=a,Y(a)=1)=P(A=a)\,P(Y(a)=1).
\]
By \cref{obj:ass:balanced-randomization},
\[
P(A=1)=\frac12,
\qquad
P(A=0)=1-P(A=1)=\frac12.
\]
The assigned-arm outcome in \cref{obj:synth_4} satisfies \(Y=Y(a)\) on \(\{A=a\}\). Consequently,
\[
\begin{aligned}
2P(A=a,Y=1)
&=2P(A=a,Y(a)=1)\\
&=2P(A=a)\,P(Y(a)=1)
=P(Y(a)=1).
\end{aligned}
\]
Applying this identity first to \(a=1\) and then to \(a=0\) gives
\[
2P(A=1,Y=1)=P(Y(1)=1),
\qquad
2P(A=0,Y=1)=P(Y(0)=1).
\] % lean: CausalSmith.Stat.MarRareqLogfrontier.randomized_armOutcomeMass

\item Subtracting the arm-zero identity from the arm-one identity and substituting into the expressions for \(\tau(P)\) and \(\Phi(P)\) proves
\[
\tau(P)
=\Phi(P)
=2\sum_{(a,x,s)\in\mathcal J_d}g(a)c_{axs}(P),
\]
where the final equality is \cref{obj:def:cell-functional}. % lean: CausalSmith.Stat.MarRareqLogfrontier.unrestricted_cell_identification
\end{enumerate}
\end{proof}

The identity makes the observed-cell functional an exact population target throughout the unrestricted-arrival class. The surrogate enters through the cells on which arrival is conditioned, and the signed sum aggregates their contributions into the treatment contrast. This is the identification relation used by the unrestricted risk bound in \cref{obj:thm:unrestricted-mixed-count-upper}.

For the rare-arrival class, the same relation is recorded under the model conditions governing its minimax analysis.

\begin{lemmav}{lem:cell-identification-background}[Identification by cell contributions]
Let \(n,d\in\mathbb N\), \(q\in\mathbb R\), and let \(P\in\mathcal M_{n,d,q}\) satisfy the model conditions in \cref{obj:def:model-class}. The average treatment effect \(\tau(P)\) of \cref{obj:def:ate-functional} equals the total signed observed-cell functional \(\Phi(P)\) of \cref{obj:def:cell-functional}:
\[
\tau(P)=\Phi(P)
=2\sum_{(a,x,s)\in\mathcal J_d}g(a)c_{axs}(P).
\]
Here the cell contributions and their null-cell convention are as defined in \cref{obj:def:cell-functional}.
\end{lemmav}

\begin{proof}[Proof of \cref{obj:lem:cell-identification-background}]
\begin{enumerate}
\item By \cref{obj:lem:unrestricted-cell-identification}, it suffices to establish
\[
P\in\mathcal M^{+}_{n,d,q},
\]
where \(\mathcal M^{+}_{n,d,q}\) is the arrival-floor model of \cref{obj:def:unrestricted-arrival-model-class}. For such a law, the cited result supplies
\[
\tau(P)=\Phi(P)
=2\sum_{(a,x,s)\in\mathcal J_d}g(a)c_{axs}(P).
\]
% lean: CausalSmith.Stat.MarRareqLogfrontier.unrestricted_cell_identification

\item Since \(P\in\mathcal M_{n,d,q}\), \cref{obj:def:model-class} gives
\[
n\ge1,\qquad d\ge1,\qquad 0<q\le1,
\]
together with
\cref{obj:ass:randomized-independence,obj:ass:balanced-randomization},
\cref{obj:ass:surrogate-consistency,obj:ass:outcome-consistency}, and
\cref{obj:ass:arrival-mar,obj:ass:occupied-cell-arrival}.
These are precisely the conditions defining the arrival-floor model. Thus
\[
P\in\mathcal M^{+}_{n,d,q},
\]
and the identity in the preceding step follows.
% lean: CausalSmith.Stat.MarRareqLogfrontier.cell_identification_background
\end{enumerate}
\end{proof}

The rare-arrival restriction in \cref{obj:ass:rare-arrival-slice} specifies the regime for the risk comparison, while this identity fixes its estimand. In particular, the null-cell convention gives a total cell functional across laws with different occupied categories. The mixed-count analysis underlying \cref{obj:thm:matched-minimax-frontier} therefore targets the same average treatment effect throughout the rare-arrival model.

\subsection{Synthetic assignment and experiment comparison}

The observational class in \cref{obj:def:zeng-discrete-ate-class} supplies an experiment that can be mapped into the unrestricted randomized model. The map assigns an independent fair binary arm to each observational record, records outcome arrival when that arm matches the observational treatment, and sets the surrogate to zero. The following statement establishes a parameter-independent observation map together with a full-data completion preserving the target.

\begin{lemmav}{lem:synthetic-randomization-kernel}[Synthetic randomization kernel]
Let the inputs satisfy:
\begin{itemize}
\item \textbf{(Sample size.)} \(n,d\in\mathbb N\) and \(n\ge1\).
\item \textbf{(Arrival floor.)} \(0<q<1/2\).
\item \textbf{(Observational law.)} \(P_Z\in\mathcal D_{d,q}\), the discrete positivity class of \cref{obj:def:zeng-discrete-ate-class}.
\end{itemize}
Index observational baseline categories by \(x\in\{1,\ldots,d\}\) and full-data baseline categories by their corresponding zero-based indices \(x-1\in\{0,\ldots,d-1\}\). For an observational record \((x,b,y)\), define
\[
o(x,b,y;u)
=
\left(x-1,u,0,\mathbf 1\{b=u\},\mathbf 1\{b=u\}y\right),
\qquad u\in\{0,1\}.
\]
Applying this map independently to each input record with independent fair binary draws defines a Markov kernel independent of \(P_Z\).

There exists a full-data law \(P_Z^*\in\mathcal M^{+}_{n,d,q}\), as defined in \cref{obj:def:unrestricted-arrival-model-class}, such that applying this kernel to \(\mathbf Z_n\), an i.i.d.\ sample of \(n\) observational records from \(P_Z\), yields exactly the \(n\)-fold product of the observed-record law induced by \(P_Z^*\). The resulting observed sample \(\mathbf O_n\) satisfies \cref{obj:ass:iid-sampling} under \(P_Z^*\), and
\[
\tau(P_Z^*)=\theta_Z(P_Z),
\]
with the functionals defined in \cref{obj:def:ate-functional,obj:def:zeng-ate-functional}.
\end{lemmav}

\begin{proof}[Proof of \cref{obj:lem:synthetic-randomization-kernel}]
\begin{enumerate}
\item Construct the full-data completion. For observational categories \(x\in\{1,\ldots,d\}\) and \(b\in\{0,1\}\), use the following total versions of the quantities in \cref{obj:def:zeng-discrete-ate-class}:
\[
p_x:=P_Z(X=x),\qquad
\pi_x:=
\begin{cases}
P_Z(X=x,B^{\mathrm{obs}}=1)/p_x,&p_x>0,\\
0,&p_x=0,
\end{cases}
\]
and
\[
\mu_{bx}:=
\begin{cases}
\dfrac{P_Z(X=x,B^{\mathrm{obs}}=b,Y=1)}
      {P_Z(X=x,B^{\mathrm{obs}}=b)},
&P_Z(X=x,B^{\mathrm{obs}}=b)>0,\\[6pt]
0,&P_Z(X=x,B^{\mathrm{obs}}=b)=0.
\end{cases}
\]
The two observational arm probabilities sum to \(p_x\). Positivity therefore gives, when \(p_x>0\),
\[
P_Z(X=x,B^{\mathrm{obs}}=1)\ge qp_x>0,\qquad
P_Z(X=x,B^{\mathrm{obs}}=0)\ge qp_x>0.
\]
Thus these means agree with those in \cref{obj:def:zeng-discrete-ate-class} on occupied categories. Each ratio lies in \([0,1]\), since its numerator is the probability of a subset of the event in its denominator.

Define the binary mass functions
\[
\beta^{\mathrm{mark}}_{bx}:=
\begin{cases}
1-\pi_x,&b=0,\\
\pi_x,&b=1,
\end{cases}
\qquad
\beta^{\mathrm{out}}_{bxy}:=
\begin{cases}
1-\mu_{bx},&y=0,\\
\mu_{bx},&y=1,
\end{cases}
\quad
(x\in\{1,\ldots,d\},\ b,y\in\{0,1\}).
\]
Their masses are nonnegative and satisfy
\[
\sum_{b=0}^1\beta^{\mathrm{mark}}_{bx}=1,
\qquad
\sum_{y=0}^1\beta^{\mathrm{out}}_{bxy}=1.
\]
Let \(\mathbb P_{\mathrm{lat}}\) be the probability law on
\(\{0,\ldots,d-1\}\times\{0,1\}^4\), with coordinates
\((X,B^{\mathrm{lat}},Y(0),Y(1),A)\), specified by
\[
\mathbb P_{\mathrm{lat}}
\bigl(X=x-1,B^{\mathrm{lat}}=b,Y(0)=y_0,Y(1)=y_1,A=a\bigr)
=
\frac{p_x}{2}\,
\beta^{\mathrm{mark}}_{bx}
\beta^{\mathrm{out}}_{0xy_0}
\beta^{\mathrm{out}}_{1xy_1},
\]
for \(x\in\{1,\ldots,d\}\) and \(a,b,y_0,y_1\in\{0,1\}\).
The preceding normalization identities and \(\sum_xp_x=1\) show that this is a probability law.

On this probability space set
\[
S(0):=S(1):=0,\qquad
S:=S(A),\qquad
Y:=Y(A),\qquad
R:=\mathbf 1\{B^{\mathrm{lat}}=A\},
\]
and define the full-data law
\[
P_Z^*
:=
\mathcal L_{\mathbb P_{\mathrm{lat}}}
\bigl(X,A,S(0),S(1),Y(0),Y(1),S,Y,R\bigr).
\]
Here \(\mathcal L\) denotes the distribution of the displayed coordinates.
% lean: CausalSmith.Stat.MarRareqLogfrontier.SyntheticCompletion.fullLaw

\item Verify membership in the arrival-floor model. The product factor \(1/2\) in the construction gives
\[
A\perp\!\!\!\perp
\bigl(X,S(0),S(1),Y(0),Y(1)\bigr),
\qquad
P_Z^*(A=1)=\frac12.
\]
Together with \(S=S(A)\) and \(Y=Y(A)\), these establish
\cref{obj:ass:randomized-independence,obj:ass:balanced-randomization,obj:ass:surrogate-consistency,obj:ass:outcome-consistency}.

For \(a,r_{\mathrm{arr}}\in\{0,1\}\) and \(x\in\{1,\ldots,d\}\), define the arrival mass
\[
\beta^{\mathrm{arr}}_{axr_{\mathrm{arr}}}
:=
\begin{cases}
\beta^{\mathrm{mark}}_{ax},&r_{\mathrm{arr}}=1,\\
\beta^{\mathrm{mark}}_{1-a,x},&r_{\mathrm{arr}}=0.
\end{cases}
\]
Fix also \(y\in\{0,1\}\). The event \(R=r_{\mathrm{arr}}\), \(A=a\) fixes the latent mark to \(a\) when \(r_{\mathrm{arr}}=1\), and to \(1-a\) when \(r_{\mathrm{arr}}=0\). Summing over the unselected potential outcome gives
\begin{equation}\label{eq:synthetic-randomization-joint-mass}
P_Z^*(A=a,X=x-1,S=0,R=r_{\mathrm{arr}},Y=y)
=
\frac{p_x}{2}\,
\beta^{\mathrm{arr}}_{axr_{\mathrm{arr}}}
\beta^{\mathrm{out}}_{axy}.
\end{equation}
Indeed, the mass of the unselected potential outcome sums to one. Summing this identity over the binary outcome and arrival coordinates yields
\[
\begin{aligned}
p_{a,x-1,0}(P_Z^*)&=\frac{p_x}{2},\\
P_Z^*(A=a,X=x-1,S=0,R=r_{\mathrm{arr}})
&=\frac{p_x}{2}\beta^{\mathrm{arr}}_{axr_{\mathrm{arr}}},\\
P_Z^*(A=a,X=x-1,S=0,Y=y)
&=\frac{p_x}{2}\beta^{\mathrm{out}}_{axy}.
\end{aligned}
\]
Consequently,
\[
\begin{aligned}
&P_Z^*(A=a,X=x-1,S=0,R=r_{\mathrm{arr}},Y=y)\,
p_{a,x-1,0}(P_Z^*)\\
&\quad=
P_Z^*(A=a,X=x-1,S=0,R=r_{\mathrm{arr}})
P_Z^*(A=a,X=x-1,S=0,Y=y).
\end{aligned}
\]
For \(S=1\), every probability in this identity is zero, since \(S=0\) throughout the construction. This proves \cref{obj:ass:arrival-mar}, including cells of zero mass.

The same calculation gives
\[
P_Z^*(A=a,X=x-1,S=0,R=1)
=\frac{p_x}{2}\beta^{\mathrm{mark}}_{ax},
\qquad
p_{a,x-1,1}(P_Z^*)=0.
\]
Thus an occupied cell has \(S=0\) and \(p_x>0\). By the observational positivity bounds,
\[
\begin{aligned}
\rho_{0,x-1,0}(P_Z^*)&=1-\pi_x\ge q,\\
\rho_{1,x-1,0}(P_Z^*)&=\pi_x\ge q.
\end{aligned}
\]
This establishes \cref{obj:ass:occupied-cell-arrival}. Finally, \(n\ge1\), membership of \(P_Z\) in the observational class supplies \(d\ge1\), and \(0<q<1/2\) implies \(q\le1\). Hence
\[
P_Z^*\in\mathcal M^+_{n,d,q}
\]
by \cref{obj:def:unrestricted-arrival-model-class}.
% lean: CausalSmith.Stat.MarRareqLogfrontier.SyntheticCompletion.fullLaw_randomized
% lean: CausalSmith.Stat.MarRareqLogfrontier.SyntheticCompletion.fullLaw_balanced
% lean: CausalSmith.Stat.MarRareqLogfrontier.SyntheticCompletion.fullLaw_surrogate
% lean: CausalSmith.Stat.MarRareqLogfrontier.SyntheticCompletion.fullLaw_outcome
% lean: CausalSmith.Stat.MarRareqLogfrontier.SyntheticCompletion.fullLaw_mar
% lean: CausalSmith.Stat.MarRareqLogfrontier.SyntheticCompletion.fullLaw_arrival

\item Record the sampling property of the canonical observed product law. Define the observed-record alphabet and its one-record law by
\[
\mathcal O_d:=\{0,\ldots,d-1\}\times\{0,1\}^4,
\qquad
\nu_{\mathrm{obs}}
:=\mathcal L_{P_Z^*}(X,A,S,R,RY).
\]
For \(i\in\{1,\ldots,n\}\), let
\[
\operatorname{pr}_i:\mathcal O_d^n\longrightarrow\mathcal O_d,
\qquad
\operatorname{pr}_i(\mathbf o):=\mathbf o_i.
\]
For arbitrary subsets \(E_i\subseteq\mathcal O_d\), the defining property of the product measure gives
\[
\nu_{\mathrm{obs}}^{\otimes n}
\left(\bigcap_{i=1}^n\operatorname{pr}_i^{-1}(E_i)\right)
=
\prod_{i=1}^n\nu_{\mathrm{obs}}(E_i).
\]
The coordinate projections are measurable on these finite spaces. Taking all but one \(E_i\) to be the entire alphabet also gives the common marginal \(\nu_{\mathrm{obs}}\). Thus this product experiment satisfies \cref{obj:ass:iid-sampling}.
% lean: CausalSmith.Stat.MarRareqLogfrontier.sampleLaw_iid

\item Identify the kernel output law. First, the selected latent record has exactly the observational law:
\[
\mathcal L_{\mathbb P_{\mathrm{lat}}}
\bigl(X+1,B^{\mathrm{lat}},Y(B^{\mathrm{lat}})\bigr)=P_Z.
\]
To verify this identity, sum over \(A\) and the unselected potential outcome. For each \(x\in\{1,\ldots,d\}\) and \(b,y\in\{0,1\}\), the resulting singleton mass is
\[
\mathbb P_{\mathrm{lat}}
\bigl(X=x-1,B^{\mathrm{lat}}=b,Y(b)=y\bigr)
=
p_x\beta^{\mathrm{mark}}_{bx}\beta^{\mathrm{out}}_{bxy}.
\]
When \(p_x>0\),
\[
p_x\beta^{\mathrm{mark}}_{bx}
=P_Z(X=x,B^{\mathrm{obs}}=b).
\]
For \(y=1\), multiplying by \(\mu_{bx}\) gives the corresponding observational singleton probability by its defining ratio. For \(y=0\), subtracting that singleton probability from the arm probability gives the same conclusion:
\[
p_x\beta^{\mathrm{mark}}_{bx}\beta^{\mathrm{out}}_{bxy}
=P_Z(X=x,B^{\mathrm{obs}}=b,Y=y).
\]
When \(p_x=0\), both sides vanish because the singleton event is contained in the null baseline category. Equality of all singleton masses proves the selected-record identity.

Define the fair binary law and one-record kernel by
\[
\nu_{\mathrm{coin}}:=\frac12\delta_0+\frac12\delta_1,
\qquad
\kappa_{\mathrm{syn}}\bigl((x,b,y),\cdot\bigr)
:=
\frac12\delta_{o(x,b,y;0)}
+\frac12\delta_{o(x,b,y;1)}.
\]
This kernel is a measurable probability kernel on finite alphabets and depends solely on the input record. Under the completion, the observed tuple satisfies the pointwise identity
\[
(X,A,S,R,RY)
=
o\bigl(X+1,B^{\mathrm{lat}},Y(B^{\mathrm{lat}});A\bigr).
\]
For a matching assignment, the selected outcome equals \(Y(A)\); for a mismatching assignment, both arrived-outcome coordinates are zero. Moreover, \(A\) is independent of the selected latent record. The preceding distributional identity therefore gives
\[
\nu_{\mathrm{obs}}
=
P_Z\kappa_{\mathrm{syn}}
:=
\int\kappa_{\mathrm{syn}}\bigl((x,b,y),\cdot\bigr)\,
P_Z(d(x,b,y)).
\]

For a deterministic input sample
\[
\mathbf z=((x_i,b_i,y_i))_{i=1}^n
\in\bigl(\{1,\ldots,d\}\times\{0,1\}^2\bigr)^n
\]
and \(\mathbf u=(u_i)_{i=1}^n\in\{0,1\}^n\), define
\[
\mathsf o_n(\mathbf z,\mathbf u)
:=
\bigl(o(x_i,b_i,y_i;u_i)\bigr)_{i=1}^n,
\qquad
\kappa_{\mathrm{syn},n}(\mathbf z,\cdot)
:=
\bigotimes_{i=1}^n
\kappa_{\mathrm{syn}}\bigl((x_i,b_i,y_i),\cdot\bigr).
\]
Independent fair draws give
\[
\kappa_{\mathrm{syn},n}(\mathbf z,\cdot)
=
\bigl(\mathsf o_n(\mathbf z,\cdot)\bigr)_{\#}
\nu_{\mathrm{coin}}^{\otimes n},
\]
where the subscript \(\#\) denotes pushforward. Integrating over the input sample, regrouping the product coordinates into record--coin pairs, and then applying the record map coordinatewise gives
\[
\begin{aligned}
\int\kappa_{\mathrm{syn},n}(\mathbf z,\cdot)\,
P_Z^{\otimes n}(d\mathbf z)
&=
(\mathsf o_n)_{\#}
\bigl(P_Z^{\otimes n}\otimes\nu_{\mathrm{coin}}^{\otimes n}\bigr)\\
&=
\bigotimes_{i=1}^n
\left[
\int\kappa_{\mathrm{syn}}\bigl((x,b,y),\cdot\bigr)\,
P_Z(d(x,b,y))
\right]\\
&=
\nu_{\mathrm{obs}}^{\otimes n}.
\end{aligned}
\]
The regrouping is justified by the product identity
\[
\left(\prod_{i=1}^nP_Z(dz_i)\right)
\left(\prod_{i=1}^n\nu_{\mathrm{coin}}(du_i)\right)
=
\prod_{i=1}^n
\left(P_Z(dz_i)\nu_{\mathrm{coin}}(du_i)\right),
\]
and coordinatewise pushforward preserves this product form. Thus the kernel produces exactly the observed product law identified in the preceding step, and its output satisfies \cref{obj:ass:iid-sampling}.
% lean: CausalSmith.Stat.MarRareqLogfrontier.SyntheticCompletion.selected_eq
% lean: CausalSmith.Stat.MarRareqLogfrontier.SyntheticCompletion.observedLaw_eq_bind
% lean: CausalSmith.Stat.MarRareqLogfrontier.syntheticProductLaw_eq_coin_map
% lean: CausalSmith.Stat.MarRareqLogfrontier.syntheticSample_bind_eq_pi

\item Compute the target through cell contributions. Since \(P_Z^*\) belongs to the arrival-floor model, \cref{obj:lem:unrestricted-cell-identification} supplies
\[
\tau(P_Z^*)=\Phi(P_Z^*)
=
2\sum_{(a,x,s)\in\mathcal J_d}g(a)c_{axs}(P_Z^*).
\]
Here the contributions and their zero convention are given by \cref{obj:def:cell-functional}, and \cref{obj:synth_19} gives \(g(0)=-1\), \(g(1)=1\).

For \(p_x>0\), the arrival mass in either arm is strictly positive:
\[
P_Z^*(A=a,X=x-1,S=0,R=1)
=
\frac{p_x}{2}\beta^{\mathrm{mark}}_{ax}
\ge\frac{qp_x}{2}>0.
\]
The joint-mass identity in \cref{eq:synthetic-randomization-joint-mass} therefore yields
\[
\mu_{a,x-1,0}(P_Z^*)
=
\frac{\frac{p_x}{2}\beta^{\mathrm{mark}}_{ax}\mu_{ax}}
     {\frac{p_x}{2}\beta^{\mathrm{mark}}_{ax}}
=\mu_{ax},
\qquad
c_{a,x-1,0}(P_Z^*)=\frac{p_x}{2}\mu_{ax}.
\]
When \(p_x=0\), the arrived-cell probability is zero, so the defining convention gives the same contribution identity. Every cell with \(S=1\) also has zero arrived-cell probability. Hence, for every category and arm,
\[
c_{a,x-1,0}(P_Z^*)=\frac{p_x}{2}\mu_{ax},
\qquad
c_{a,x-1,1}(P_Z^*)=0.
\]
Substituting these identities into the signed cell sum gives
\[
\begin{aligned}
\tau(P_Z^*)
&=2\sum_{x=1}^d
\left(\frac{p_x}{2}\mu_{1x}
-\frac{p_x}{2}\mu_{0x}\right)\\
&=\sum_{x:p_x>0}p_x(\mu_{1x}-\mu_{0x})
=\theta_Z(P_Z),
\end{aligned}
\]
where the last equality uses \cref{obj:def:zeng-ate-functional}; null categories contribute zero.
% lean: CausalSmith.Stat.MarRareqLogfrontier.SyntheticCompletion.full_cellContribution_false
% lean: CausalSmith.Stat.MarRareqLogfrontier.SyntheticCompletion.full_cellContribution_true
% lean: CausalSmith.Stat.MarRareqLogfrontier.SyntheticCompletion.fullLaw_ate
\end{enumerate}
\end{proof}

The map preserves both the product observation law and the target at the corresponding full-data completion. Its parameter independence permits composition with any randomized estimator in \cref{obj:def:unrestricted-minimax-risk}. This fixes the direction of the experiment comparison used in \cref{obj:thm:unrestricted-fixed-q-minimax-frontier}: observational minimax lower bounds transfer to the unrestricted randomized model at the same arrival floor.

For construction, the same map supplies the synthetic records used by the observational procedure in \cref{obj:thm:zeng-fixed-positivity-polynomial-frontier}. The target equality connects its squared loss to that of the randomized estimator, and \cref{obj:lem:randomized-kernel-barycenter} supplies the risk comparison for averaging over the synthetic assignments. Thus the observation map supports both the fixed-arrival lower-bound comparison and the constructive observational extension.


\section{Risk certificates for the constructive estimator}

The risk analysis separates estimation of cell membership probabilities from estimation of arrived-cell outcome means. Independence holds in the uncapped Poisson comparison experiment used for this analysis. The reported finite-sample estimator instead averages assignments of a capped prefix, and the prefix-transfer argument below connects its risk to that independent-stream comparison. The polynomial branch controls the contribution of sparsely observed cells. The resulting certificate bounds squared error for the observed-cell functional in \cref{obj:def:cell-functional}; identification then connects that certificate to the causal upper bound in \cref{obj:thm:unrestricted-mixed-count-upper}.

For the deterministic mixed-count estimator \leanref{sym:\widehat\tau^{\mathrm{mix}}_{n,d,q}}{\ensuremath{\widehat\tau^{\mathrm{mix}}_{n,d,q}}} of \cref{obj:def:mixed-count-estimator}, write \(N=nq\) for the effective sample size, \(m=n/6\) for the intensity of each auxiliary stream, and \(N_0=mq=N/6\) for the effective stream size. The logarithmic scale \(\ell\), degree \(k\), and threshold \(B\) are specified in \cref{obj:synth_7,obj:synth_8,obj:synth_18}. Define the prefix-tail constant by
\[
\gamma_0=\log 2-\frac12.
\]
On the polynomial branch, where \(N>1\), \(\ell\ge128\), and \(d\ge\sqrt N\), introduce the certificate quantities
\[
b_{n,d,q}
=2\left\{\frac{4dB}{N_0k^2}+3e^{-16\ell}\right\}
\]
and
\[
v_{n,d,q}
=8\left(\frac4{N_0}+\frac1m\right)
+64d(e^{\ell/8}+1)
\left(\frac{B^2}{N_0^2}+\frac{B}{mN_0}\right)
+32e^{-32\ell}\left(1+\frac1m\right).
\]
Here \(b_{n,d,q}\) is the bias allowance and \(v_{n,d,q}\) is the additive second-moment allowance. These quantities summarize the polynomial approximation, marked-count fluctuations, and estimation of membership weights.

The deterministic squared-error envelope \leanref{sym:\mathfrak U_{n,d,q}}{\ensuremath{\mathfrak U_{n,d,q}}} combines these allowances with the elementary branch and the prefix-overflow contribution. The following statement gives its full formula and compares it with the rate \leanref{sym:r_{n,d,q}}{\ensuremath{r_{n,d,q}}} of \cref{obj:def:frontier-rate} throughout the unrestricted-arrival class \leanref{sym:\mathcal M^{+}_{n,d,q}}{\ensuremath{\mathcal M^{+}_{n,d,q}}} of \cref{obj:def:unrestricted-arrival-model-class}.

\begin{theoremv}{thm:unrestricted-uniform-mixed-count-certificate}[Uniform mixed-count risk certificate]
There exists a universal constant \(0<\overline C<\infty\) such that the following holds.
\begin{itemize}
\item \textbf{(Domain.)} Let \(n,d\ge1\) be integers and let \(0<q\le1\).
\item \textbf{(Model.)} Let \(P\in\mathcal M^{+}_{n,d,q}\), the unrestricted-arrival model class of \cref{obj:def:unrestricted-arrival-model-class}.
\item \textbf{(Sampling.)} Let \(\mathbf O_n\), the sample of observed records, satisfy \cref{obj:ass:iid-sampling} under \(P\).
\end{itemize}
Use the mixed-count estimator of \cref{obj:def:mixed-count-estimator}. Let \(\ell\), the logarithmic scale, be defined by \cref{obj:synth_7}; let \(k\), the polynomial degree, be defined by \cref{obj:synth_8}; and let \(B\), the threshold, be defined by \cref{obj:synth_18}. Put
\[
N=nq,\qquad m=\frac n6,\qquad N_0=mq=\frac N6,
\qquad \gamma_0=\log 2-\frac12.
\]
Define the deterministic envelope \(\mathfrak U_{n,d,q}\) as follows. If \(N\le1\), set
\[
\mathfrak U_{n,d,q}=1.
\]
If \(N>1\) and either \(\ell<128\) or \(d<\sqrt N\), set
\[
\mathfrak U_{n,d,q}
=\min\left\{4,
\left(\frac{8d}{eN_0}\right)^2
+4\left(\frac4{N_0}+\frac1m\right)
+4e^{-n\gamma_0}\right\}.
\]
If \(N>1\), \(\ell\ge128\), and \(d\ge\sqrt N\), put
\[
b_{n,d,q}
=2\left\{\frac{4dB}{N_0k^2}+3e^{-16\ell}\right\},
\]
\[
v_{n,d,q}
=8\left(\frac4{N_0}+\frac1m\right)
+64d(e^{\ell/8}+1)
\left(\frac{B^2}{N_0^2}+\frac{B}{mN_0}\right)
+32e^{-32\ell}\left(1+\frac1m\right),
\]
and set
\[
\mathfrak U_{n,d,q}
=\min\left\{4,b_{n,d,q}^2+v_{n,d,q}
+4e^{-n\gamma_0}\right\}.
\]
Then, with \(\Phi(P)\) the cell functional of \cref{obj:def:cell-functional} and \(r_{n,d,q}\) the rate of \cref{obj:def:frontier-rate}, expectation under the canonical i.i.d. observed-sample law satisfies
\[
E_P\!\left[
\left(\widehat\tau^{\mathrm{mix}}_{n,d,q}(\mathbf O_n)-\Phi(P)\right)^2
\right]
\le \mathfrak U_{n,d,q}
\le \overline C\,r_{n,d,q}.
\]
\end{theoremv}

The certificate distinguishes information from arrived outcomes, measured by \(N_0\), from information about cell membership, measured by \(m\). On the polynomial branch, the squared bias allowance captures the aggregate effect of estimating means in sparsely observed cells, while the additive allowance accounts for fluctuations in both components of the cell contribution. The elementary branch supplies the corresponding bound in its stated parameter range. Together, the branches give a deterministic precision benchmark depending on \(n,d,q\), with a universal comparison constant.

\subsection{Polynomial bias and marked-count second moments}

The polynomial statistic \(H_j\) and the pilot-selected statistic \(G_j\) are specified in \cref{obj:synth_14,obj:synth_15}. Their mathematical role is to replace the elementary cell statistic where the pilot count indicates sparse outcome arrival. The degree and threshold jointly determine the approximation contribution \(4dB/(N_0k^2)\) in the bias allowance. The exponentially small term in that allowance accounts for branch control. The certificate combines these contributions before squaring, preserving the distinction between bias and second-moment terms.

Marked counts retain the association between an arrived outcome and its cell. Consequently, the second-moment analysis concerns the marked polynomial statistic, together with the pilot-selected branch, rather than an unmarked count alone. The expression \(v_{n,d,q}\) records the resulting allowance: its terms involve the effective arrived-outcome intensity, the membership-stream intensity, the alphabet size, and the polynomial tuning parameters. This decomposition explains why the degree enters the bias allowance while the coefficient growth associated with the polynomial construction enters the additive allowance.

\subsection{Membership weights and branch control}

The membership weight \(W_j\) in \cref{obj:synth_11} is formed from the membership-labelled records. In the uncapped Poisson comparison experiment, that stream is independent of the pilot and estimation streams. The finite-sample estimator uses the capped prefix and conditional averaging in \cref{obj:def:mixed-count-estimator}; the overflow allowance and prefix-transfer comparison connect it to the independent-stream analysis. The weight supplies the population mass attached to each arrived-cell mean while retaining information from all observed memberships.

The two sample scales appear explicitly in the certificate through \(1/N_0\), \(1/m\), and the mixed term \(B/(mN_0)\). These terms retain the cost of estimating membership weights alongside the cost of estimating arrived-cell means. The pilot-selected statistic in \cref{obj:synth_15} organizes the latter task within the polynomial branch, whereas the elementary ratio in \cref{obj:synth_12} supplies the estimator on the elementary branch. The branch conditions in \cref{obj:thm:unrestricted-uniform-mixed-count-certificate} specify where each deterministic expression applies.

\subsection{Transfer to the fixed observed sample and causal target}

The auxiliary construction in \cref{obj:def:mixed-count-estimator} uses an independent Poisson prefix and uniform stream assignments. It assigns zero output when the prefix exceeds the available \(n\) records and clips the output to \([-1,1]\). The term \(4e^{-n\gamma_0}\) in \cref{obj:thm:unrestricted-uniform-mixed-count-certificate} is the prefix-overflow allowance. Its dependence on the full record count \(n\) reflects the supply of records for the auxiliary construction, while \(N=nq\) governs the effective outcome information. Clipping supplies the bounded-loss allowance represented by the minimum with \(4\); the zero estimator on the small-effective-sample branch supplies the bound \(1\).

Conditional averaging over the auxiliary randomization produces the deterministic estimator in \cref{obj:def:mixed-count-estimator}. The squared-loss comparison for such averaging is supplied by \cref{obj:lem:randomized-kernel-barycenter}. Thus the certificate concerns the estimator evaluated on the fixed observed sample, with the auxiliary construction serving its risk analysis.

For the causal interpretation, \cref{obj:lem:unrestricted-cell-identification} identifies the observed-cell functional with the average treatment effect under the unrestricted model conditions. Substituting that identity into \cref{obj:thm:unrestricted-uniform-mixed-count-certificate} gives the causal risk control used in \cref{obj:thm:unrestricted-mixed-count-upper}. The same identification step applies throughout the rare-arrival class \leanref{sym:\mathcal M_{n,d,q}}{\ensuremath{\mathcal M_{n,d,q}}} of \cref{obj:def:model-class}. The following specialization records the resulting envelope for that class.

\begin{theoremv}{thm:uniform-mixed-count-certificate}[Uniform mixed-count risk certificate]
There exists a universal constant \(0<\overline C<\infty\) such that the following holds for every \(n,d\in\mathbb N\), \(q\in\mathbb R\), and full-data law \(P\in\mathcal M_{n,d,q}\) satisfying \cref{obj:def:model-class}, including \(n\ge1\), \(d\ge1\), and \(0<q\le1\).

Let \(\widehat\tau^{\mathrm{mix}}_{n,d,q}\) be the estimator in \cref{obj:def:mixed-count-estimator}, and let \(r_{n,d,q}\) be the rate in \cref{obj:def:frontier-rate}. Let \(\ell\) be the logarithmic scale defined in \cref{obj:synth_7}, \(k\) the polynomial degree defined in \cref{obj:synth_8}, and \(B\) the threshold defined in \cref{obj:synth_18}. Put
\[
N=nq,\qquad m=\frac n6,\qquad N_0=mq=\frac N6,\qquad
\gamma_0=\log 2-\frac12,
\]
where \(N\) is the effective sample size, \(m\) the stream intensity, \(N_0\) the effective stream size, and \(\gamma_0\) the prefix-tail constant. Define the deterministic envelope \(\mathfrak U_{n,d,q}\) by the following branches:
\begin{itemize}
\item \textbf{(Small effective sample.)} If \(N\le1\), set
\[
\mathfrak U_{n,d,q}=1.
\]
\item \textbf{(Elementary branch.)} If \(N>1\) and either \(\ell<128\) or \(d<\sqrt N\), set
\[
\mathfrak U_{n,d,q}
=\min\left\{4,
\left(\frac{8d}{eN_0}\right)^2
+4\left(\frac4{N_0}+\frac1m\right)
+4e^{-n\gamma_0}\right\}.
\]
\item \textbf{(Polynomial branch.)} If \(N>1\), \(\ell\ge128\), and \(d\ge\sqrt N\), define the certificate quantities
\[
b_{n,d,q}
=2\left(\frac{4dB}{N_0k^2}+3e^{-16\ell}\right),
\]
\[
v_{n,d,q}
=8\left(\frac4{N_0}+\frac1m\right)
+64d\left(e^{\ell/8}+1\right)
\left(\frac{B^2}{N_0^2}+\frac{B}{mN_0}\right)
+32e^{-32\ell}\left(1+\frac1m\right),
\]
and set
\[
\mathfrak U_{n,d,q}
=\min\left\{4,b_{n,d,q}^2+v_{n,d,q}+4e^{-n\gamma_0}\right\}.
\]
\end{itemize}
Then, writing \(\mathbf O_n\) for the observed sample induced by \(P\) and \(\tau(P)\) for the average treatment effect in \cref{obj:def:ate-functional},
\[
E_P\!\left[
\left(\widehat\tau^{\mathrm{mix}}_{n,d,q}(\mathbf O_n)-\tau(P)\right)^2
\right]
\le \mathfrak U_{n,d,q}
\le \overline C\,r_{n,d,q}.
\]
\end{theoremv}

The rare-arrival restriction places the model within the domain of the unrestricted certificate, and \cref{obj:lem:cell-identification-background} fixes the causal target there. Accordingly, the specialization retains the same numerical envelope. Its polynomial bias, second-moment, and prefix-transfer components are those of \cref{obj:thm:unrestricted-uniform-mixed-count-certificate}; model inclusion and identification supply the specialization step.

The next statement records this causal upper bound together with the squared-loss comparison between the deterministic estimator and its auxiliary randomized output. It makes explicit the role of conditional averaging in the final estimator.

\begin{theoremv}{thm:mixed-count-upper}[Mixed-count risk upper bound]
There exists a universal constant \(0<\overline C<\infty\) such that the following holds for every choice of parameters and law satisfying:
\begin{itemize}
\item \textbf{(Sample size and dimension.)} \(n,d\in\mathbb N\) satisfy \(n,d\ge1\).
\item \textbf{(Arrival floor.)} \(q\in\mathbb R\) satisfies \(0<q\le1\).
\item \textbf{(Model.)} The full-data law \(P\) belongs to \(\mathcal M_{n,d,q}\) in \cref{obj:def:model-class}.
\end{itemize}
Let \(\widehat\tau^{\mathrm{mix}}_{n,d,q}\) and \(\widetilde\tau(\mathbf O_n;\omega)\) be the deterministic estimator and auxiliary randomized output of \cref{obj:def:mixed-count-estimator}, where \(\mathbf O_n\) is the observed sample and \(\omega\) is the auxiliary randomization. Let \(\ell\) be the logarithmic scale defined in \cref{obj:synth_7}, \(k\) the polynomial degree defined in \cref{obj:synth_8}, and \(B=256\ell\) the threshold defined in \cref{obj:synth_18}. Use the elementary and polynomial branch rules of \cref{obj:def:mixed-count-estimator}. Define the effective sample size \(N\), stream intensity \(m\), effective stream size \(N_0\), and prefix-tail constant \(\gamma_0\) by
\[
N=nq,\qquad m=\frac n6,\qquad N_0=mq=\frac N6,
\qquad \gamma_0=\log 2-\frac12.
\]
Define the certificate quantities
\[
b_{n,d,q}
=
2\left(\frac{4dB}{N_0k^2}+3e^{-16\ell}\right)
\]
and
\[
v_{n,d,q}
=
8\left(\frac4{N_0}+\frac1m\right)
+
64d\left(e^{\ell/8}+1\right)
\left(\frac{B^2}{N_0^2}+\frac{B}{mN_0}\right)
+
32e^{-32\ell}\left(1+\frac1m\right).
\]
The deterministic risk envelope \(\mathfrak U_{n,d,q}\) is defined as follows. For \(N\le1\), set
\[
\mathfrak U_{n,d,q}=1.
\]
For \(N>1\) on the elementary branch, set
\[
\mathfrak U_{n,d,q}
=
\min\left\{
4,\,
\left(\frac{8d}{eN_0}\right)^2
+4\left(\frac4{N_0}+\frac1m\right)
+4e^{-n\gamma_0}
\right\}.
\]
For \(N>1\) on the polynomial branch, set
\[
\mathfrak U_{n,d,q}
=
\min\left\{
4,\,
b_{n,d,q}^2+v_{n,d,q}+4e^{-n\gamma_0}
\right\}.
\]
Then, with \(\tau(P)\), the average treatment effect, defined in \cref{obj:def:ate-functional}, and \(r_{n,d,q}\), the frontier rate, defined in \cref{obj:def:frontier-rate},
\[
E_P\!\left[
\left(
\widehat\tau^{\mathrm{mix}}_{n,d,q}(\mathbf O_n)-\tau(P)
\right)^2
\right]
\le
\mathfrak U_{n,d,q}
\le
\overline C\,r_{n,d,q}.
\]
The deterministic estimator also satisfies
\[
E_P\!\left[
\left(
\widehat\tau^{\mathrm{mix}}_{n,d,q}(\mathbf O_n)-\tau(P)
\right)^2
\right]
\le
E_{P,\omega}\!\left[
\left(
\widetilde\tau(\mathbf O_n;\omega)-\tau(P)
\right)^2
\right],
\]
where the auxiliary expectation uses the independent Poisson prefix and uniform stream assignments, including the zero output when the prefix exceeds \(n\), specified in \cref{obj:def:mixed-count-estimator}.
\end{theoremv}

Averaging over the prefix and stream assignments removes auxiliary variation under squared loss. The first inequality in \cref{obj:thm:mixed-count-upper} is the certificate of \cref{obj:thm:uniform-mixed-count-certificate}, and the final comparison is the conditional-averaging relation of \cref{obj:lem:randomized-kernel-barycenter} applied to the specified auxiliary output. These references preserve the same construction and envelope across the two statements. For a prospective choice of \(n,d,q\), the envelope therefore supplies a deterministic causal risk benchmark that incorporates cell complexity, outcome arrival, and the fixed-sample implementation.


\section{Lower bounds with full observed memberships}

The lower comparisons concern the observed experiment in \cref{obj:def:observed-experiment}, which retains baseline category, assignment, surrogate, arrival status, and arrived outcome for every sampled individual. Two constructions account for the two components of the rate in \cref{obj:def:frontier-rate}: a one-category family isolates the effective sample-size contribution, while an activated many-category experiment supplies the category-complexity comparison. We first introduce the \leanref{S-3}{activated moment-matching construction}, then state its approximation inputs and the one-category bounds. A separate comparison addresses the ascertainment deficit in the unrestricted large-alphabet regime.

\subsection{The activated comparison experiment}

The construction in \cref{obj:def:activation-lower-handle} augments each observed record with an activation flag whose distribution is common to the two hypotheses. Conditional on a latent vector drawn once for the sample, it generates exactly \(n\) independent records with balanced treatment assignment. The following algorithm specifies both augmented laws and the conditions on their finite priors.

\begin{algorithmv}{def:activation-lower-handle}[Activated fuzzy-hypothesis experiment]
The activated handle \(\mathfrak H_{\mathrm{act}}\) is the ordered pair of augmented \(n\)-record laws constructed below, with the following inputs:
\begin{itemize}
\item \textbf{(Parameters.)} A constant \(\eta>0\), fixed before integers \(n,d\ge1\) and an arrival floor \(0<q\le1\), satisfying \cref{obj:ass:rare-arrival-slice}. Define the effective sample size, logarithmic scale, moment-matching degree, support endpoint, nominal rare-cell mass, and rare-cell count by
\[
N=nq,\qquad
\ell=\log(e+N),\qquad
K=\lceil\ell\rceil,\qquad
H=K^2,\qquad
b_0=\frac{\eta}{N\ell},\qquad
J=\min\left\{d-1,\left\lfloor\frac{1}{2b_0}\right\rfloor\right\}.
\]
\item \textbf{(Finite priors.)} An ordered pair \((\Pi_0,\Pi_1)\) of probability laws, each assigning probability one to a finite subset of \([1,H]\).
\item \textbf{(Moment matching and separation.)} The priors satisfy
\[
\int z^v\,d\Pi_0(z)=\int z^v\,d\Pi_1(z)
\qquad\text{for every integer }0\le v\le K,
\]
and
\[
\left|\int z^{-1}\,d\Pi_1(z)-\int z^{-1}\,d\Pi_0(z)\right|
\ge\frac1{12}.
\]
\end{itemize}

For each \(i\in\{0,1\}\), independently draw a latent coordinate \(Z_x\sim\Pi_i\) for every baseline index \(x\in\{0,\ldots,d-1\}\). Conditional on this single latent vector, generate \(n\) independent records using the following common rule. Draw the baseline index with probabilities
\[
P(X=x)=
\begin{cases}
b_0,&0\le x<J,\\
1-Jb_0,&x=J,\\
0,&J<x<d.
\end{cases}
\]
The first \(J\) indices are the rare cells, and index \(J\) is the reservoir. Draw treatment independently with
\[
A\sim\operatorname{Bernoulli}(1/2),
\qquad S=0.
\]
Independently of the latent vector, baseline index, and treatment, draw and reveal the activation flag
\[
\mathsf B_{\mathrm{act}}\sim\operatorname{Bernoulli}(qH).
\]
For a rare cell \(x\), the same coordinate \(Z_x\) governs both treatment arms and all records in that cell.

If \(\mathsf B_{\mathrm{act}}=0\), set \((R,RY)=(0,0)\). Conditional on \(\mathsf B_{\mathrm{act}}=1\), draw \((R,RY)\) according to
\[
\begin{array}{c|ccc}
 & (R,RY)=(1,1) & (R,RY)=(1,0) & (R,RY)=(0,0)\\ \hline
 A=1,\ X=x<J & 1/H & (Z_x-1)/H & 1-Z_x/H\\
 A=0,\ X=x<J & 0 & Z_x/H & 1-Z_x/H\\
 X\ge J,\ \text{either arm} & 0 & 1 & 0
\end{array}
\]
Each augmented record is
\[
\bigl(\mathsf B_{\mathrm{act}},(X,A,S,R,RY)\bigr).
\]
The \(i\)-th component of \(\mathfrak H_{\mathrm{act}}\) is the law of these \(n\) augmented records after integrating over the product prior \(\Pi_i^{\otimes d}\). Thus the latent vector is drawn once for the entire sample, and the records are independent conditional on that vector. The handle is defined for every prior pair satisfying the stated conditions.
\end{algorithmv}

The augmented pair \leanref{sym:\mathfrak H_{\mathrm{act}}}{\ensuremath{\mathfrak H_{\mathrm{act}}}} in \cref{obj:def:activation-lower-handle} preserves all original observations. Its revealed activation flag \leanref{sym:\mathsf B_{\mathrm{act}}}{\ensuremath{\mathsf B_{\mathrm{act}}}} has a distribution depending on the experiment's parameters and independent of the latent coordinates and prior index. Projection onto the original record coordinates is therefore a common, parameter-independent observation map. The total-variation comparison in \cref{obj:lem:randomized-kernel-tv-comparison} supplies the decision-theoretic connection between an observed-law comparison and any common randomized procedure.

The tuning constant \leanref{sym:\eta}{\ensuremath{\eta}}, moment-matching degree \leanref{sym:K}{\ensuremath{K}}, support endpoint \leanref{sym:H}{\ensuremath{H}}, nominal rare-cell mass \leanref{sym:b_0}{\ensuremath{b_0}}, and rare-cell count \leanref{sym:J}{\ensuremath{J}} are specified in \cref{obj:def:activation-lower-handle}. The mass cap reserves at least half of the baseline probability for the reservoir. Each rare cell has a single latent coordinate \leanref{sym:Z}{\ensuremath{Z_x}}, shared by its two treatment arms and all its sampled records. In the table, this coordinate governs arrival, while its reciprocal governs the treated arrived-outcome mean; the control outcome is zero. This arrangement places reciprocal separation and polynomial moment matching in the same experiment.

\subsection{Reciprocal approximation and finite priors}

The approximation inputs concern the interval and degree used by \cref{obj:def:activation-lower-handle}. Best uniform approximation quantifies the discrepancy between the reciprocal function and polynomials of the matching degree. The first input gives that discrepancy exactly.

\begin{lemmav}{lem:reciprocal-best-approximation}[Reciprocal best uniform approximation]
For every integer \(K\ge 2\), set \(H=K^2\). Then
\[
\inf_{\substack{p\ \mathrm{real\ polynomial}\\ \deg(p)\le K}}
\sup_{z\in[1,H]}\left|z^{-1}-p(z)\right|
=
\frac{H-1}{2H}
\left(\frac{K-1}{K+1}\right)^K.
\]
\end{lemmav}

\begin{proof}[Proof of \cref{obj:lem:reciprocal-best-approximation}]
\begin{enumerate}
\item Fix an integer \(K\ge 2\) and set
\[
H:=K^2>1.
\]
We first establish the approximation formula on a general positive interval. Let
\[
0<a_{\mathrm{app}}<b_{\mathrm{app}},
\qquad m_{\mathrm{app}}\in\mathbb N_0,
\]
and define
\[
\Delta_{\mathrm{app}}:=b_{\mathrm{app}}-a_{\mathrm{app}},
\qquad
v_{\mathrm{app}}:=\frac{a_{\mathrm{app}}+b_{\mathrm{app}}}
{\Delta_{\mathrm{app}}},
\qquad
s_{\mathrm{app}}:=\sqrt{v_{\mathrm{app}}^2-1},
\]
\[
\rho_{\mathrm{app}}:=v_{\mathrm{app}}-s_{\mathrm{app}},
\qquad
E_{\mathrm{app}}:=
\frac{2\rho_{\mathrm{app}}^{m_{\mathrm{app}}}}
{\Delta_{\mathrm{app}}(v_{\mathrm{app}}^2-1)}.
\]
Since \(\Delta_{\mathrm{app}}>0\) and
\(a_{\mathrm{app}}+b_{\mathrm{app}}>\Delta_{\mathrm{app}}\), we have
\(v_{\mathrm{app}}>1\). Consequently,
\[
0\le s_{\mathrm{app}}<v_{\mathrm{app}},
\qquad
s_{\mathrm{app}}^2=v_{\mathrm{app}}^2-1,
\qquad
(v_{\mathrm{app}}-s_{\mathrm{app}})
(v_{\mathrm{app}}+s_{\mathrm{app}})=1.
\]
Thus \(\rho_{\mathrm{app}}>0\), and
\(v_{\mathrm{app}}+s_{\mathrm{app}}>1\) gives
\[
0<\rho_{\mathrm{app}}<1,
\qquad
\rho_{\mathrm{app}}^{-1}=v_{\mathrm{app}}+s_{\mathrm{app}},
\qquad
\rho_{\mathrm{app}}+\rho_{\mathrm{app}}^{-1}=2v_{\mathrm{app}}.
\]
In particular,
\[
\rho_{\mathrm{app}}^2-2v_{\mathrm{app}}\rho_{\mathrm{app}}+1=0,
\qquad E_{\mathrm{app}}>0.
\]
% lean: Causalean.Mathlib.Analysis.Approximation.Chebyshev.Reciprocal.intervalDecay_add_inv

\item Define the Chebyshev polynomials, with integer indices, by
\[
T_0(t_{\mathrm{app}})=1,\qquad
T_1(t_{\mathrm{app}})=t_{\mathrm{app}},\qquad
T_{\nu+2}(t_{\mathrm{app}})
=2t_{\mathrm{app}}T_{\nu+1}(t_{\mathrm{app}})
-T_\nu(t_{\mathrm{app}})
\quad(\nu\in\mathbb N_0),
\]
\[
T_{-\nu}=T_\nu\quad(\nu\in\mathbb N_0).
\]
Here \(t_{\mathrm{app}}\in\mathbb R\). These polynomials satisfy
\[
\deg T_\nu=|\nu|,
\qquad
T_\nu(\cos\theta_{\mathrm{app}})
=\cos(\nu\theta_{\mathrm{app}})
\quad(\nu\in\mathbb Z,\ \theta_{\mathrm{app}}\in\mathbb R).
\]
For every \(\nu\in\mathbb N_0\), define the residual polynomial
\[
\mathcal R_{\mathrm{app},\nu}(t_{\mathrm{app}})
:=
\rho_{\mathrm{app}}^{-1}T_{\nu+1}(t_{\mathrm{app}})
-2T_\nu(t_{\mathrm{app}})
+\rho_{\mathrm{app}}T_{\nu-1}(t_{\mathrm{app}}).
\]
The degree bounds for its three terms give
\[
\deg\mathcal R_{\mathrm{app},\nu}\le \nu+1.
\]
The Chebyshev recurrence also gives
\[
\mathcal R_{\mathrm{app},\nu+2}(t_{\mathrm{app}})
=
2t_{\mathrm{app}}\mathcal R_{\mathrm{app},\nu+1}(t_{\mathrm{app}})
-\mathcal R_{\mathrm{app},\nu}(t_{\mathrm{app}}).
\]
We claim that
\[
\mathcal R_{\mathrm{app},\nu}(v_{\mathrm{app}})
=
\frac{2(v_{\mathrm{app}}^2-1)}
{\rho_{\mathrm{app}}^\nu}
\qquad(\nu\in\mathbb N_0).
\]
For the two initial indices, the identities from the preceding step yield
\[
\begin{aligned}
\mathcal R_{\mathrm{app},0}(v_{\mathrm{app}})
&=(\rho_{\mathrm{app}}^{-1}+\rho_{\mathrm{app}})
v_{\mathrm{app}}-2
=2(v_{\mathrm{app}}^2-1),\\
\mathcal R_{\mathrm{app},1}(v_{\mathrm{app}})
&=\rho_{\mathrm{app}}^{-1}(2v_{\mathrm{app}}^2-1)
-2v_{\mathrm{app}}+\rho_{\mathrm{app}}
=\frac{2(v_{\mathrm{app}}^2-1)}{\rho_{\mathrm{app}}}.
\end{aligned}
\]
If the claim holds at \(\nu\) and \(\nu+1\), the recurrence gives
\[
\begin{aligned}
\mathcal R_{\mathrm{app},\nu+2}(v_{\mathrm{app}})
&=
2(v_{\mathrm{app}}^2-1)
\left(
\frac{2v_{\mathrm{app}}}{\rho_{\mathrm{app}}^{\nu+1}}
-\frac{1}{\rho_{\mathrm{app}}^\nu}
\right)\\
&=
\frac{2(v_{\mathrm{app}}^2-1)}
{\rho_{\mathrm{app}}^{\nu+2}},
\end{aligned}
\]
because \(2v_{\mathrm{app}}\rho_{\mathrm{app}}
-\rho_{\mathrm{app}}^2=1\). This proves the claim by induction, including the degree-zero case.
% lean: Causalean.Mathlib.Analysis.Approximation.Chebyshev.Reciprocal.reciprocalResidualPoly_at_shape

\item Define the affine map
\[
t_{\mathrm{app}}(z):=
v_{\mathrm{app}}-\frac{2z}{\Delta_{\mathrm{app}}}
\qquad(z\in\mathbb R).
\]
Since \(\mathcal R_{\mathrm{app},m_{\mathrm{app}}}
(v_{\mathrm{app}})>0\), the polynomial
\[
1-
\frac{\mathcal R_{\mathrm{app},m_{\mathrm{app}}}
(t_{\mathrm{app}}(z))}
{\mathcal R_{\mathrm{app},m_{\mathrm{app}}}(v_{\mathrm{app}})}
\]
is well defined and vanishes at \(z=0\). It therefore determines a unique real polynomial \(p_{\mathrm{app}}\) through
\[
z\,p_{\mathrm{app}}(z)
=
1-
\frac{\mathcal R_{\mathrm{app},m_{\mathrm{app}}}
(t_{\mathrm{app}}(z))}
{\mathcal R_{\mathrm{app},m_{\mathrm{app}}}(v_{\mathrm{app}})}
\qquad(z\in\mathbb R).
\]
The numerator has degree at most \(m_{\mathrm{app}}+1\), so division by \(z\) gives
\[
\deg p_{\mathrm{app}}\le m_{\mathrm{app}}.
\]
For every \(z\ne0\), rearranging the defining identity yields
\[
z^{-1}-p_{\mathrm{app}}(z)
=
\frac{\mathcal R_{\mathrm{app},m_{\mathrm{app}}}
(t_{\mathrm{app}}(z))}
{z\,\mathcal R_{\mathrm{app},m_{\mathrm{app}}}
(v_{\mathrm{app}})}.
\]
% lean: Causalean.Mathlib.Analysis.Approximation.Chebyshev.Reciprocal.reciprocalApproxPoly_residual

\item For \(\theta_{\mathrm{app}}\in\mathbb R\), define
\[
\alpha_{\mathrm{app},+}(\theta_{\mathrm{app}})
:=\frac{(m_{\mathrm{app}}+1)\theta_{\mathrm{app}}}{2},
\qquad
\alpha_{\mathrm{app},-}(\theta_{\mathrm{app}})
:=\frac{(m_{\mathrm{app}}-1)\theta_{\mathrm{app}}}{2},
\]
\[
F_{\mathrm{app}}^{\sin}(\theta_{\mathrm{app}})
:=
\sin\alpha_{\mathrm{app},+}(\theta_{\mathrm{app}})
-\rho_{\mathrm{app}}
\sin\alpha_{\mathrm{app},-}(\theta_{\mathrm{app}}),
\]
\[
F_{\mathrm{app}}^{\cos}(\theta_{\mathrm{app}})
:=
\cos\alpha_{\mathrm{app},+}(\theta_{\mathrm{app}})
-\rho_{\mathrm{app}}
\cos\alpha_{\mathrm{app},-}(\theta_{\mathrm{app}}).
\]
The Chebyshev cosine identity gives
\[
\begin{aligned}
\mathcal R_{\mathrm{app},m_{\mathrm{app}}}
(\cos\theta_{\mathrm{app}})
={}&
\rho_{\mathrm{app}}^{-1}
\cos((m_{\mathrm{app}}+1)\theta_{\mathrm{app}})
-2\cos(m_{\mathrm{app}}\theta_{\mathrm{app}})\\
&+\rho_{\mathrm{app}}
\cos((m_{\mathrm{app}}-1)\theta_{\mathrm{app}}).
\end{aligned}
\]
To factor this expression, use
\[
\cos((m_{\mathrm{app}}\pm1)\theta_{\mathrm{app}})
=
1-2\sin^2\alpha_{\mathrm{app},\pm}(\theta_{\mathrm{app}})
=
2\cos^2\alpha_{\mathrm{app},\pm}(\theta_{\mathrm{app}})-1,
\]
\[
\begin{aligned}
\cos(m_{\mathrm{app}}\theta_{\mathrm{app}})
-\cos\theta_{\mathrm{app}}
&=-2\sin\alpha_{\mathrm{app},+}(\theta_{\mathrm{app}})
\sin\alpha_{\mathrm{app},-}(\theta_{\mathrm{app}}),\\
\cos(m_{\mathrm{app}}\theta_{\mathrm{app}})
+\cos\theta_{\mathrm{app}}
&=2\cos\alpha_{\mathrm{app},+}(\theta_{\mathrm{app}})
\cos\alpha_{\mathrm{app},-}(\theta_{\mathrm{app}}).
\end{aligned}
\]
Substitution, together with
\(\rho_{\mathrm{app}}+\rho_{\mathrm{app}}^{-1}
=2v_{\mathrm{app}}\), produces the two identities
\[
\begin{aligned}
\mathcal R_{\mathrm{app},m_{\mathrm{app}}}
(\cos\theta_{\mathrm{app}})
&=
2(v_{\mathrm{app}}-\cos\theta_{\mathrm{app}})
-\frac{2}{\rho_{\mathrm{app}}}
\bigl(F_{\mathrm{app}}^{\sin}(\theta_{\mathrm{app}})\bigr)^2,\\
\mathcal R_{\mathrm{app},m_{\mathrm{app}}}
(\cos\theta_{\mathrm{app}})
&=
-2(v_{\mathrm{app}}-\cos\theta_{\mathrm{app}})
+\frac{2}{\rho_{\mathrm{app}}}
\bigl(F_{\mathrm{app}}^{\cos}(\theta_{\mathrm{app}})\bigr)^2.
\end{aligned}
\]
For \(z\in[a_{\mathrm{app}},b_{\mathrm{app}}]\),
\[
t_{\mathrm{app}}(z)
=\frac{a_{\mathrm{app}}+b_{\mathrm{app}}-2z}
{\Delta_{\mathrm{app}}}
\in[-1,1].
\]
Choose \(\theta_{\mathrm{app}}\in[0,\pi]\) with
\(\cos\theta_{\mathrm{app}}=t_{\mathrm{app}}(z)\).
The nonnegative squares in the two identities imply
\[
-2(v_{\mathrm{app}}-t_{\mathrm{app}}(z))
\le
\mathcal R_{\mathrm{app},m_{\mathrm{app}}}(t_{\mathrm{app}}(z))
\le
2(v_{\mathrm{app}}-t_{\mathrm{app}}(z)).
\]
Hence
\[
\left|
\mathcal R_{\mathrm{app},m_{\mathrm{app}}}(t_{\mathrm{app}}(z))
\right|
\le
2(v_{\mathrm{app}}-t_{\mathrm{app}}(z))
=\frac{4z}{\Delta_{\mathrm{app}}}.
\]
% lean: Causalean.Mathlib.Analysis.Approximation.Chebyshev.Reciprocal.reciprocalResidualPoly_bound

\item Every \(z\in[a_{\mathrm{app}},b_{\mathrm{app}}]\) is positive. The residual identity, the preceding bound, and the value of the residual polynomial at \(v_{\mathrm{app}}\) therefore give
\[
\begin{aligned}
|z^{-1}-p_{\mathrm{app}}(z)|
&=
\frac{
|\mathcal R_{\mathrm{app},m_{\mathrm{app}}}(t_{\mathrm{app}}(z))|
}{
z\,\mathcal R_{\mathrm{app},m_{\mathrm{app}}}(v_{\mathrm{app}})
}\\
&\le
\frac{4z/\Delta_{\mathrm{app}}}
{z\,\mathcal R_{\mathrm{app},m_{\mathrm{app}}}(v_{\mathrm{app}})}
=
\frac{2\rho_{\mathrm{app}}^{m_{\mathrm{app}}}}
{\Delta_{\mathrm{app}}(v_{\mathrm{app}}^2-1)}
=E_{\mathrm{app}}.
\end{aligned}
\]
Taking the supremum yields
\[
\sup_{z\in[a_{\mathrm{app}},b_{\mathrm{app}}]}
|z^{-1}-p_{\mathrm{app}}(z)|
\le E_{\mathrm{app}}.
\]
% lean: Causalean.Mathlib.Analysis.Approximation.Chebyshev.Reciprocal.reciprocalApproxPoly_upper

\item We next construct alternating equality points. Define the grid
\[
\gamma_{\mathrm{app},\iota}
:=
\frac{\iota\pi}{m_{\mathrm{app}}+1}
\qquad
(\iota\in\{0,\ldots,m_{\mathrm{app}}+1\}).
\]
It is strictly increasing from \(0\) to \(\pi\). At its points,
\[
\alpha_{\mathrm{app},+}(\gamma_{\mathrm{app},\iota})
=\frac{\iota\pi}{2},
\qquad
\alpha_{\mathrm{app},-}(\gamma_{\mathrm{app},\iota})
=\frac{\iota\pi}{2}-\gamma_{\mathrm{app},\iota}.
\]
The sine and cosine subtraction formulas consequently give, for
\(\eta_{\mathrm{app}}\in\mathbb N_0\) whenever the displayed grid indices are in range,
\[
\begin{aligned}
F_{\mathrm{app}}^{\cos}(\gamma_{\mathrm{app},2\eta_{\mathrm{app}}})
&=(-1)^{\eta_{\mathrm{app}}}
\bigl(1-\rho_{\mathrm{app}}
\cos\gamma_{\mathrm{app},2\eta_{\mathrm{app}}}\bigr),\\
F_{\mathrm{app}}^{\sin}(\gamma_{\mathrm{app},2\eta_{\mathrm{app}}})
&=\rho_{\mathrm{app}}(-1)^{\eta_{\mathrm{app}}}
\sin\gamma_{\mathrm{app},2\eta_{\mathrm{app}}},\\
F_{\mathrm{app}}^{\cos}(\gamma_{\mathrm{app},2\eta_{\mathrm{app}}+1})
&=-\rho_{\mathrm{app}}(-1)^{\eta_{\mathrm{app}}}
\sin\gamma_{\mathrm{app},2\eta_{\mathrm{app}}+1},\\
F_{\mathrm{app}}^{\sin}(\gamma_{\mathrm{app},2\eta_{\mathrm{app}}+1})
&=(-1)^{\eta_{\mathrm{app}}}
\bigl(1-\rho_{\mathrm{app}}
\cos\gamma_{\mathrm{app},2\eta_{\mathrm{app}}+1}\bigr).
\end{aligned}
\]
The magnitudes controlling the left and right endpoints satisfy
\[
1-\rho_{\mathrm{app}}\cos\theta_{\mathrm{app}}
\ge1-\rho_{\mathrm{app}}>0
\quad(\theta_{\mathrm{app}}\in\mathbb R),
\qquad
\rho_{\mathrm{app}}\sin\theta_{\mathrm{app}}\ge0
\quad(\theta_{\mathrm{app}}\in[0,\pi]).
\]
Thus, on a cell with odd right index, the cosine factor has a nonzero left value and a right value of the opposite weak sign. On a cell with positive even right index, the sine factor has this property, since
\[
F_{\mathrm{app}}^{\sin}
(\gamma_{\mathrm{app},2(\eta_{\mathrm{app}}+1)})
=
-(-1)^{\eta_{\mathrm{app}}}\rho_{\mathrm{app}}
\sin\gamma_{\mathrm{app},2(\eta_{\mathrm{app}}+1)}.
\]
Continuity and the intermediate value theorem give choices satisfying
\[
\theta_{\mathrm{app},0}:=0,
\qquad
\gamma_{\mathrm{app},\iota-1}
<\theta_{\mathrm{app},\iota}
\le\gamma_{\mathrm{app},\iota}
\quad(1\le\iota\le m_{\mathrm{app}}+1),
\]
\[
\begin{cases}
F_{\mathrm{app}}^{\sin}(\theta_{\mathrm{app},\iota})=0,
&\iota\ \text{even},\\
F_{\mathrm{app}}^{\cos}(\theta_{\mathrm{app},\iota})=0,
&\iota\ \text{odd}.
\end{cases}
\]
The left endpoints are excluded because their factor values are nonzero. Also
\(F_{\mathrm{app}}^{\sin}(0)=0\). The cell bounds therefore imply
\[
0=\theta_{\mathrm{app},0}
<\theta_{\mathrm{app},1}<\cdots
<\theta_{\mathrm{app},m_{\mathrm{app}}+1}\le\pi.
\]
This construction applies also when \(m_{\mathrm{app}}=0\).
% lean: Causalean.Mathlib.Analysis.Approximation.Chebyshev.Reciprocal.exists_reciprocalAlternatingAngles

\item Define the nodes
\[
z_{\mathrm{app},\iota}:=
\frac{a_{\mathrm{app}}+b_{\mathrm{app}}
-\Delta_{\mathrm{app}}\cos\theta_{\mathrm{app},\iota}}{2}
\qquad(0\le\iota\le m_{\mathrm{app}}+1).
\]
Because cosine is strictly decreasing on \([0,\pi]\),
\[
a_{\mathrm{app}}\le z_{\mathrm{app},0}
<z_{\mathrm{app},1}<\cdots
<z_{\mathrm{app},m_{\mathrm{app}}+1}\le b_{\mathrm{app}},
\qquad
t_{\mathrm{app}}(z_{\mathrm{app},\iota})
=\cos\theta_{\mathrm{app},\iota}.
\]
For even \(\iota\), the sine-square identity has zero square; for odd \(\iota\), the cosine-square identity has zero square. Hence
\[
\mathcal R_{\mathrm{app},m_{\mathrm{app}}}
(t_{\mathrm{app}}(z_{\mathrm{app},\iota}))
=
(-1)^\iota
2\bigl(v_{\mathrm{app}}-t_{\mathrm{app}}(z_{\mathrm{app},\iota})\bigr)
=
(-1)^\iota\frac{4z_{\mathrm{app},\iota}}{\Delta_{\mathrm{app}}}.
\]
Substituting into the residual identity gives
\[
z_{\mathrm{app},\iota}^{-1}
-p_{\mathrm{app}}(z_{\mathrm{app},\iota})
=
(-1)^\iota
\frac{4}
{\Delta_{\mathrm{app}}\,
\mathcal R_{\mathrm{app},m_{\mathrm{app}}}(v_{\mathrm{app}})}
=
(-1)^\iota E_{\mathrm{app}}.
\]
% lean: Causalean.Mathlib.Analysis.Approximation.Chebyshev.Reciprocal.reciprocalApproxPoly_alternates

\item To certify optimality, define the signed Lagrange weights
\[
\beta_{\mathrm{app},\iota}:=
\prod_{\substack{0\le\jmath\le m_{\mathrm{app}}+1\\\jmath\ne\iota}}
(z_{\mathrm{app},\iota}-z_{\mathrm{app},\jmath})^{-1},
\qquad
\Lambda_{\mathrm{app}}:=
\sum_{\jmath=0}^{m_{\mathrm{app}}+1}
|\beta_{\mathrm{app},\jmath}|,
\qquad
w_{\mathrm{app},\iota}:=
\frac{\beta_{\mathrm{app},\iota}}{\Lambda_{\mathrm{app}}}
\quad(0\le\iota\le m_{\mathrm{app}}+1).
\]
Distinctness of the nodes makes every \(\beta_{\mathrm{app},\iota}\) nonzero, so \(\Lambda_{\mathrm{app}}>0\). There are exactly
\(m_{\mathrm{app}}+1-\iota\) negative factors in the product defining
\(\beta_{\mathrm{app},\iota}\). Therefore
\[
\sum_{\iota=0}^{m_{\mathrm{app}}+1}|w_{\mathrm{app},\iota}|=1,
\qquad
w_{\mathrm{app},\iota}
=
(-1)^{m_{\mathrm{app}}+1-\iota}
|w_{\mathrm{app},\iota}|.
\]
For any real polynomial \(q_{\mathrm{app}}\) of degree at most
\(m_{\mathrm{app}}\), Lagrange interpolation gives the polynomial identity
\[
q_{\mathrm{app}}(z)
=
\sum_{\iota=0}^{m_{\mathrm{app}}+1}
\beta_{\mathrm{app},\iota}\,
q_{\mathrm{app}}(z_{\mathrm{app},\iota})
\prod_{\substack{0\le\jmath\le m_{\mathrm{app}}+1\\\jmath\ne\iota}}
(z-z_{\mathrm{app},\jmath}).
\]
Comparing coefficients of \(z^{m_{\mathrm{app}}+1}\), and then dividing by
\(\Lambda_{\mathrm{app}}\), yields
\[
\sum_{\iota=0}^{m_{\mathrm{app}}+1}
\beta_{\mathrm{app},\iota}
q_{\mathrm{app}}(z_{\mathrm{app},\iota})=0,
\qquad
\sum_{\iota=0}^{m_{\mathrm{app}}+1}
w_{\mathrm{app},\iota}
q_{\mathrm{app}}(z_{\mathrm{app},\iota})=0.
\]
This cancellation applies to \(p_{\mathrm{app}}\) as well. Combining the weight signs with the alternating residuals gives
\[
\begin{aligned}
\sum_{\iota=0}^{m_{\mathrm{app}}+1}
w_{\mathrm{app},\iota}
\bigl(z_{\mathrm{app},\iota}^{-1}
-p_{\mathrm{app}}(z_{\mathrm{app},\iota})\bigr)
&=
(-1)^{m_{\mathrm{app}}+1}E_{\mathrm{app}}
\sum_{\iota=0}^{m_{\mathrm{app}}+1}
|w_{\mathrm{app},\iota}|\\
&=(-1)^{m_{\mathrm{app}}+1}E_{\mathrm{app}}.
\end{aligned}
\]
The cancellation identities consequently imply, for every such
\(q_{\mathrm{app}}\),
\[
\sum_{\iota=0}^{m_{\mathrm{app}}+1}
w_{\mathrm{app},\iota}
\bigl(z_{\mathrm{app},\iota}^{-1}
-q_{\mathrm{app}}(z_{\mathrm{app},\iota})\bigr)
=
(-1)^{m_{\mathrm{app}}+1}E_{\mathrm{app}}.
\]
Taking absolute values and using the unit absolute mass,
\[
\begin{aligned}
E_{\mathrm{app}}
&\le
\sum_{\iota=0}^{m_{\mathrm{app}}+1}
|w_{\mathrm{app},\iota}|
\left|z_{\mathrm{app},\iota}^{-1}
-q_{\mathrm{app}}(z_{\mathrm{app},\iota})\right|\\
&\le
\sup_{z\in[a_{\mathrm{app}},b_{\mathrm{app}}]}
|z^{-1}-q_{\mathrm{app}}(z)|.
\end{aligned}
\]
The residuals are continuous on this compact positive interval, so these suprema are finite. The lower bound for every admissible polynomial and the upper bound attained by \(p_{\mathrm{app}}\) establish
\[
\inf_{\substack{q_{\mathrm{app}}\ \mathrm{real\ polynomial}\\
\deg(q_{\mathrm{app}})\le m_{\mathrm{app}}}}
\sup_{z\in[a_{\mathrm{app}},b_{\mathrm{app}}]}
|z^{-1}-q_{\mathrm{app}}(z)|
=
\frac{2\rho_{\mathrm{app}}^{m_{\mathrm{app}}}}
{\Delta_{\mathrm{app}}(v_{\mathrm{app}}^2-1)}.
\]
% lean: Causalean.Mathlib.Analysis.Approximation.Chebyshev.Reciprocal.bestUniformApproxError_reciprocal

\item Specialize to
\[
a_{\mathrm{app}}=1,\qquad b_{\mathrm{app}}=H=K^2,
\qquad m_{\mathrm{app}}=K.
\]
These parameters meet the interval and degree conditions because \(K\ge2\).
Now
\[
v_{\mathrm{app}}=\frac{1+K^2}{K^2-1},
\qquad
(1+K^2)^2-(K^2-1)^2=4K^2,
\qquad
v_{\mathrm{app}}^2-1=\frac{4K^2}{(K^2-1)^2}.
\]
Both \(K\) and \(K^2-1\) are positive, so
\[
\sqrt{v_{\mathrm{app}}^2-1}
=\frac{2K}{K^2-1},
\qquad
\rho_{\mathrm{app}}
=
\frac{1+K^2-2K}{K^2-1}
=
\frac{K-1}{K+1}.
\]
Finally,
\[
\frac{2}{\Delta_{\mathrm{app}}(v_{\mathrm{app}}^2-1)}
=
\frac{K^2-1}{2K^2}
=
\frac{H-1}{2H}.
\]
Substitution into the established interval formula gives
\[
\inf_{\substack{p\ \mathrm{real\ polynomial}\\ \deg(p)\le K}}
\sup_{z\in[1,H]}|z^{-1}-p(z)|
=
\frac{H-1}{2H}
\left(\frac{K-1}{K+1}\right)^K,
\]
as required.
% lean: Causalean.Mathlib.Analysis.Approximation.Chebyshev.Reciprocal.bestUniformApproxError_reciprocal_one_sq
\end{enumerate}
\end{proof}

\Cref{obj:lem:reciprocal-best-approximation} identifies the approximation scale underlying the prior construction. Its mathematical role is to distinguish the reciprocal functional from the polynomial moments used to compare the two experiments. The next input supplies finite probability laws with exactly matching moments and a fixed reciprocal-mean separation.

\begin{lemmav}{lem:moment-matched-reciprocal-priors}[Moment-matched reciprocal priors]
For every integer \(K\ge 2\), set \(H=K^2\). There exist finitely supported probability laws \(\Pi_0\) and \(\Pi_1\) on \([1,H]\) such that
\[
\int z^v\,d\Pi_0(z)=\int z^v\,d\Pi_1(z)
\qquad\text{for every integer }v\in\{0,\ldots,K\},
\]
and
\[
\left|\int z^{-1}\,d\Pi_1(z)-\int z^{-1}\,d\Pi_0(z)\right|
\ge \frac{1}{12}.
\]
\end{lemmav}

\begin{proof}[Proof of \cref{obj:lem:moment-matched-reciprocal-priors}]
\begin{enumerate}
\item Fix \(K\ge2\), and define
\[
H=K^2,\qquad
\mathcal I_\star=[1,H],\qquad
f_\star:\mathcal I_\star\to\mathbb R,\quad f_\star(z)=z^{-1}.
\]
Since \(H>1\) and every point of \(\mathcal I_\star\) is positive, \(f_\star\) is continuous. For a continuous function \(\varphi:\mathcal I_\star\to\mathbb R\), write
\[
\|\varphi\|_\infty=\max_{z\in\mathcal I_\star}|\varphi(z)|,
\qquad
E_K=\inf_{\substack{\mathfrak p\in\mathbb R[z]\\ \deg\mathfrak p\le K}}
\|f_\star-\mathfrak p\|_\infty.
\]
In particular, \(\|f_\star\|_\infty=1\).

To obtain a best approximant, consider the finite-dimensional space
\[
\mathcal V_K
=\{\mathfrak p|_{\mathcal I_\star}:
  \mathfrak p\in\mathbb R[z],\ \deg\mathfrak p\le K\},
\qquad
\mathcal B_\star
=\{\varphi\in\mathcal V_K:\|\varphi\|_\infty\le3\}.
\]
The closed ball \(\mathcal B_\star\) is compact. The continuous error function therefore attains its minimum there at the restriction of a polynomial \(p_\star\) of degree at most \(K\). Comparison with the zero polynomial gives
\[
\|f_\star-p_\star\|_\infty\le1.
\]
For every \(\varphi\in\mathcal V_K\setminus\mathcal B_\star\), the reverse triangle inequality gives
\[
\|f_\star-\varphi\|_\infty
\ge\|\varphi\|_\infty-\|f_\star\|_\infty
>3-1=2
\ge\|f_\star-p_\star\|_\infty.
\]
Thus this minimum is global:
\[
\deg p_\star\le K,\qquad
\|f_\star-p_\star\|_\infty=E_K\ge0.
\]
% lean: Causalean.Mathlib.Analysis.FinitePolynomialAlternationDuality.exists_bestPolynomial

\item Define the residual
\[
e_\star(z)=f_\star(z)-p_\star(z)
\qquad(z\in\mathcal I_\star).
\]
We establish the existence of points and an orientation satisfying
\[
1\le z_0<z_1<\cdots<z_{K+1}\le H,\qquad
\eta_\star\in\{-1,1\},\qquad
e_\star(z_\iota)=\eta_\star(-1)^\iota E_K
\quad(0\le\iota\le K+1).
\]
If \(E_K=0\), the residual vanishes throughout the interval, and the choices
\[
z_\iota=1+\frac{\iota(H-1)}{K+1}
\quad(0\le\iota\le K+1),
\qquad
\eta_\star=1
\]
have these properties.

Suppose now that \(E_K>0\). Define the compact extremal sets and the residual sign by
\[
\begin{aligned}
\mathcal Z_\star
&=\{z\in\mathcal I_\star:|e_\star(z)|=E_K\},\\
\mathcal Z_\star^+
&=\{z\in\mathcal Z_\star:e_\star(z)=E_K\},\\
\mathcal Z_\star^-
&=\{z\in\mathcal Z_\star:e_\star(z)=-E_K\},\\
\chi_\star(z)
&=
\begin{cases}
1,&e_\star(z)>0,\\
0,&e_\star(z)=0,\\
-1,&e_\star(z)<0
\end{cases}
\qquad(z\in\mathcal I_\star).
\end{aligned}
\]
The set \(\mathcal Z_\star\) is nonempty because the maximum residual magnitude is attained.

Assume, for a contradiction, that the required \(K+2\) alternating extrema do not exist. We first construct a polynomial \(p_{\rm sign}\) such that
\[
\deg p_{\rm sign}\le K,\qquad
e_\star(z)p_{\rm sign}(z)>0
\quad(z\in\mathcal Z_\star).
\]
If \(\mathcal Z_\star^+=\varnothing\), take \(p_{\rm sign}(z)=-1\); if \(\mathcal Z_\star^-=\varnothing\), take \(p_{\rm sign}(z)=1\). These choices give
\[
e_\star(z)p_{\rm sign}(z)=E_K>0
\quad(z\in\mathcal Z_\star).
\]

When both sign sets are nonempty, their disjointness and compactness give a separation \(\delta_{\rm sep}>0\) with
\[
\delta_{\rm sep}<|z-z'|
\quad(z\in\mathcal Z_\star^+,\ z'\in\mathcal Z_\star^-).
\]
Set
\[
\varepsilon_{\rm sep}=\frac{\delta_{\rm sep}}5.
\]
Choose finitely many distinct centers in \(\mathcal Z_\star\), ordered as
\[
c_0^{\rm ctr}<\cdots<c_{n_{\rm ctr}-1}^{\rm ctr},
\qquad n_{\rm ctr}\ge1,
\]
whose \(\varepsilon_{\rm sep}\)-neighborhoods cover \(\mathcal Z_\star\). Thus
\[
\forall z\in\mathcal Z_\star\ \exists\,\iota\in\{0,\ldots,n_{\rm ctr}-1\}:
\quad |z-c_\iota^{\rm ctr}|<\varepsilon_{\rm sep}.
\]
Define the consecutive sign-change indices by
\[
\mathcal D_{\rm ctr}
=\{\iota\in\{0,\ldots,n_{\rm ctr}-2\}:
  \chi_\star(c_\iota^{\rm ctr})
  \ne\chi_\star(c_{\iota+1}^{\rm ctr})\}.
\]
Selecting one center from each successive sign block produces an alternating subsequence of length \(|\mathcal D_{\rm ctr}|+1\). The assumed absence of \(K+2\) alternating extrema therefore implies
\[
|\mathcal D_{\rm ctr}|\le K.
\]
Define
\[
p_{\rm sign}(z)
=\chi_\star(c_{n_{\rm ctr}-1}^{\rm ctr})
 \prod_{\iota\in\mathcal D_{\rm ctr}}
 \left(z-\frac{c_\iota^{\rm ctr}+c_{\iota+1}^{\rm ctr}}2\right),
\]
with an empty product equal to one. Its degree is \(|\mathcal D_{\rm ctr}|\le K\). Each factor changes sign exactly between consecutive blocks, so its sign at every center equals the residual sign there:
\[
\chi_\star(c_\iota^{\rm ctr})p_{\rm sign}(c_\iota^{\rm ctr})>0.
\]
Oppositely signed centers are more than \(5\varepsilon_{\rm sep}\) apart. Each midpoint appearing as a root is therefore more than \(2\varepsilon_{\rm sep}\) from every center. Consequently,
\[
\chi_\star(c_\iota^{\rm ctr})p_{\rm sign}(z)>0
\quad\text{whenever }|z-c_\iota^{\rm ctr}|<\varepsilon_{\rm sep}.
\]
An extremal point and a center within \(\varepsilon_{\rm sep}\) have the same residual sign, by the separation of \(\mathcal Z_\star^+\) and \(\mathcal Z_\star^-\). Hence
\[
e_\star(z)p_{\rm sign}(z)
=E_K\chi_\star(z)p_{\rm sign}(z)>0
\quad(z\in\mathcal Z_\star).
\]
% lean: Causalean.Mathlib.Analysis.FinitePolynomialAlternationDuality.exists_signPolynomial_of_no_alternatingExtrema

We next turn this sign agreement into a strict improvement. Define
\[
\mathcal B_{\rm bad}
=\{z\in\mathcal I_\star:e_\star(z)p_{\rm sign}(z)\le0\},
\qquad
c_{\rm cut}
=
\begin{cases}
\displaystyle
\frac{E_K+\max_{z\in\mathcal B_{\rm bad}}|e_\star(z)|}{2},
&\mathcal B_{\rm bad}\ne\varnothing,\\[6pt]
E_K/2,&\mathcal B_{\rm bad}=\varnothing.
\end{cases}
\]
When nonempty, \(\mathcal B_{\rm bad}\) is compact and contains no extremal point. Its maximum residual magnitude is therefore strictly below \(E_K\). In either case,
\[
c_{\rm cut}<E_K,\qquad
|e_\star(z)|\ge c_{\rm cut}
\ \Longrightarrow\
e_\star(z)p_{\rm sign}(z)>0.
\]
Define
\[
\begin{aligned}
\mathcal U_{\rm high}
&=\{z\in\mathcal I_\star:|e_\star(z)|\ge c_{\rm cut}\},\\
\mu_{\rm sign}
&=\min_{z\in\mathcal U_{\rm high}}
  e_\star(z)p_{\rm sign}(z)>0,\\
M_{\rm sign}
&=\|p_{\rm sign}\|_\infty>0,\\
t_\star
&=\min\left\{
\frac{\mu_{\rm sign}}{M_{\rm sign}^2},
\frac{E_K-c_{\rm cut}}{2M_{\rm sign}}
\right\}>0.
\end{aligned}
\]
The minimum defining \(\mu_{\rm sign}\) exists because \(\mathcal U_{\rm high}\) is compact and contains \(\mathcal Z_\star\). The positivity of \(M_{\rm sign}\) follows from the strict sign agreement at an extremal point.

For \(z\in\mathcal U_{\rm high}\),
\[
\begin{aligned}
\bigl(e_\star(z)-t_\star p_{\rm sign}(z)\bigr)^2
&=e_\star(z)^2
 -2t_\star e_\star(z)p_{\rm sign}(z)
 +t_\star^2p_{\rm sign}(z)^2\\
&\le E_K^2-2t_\star\mu_{\rm sign}
       +t_\star^2M_{\rm sign}^2\\
&\le E_K^2-t_\star\mu_{\rm sign}
<E_K^2.
\end{aligned}
\]
For \(z\in\mathcal I_\star\setminus\mathcal U_{\rm high}\),
\[
\begin{aligned}
|e_\star(z)-t_\star p_{\rm sign}(z)|
&\le |e_\star(z)|+t_\star M_{\rm sign}\\
&<c_{\rm cut}+\frac{E_K-c_{\rm cut}}2
=\frac{E_K+c_{\rm cut}}2<E_K.
\end{aligned}
\]
Continuity and compactness now give
\[
\|f_\star-(p_\star+t_\star p_{\rm sign})\|_\infty<E_K,
\qquad
\deg(p_\star+t_\star p_{\rm sign})\le K,
\]
contradicting the definition of \(E_K\). Thus the required alternating extrema exist.
% lean: Causalean.Mathlib.Analysis.FinitePolynomialAlternationDuality.exists_equioscillationWitness

\item Using these nodes, define the index set, nodal weights, normalization, and signed weights by
\[
\begin{aligned}
\mathcal N_K&=\{0,\ldots,K+1\},\\
b_\iota^{\rm lag}
&=\left(\prod_{\substack{\iota'\in\mathcal N_K\\\iota'\ne\iota}}
(z_\iota-z_{\iota'})\right)^{-1},\\
Z_{\rm lag}&=\sum_{\iota\in\mathcal N_K}|b_\iota^{\rm lag}|,\\
w_\iota&=\frac{b_\iota^{\rm lag}}{Z_{\rm lag}}
\qquad(\iota\in\mathcal N_K).
\end{aligned}
\]
All nodal weights are nonzero, so \(Z_{\rm lag}>0\). There are exactly \(K+1-\iota\) negative factors in the product defining \(b_\iota^{\rm lag}\). It follows that
\[
\sum_{\iota\in\mathcal N_K}|w_\iota|=1,
\qquad
w_\iota=(-1)^{K+1-\iota}|w_\iota|.
\]

For any polynomial \(\mathfrak p\) of degree at most \(K\), Lagrange interpolation gives
\[
\mathfrak p(z)
=\sum_{\iota\in\mathcal N_K}
\mathfrak p(z_\iota)b_\iota^{\rm lag}
\prod_{\substack{\iota'\in\mathcal N_K\\\iota'\ne\iota}}
(z-z_{\iota'}).
\]
Comparing coefficients of \(z^{K+1}\), whose coefficient on the left is zero, and dividing by \(Z_{\rm lag}\), yields
\[
\sum_{\iota\in\mathcal N_K}w_\iota\mathfrak p(z_\iota)=0.
\]
In particular,
\[
\sum_{\iota\in\mathcal N_K}w_\iota z_\iota^v=0
\quad(0\le v\le K),
\qquad
\sum_{\iota\in\mathcal N_K}w_\iota p_\star(z_\iota)=0.
\]
Define the signed reciprocal evaluation by
\[
d_\star=\sum_{\iota\in\mathcal N_K}w_\iota z_\iota^{-1}.
\]
Polynomial cancellation, alternation, and normalization give
\[
\begin{aligned}
d_\star
&=\sum_{\iota\in\mathcal N_K}
  w_\iota\bigl(f_\star(z_\iota)-p_\star(z_\iota)\bigr)\\
&=\sum_{\iota\in\mathcal N_K}
  (-1)^{K+1-\iota}|w_\iota|
  \eta_\star(-1)^\iota E_K\\
&=\eta_\star(-1)^{K+1}E_K
  \sum_{\iota\in\mathcal N_K}|w_\iota|
=\eta_\star(-1)^{K+1}E_K.
\end{aligned}
\]
Since \(E_K\ge0\) and \(|\eta_\star|=1\), this proves
\[
|d_\star|=E_K.
\]
% lean: Causalean.Mathlib.Analysis.FinitePolynomialAlternationDuality.exists_finiteMomentDual

\item Define two nonnegative atomic measures by
\[
\Pi^+
=2\sum_{\iota\in\mathcal N_K}\max(w_\iota,0)\,\delta_{z_\iota},
\qquad
\Pi^-
=2\sum_{\iota\in\mathcal N_K}\max(-w_\iota,0)\,\delta_{z_\iota},
\]
where the unit point mass is specified, for every Borel set \(\mathcal T\subseteq\mathbb R\), by
\[
\delta_z(\mathcal T)
=
\begin{cases}
1,&z\in\mathcal T,\\
0,&z\notin\mathcal T.
\end{cases}
\]
The moment identity at \(v=0\) gives \(\sum_\iota w_\iota=0\). Applying
\[
|w_\iota|=2\max(w_\iota,0)-w_\iota
         =2\max(-w_\iota,0)+w_\iota
\]
and summing, we obtain
\[
1=2\sum_{\iota\in\mathcal N_K}\max(w_\iota,0)
 =2\sum_{\iota\in\mathcal N_K}\max(-w_\iota,0).
\]
Thus \(\Pi^+\) and \(\Pi^-\) are probability measures.

Orient the pair according to the reciprocal evaluation:
\[
(\Pi_0,\Pi_1)
=
\begin{cases}
(\Pi^-,\Pi^+),&d_\star\ge0,\\
(\Pi^+,\Pi^-),&d_\star<0.
\end{cases}
\]
Both oriented measures are therefore probability measures. Define their common finite supporting set by
\[
\mathcal S_\star=\{z_\iota:\iota\in\mathcal N_K\}.
\]
Every atom belongs to \(\mathcal S_\star\subseteq[1,H]\), and hence
\[
\Pi^+(\mathcal S_\star)=\Pi^-(\mathcal S_\star)=1,
\qquad
\Pi_0(\mathcal S_\star)=\Pi_1(\mathcal S_\star)=1.
\]
% lean: CausalSmith.Stat.MarRareqLogfrontier.moment_matched_reciprocal_priors

\item For any real-valued function \(\varphi_{\rm test}:\mathbb R\to\mathbb R\), the finite atomic formulas and
\(\max(w_\iota,0)-\max(-w_\iota,0)=w_\iota\) give
\[
\begin{aligned}
\int\varphi_{\rm test}\,d\Pi^+
-\int\varphi_{\rm test}\,d\Pi^-
&=2\sum_{\iota\in\mathcal N_K}
 \bigl(\max(w_\iota,0)-\max(-w_\iota,0)\bigr)
 \varphi_{\rm test}(z_\iota)\\
&=2\sum_{\iota\in\mathcal N_K}
 w_\iota\varphi_{\rm test}(z_\iota).
\end{aligned}
\]
Taking \(\varphi_{\rm test}(z)=z^v\), for an integer \(0\le v\le K\), yields
\[
\int z^v\,d\Pi^+(z)-\int z^v\,d\Pi^-(z)
=2\sum_{\iota\in\mathcal N_K}w_\iota z_\iota^v=0.
\]
Either orientation therefore satisfies
\[
\int z^v\,d\Pi_0(z)=\int z^v\,d\Pi_1(z)
\qquad(0\le v\le K).
\]
% lean: Causalean.Stat.Minimax.MomentMatchedMixture.FiniteSignedMomentMarkedPoissonMixture.NormalizedFiniteSignedMomentCertificate.orientedPriors_moments_eq

For the reciprocal function, the orientation gives
\[
\int z^{-1}\,d\Pi_1(z)-\int z^{-1}\,d\Pi_0(z)
=
\begin{cases}
2d_\star,&d_\star\ge0,\\
-2d_\star,&d_\star<0
\end{cases}
=2|d_\star|=2E_K.
\]
% lean: Causalean.Stat.Minimax.MomentMatchedMixture.FiniteSignedMomentMarkedPoissonMixture.NormalizedFiniteSignedMomentCertificate.orientedPrior_target_separation

\item It remains to bound \(2E_K\). By \cref{obj:lem:reciprocal-best-approximation},
\[
E_K=\frac{K^2-1}{2K^2}
\left(\frac{K-1}{K+1}\right)^K.
\]
If \(K=2\), this gives
\[
E_2=\frac1{24},
\qquad
2E_2=\frac1{12}.
\]

Suppose \(K\ge3\). First,
\[
\frac{K^2-1}{2K^2}-\frac49
=\frac{K^2-9}{18K^2}\ge0.
\]
To bound the power term, define
\[
x_{\rm rec}=\frac1K\in(0,1/3].
\]
The logarithmic remainder estimate at order zero gives
\[
\frac12\log\frac{1+x_{\rm rec}}{1-x_{\rm rec}}
\le\frac{x_{\rm rec}}{1-x_{\rm rec}^2}.
\]
Since
\[
\frac{1+x_{\rm rec}}{1-x_{\rm rec}}
=\left(\frac{K-1}{K+1}\right)^{-1},
\qquad
8K^2\le9(K^2-1),
\]
we obtain
\[
-K\log\frac{K-1}{K+1}
\le K\frac{2x_{\rm rec}}{1-x_{\rm rec}^2}
=\frac{2K^2}{K^2-1}
\le\frac94.
\]
Exponentiating, using the positivity of \((K-1)/(K+1)\), gives
\[
e^{-9/4}\le\left(\frac{K-1}{K+1}\right)^K.
\]

For the numerical calibration, the standard numerical exponential bound and the rational exponential inequality give
\[
e<2.7182818286<\frac{11}{4},
\qquad
e^{1/4}\le\frac97.
\]
Here the latter estimate is the specialization at \(1/4\) of
\[
e^y\le\frac{2+y}{2-y}\qquad(0\le y<2).
\]
Consequently,
\[
e^{9/4}=e^2e^{1/4}
\le\left(\frac{11}{4}\right)^2\frac97
=\frac{1089}{112}
\le\frac{32}{3}.
\]
Taking reciprocals yields
\[
\frac3{32}\le e^{-9/4}
\le\left(\frac{K-1}{K+1}\right)^K.
\]
% lean: CausalSmith.Stat.MarRareqLogfrontier.reciprocal_ratio_power_lower

Combining the coefficient and power bounds,
\[
2E_K
=2\frac{K^2-1}{2K^2}
 \left(\frac{K-1}{K+1}\right)^K
\ge2\cdot\frac49\cdot\frac3{32}
=\frac1{12}.
\]
Thus, in both cases,
\[
\left|\int z^{-1}\,d\Pi_1(z)-\int z^{-1}\,d\Pi_0(z)\right|
=|2E_K|
\ge2E_K
\ge\frac1{12}.
\]
Together with the finite support and moment identities established above, this proves the result.
% lean: CausalSmith.Stat.MarRareqLogfrontier.reciprocal_best_error_lower
\end{enumerate}
\end{proof}

The finite priors \leanref{sym:\Pi_a}{\ensuremath{\Pi_0,\Pi_1}} supplied by \cref{obj:lem:moment-matched-reciprocal-priors} satisfy the prior requirements of \cref{obj:def:activation-lower-handle}. Matching includes the zeroth moment, and reciprocal separation remains bounded below by an absolute constant. Finite support makes the augmented laws finite mixtures of experiments generated conditional on a fixed latent vector.

\subsection{A one-category family}

For the effective sample-size component, the local family in \cref{obj:lem:one-cell-testing-family} places the entire baseline mass on category zero. Both potential surrogates and the control potential outcome are zero; the treated potential outcome is Bernoulli with mean \(u\in[1/4,3/4]\). Treatment is an independent fair coin, and arrival is an independent Bernoulli variable with probability \(q\). Thus \(P_u\), the full-data law indexed by the treated outcome mean and defined through \cref{obj:lem:one-cell-testing-family}, has average treatment effect \(u\). The following statement records both squared-error and honest-interval comparisons for this family.

\begin{lemmav}{lem:one-cell-testing-family}[One-cell testing family bounds]
There exist a universal constant \(\underline c>0\) and constants \(\underline c_\alpha>0\) for every \(\alpha\in(0,1)\), independent of \(n,d,q\), such that the following holds whenever:
\begin{itemize}
\item \textbf{(Sample and alphabet sizes.)} \(n,d\) are integers with \(n\ge1\) and \(d\ge1\).
\item \textbf{(Arrival probability.)} \(0<q\le1\).
\item \textbf{(Rare-arrival slice.)} The restriction in \cref{obj:ass:rare-arrival-slice} holds.
\end{itemize}
Write \(N=nq\) for the effective sample size. There exists a family of full-data laws \(\{P_u:u\in[1/4,3/4]\}\), chosen independently of \(\alpha\), such that, for every \(u\in[1/4,3/4]\),
\[
P_u\in\mathcal M_{n,d,q},\qquad
P_u\!\left(X=0,\ S(0)=S(1)=0,\ Y(0)=0\right)=1,
\qquad
P_u(R=1)=q,\qquad
\tau(P_u)=u.
\]
Here \(X=0\) denotes the first baseline category in the alphabet \(\{0,\ldots,d-1\}\); the model class \(\mathcal M_{n,d,q}\) and average treatment effect \(\tau(P)\) are defined in \cref{obj:def:model-class,obj:def:ate-functional}.

There exist distinct \(u_0,u_1\in[1/4,3/4]\) such that every possibly randomized estimator \(\mathsf T\in\mathcal T_n\), as defined in \cref{obj:def:minimax-risk}, satisfies
\[
\underline c\min\{1,N^{-1}\}
\le
\max\left\{
\mathbb E_{P_{u_0}}\!\left[(\mathsf T-u_0)^2\right],
\mathbb E_{P_{u_1}}\!\left[(\mathsf T-u_1)^2\right]
\right\}.
\]
Moreover, the minimax squared-error risk in \cref{obj:def:minimax-risk} satisfies
\[
\mathfrak R_{n,d,q}\ge
\underline c\min\{1,N^{-1}\}.
\]

For every \(\alpha\in(0,1)\), there exist an integer \(M\ge1\) and pairwise distinct grid points \(u_0,\ldots,u_M\in[1/4,3/4]\) from this same family such that every uniformly honest connected-interval procedure \(\mathsf I\in\mathcal C_{n,d,q,\alpha}\), as defined in \cref{obj:def:interval-length-risk}, satisfies
\[
\underline c_\alpha\min\{1,N^{-1/2}\}
\le
\max_{i=0,\ldots,M}\mathbb E_{P_{u_i}}[u-l],
\]
where \((l,u)\) are the random endpoints returned by \(\mathsf I\). Moreover, the minimax expected interval length in \cref{obj:def:interval-length-risk} satisfies
\[
\mathfrak L_{n,d,q,\alpha}\ge
\underline c_\alpha\min\{1,N^{-1/2}\}.
\]
All expectations include the observed sample and the procedure's randomization.
\end{lemmav}

\begin{proof}[Proof of \cref{obj:lem:one-cell-testing-family}]
\begin{enumerate}
\item \textbf{Construction of the family.}
Choose
\[
\underline c:=\frac1{256},
\qquad
\underline c_\alpha:=\frac{3(1-\alpha)^2}{512}
\quad(\alpha\in\mathbb R).
\]
These constants have the required positivity. Fix admissible \(n,d,q\), and write \(N=nq>0\). Define
\[
\operatorname{clip}_{[0,1]}(u):=\max\{0,\min\{1,u\}\}
\quad(u\in\mathbb R).
\]
For every \(u\in\mathbb R\), define \(P_u\) by
\[
X=0,\qquad S(0)=S(1)=Y(0)=0,
\]
\[
A\sim\operatorname{Bernoulli}(1/2),\qquad
Y(1)\sim\operatorname{Bernoulli}
       \bigl(\operatorname{clip}_{[0,1]}(u)\bigr),\qquad
R\sim\operatorname{Bernoulli}(q),
\]
with \(A,Y(1),R\) mutually independent. The first baseline category exists because \(d\ge1\). All atom weights are nonnegative, and their sum is
\[
\left(\frac12+\frac12\right)
\left(\operatorname{clip}_{[0,1]}(u)
      +1-\operatorname{clip}_{[0,1]}(u)\right)
\left(q+1-q\right)=1.
\]
Thus \(P_u\) is a full-data probability law for every real \(u\).

For \(u\in[1/4,3/4]\), clipping leaves \(u\) unchanged. Set the realized variables to \(S=S(A)\) and \(Y=Y(A)\), and use independent observed records with their induced marginal law. This gives the sampling and consistency conditions. Independence and the fair treatment coin give
\[
A\perp\!\!\!\perp
  \{X,S(0),S(1),Y(0),Y(1)\},
\qquad P_u(A=1)=\frac12.
\]
Moreover, for \(a,y,r\in\{0,1\}\),
\[
P_u(R=r,Y=y\mid X=0,A=a,S=0)
=
P_u(R=r)\,P_u(Y=y\mid A=a),
\]
because \(R\) is independent of \((A,Y(1))\) and \(Y=AY(1)\). This verifies \cref{obj:ass:arrival-mar} on both occupied cells; all joint probabilities on the other cells vanish. Direct summation also gives
\[
p_{axs}(P_u)=
\begin{cases}
1/2,&x=0,\ s=0,\\
0,&\text{otherwise},
\end{cases}
\qquad
P_u(A=a,X=x,S=s,R=1)=q\,p_{axs}(P_u).
\]
Consequently, \(\rho_{a00}(P_u)=q\), so
\cref{obj:ass:occupied-cell-arrival} holds. The fixed parameters satisfy
\cref{obj:ass:rare-arrival-slice}; hence
\[
P_u\in\mathcal M_{n,d,q}.
\]
Finally, the construction yields
\[
P_u\!\left(X=0,S(0)=S(1)=0,Y(0)=0\right)=1,
\qquad P_u(R=1)=q,
\qquad
\tau(P_u)=E_{P_u}[Y(1)-Y(0)]=u.
\]
The family is chosen independently of \(\alpha\).
% lean: CausalSmith.Stat.MarRareqLogfrontier.oneCellFamilyLaw_model
% lean: CausalSmith.Stat.MarRareqLogfrontier.oneCellFamilyLaw_support
% lean: CausalSmith.Stat.MarRareqLogfrontier.oneCellFamilyLaw_arrival_mass
% lean: CausalSmith.Stat.MarRareqLogfrontier.oneCellFamilyLaw_ate
% lean: CausalSmith.Stat.MarRareqLogfrontier.one_cell_testing_family

\item \textbf{Observed laws and chi-square divergence.}
Define the observed-record space, its one-record laws, and its sample laws by
\[
\mathcal O_d:=\{0,\ldots,d-1\}\times\{0,1\}^4,
\qquad
Q_u:=\mathcal L_{P_u}(X,A,S,R,RY),
\qquad
\mathbb Q_u:=Q_u^{\otimes n}.
\]
For \(u\in[1/4,3/4]\), the potentially nonzero atom probabilities are
\[
\begin{array}{c|c}
o=(X,A,S,R,RY)&Q_u(\{o\})\\ \hline
(0,0,0,0,0)&(1-q)/2\\
(0,1,0,0,0)&(1-q)/2\\
(0,0,0,1,0)&q/2\\
(0,1,0,1,0)&q(1-u)/2\\
(0,1,0,1,1)&qu/2
\end{array}
\]
and all remaining atoms have probability zero. Each full-data atom of positive \(P_u\)-mass has positive \(P_{1/2}\)-mass: the treatment weights are identical, both central outcome weights are positive, and a zero arrival weight vanishes under both laws. Passing to the observed marginal and then to products gives
\[
Q_u\ll Q_{1/2},
\qquad
\mathbb Q_u\ll\mathbb Q_{1/2}.
\]

For these finite laws, define
\[
\chi^2(Q_u\Vert Q_{1/2})
:=
\sum_{\substack{o\in\mathcal O_d\\Q_{1/2}(\{o\})>0}}
\frac{\bigl(Q_u(\{o\})-Q_{1/2}(\{o\})\bigr)^2}
     {Q_{1/2}(\{o\})}.
\]
Expanding the square and using the unit total masses gives
\[
1+\chi^2(Q_u\Vert Q_{1/2})
=
\sum_{\substack{o\in\mathcal O_d\\Q_{1/2}(\{o\})>0}}
\frac{Q_u(\{o\})^2}{Q_{1/2}(\{o\})}.
\]
If \(0<q<1\), summing the five displayed atoms yields
\[
1+\chi^2(Q_u\Vert Q_{1/2})
=
(1-q)+\frac q2+q(1-u)^2+qu^2
=
1+2q\left(u-\frac12\right)^2.
\]
If \(q=1\), the two nonarrival atoms have zero mass and are omitted; the remaining three terms give the same identity. Thus
\[
\chi^2(Q_u\Vert Q_{1/2})
=
2q\left(u-\frac12\right)^2.
\]

To record the product calculation, define
\[
\ell_u(o):=
\begin{cases}
Q_u(\{o\})/Q_{1/2}(\{o\}),&Q_{1/2}(\{o\})>0,\\
0,&Q_{1/2}(\{o\})=0,
\end{cases}
\qquad
L_{u,n}(\boldsymbol o):=\prod_{i=1}^{n}\ell_u(o_i)
\quad(\boldsymbol o\in\mathcal O_d^n).
\]
These are likelihood ratios almost surely under the corresponding reference laws. Independence gives
\[
E_{\mathbb Q_{1/2}}L_{u,n}=1,
\qquad
E_{\mathbb Q_{1/2}}L_{u,n}^2
=
\left(E_{Q_{1/2}}\ell_u^2\right)^n.
\]
Defining the sample chi-square divergence by
\[
\chi^2(\mathbb Q_u\Vert\mathbb Q_{1/2})
:=
E_{\mathbb Q_{1/2}}(L_{u,n}-1)^2,
\]
we obtain
\[
1+\chi^2(\mathbb Q_u\Vert\mathbb Q_{1/2})
=
\left(1+\chi^2(Q_u\Vert Q_{1/2})\right)^n
\le
\exp\!\left(n\chi^2(Q_u\Vert Q_{1/2})\right).
\]
The final inequality follows from \(1+x\le e^x\) for \(x\ge0\). With the total-variation convention of
\cref{obj:lem:randomized-kernel-tv-comparison}, the finite-space formula and Cauchy--Schwarz give
\[
\operatorname{TV}(\mathbb Q_u,\mathbb Q_{1/2})
=
\frac12E_{\mathbb Q_{1/2}}|L_{u,n}-1|
\le
\frac12
\sqrt{\chi^2(\mathbb Q_u\Vert\mathbb Q_{1/2})}.
\]
% lean: CausalSmith.Stat.MarRareqLogfrontier.oneCellFamilyLaw_obs_ac
% lean: CausalSmith.Stat.MarRareqLogfrontier.oneCellFamilyLaw_obs_atom
% lean: CausalSmith.Stat.MarRareqLogfrontier.oneCellFamilyLaw_obs_chi
% lean: Causalean.Stat.one_add_chiSqDiv_pi_iid_general
% lean: Causalean.Stat.tvDist_le_half_sqrt_chiSqDiv

\item \textbf{A two-point testing bound.}
Define
\[
\Delta_{\mathrm{pt}}
:=
\min\left\{\frac14,\frac1{4\sqrt N}\right\},
\qquad
u_0^{\mathrm{pt}}:=\frac12,
\qquad
u_1^{\mathrm{pt}}:=\frac12+\Delta_{\mathrm{pt}}.
\]
Since \(N>0\),
\[
0<\Delta_{\mathrm{pt}}\le\frac14.
\]
The two means are therefore distinct and belong to \([1/4,3/4]\). Also,
\[
\Delta_{\mathrm{pt}}^2\le\frac1{16N},
\qquad
\chi^2(Q_{u_1^{\mathrm{pt}}}\Vert Q_{1/2})
=
2q\Delta_{\mathrm{pt}}^2
\le\frac1{8n}.
\]
The product calculation above yields
\[
\chi^2(\mathbb Q_{u_1^{\mathrm{pt}}}\Vert\mathbb Q_{1/2})
\le e^{1/8}-1<1.
\]
For the numerical bound, use
\[
(e^{1/2})^2=e<3<4,
\qquad
e^{1/8}\le e^{1/2}<2.
\]
Consequently,
\[
\operatorname{TV}
 \bigl(\mathbb Q_{u_1^{\mathrm{pt}}},\mathbb Q_{1/2}\bigr)
\le\frac12.
\]
The testing inequality in
\cref{obj:lem:randomized-kernel-tv-comparison}, applied with the identity kernel, now gives, for every measurable
\(\mathscr A\subseteq\mathcal O_d^n\),
\[
\mathbb Q_{u_1^{\mathrm{pt}}}(\mathscr A^c)
+\mathbb Q_{u_0^{\mathrm{pt}}}(\mathscr A)
\ge
1-\operatorname{TV}
 \bigl(\mathbb Q_{u_1^{\mathrm{pt}}},
       \mathbb Q_{u_0^{\mathrm{pt}}}\bigr)
\ge\frac12.
\]
% lean: CausalSmith.Stat.MarRareqLogfrontier.oneCellFamilyLaw_obs_chi_step
% lean: CausalSmith.Stat.MarRareqLogfrontier.complete_arrival_exp_eighth_lt_two
% lean: CausalSmith.Stat.MarRareqLogfrontier.complete_arrival_product_testing_half
% lean: CausalSmith.Stat.MarRareqLogfrontier.oneCell_product_testing_half

\item \textbf{From testing to squared risk.}
Fix any \(\mathsf T\in\mathcal T_n\), and define its barycenter by
\[
b_{\mathsf T}(\boldsymbol o)
:=
\int_{[-1,1]}h\,\mathsf T(\boldsymbol o,dh)
\quad(\boldsymbol o\in\mathcal O_d^n).
\]
By \cref{obj:lem:randomized-kernel-barycenter}, this is a measurable function taking values in \([-1,1]\), and, for \(j\in\{0,1\}\),
\[
E_{\mathbb Q_{u_j^{\mathrm{pt}}}}
 \left[(b_{\mathsf T}-u_j^{\mathrm{pt}})^2\right]
\le
E_{P_{u_j^{\mathrm{pt}}}}
 \left[(\mathsf T-u_j^{\mathrm{pt}})^2\right].
\]
Define the sample events
\[
\mathscr E_j:=
\left\{\boldsymbol o:
 |b_{\mathsf T}(\boldsymbol o)-u_j^{\mathrm{pt}}|
 \ge\frac{\Delta_{\mathrm{pt}}}{2}\right\}
\quad(j=0,1),
\]
\[
\mathscr A_{\mathrm{pt}}:=
\left\{\boldsymbol o:
 |b_{\mathsf T}(\boldsymbol o)-u_1^{\mathrm{pt}}|
 \le
 |b_{\mathsf T}(\boldsymbol o)-u_0^{\mathrm{pt}}|\right\}.
\]
The triangle inequality gives, at every sample,
\[
\Delta_{\mathrm{pt}}
\le
|b_{\mathsf T}-u_0^{\mathrm{pt}}|
+
|b_{\mathsf T}-u_1^{\mathrm{pt}}|.
\]
Thus
\[
\mathscr A_{\mathrm{pt}}\subseteq\mathscr E_0,
\qquad
\mathscr A_{\mathrm{pt}}^c\subseteq\mathscr E_1,
\qquad
\mathbb Q_{u_0^{\mathrm{pt}}}(\mathscr E_0)
+\mathbb Q_{u_1^{\mathrm{pt}}}(\mathscr E_1)
\ge\frac12.
\]
Since squared error is at least \(\Delta_{\mathrm{pt}}^2/4\) on the corresponding error event,
\[
\begin{aligned}
\max_{j=0,1}
E_{P_{u_j^{\mathrm{pt}}}}
 \left[(\mathsf T-u_j^{\mathrm{pt}})^2\right]
&\ge
\frac{\Delta_{\mathrm{pt}}^2}{4}
\max_{j=0,1}
\mathbb Q_{u_j^{\mathrm{pt}}}(\mathscr E_j)\\
&\ge\frac{\Delta_{\mathrm{pt}}^2}{16}.
\end{aligned}
\]

For the scale of this bound, first suppose
\[
\frac14\le\frac1{4\sqrt N}.
\]
Then \(\Delta_{\mathrm{pt}}=1/4\), and
\[
\Delta_{\mathrm{pt}}^2=\frac1{16}
\ge\frac1{16}\min\{1,N^{-1}\}.
\]
In the complementary case,
\(\Delta_{\mathrm{pt}}=1/(4\sqrt N)\), so
\[
\Delta_{\mathrm{pt}}^2=\frac1{16N}
\ge\frac1{16}\min\{1,N^{-1}\}.
\]
Both cases therefore give
\[
\frac1{256}\min\{1,N^{-1}\}
\le
\max_{j=0,1}
E_{P_{u_j^{\mathrm{pt}}}}
 \left[(\mathsf T-u_j^{\mathrm{pt}})^2\right].
\]

Both testing laws belong to the model. The estimator class is nonempty, as witnessed by the constant-zero estimator. For every full-data law,
\[
\tau(P)=P(Y(1)=1)-P(Y(0)=1)\in[-1,1],
\]
and hence every estimator in \(\mathcal T_n\) has risk between zero and four:
\[
0\le E_P[(\mathsf T-\tau(P))^2]\le4.
\]
Thus the model supremum is a finite real number and dominates the displayed two-point maximum. Taking the infimum over estimators, as in
\cref{obj:def:minimax-risk}, proves
\[
\mathfrak R_{n,d,q}
\ge\frac1{256}\min\{1,N^{-1}\}.
\]
% lean: CausalSmith.Stat.MarRareqLogfrontier.complete_arrival_two_point_decision_reduction
% lean: CausalSmith.Stat.MarRareqLogfrontier.oneCellStep_sq_lower
% lean: CausalSmith.Stat.MarRareqLogfrontier.oneCell_two_point_risk_floor
% lean: CausalSmith.Stat.MarRareqLogfrontier.unrestricted_squaredRisk_bounds
% lean: CausalSmith.Stat.MarRareqLogfrontier.oneCell_minimax_risk_floors

\item \textbf{Construction and scale of the interval grid.}
Fix \(\alpha\in(0,1)\), and define
\[
\beta_\alpha:=1-\alpha,
\qquad
\varepsilon_\alpha:=\frac{\beta_\alpha}{8}.
\]
Choose an integer
\[
M\ge\frac8{\beta_\alpha}.
\]
Then \(M\ge1\), and \(0<\varepsilon_\alpha\le1\). Define
\[
w_{\mathrm{grid}}
:=
\min\left\{\frac14,
           \frac{\varepsilon_\alpha}{4\sqrt N}\right\},
\qquad
\Delta_{\mathrm{grid}}:=\frac{w_{\mathrm{grid}}}{M},
\qquad
u_j^{\mathrm{grid}}
:=
\frac12+j\Delta_{\mathrm{grid}}
\quad(j\in\{0,\ldots,M\}).
\]
These quantities satisfy
\[
0<w_{\mathrm{grid}}\le\frac14,
\qquad
\Delta_{\mathrm{grid}}>0,
\qquad
2Nw_{\mathrm{grid}}^2
\le
2N\left(\frac{\varepsilon_\alpha}{4\sqrt N}\right)^2
=
\frac{\varepsilon_\alpha^2}{8}.
\]
For every grid index,
\[
0\le j\Delta_{\mathrm{grid}}
=\frac{j}{M}w_{\mathrm{grid}}
\le w_{\mathrm{grid}},
\qquad
\frac12\le u_j^{\mathrm{grid}}\le\frac34.
\]
Positive spacing makes the grid points pairwise distinct.

We also need
\[
\frac{\varepsilon_\alpha}{4}
\min\{1,N^{-1/2}\}
\le w_{\mathrm{grid}}.
\]
If \(1/4\le\varepsilon_\alpha/(4\sqrt N)\), then
\[
w_{\mathrm{grid}}=\frac14,
\qquad
\frac{\varepsilon_\alpha}{4}
\min\{1,N^{-1/2}\}
\le\frac{\varepsilon_\alpha}{4}\le\frac14.
\]
In the complementary case,
\[
w_{\mathrm{grid}}
=\frac{\varepsilon_\alpha}{4\sqrt N},
\qquad
\frac{\varepsilon_\alpha}{4}
\min\{1,N^{-1/2}\}
\le\frac{\varepsilon_\alpha}{4\sqrt N}.
\]
This proves the required scale inequality in both cases.
% lean: CausalSmith.Stat.MarRareqLogfrontier.oneCell_interval_minimax_floor

\item \textbf{Coverage and comparison of interval-output laws.}
Let \(\mathsf I\in\mathcal C_{n,d,q,\alpha}\). Define its endpoint-output laws by
\[
\mu_{\mathrm{ref}}
:=
\mathbb Q_{1/2}\mathsf I,
\qquad
\mu_j:=\mathbb Q_{u_j^{\mathrm{grid}}}\mathsf I
\quad(j=0,\ldots,M),
\]
meaning explicitly that, for every measurable endpoint set \(\mathscr A_{\mathrm{end}}\),
\[
\mu_{\mathrm{ref}}(\mathscr A_{\mathrm{end}})
=
\int\mathsf I(\boldsymbol o,\mathscr A_{\mathrm{end}})\,
     \mathbb Q_{1/2}(d\boldsymbol o),
\qquad
\mu_j(\mathscr A_{\mathrm{end}})
=
\int\mathsf I(\boldsymbol o,\mathscr A_{\mathrm{end}})\,
     \mathbb Q_{u_j^{\mathrm{grid}}}(d\boldsymbol o).
\]
Write the returned endpoint pair as
\[
(l,u_{\mathrm{end}}),\qquad -1\le l\le u_{\mathrm{end}}\le1.
\]
Define the coverage events
\[
B_j:=
\{(l,u_{\mathrm{end}}):-1\le l\le u_{\mathrm{end}}\le1,\ 
          l\le u_j^{\mathrm{grid}}\le u_{\mathrm{end}}\}.
\]
All grid laws belong to the model and have the corresponding grid means as their average treatment effects. Honesty therefore gives
\[
\mu_j(B_j)\ge\beta_\alpha
\qquad(j=0,\ldots,M).
\]

For the comparison with the reference sample law, define
\[
\xi_j:=2N(j\Delta_{\mathrm{grid}})^2
\quad(j=0,\ldots,M).
\]
The width budget implies
\[
0\le\xi_j
\le2Nw_{\mathrm{grid}}^2
\le\frac{\varepsilon_\alpha^2}{8}
\le\frac18<1.
\]
The one-record identity and product calculation give
\[
\chi^2\!\left(
 \mathbb Q_{u_j^{\mathrm{grid}}}\Vert\mathbb Q_{1/2}
 \right)
\le e^{\xi_j}-1.
\]
For \(0\le\xi_j<1\), use the exponential bound
\(e^{\xi_j}\le(1-\xi_j)^{-1}\). Furthermore,
\[
\xi_j(1+\varepsilon_\alpha^2)
\le2\xi_j
\le\frac{\varepsilon_\alpha^2}{4}
\le\varepsilon_\alpha^2.
\]
Since \(1-\xi_j>0\), this implies
\[
e^{\xi_j}-1
\le\frac1{1-\xi_j}-1
=\frac{\xi_j}{1-\xi_j}
\le\varepsilon_\alpha^2.
\]
Consequently,
\[
\operatorname{TV}
 \bigl(\mathbb Q_{1/2},
       \mathbb Q_{u_j^{\mathrm{grid}}}\bigr)
\le
\frac12
\sqrt{\chi^2\!\left(
 \mathbb Q_{u_j^{\mathrm{grid}}}\Vert\mathbb Q_{1/2}
 \right)}
\le\frac{\varepsilon_\alpha}{2}
\le\varepsilon_\alpha.
\]
Applying \cref{obj:lem:randomized-kernel-tv-comparison} to the common interval kernel yields
\[
\operatorname{TV}(\mu_{\mathrm{ref}},\mu_j)
\le\varepsilon_\alpha.
\]
Thus, by the defining event bound for total variation,
\[
\mu_{\mathrm{ref}}(B_j)
\ge\mu_j(B_j)-\varepsilon_\alpha
\ge\beta_\alpha-\varepsilon_\alpha
=\frac{7\beta_\alpha}{8}.
\]
% lean: CausalSmith.Stat.MarRareqLogfrontier.oneCell_product_tv_le
% lean: Causalean.Stat.tvDist_bind_le
% lean: CausalSmith.Stat.MarRareqLogfrontier.oneCell_interval_minimax_floor

\item \textbf{Counting covered grid points.}
For every ordered endpoint pair, define
\[
C_{\mathrm{grid}}(l,u_{\mathrm{end}})
:=
\sum_{j=0}^{M}
\mathbf1\{l\le u_j^{\mathrm{grid}}\le u_{\mathrm{end}}\}
\quad(-1\le l\le u_{\mathrm{end}}\le1).
\]
We claim the pointwise inequality
\[
\Delta_{\mathrm{grid}}
\bigl(C_{\mathrm{grid}}(l,u_{\mathrm{end}})-1\bigr)
\le u_{\mathrm{end}}-l.
\]
If the covered index set is empty, its left-hand side equals
\(-\Delta_{\mathrm{grid}}\le0\le u_{\mathrm{end}}-l\). Otherwise define its extreme indices by
\[
j_-:=\min\{j\in\{0,\ldots,M\}:
                 l\le u_j^{\mathrm{grid}}\le u_{\mathrm{end}}\},
\qquad
j_+:=\max\{j\in\{0,\ldots,M\}:
                 l\le u_j^{\mathrm{grid}}\le u_{\mathrm{end}}\}.
\]
Every covered index lies between \(j_-\) and \(j_+\), so
\[
C_{\mathrm{grid}}(l,u_{\mathrm{end}})\le j_+-j_-+1.
\]
The two extreme grid points lie in the interval, and hence
\[
\Delta_{\mathrm{grid}}
 \bigl(C_{\mathrm{grid}}(l,u_{\mathrm{end}})-1\bigr)
\le
(j_+-j_-)\Delta_{\mathrm{grid}}
=
u_{j_+}^{\mathrm{grid}}-u_{j_-}^{\mathrm{grid}}
\le u_{\mathrm{end}}-l.
\]
This proves the claim, including the empty and singleton cases.

Summing coverage probabilities under the reference output law gives
\[
\int C_{\mathrm{grid}}(l,u_{\mathrm{end}})\,
     \mu_{\mathrm{ref}}(d(l,u_{\mathrm{end}}))
=
\sum_{j=0}^{M}\mu_{\mathrm{ref}}(B_j)
\ge(M+1)(\beta_\alpha-\varepsilon_\alpha).
\]
Integrating the pointwise length bound therefore yields
\[
\int(u_{\mathrm{end}}-l)\,\mu_{\mathrm{ref}}(d(l,u_{\mathrm{end}}))
\ge
\Delta_{\mathrm{grid}}
\left((M+1)\frac{7\beta_\alpha}{8}-1\right).
\]
The choice of \(M\) gives \(M\beta_\alpha\ge8\), and
\[
(M+1)\frac{7\beta_\alpha}{8}-1
-\frac{3\beta_\alpha M}{4}
=
\frac{M\beta_\alpha}{8}
+\frac{7\beta_\alpha}{8}-1
\ge0.
\]
Since \(M\Delta_{\mathrm{grid}}=w_{\mathrm{grid}}\), we conclude
\[
\int(u_{\mathrm{end}}-l)\,\mu_{\mathrm{ref}}(d(l,u_{\mathrm{end}}))
\ge\frac{3\beta_\alpha}{4}\,w_{\mathrm{grid}}.
\]
The endpoint length is measurable and bounded by two. Integrating first over the sample and then over the kernel therefore equals integration under its output law, giving
\[
E_{P_{1/2}}[u_{\mathrm{end}}-l]
=
\int(u_{\mathrm{end}}-l)\,\mu_{\mathrm{ref}}(d(l,u_{\mathrm{end}}))
\ge\frac{3\beta_\alpha}{4}\,w_{\mathrm{grid}}.
\]
All randomization is included in this expectation.
% lean: Causalean.Stat.intervalLength_ge_gridCount
% lean: Causalean.Stat.integral_gridCount
% lean: Causalean.Stat.honest_interval_length_lower_of_grid_tv_positive
% lean: CausalSmith.Stat.MarRareqLogfrontier.intervalRisk_bind_lintegral
% lean: CausalSmith.Stat.MarRareqLogfrontier.oneCell_interval_minimax_floor

\item \textbf{Minimax length and constant calibration.}
The honest interval class is nonempty. Indeed, for every full-data law,
\[
\tau(P)=P(Y(1)=1)-P(Y(0)=1)\in[-1,1],
\]
so the constant procedure returning \([-1,1]\) has coverage one. Every admissible interval procedure also satisfies, for every full-data law,
\[
0\le E_P[u_{\mathrm{end}}-l]\le2,
\]
because \(0\le u_{\mathrm{end}}-l\le2\) for every returned endpoint pair. Thus the real-valued model suprema are finite. Since \(P_{1/2}\) belongs to the model, the preceding bound implies
\[
\sup_{P\in\mathcal M_{n,d,q}}E_P[u_{\mathrm{end}}-l]
\ge E_{P_{1/2}}[u_{\mathrm{end}}-l]
\ge\frac{3\beta_\alpha}{4}\,w_{\mathrm{grid}}.
\]
Taking the infimum over honest procedures in
\cref{obj:def:interval-length-risk} gives
\[
\mathfrak L_{n,d,q,\alpha}
\ge\frac{3\beta_\alpha}{4}\,w_{\mathrm{grid}}.
\]

The width inequality and
\(\varepsilon_\alpha=\beta_\alpha/8\) give the required calibration:
\[
\begin{aligned}
\underline c_\alpha\min\{1,N^{-1/2}\}
&=
\frac{3\beta_\alpha^2}{512}
\min\{1,N^{-1/2}\}\\
&\le
\frac{3\beta_\alpha\varepsilon_\alpha}{16}
\min\{1,N^{-1/2}\}\\
&\le
\frac{3\beta_\alpha}{4}\,w_{\mathrm{grid}}.
\end{aligned}
\]
Finally, \(u_0^{\mathrm{grid}}=1/2\). Hence every honest procedure satisfies
\[
\underline c_\alpha\min\{1,N^{-1/2}\}
\le
E_{P_{u_0^{\mathrm{grid}}}}[u_{\mathrm{end}}-l]
\le
\max_{j=0,\ldots,M}E_{P_{u_j^{\mathrm{grid}}}}[u_{\mathrm{end}}-l],
\]
and the minimax length satisfies the same lower bound. These are grid points from the family fixed in the first step, completing all assertions.
% lean: CausalSmith.Stat.MarRareqLogfrontier.ate_mem_unit_interval
% lean: CausalSmith.Stat.MarRareqLogfrontier.intervalRisk_bounds
% lean: CausalSmith.Stat.MarRareqLogfrontier.oneCell_interval_minimax_floor
% lean: CausalSmith.Stat.MarRareqLogfrontier.one_cell_testing_family
\end{enumerate}
\end{proof}

The squared-error comparison already applies when all individuals share a baseline category and both potential surrogates are constant. It isolates the information supplied by arrived treated outcomes, whose frequency is governed by balanced assignment and the arrival probability. The same family supports the interval comparison: the family is common across coverage levels, while the finite grid and interval lower constant may depend on the chosen coverage level. Both conclusions retain the randomized decision classes of \cref{obj:def:minimax-risk,obj:def:interval-length-risk}.

\subsection{The many-category comparison in the full experiment}

The many-category component uses the augmented laws of \cref{obj:def:activation-lower-handle}, with the finite priors supplied by \cref{obj:lem:moment-matched-reciprocal-priors}. Their membership distribution, treatment allocation, and activation distribution are common to the two hypotheses. The comparison therefore concerns the full fixed-sample experiment, including the information in repeated memberships and in both randomized arms.

Two mathematical inputs have distinct roles in this comparison. Moment matching governs comparison of the integrated observation laws, while reciprocal-mean separation governs separation of the treatment-effect functional across the prior constructions. The common activation mechanism connects these roles while preserving the original record coordinates. Together with the one-category component in \cref{obj:lem:one-cell-testing-family}, these inputs support the following lower bound for the entire decision class.

\begin{theoremv}{thm:full-experiment-lower}[Full-experiment minimax lower bound]
There exists a universal lower constant \(\underline c>0\) such that, for every \(n,d\in\mathbb N\) and \(q\in\mathbb R\) satisfying
\begin{itemize}
\item \textbf{(Sample size and dimension.)} \(n\ge1\) and \(d\ge1\).
\item \textbf{(Arrival parameter.)} \(0<q\le1\).
\item \textbf{(Rare-arrival slice.)} The restriction in \cref{obj:ass:rare-arrival-slice} holds.
\end{itemize}
the minimax squared-error risk \(\mathfrak R_{n,d,q}\) defined in \cref{obj:def:minimax-risk} obeys
\[
\mathfrak R_{n,d,q}\ge \underline c\,r_{n,d,q},
\]
where \(r_{n,d,q}\) is the frontier rate defined in \cref{obj:def:frontier-rate}.
\end{theoremv}

The bound combines effective sample-size uncertainty with uncertainty across many categories. The one-category family accounts for the reciprocal effective sample-size term, and the activated comparison accounts for the squared category term in \cref{obj:def:frontier-rate}. Full membership observations remain available throughout the experiment. Consequently, the lower bound describes the difficulty of estimating the treatment contrast with all recorded baseline, assignment, surrogate, and arrival information under the stated rare-arrival conditions.

\subsection{Large alphabets and the ascertainment deficit}

The separate unrestricted comparison is \cref{obj:thm:unrestricted-large-alphabet-deficit-lower}. Its domain is \(n\ge1\), \(d\ge1024n^2\), and \(0<q\le1\), with the model and risk defined in \cref{obj:def:unrestricted-arrival-model-class,obj:def:unrestricted-minimax-risk}. The observed-law comparison concerns all \(n\) original records, including their category memberships. Target separation supplies the ascertainment-deficit component, quantified by the lower bound \((1-q)^2/1024\); the sampling component is bounded below by \(1/(256n)\).

These components have complementary interpretations. The sampling term remains present at complete ascertainment. The squared deficit measures the additional difficulty captured by the large-alphabet comparison when the arrival floor is below one. The matching upper bound \(4/n+(1-q)^2\) in \cref{obj:thm:unrestricted-large-alphabet-deficit-lower} comes from the observed-outcome envelope of \cref{obj:thm:unrestricted-near-complete-arrival-envelope}. Thus the large-alphabet result supplies a comparison with universal constants for the joint sampling-and-deficit scale.

\subsection{Squared-error comparisons and their consequences}

The lower bound in \cref{obj:thm:full-experiment-lower} and the constructive upper comparison in \cref{obj:thm:unrestricted-mixed-count-upper} provide the two sides of \cref{obj:thm:matched-minimax-frontier}. Under \cref{obj:ass:rare-arrival-slice}, the minimax squared-error risk is therefore comparable, with universal constants, to the rate in \cref{obj:def:frontier-rate}. Conditional averaging preserves the upper guarantee by \cref{obj:lem:randomized-kernel-barycenter}, while the lower comparison covers the full class of possibly randomized estimators.

For a fixed alphabet and growing effective sample size, \cref{obj:prop:fixed-d-effective-sample-reduction} records risk of order \(N^{-1}\), including within the constant-surrogate submodel. Along growing-alphabet sequences satisfying the rare-arrival restriction and \(N\to\infty\), \cref{obj:prop:phase-boundaries} gives consistency exactly when \(d=o\{N\log(e+N)\}\), and risk of order \(N^{-1}\) exactly when \(d=O\{\sqrt N\log(e+N)\}\). These consequences distinguish the category-growth range permitting consistency from the range preserving effective sample-size precision.

For the unrestricted model, \cref{obj:prop:unrestricted-uniform-near-complete-elbow} combines the large-alphabet lower comparison with the dimension-free upper envelope. An eventual risk bound of order \(n^{-1}\), simultaneous over every alphabet size, is equivalent to an ascertainment deficit \(1-q_n=O(n^{-1/2})\). The numerical comparison along arbitrary alphabet-size sequences is given by \cref{obj:prop:unrestricted-near-complete-phase-boundary}. At complete arrival, \cref{obj:thm:unrestricted-complete-arrival-frontier,obj:prop:unrestricted-complete-arrival-phase-boundary} establish the same sampling order uniformly over category growth.

Finally, the fixed-floor comparison in \cref{obj:thm:unrestricted-fixed-q-minimax-frontier} concerns each fixed \(q\in(0,1/2)\), with constants depending on that floor. Its lower comparison invokes the discrete observational result of \citet[Theorem~2, Section~3.2, p.~9, and its proof in Section~C.5, pp.~39--44]{ZengBalakrishnanHanKennedy2026}, including the submodel with zero control regression. The resulting category-growth consequences are stated in \cref{obj:prop:unrestricted-fixed-q-phase-boundaries}: consistency corresponds to \(d=o\{n\log(e+n)\}\), and parametric squared-error order corresponds to \(d=O\{\sqrt n\log(e+n)\}\). These comparisons preserve the distinct parameter domains and constant dependencies of the rare-arrival, fixed-floor, and near-complete ascertainment results.


\section{Interval and observational proofs}

The interval and observational comparisons use two features of the constructive analysis: a deterministic bound on squared-error risk and a parameter-independent transformation preserving the target. The decision-theoretic results in \cref{obj:lem:randomized-kernel-barycenter,obj:lem:randomized-kernel-tv-comparison} supply the corresponding comparisons for procedures using auxiliary randomization.

\subsection{Coverage and expected interval length}

Fix \(\alpha\in(0,1)\), and suppose that \(n,d\ge1\), \(0<q\le1\), and \cref{obj:ass:rare-arrival-slice} holds. The deterministic envelope \leanref{sym:\mathfrak U_{n,d,q}}{\ensuremath{\mathfrak U_{n,d,q}}} in \cref{obj:thm:unrestricted-mixed-count-upper} controls squared-error risk uniformly over the model in \cref{obj:def:model-class}. Its role in \cref{obj:def:risk-envelope-interval} is to determine the radius before observing the sample. Markov's inequality converts this risk control into coverage at least \(1-\alpha\). Because the target belongs to \([-1,1]\), intersecting the interval with that range preserves coverage.

The mixed-count estimator \leanref{sym:\widehat\tau^{\mathrm{mix}}_{n,d,q}}{\ensuremath{\widehat\tau^{\mathrm{mix}}_{n,d,q}}} is defined in \cref{obj:def:mixed-count-estimator}. The resulting interval \leanref{sym:I^{\mathrm{mix}}_\alpha}{\ensuremath{I^{\mathrm{mix}}_\alpha}} has length bounded by
\[
\operatorname{length}(I^{\mathrm{mix}}_\alpha)
\le
\min\left\{2,\,
2\sqrt{\frac{\mathfrak U_{n,d,q}}{\alpha}}\right\}.
\]
Thus the envelope comparison with the rate \leanref{sym:r_{n,d,q}}{\ensuremath{r_{n,d,q}}} in \cref{obj:def:frontier-rate} supplies the upper expected-length comparison in \cref{obj:thm:connected-interval-frontier}. The upper constant \leanref{sym:\overline C_\alpha}{\ensuremath{\overline C_\alpha}} may depend on the fixed miscoverage level and is uniform over the sample sizes, dimensions, and arrival floors satisfying these conditions.

The lower comparison concerns every endpoint procedure admitted by \cref{obj:def:interval-length-risk}. Its sampling component is supplied directly by \cref{obj:lem:one-cell-testing-family}. For each fixed \(\alpha\), that result supplies a finite grid of distinct target values within the same one-category family and bounds maximal expected length below at order
\[
\min\{1,(nq)^{-1/2}\}.
\]
The underlying family is common across coverage levels; the grid and its comparison constant may depend on \(\alpha\). This finite-family formulation supports every coverage level \(1-\alpha\in(0,1)\).

The many-category component of \cref{obj:thm:connected-interval-frontier} uses the separation and observation-law comparison associated with \cref{obj:def:activation-lower-handle,obj:lem:moment-matched-reciprocal-priors}. Moment matching controls statistical distinguishability, while reciprocal separation determines the target scale. Connectedness supplies the geometric link to expected length: an interval containing separated target values spans their separation. Finite testing families at the specified coverage level retain this link throughout the coverage range in the theorem.

Auxiliary endpoint randomization enters through the joint observation--action laws defined in \cref{obj:lem:randomized-kernel-tv-comparison}. Apply that result with the action space equal to the ordered endpoint pairs and the common decision rule equal to the interval procedure. The joint-law total variation equals the total variation of the observation laws. Consequently, coverage events can be compared while retaining both the observed sample and the reported endpoints. Expected length is evaluated under this same joint law, exactly as required by \cref{obj:def:interval-length-risk}.

Together, these components give the lower constant \leanref{sym:\underline c_\alpha}{\ensuremath{\underline c_\alpha}} in \cref{obj:thm:connected-interval-frontier}. The minimax expected length \leanref{sym:\mathfrak L_{n,d,q,\alpha}}{\ensuremath{\mathfrak L_{n,d,q,\alpha}}} is therefore comparable to \(\sqrt{r_{n,d,q}}\), with constants depending on \(\alpha\). The square-root scale reflects both sources of uncertainty: effective sample size and the number of covariate categories. Along sequences satisfying the rare-arrival restriction, at fixed \(\alpha\), optimal expected length tends to zero exactly when \(r_{n,d,q}\) tends to zero.

\subsection{Synthetic assignment and observational risk}

Fix a strict positivity margin \(0<\epsilon<1/2\). The observational class \leanref{sym:\mathcal D_{d,\epsilon}}{\ensuremath{\mathcal D_{d,\epsilon}}} is defined in \cref{obj:def:zeng-discrete-ate-class}, and its target is the contrast in \cref{obj:def:zeng-ate-functional}. The transformation in \cref{obj:lem:synthetic-randomization-kernel}, with arrival floor \(q=\epsilon\), assigns an independent fair arm to each observational record and retains the outcome when that arm agrees with the observational treatment. Positivity supplies the arrival floor in both arms.

The induced full-data law \(P_Z^*\), defined by \cref{obj:lem:synthetic-randomization-kernel}, belongs to the unrestricted randomized class and preserves the target:
\[
\tau(P_Z^*)=\theta_Z(P_Z).
\]
The same statement identifies the synthetic sample law with the product observed-record law of this completion. These two properties make the unrestricted risk certificate in \cref{obj:thm:unrestricted-mixed-count-upper} applicable to the mixed-count estimator evaluated on the synthetic records.

The observational estimator \leanref{sym:\widehat\theta^{\mathrm{poly}}_{n,d,\epsilon}}{\ensuremath{\widehat\theta^{\mathrm{poly}}_{n,d,\epsilon}}} in \cref{obj:thm:zeng-fixed-positivity-polynomial-frontier} averages over the independent synthetic assignments conditional on the observational sample. The barycenter comparison in \cref{obj:lem:randomized-kernel-barycenter} weakly reduces squared-error risk under this averaging. Accordingly, the constructive bound has a universal upper constant and retains the positivity dependence explicitly through
\[
r_{n,d,\epsilon}
=
\min\left\{
1,\frac{1}{n\epsilon}
+
\left[
\frac{d}{n\epsilon\log(e+n\epsilon)}
\right]^2
\right\}.
\]
Synthetic assignment converts observational treatment imbalance into selective outcome arrival; conditional averaging then removes the variation introduced by that assignment. Target preservation connects both operations to the same observational contrast.

The matching lower comparison invokes \citet[Theorem~2, Section~3.2, p.~9, together with its proof in Section~C.5, pp.~39--44]{ZengBalakrishnanHanKennedy2026}. The relevant published conclusion concerns fixed strict positivity, sufficiently large \(n\), and \(1\le d\le c_0n\log n\), where \(c_0>0\) is a constant in that conclusion. Its proof uses the submodel with zero control regression on occupied categories and supplies the lower order
\[
\frac1n+\left(\frac{d}{n\log n}\right)^2.
\]
This conclusion supplies the observational lower-bound input. The matched comparison in \cref{obj:thm:zeng-fixed-positivity-polynomial-frontier} covers every \(n,d\ge1\), incorporates bounded-risk truncation, and uses \(\log(e+n)\).

For each fixed \(\epsilon\), the observational minimax risk \leanref{sym:\mathfrak R^Z_{n,d,\epsilon}}{\ensuremath{\mathfrak R^Z_{n,d,\epsilon}}} consequently has the two-sided comparison stated in \cref{obj:thm:zeng-fixed-positivity-polynomial-frontier}, with positive finite constants \(c_\epsilon\) and \(C_\epsilon\) depending on that margin. This dependence is distinct from the universal constant in the constructive bound expressed through \(r_{n,d,\epsilon}\).

\subsection{Fixed arrival and category-growth consequences}

The same transformation determines the direction of the fixed-arrival experiment comparison. Fix \(q\in(0,1/2)\). For every observational law in \(\mathcal D_{d,q}\), \cref{obj:lem:synthetic-randomization-kernel} supplies a target-preserving completion in the unrestricted class \leanref{sym:\mathcal M^{+}_{n,d,q}}{\ensuremath{\mathcal M^{+}_{n,d,q}}} of \cref{obj:def:unrestricted-arrival-model-class}. Composing any randomized estimation procedure with the synthetic assignment gives an observational procedure with the same risk at the corresponding completion. Hence the unrestricted randomized minimax risk \leanref{sym:\mathfrak R^{+}_{n,d,q}}{\ensuremath{\mathfrak R^{+}_{n,d,q}}} dominates the observational minimax risk at positivity margin \(q\). The published observational lower comparison therefore transfers to this randomized class.

Combined with the constructive upper comparison, this transfer supports \cref{obj:thm:unrestricted-fixed-q-minimax-frontier}. Its constants \(\underline c(q)\) and \(\overline C(q)\) depend only on the fixed arrival floor and apply uniformly over \(n,d\ge1\). The resulting category-growth conclusions are recorded in \cref{obj:prop:unrestricted-fixed-q-phase-boundaries}: consistency is equivalent to
\[
d_n=o\{n\log(e+n)\},
\]
and parametric squared-error order is equivalent to
\[
d_n=O\{\sqrt n\log(e+n)\}.
\]
For an eventual bound \(d_n\le M\sqrt n\log(e+n)\), the lower risk constant may depend on \(q\), while the upper constant may depend on both \(q\) and \(M\). The eventual sample-size threshold may depend on the category sequence. These dependencies specify the fixed-arrival scope of the matched comparison and its sequence-level consequences.


\section{Verification note}

The theorem and lemma statements and proofs attached to the frozen results are machine-checked in Lean 4 under their stated hypotheses and cited dependencies. The checked conclusions cover the identification identities in \cref{obj:lem:unrestricted-cell-identification,obj:lem:cell-identification-background}, the decision-rule comparisons in \cref{obj:lem:randomized-kernel-barycenter,obj:lem:randomized-kernel-tv-comparison}, and the synthetic-assignment construction in \cref{obj:lem:synthetic-randomization-kernel}. They also cover the estimator risk certificates in \cref{obj:thm:unrestricted-uniform-mixed-count-certificate,obj:thm:uniform-mixed-count-certificate,obj:thm:mixed-count-upper}, the approximation and testing conclusions in \cref{obj:lem:reciprocal-best-approximation,obj:lem:moment-matched-reciprocal-priors,obj:lem:one-cell-testing-family}, and the lower bound in \cref{obj:thm:full-experiment-lower}. The point-risk comparisons and their stated phase-boundary consequences, the connected-interval comparison in \cref{obj:thm:connected-interval-frontier}, and the observational comparison in \cref{obj:thm:zeng-fixed-positivity-polynomial-frontier} are checked with the conditions and constant dependencies recorded in their statements.

The model restrictions enter as assumptions. These include independent sampling, randomized assignment and balance, surrogate and outcome consistency, conditional independence of outcome arrival, and the occupied-cell arrival floor in \cref{obj:ass:iid-sampling,obj:ass:randomized-independence,obj:ass:balanced-randomization,obj:ass:surrogate-consistency,obj:ass:outcome-consistency,obj:ass:arrival-mar,obj:ass:occupied-cell-arrival}. Results for rare arrival additionally use \cref{obj:ass:rare-arrival-slice}; the observational comparison uses the positivity class specified in \cref{obj:def:zeng-discrete-ate-class}. Presentation definitions specify the laws, targets, decision classes, and procedures to which the mathematical assertions apply.

Published mathematical conclusions used through theorem-local cited dependencies remain part of the assumptions of the corresponding checks. The generated local verification disclosures identify those conclusions, preserve their source locators, and distinguish them from proofs formalized in Lean. Together, the stated hypotheses and those disclosures specify the scope of the machine-checked conditional conclusions.


\section{Proofs of the main results}\label{sec:deferred-proofs}

\begin{proof}[Proof of \cref{obj:thm:unrestricted-near-complete-arrival-envelope}]
\begin{enumerate}
\item Choose the universal constant \(\overline C>0\) supplied by
\cref{obj:thm:unrestricted-mixed-count-upper}. For every admissible law \(P\) and the specified sampling law, that result gives
\[
E_P\!\left[
\left(\widehat\tau^{\mathrm{mix}}_{n,d,q}(\mathbf O_n)-\tau(P)\right)^2
\right]
\le \mathfrak U_{n,d,q}
\le \overline C\,r_{n,d,q},
\]
and
\[
E_P\!\left[
\left(\widehat\tau^{\mathrm{mix}}_{n,d,q}(\mathbf O_n)-\tau(P)\right)^2
\right]
\le
E_P\!\left[
E_\omega\!\left[
\left(\widetilde\tau(\mathbf O_n;\omega)-\tau(P)\right)^2
\,\middle|\,\mathbf O_n
\right]\right].
\]
Its envelope is exactly the branchwise envelope in the statement, and its auxiliary output includes the zero output on the Poisson tail. We retain this same constant throughout.
% lean: CausalSmith.Stat.MarRareqLogfrontier.unrestricted_mixed_count_upper

\item We first record how a deterministic risk bound yields a minimax bound. Write
\[
\mathcal O_d:=\{0,\ldots,d-1\}\times\{0,1\}^{4}
\]
for the observed-record space. For a bounded estimator kernel \(\mathsf T\), define its risk by
\[
\mathcal R_P(\mathsf T)
:=
\int_{\mathcal O_d^n}\int_{[-1,1]}
(h-\tau(P))^2\,\mathsf T(\mathbf o,dh)\,
\mathcal L_P(\mathbf O_n)(d\mathbf o).
\]
These risks are nonnegative.

Let \(f_{\mathrm{det}}:\mathcal O_d^n\to[-1,1]\) be measurable, and let \(b_{\mathrm{up}}\ge0\). Suppose
\[
E_P\!\left[
\left(f_{\mathrm{det}}(\mathbf O_n)-\tau(P)\right)^2
\right]\le b_{\mathrm{up}}
\qquad
\text{for every }P\in\mathcal M^{+}_{n,d,q}.
\]
The parameter-independent kernel
\[
\mathsf T_{\mathrm{det}}(\mathbf o,\cdot)
:=\operatorname{Dirac}\!\left(f_{\mathrm{det}}(\mathbf o)\right)
\]
belongs to the decision class in \cref{obj:def:unrestricted-minimax-risk}, and satisfies
\[
\mathcal R_P(\mathsf T_{\mathrm{det}})
=
E_P\!\left[
\left(f_{\mathrm{det}}(\mathbf O_n)-\tau(P)\right)^2
\right].
\]
Consequently,
\[
\mathfrak R^{+}_{n,d,q}
\le
\sup_{P\in\mathcal M^{+}_{n,d,q}}
E_P\!\left[
\left(f_{\mathrm{det}}(\mathbf O_n)-\tau(P)\right)^2
\right]
\le b_{\mathrm{up}}.
\]
For a nonempty model class, the last inequality follows by taking the supremum of the pointwise bounds. For an empty class, the real supremum is defined to be zero, and the inequality follows from \(b_{\mathrm{up}}\ge0\).
% lean: CausalSmith.Stat.MarRareqLogfrontier.unrestricted_deterministic_upper

\item Fix \(n,d\ge1\) and \(0<q\le1\). Binary potential outcomes and
\cref{obj:def:ate-functional} give
\[
\tau(P)=P(Y(1)=1)-P(Y(0)=1)\in[-1,1].
\]
Indeed, each probability belongs to \([0,1]\). The zero map is measurable and takes values in \([-1,1]\), and its risk under the canonical product probability law is
\[
E_P\!\left[(0-\tau(P))^2\right]=\tau(P)^2\le1.
\]
Applying the deterministic comparison just established yields
\[
\mathfrak R^{+}_{n,d,q}\le1.
\]
% lean: CausalSmith.Stat.MarRareqLogfrontier.ate_mem_unit_interval

\item The rate in \cref{obj:def:frontier-rate} satisfies
\[
r_{n,d,q}
=
\min\left\{
1,\,
\frac1{nq}
+
\left[\frac{d}{nq\log(e+nq)}\right]^2
\right\}\ge0,
\qquad
\overline C\,r_{n,d,q}\ge0.
\]
The estimator in \cref{obj:def:mixed-count-estimator} is measurable. Its auxiliary output lies in \([-1,1]\) for every sample and every randomization, so averaging gives
\[
-1
\le
E_\omega\!\left[
\widetilde\tau(\mathbf o;\omega)\mid\mathbf o
\right]
=
\widehat\tau^{\mathrm{mix}}_{n,d,q}(\mathbf o)
\le1.
\]
Under the canonical product law, the coordinate observations are measurable and independent, with
\[
\mathcal L_P(O_i)=\mathcal L_P(O_1),
\qquad i=1,\ldots,n,
\]
so they satisfy \cref{obj:ass:iid-sampling}. The first step therefore applies to every law in the model class. The deterministic comparison, with \(b_{\mathrm{up}}=\overline C\,r_{n,d,q}\), gives
\[
\mathfrak R^{+}_{n,d,q}\le\overline C\,r_{n,d,q}.
\]
% lean: CausalSmith.Stat.MarRareqLogfrontier.unrestricted_near_complete_arrival_envelope

\item We next establish the observed-contrast bound. Fix
\(P\in\mathcal M^{+}_{n,d,q}\). For an observed record, write its coordinates as
\[
o=(X(o),A(o),S(o),R(o),(RY)(o)),
\]
and define
\[
\xi_{\mathrm{obs}}(o)
:=
2(2A(o)-1)\mathbf1\{R(o)=1,\ (RY)(o)=1\}.
\]
For a sample \(\mathbf o=(o_1,\ldots,o_n)\in\mathcal O_d^n\), define
\[
\overline\xi_n(\mathbf o)
:=
\frac1n\sum_{i=1}^n\xi_{\mathrm{obs}}(o_i).
\]
The canonical product law gives
\[
E_P[\overline\xi_n(\mathbf O_n)]
=E_P[\xi_{\mathrm{obs}}(O_1)],
\qquad
\operatorname{Var}_P(\overline\xi_n(\mathbf O_n))
=\frac1n\operatorname{Var}_P(\xi_{\mathrm{obs}}(O_1)).
\]
Since every observed coordinate is binary,
\[
-2\le\xi_{\mathrm{obs}}(o)\le2,
\qquad
\operatorname{Var}_P(\xi_{\mathrm{obs}}(O_1))
\le E_P[\xi_{\mathrm{obs}}(O_1)^2]\le4.
\]
Expanding the squared error around the mean now gives
\[
E_P\!\left[
\left(\overline\xi_n(\mathbf O_n)-\tau(P)\right)^2
\right]
=
\frac{\operatorname{Var}_P(\xi_{\mathrm{obs}}(O_1))}{n}
+
\left(E_P[\xi_{\mathrm{obs}}(O_1)]-\tau(P)\right)^2.
\]
By \cref{obj:def:complete-arrival-ht-estimator},
\[
\widehat\tau^{\mathrm{HT}}_n(\mathbf o)
=
\operatorname{clip}_{[-1,1]}(\overline\xi_n(\mathbf o)).
\]
Because \(\tau(P)\in[-1,1]\), considering whether
\(\overline\xi_n(\mathbf o)<-1\), belongs to \([-1,1]\), or exceeds \(1\), gives
\[
\left(
\widehat\tau^{\mathrm{HT}}_n(\mathbf o)-\tau(P)
\right)^2
\le
\left(\overline\xi_n(\mathbf o)-\tau(P)\right)^2.
\]
% lean: CausalSmith.Stat.MarRareqLogfrontier.observed_contrast_risk_le

\item To bound the bias, define the missing-outcome masses for both arms by
\[
\Delta_a(P)
:=
2P(A=a,R=0,Y=1),
\qquad a\in\{0,1\}.
\]
For every cell indexed by \(a,x,s\), the occupied-cell arrival condition implies
\[
q\,p_{axs}(P)
\le P(A=a,X=x,S=s,R=1).
\]
For an occupied cell this is the arrival-floor inequality in
\cref{obj:ass:occupied-cell-arrival}; for a null cell its left side is zero and its right side is nonnegative. Partitioning a cell according to \(R\) therefore gives
\[
\begin{aligned}
P(A=a,X=x,S=s,R=0)
&=
p_{axs}(P)-P(A=a,X=x,S=s,R=1)\\
&\le(1-q)p_{axs}(P).
\end{aligned}
\]
Balanced assignment gives \(P(A=1)=1/2\) and
\(P(A=0)=1-P(A=1)=1/2\). Thus, for each \(a\in\{0,1\}\),
\[
\sum_{x=0}^{d-1}\sum_{s=0}^{1}p_{axs}(P)=\frac12,
\]
and the finite cell partition yields
\[
\begin{aligned}
P(A=a,R=0)
&=
\sum_{x=0}^{d-1}\sum_{s=0}^{1}
P(A=a,X=x,S=s,R=0)\\
&\le
(1-q)\sum_{x=0}^{d-1}\sum_{s=0}^{1}p_{axs}(P)
=\frac{1-q}{2}.
\end{aligned}
\]
The event \(\{A=a,R=0,Y=1\}\) is contained in \(\{A=a,R=0\}\), so
\[
0\le\Delta_a(P)\le2P(A=a,R=0)\le1-q.
\]

Outcome consistency, randomized independence, and balanced assignment give
\[
\begin{aligned}
P(A=a,Y=1)
&=P(A=a,Y(a)=1)\\
&=P(A=a)P(Y(a)=1)
=\frac12P(Y(a)=1).
\end{aligned}
\]
Using the observed-record map and partitioning the outcome-one event according to arrival, we obtain
\[
\begin{aligned}
E_P[\xi_{\mathrm{obs}}(O_1)]
&=
2P(A=1,R=1,Y=1)-2P(A=0,R=1,Y=1)\\
&=
2P(A=1,Y=1)-2P(A=0,Y=1)
-\Delta_1(P)+\Delta_0(P)\\
&=\tau(P)-\Delta_1(P)+\Delta_0(P).
\end{aligned}
\]
The bounds on the two deficits imply
\[
-(1-q)
\le
E_P[\xi_{\mathrm{obs}}(O_1)]-\tau(P)
\le1-q,
\]
and hence, since \(1-q\ge0\),
\[
\left(E_P[\xi_{\mathrm{obs}}(O_1)]-\tau(P)\right)^2
\le(1-q)^2.
\]
Together with the preceding variance and clipping bounds, this gives
\[
\begin{aligned}
E_P\!\left[
\left(\widehat\tau^{\mathrm{HT}}_n(\mathbf O_n)-\tau(P)\right)^2
\right]
&\le
E_P\!\left[
\left(\overline\xi_n(\mathbf O_n)-\tau(P)\right)^2
\right]\\
&\le\frac4n+(1-q)^2.
\end{aligned}
\]
The clipped contrast is measurable on the finite sample space and takes values in \([-1,1]\). Since \(4/n+(1-q)^2\ge0\), the deterministic comparison gives both
\[
\mathfrak R^{+}_{n,d,q}\le\frac4n+(1-q)^2
\]
and
\[
\sup_{P\in\mathcal M^{+}_{n,d,q}}
E_P\!\left[
\left(\widehat\tau^{\mathrm{HT}}_n(\mathbf O_n)-\tau(P)\right)^2
\right]
\le\frac4n+(1-q)^2.
\]
% lean: CausalSmith.Stat.MarRareqLogfrontier.observed_contrast_risk_le

\item For the lower bound, first construct a complete-arrival pair. Define
\[
\delta_{\mathrm{gap}}:=\frac1{4\sqrt n},
\qquad
0<\delta_{\mathrm{gap}}\le\frac14,
\qquad
\delta_{\mathrm{gap}}^2=\frac1{16n}.
\]
For \(\iota\in\{0,1\}\), let \(P_{\mathrm{test},\iota}\) be the full-data law specified by
\[
\begin{gathered}
X=0,\qquad S(0)=S(1)=0,\qquad Y(0)=0,\qquad R=1,\\
A\sim\operatorname{Bernoulli}(1/2),\qquad
Y(1)\sim\operatorname{Bernoulli}(1/2+\iota\delta_{\mathrm{gap}}),
\qquad A\perp\!\!\!\perp Y(1),\\
S=S(A)=0,\qquad Y=Y(A)=A\,Y(1).
\end{gathered}
\]
The baseline value \(0\) is available for every \(d\ge1\), and the outcome parameters lie in \([1/2,3/4]\). Independence and the fair treatment coin verify randomized independence and balanced assignment. The displayed definitions verify consistency. Constant arrival verifies arrival MAR and gives arrival probability one in every occupied cell. Therefore
\[
P_{\mathrm{test},0},P_{\mathrm{test},1}
\in\mathcal M^{+}_{n,d,1},
\qquad
\tau(P_{\mathrm{test},\iota})
=\frac12+\iota\delta_{\mathrm{gap}}.
\]

Define the observed laws and their sample laws by
\[
Q_{\mathrm{test},\iota}
:=\mathcal L_{P_{\mathrm{test},\iota}}(O_1),
\qquad
Q_{\mathrm{test},\iota}^{(n)}
:=Q_{\mathrm{test},\iota}^{\otimes n}.
\]
Their common support consists of the three records
\[
\begin{aligned}
o_{\mathrm{test},\mathrm{control}}&:=(0,0,0,1,0),\\
o_{\mathrm{test},\mathrm{zero}}&:=(0,1,0,1,0),\\
o_{\mathrm{test},\mathrm{one}}&:=(0,1,0,1,1),
\end{aligned}
\qquad
\mathcal Z_{\mathrm{test}}
:=
\{o_{\mathrm{test},\mathrm{control}},
o_{\mathrm{test},\mathrm{zero}},
o_{\mathrm{test},\mathrm{one}}\}.
\]
The atom probabilities are
\[
\begin{aligned}
Q_{\mathrm{test},\iota}(\{o_{\mathrm{test},\mathrm{control}}\})&=\frac12,\\
Q_{\mathrm{test},\iota}(\{o_{\mathrm{test},\mathrm{zero}}\})&=\frac14-\frac{\iota\delta_{\mathrm{gap}}}{2},\\
Q_{\mathrm{test},\iota}(\{o_{\mathrm{test},\mathrm{one}}\})&=\frac14+\frac{\iota\delta_{\mathrm{gap}}}{2},
\end{aligned}
\]
with probability zero elsewhere. In particular,
\(Q_{\mathrm{test},1}\ll Q_{\mathrm{test},0}\).

For finite laws with a positive reference law on their common support, define
\[
\chi^2(Q_{\mathrm{test},1}\Vert Q_{\mathrm{test},0})
:=
\sum_{o\in\mathcal Z_{\mathrm{test}}}
\frac{
\left(Q_{\mathrm{test},1}(\{o\})-Q_{\mathrm{test},0}(\{o\})\right)^2
}{
Q_{\mathrm{test},0}(\{o\})
}.
\]
Summing the squared probability ratios gives
\[
\begin{aligned}
1+\chi^2(Q_{\mathrm{test},1}\Vert Q_{\mathrm{test},0})
&=
\sum_{o\in\mathcal Z_{\mathrm{test}}}
\frac{Q_{\mathrm{test},1}(\{o\})^2}{Q_{\mathrm{test},0}(\{o\})}\\
&=
\frac{(1/2)^2}{1/2}
+
\frac{(1/4-\delta_{\mathrm{gap}}/2)^2}{1/4}
+
\frac{(1/4+\delta_{\mathrm{gap}}/2)^2}{1/4}\\
&=1+2\delta_{\mathrm{gap}}^2.
\end{aligned}
\]
Thus
\[
\chi^2(Q_{\mathrm{test},1}\Vert Q_{\mathrm{test},0})
=2\delta_{\mathrm{gap}}^2=\frac1{8n}.
\]
% lean: CausalSmith.Stat.MarRareqLogfrontier.completeFamilyLaw_obs_chi

\item Define the sample likelihood ratio on \(\mathcal O_d^n\) by
\[
L_{\mathrm{test},n}(\mathbf o)
:=
\begin{cases}
\displaystyle\prod_{i=1}^{n}
\frac{Q_{\mathrm{test},1}(\{o_i\})}
{Q_{\mathrm{test},0}(\{o_i\})},
&\mathbf o\in\mathcal Z_{\mathrm{test}}^n,\\
0,&\mathbf o\notin\mathcal Z_{\mathrm{test}}^n.
\end{cases}
\]
The product laws satisfy
\[
E_{Q_{\mathrm{test},0}^{(n)}}[L_{\mathrm{test},n}]=1,
\]
and independence gives
\[
\begin{aligned}
1+\chi^2(Q_{\mathrm{test},1}^{(n)}\Vert Q_{\mathrm{test},0}^{(n)})
&=
E_{Q_{\mathrm{test},0}^{(n)}}[L_{\mathrm{test},n}^2]\\
&=
\left(
1+\chi^2(Q_{\mathrm{test},1}\Vert Q_{\mathrm{test},0})
\right)^n.
\end{aligned}
\]
Here the product divergence is defined by the same finite-sum formula on
\(\mathcal Z_{\mathrm{test}}^n\). Since \(1+t\le e^t\) for \(t\ge0\),
\[
\chi^2(Q_{\mathrm{test},1}^{(n)}\Vert Q_{\mathrm{test},0}^{(n)})
\le e^{1/8}-1<1.
\]
For the numerical inequality, monotonicity gives \(e^{1/8}\le e^{1/2}\), while
\[
(e^{1/2})^2=e<3<4,
\]
so \(e^{1/8}<2\).

Define total variation for these sample laws by
\[
\operatorname{TV}(Q_{\mathrm{test},1}^{(n)},Q_{\mathrm{test},0}^{(n)})
:=
\frac12
\sum_{\mathbf o\in\mathcal Z_{\mathrm{test}}^n}
\left|
Q_{\mathrm{test},1}^{(n)}(\{\mathbf o\})
-
Q_{\mathrm{test},0}^{(n)}(\{\mathbf o\})
\right|.
\]
Cauchy--Schwarz applied to the likelihood ratio yields
\[
\begin{aligned}
\operatorname{TV}(Q_{\mathrm{test},1}^{(n)},Q_{\mathrm{test},0}^{(n)})
&=
\frac12E_{Q_{\mathrm{test},0}^{(n)}}|L_{\mathrm{test},n}-1|\\
&\le
\frac12\sqrt{
E_{Q_{\mathrm{test},0}^{(n)}}[(L_{\mathrm{test},n}-1)^2]
}\\
&=
\frac12\sqrt{
\chi^2(Q_{\mathrm{test},1}^{(n)}\Vert Q_{\mathrm{test},0}^{(n)})
}
\le\frac12.
\end{aligned}
\]
Consequently, for every measurable \(\mathcal E\subseteq\mathcal O_d^n\),
\[
\begin{aligned}
Q_{\mathrm{test},1}^{(n)}(\mathcal E^c)
+Q_{\mathrm{test},0}^{(n)}(\mathcal E)
&=
1-\left(
Q_{\mathrm{test},1}^{(n)}(\mathcal E)
-Q_{\mathrm{test},0}^{(n)}(\mathcal E)
\right)\\
&\ge
1-\operatorname{TV}(Q_{\mathrm{test},1}^{(n)},Q_{\mathrm{test},0}^{(n)})
\ge\frac12.
\end{aligned}
\]
% lean: CausalSmith.Stat.MarRareqLogfrontier.complete_arrival_product_testing_half

\item Take any parameter-independent bounded kernel \(\mathsf T\) in the decision class, and define its barycenter by
\[
b_{\mathsf T}(\mathbf o)
:=
\int_{[-1,1]}h\,\mathsf T(\mathbf o,dh).
\]
By \cref{obj:lem:randomized-kernel-barycenter}, this is a measurable map into \([-1,1]\), and
\[
E_{P_{\mathrm{test},\iota}}\!\left[
\left(b_{\mathsf T}(\mathbf O_n)-\tau(P_{\mathrm{test},\iota})\right)^2
\right]
\le
\mathcal R_{P_{\mathrm{test},\iota}}(\mathsf T),
\qquad \iota\in\{0,1\}.
\]
Define the two large-error events and the nearest-target test by
\[
\mathcal E_{\mathrm{test},\iota}
:=
\left\{\mathbf o:
\frac{\delta_{\mathrm{gap}}}{2}
\le
\left|b_{\mathsf T}(\mathbf o)-\tau(P_{\mathrm{test},\iota})\right|
\right\},
\qquad \iota\in\{0,1\},
\]
and
\[
\mathcal A_{\mathrm{test}}
:=
\left\{\mathbf o:
\left|b_{\mathsf T}(\mathbf o)-\tau(P_{\mathrm{test},1})\right|
\le
\left|b_{\mathsf T}(\mathbf o)-\tau(P_{\mathrm{test},0})\right|
\right\}.
\]
The target separation and the triangle inequality give, for every sample,
\[
\delta_{\mathrm{gap}}
\le
\left|b_{\mathsf T}(\mathbf o)-\tau(P_{\mathrm{test},0})\right|
+
\left|b_{\mathsf T}(\mathbf o)-\tau(P_{\mathrm{test},1})\right|.
\]
On \(\mathcal A_{\mathrm{test}}\), the first error is at least
\(\delta_{\mathrm{gap}}/2\); on its complement, the second error is at least
\(\delta_{\mathrm{gap}}/2\). Hence
\[
\mathcal A_{\mathrm{test}}\subseteq\mathcal E_{\mathrm{test},0},
\qquad
\mathcal A_{\mathrm{test}}^c\subseteq\mathcal E_{\mathrm{test},1}.
\]
Applying the testing bound and then event monotonicity gives
\[
\frac12
\le
Q_{\mathrm{test},1}^{(n)}(\mathcal A_{\mathrm{test}}^c)
+
Q_{\mathrm{test},0}^{(n)}(\mathcal A_{\mathrm{test}})
\le
Q_{\mathrm{test},1}^{(n)}(\mathcal E_{\mathrm{test},1})
+
Q_{\mathrm{test},0}^{(n)}(\mathcal E_{\mathrm{test},0}),
\]
so
\[
\max_{\iota\in\{0,1\}}
Q_{\mathrm{test},\iota}^{(n)}(\mathcal E_{\mathrm{test},\iota})
\ge\frac14.
\]
Integrating squared error over each large-error event yields
\[
\left(\frac{\delta_{\mathrm{gap}}}{2}\right)^2
Q_{\mathrm{test},\iota}^{(n)}(\mathcal E_{\mathrm{test},\iota})
\le
E_{P_{\mathrm{test},\iota}}\!\left[
\left(b_{\mathsf T}(\mathbf O_n)-\tau(P_{\mathrm{test},\iota})\right)^2
\right].
\]
Whether the probability for \(\iota=0\) is smaller than that for \(\iota=1\), or conversely, multiplication by the nonnegative factor
\((\delta_{\mathrm{gap}}/2)^2\) preserves their ordering. Therefore
\[
\begin{aligned}
\frac{\delta_{\mathrm{gap}}^2}{16}
&\le
\left(\frac{\delta_{\mathrm{gap}}}{2}\right)^2
\max_{\iota\in\{0,1\}}
Q_{\mathrm{test},\iota}^{(n)}(\mathcal E_{\mathrm{test},\iota})\\
&=
\max_{\iota\in\{0,1\}}
\left\{
\left(\frac{\delta_{\mathrm{gap}}}{2}\right)^2
Q_{\mathrm{test},\iota}^{(n)}(\mathcal E_{\mathrm{test},\iota})
\right\}\\
&\le
\max_{\iota\in\{0,1\}}
E_{P_{\mathrm{test},\iota}}\!\left[
\left(b_{\mathsf T}(\mathbf O_n)-\tau(P_{\mathrm{test},\iota})\right)^2
\right]\\
&\le
\max_{\iota\in\{0,1\}}
\mathcal R_{P_{\mathrm{test},\iota}}(\mathsf T).
\end{aligned}
\]
For every bounded kernel and every full-data law, bounded actions and
\(\tau(P)\in[-1,1]\) also give
\[
0\le\mathcal R_P(\mathsf T)\le4.
\]
Thus the worst-case risks are finite. The decision class is nonempty because it contains the zero estimator, and the complete-arrival class contains the constructed pair. Taking the supremum over that class and then the infimum over kernels gives
\[
\mathfrak R^{+}_{n,d,1}
\ge
\frac{\delta_{\mathrm{gap}}^2}{16}
=
\frac1{256n}.
\]
% lean: CausalSmith.Stat.MarRareqLogfrontier.complete_arrival_two_point_minimax_floor

\item Finally, every complete-arrival law belongs to the model with arrival floor \(q\). Indeed, all other conditions are shared, and for each occupied cell of such a law,
\[
q\,p_{axs}(P)
\le p_{axs}(P)
\le P(A=a,X=x,S=s,R=1).
\]
Together with \(0<q\le1\), this proves
\[
\mathcal M^{+}_{n,d,1}\subseteq\mathcal M^{+}_{n,d,q}.
\]
For every kernel \(\mathsf T\), class inclusion gives
\[
\sup_{P\in\mathcal M^{+}_{n,d,1}}\mathcal R_P(\mathsf T)
\le
\sup_{P\in\mathcal M^{+}_{n,d,q}}\mathcal R_P(\mathsf T).
\]
The risks lie in \([0,4]\), and both classes contain the complete-arrival pair. Taking infima over the same nonempty decision class therefore yields
\[
\frac1{256n}
\le
\mathfrak R^{+}_{n,d,1}
\le
\mathfrak R^{+}_{n,d,q}.
\]
Combining this with the three upper bounds proves
\[
\frac1{256n}
\le
\mathfrak R^{+}_{n,d,q}
\le
\min\left\{
1,\,
\overline C\,r_{n,d,q},\,
\frac4n+(1-q)^2
\right\}.
\]
The first step supplies the mixed-count envelope and its auxiliary-risk comparison, and the observed-contrast calculation supplies the asserted uniform witness bound.
% lean: CausalSmith.Stat.MarRareqLogfrontier.unrestricted_near_complete_arrival_envelope
\end{enumerate}
\end{proof}

\begin{proof}[Proof of \cref{obj:thm:unrestricted-large-alphabet-deficit-lower}]
\begin{enumerate}
\item Put
\[
\delta:=1-q\in[0,1).
\]
Since \(d\ge1024n^2\ge1\), the hypotheses of
\cref{obj:thm:unrestricted-near-complete-arrival-envelope} hold. Its lower bound and the last term of its upper envelope give
\[
\frac1{256n}\le\mathfrak R^{+}_{n,d,q},
\qquad
\mathfrak R^{+}_{n,d,q}\le\frac4n+\delta^2.
\]
If \(\delta^2\le n^{-1}\), then
\[
\frac{\delta^2}{1024}
\le\frac1{1024n}
\le\frac1{256n}
\le\mathfrak R^{+}_{n,d,q}.
\]
This case includes \(q=1\). For the remaining case, assume
\[
\delta^2>n^{-1},
\qquad\text{so in particular}\qquad n^{-1}\le\delta^2.
\]
We establish the deficit lower bound in this regime. % lean: CausalSmith.Stat.MarRareqLogfrontier.unrestricted_large_alphabet_deficit_lower

\item Write an ordered observed sample as
\[
\mathbf o=(o_1,\ldots,o_n)
\in\left(\{0,\ldots,d-1\}\times\{0,1\}^4\right)^n,
\]
with each record ordered as \((X,A,S,R,RY)\). The decision class contains the constant zero kernel
\[
\mathsf T_0(\mathbf o,\cdot):=\operatorname{Dirac}(0).
\]
Fix an arbitrary \(\mathsf T\in\mathcal T_n\), the bounded, parameter-independent kernel class of
\cref{obj:def:unrestricted-minimax-risk}, and define its risk under a full-data law \(P\) by
\[
\mathcal R_{\mathsf T}(P)
:=
\int\!\int_{[-1,1]}
(w-\tau(P))^2\,
\mathsf T(\mathbf o,dw)\,
\mathcal L_P(\mathbf O_n)(d\mathbf o).
\]
Define the separation between the two target centers by
\[
h_{\mathrm{sep}}
:=\frac12-\frac q{1+q}
=\frac{\delta}{2(1+q)}.
\]
Because \(1<1+q\le2\) and \(\delta\ge0\), this identity gives
\[
\frac{\delta}{4}\le h_{\mathrm{sep}}\le\frac{\delta}{2},
\qquad h_{\mathrm{sep}}\ge0.
\]

Construct a comparison law \(P_\star\) by taking a uniform baseline, zero surrogates and control outcome, and a Bernoulli treated potential outcome:
\[
X\sim\operatorname{Unif}\{0,\ldots,d-1\},
\qquad
S(0)=S(1)=Y(0)=0,
\qquad
Y(1)\mid X\sim\operatorname{Bernoulli}\!\left(\frac q{1+q}\right).
\]
Draw \(A\sim\operatorname{Bernoulli}(1/2)\) independently of the baseline and potential variables, set
\[
S=S(A),\qquad Y=Y(A),
\]
and draw the arrival indicator conditionally independently of \(Y(1)\), with
\[
P_\star(R=1\mid X,A,Y(1))
=
\begin{cases}
1,&A=0,\\
(1+q)/2,&A=1.
\end{cases}
\]
These distributions are well defined because
\[
0<\frac q{1+q}\le\frac12,
\qquad
q\le\frac{1+q}{2}\le1.
\]
Evaluating the average treatment effect of
\cref{obj:def:ate-functional} yields
\[
\tau(P_\star)
=E_{P_\star}[Y(1)-Y(0)]
=\frac q{1+q}
=\frac12-h_{\mathrm{sep}}.
\]
% lean: CausalSmith.Stat.MarRareqLogfrontier.largeAlphabet_center_separation
% lean: CausalSmith.Stat.MarRareqLogfrontier.largeAlphabetComparisonLaw_ate

\item For each
\[
\xi=(\xi_x)_{x=0}^{d-1}\in\{0,1\}^{\{0,\ldots,d-1\}},
\]
construct a law \(P_\xi\) with the same uniform baseline and independent fair treatment assignment, and with
\[
S(0)=S(1)=Y(0)=0,\qquad Y(1)=\xi_X,
\qquad S=S(A),\qquad Y=Y(A).
\]
Its arrival draw satisfies, for \(x\in\{0,\ldots,d-1\}\) and \(a\in\{0,1\}\),
\[
P_\xi(R=1\mid X=x,A=a)=1-a\delta\xi_x.
\]
Thus every arrival probability is either \(q\) or \(1\), and
\[
\tau(P_\xi)
=\frac1d\sum_{x=0}^{d-1}\xi_x.
\]
Put the product fair-coin prior on these vectors:
\[
\nu_{\mathrm{cell}}
:=\bigotimes_{x=0}^{d-1}\operatorname{Bernoulli}(1/2),
\qquad
\nu_{\mathrm{cell}}(\{\xi\})=2^{-d}.
\]
Define the sample laws and their mixture by
\[
Q_\xi:=\mathcal L_{P_\xi}(\mathbf O_n),
\qquad
Q_\star:=\mathcal L_{P_\star}(\mathbf O_n),
\qquad
Q_{\mathrm{mix}}
:=\int Q_\xi\,\nu_{\mathrm{cell}}(d\xi)
=\sum_\xi2^{-d}Q_\xi.
\]
For each fixed full-data law, the sample consists of \(n\) independent observed records as in \cref{obj:ass:iid-sampling}. The mixture draws \(\xi\) once for the entire sample.

Define the test event, the exceptional prior vectors, and the two risks by
\[
E_{\mathrm{test}}
:=
\left\{(\mathbf o,w):
w\le\frac12-\frac{h_{\mathrm{sep}}}{2}\right\},
\qquad
w\in\mathbb R,
\]
\[
F_{\mathrm{bad}}
:=
\left\{\xi:
\left|\tau(P_\xi)-\frac12\right|>
\frac{h_{\mathrm{sep}}}{4}\right\},
\]
\[
\mathcal R_{\mathrm{prior}}(\mathsf T)
:=\int\mathcal R_{\mathsf T}(P_\xi)\,\nu_{\mathrm{cell}}(d\xi),
\qquad
\mathcal R_\star(\mathsf T):=\mathcal R_{\mathsf T}(P_\star).
\]
The notation \(Q\ltimes\mathsf T\) denotes the joint sample--action law used in
\cref{obj:lem:randomized-kernel-tv-comparison}.

If \(\xi\notin F_{\mathrm{bad}}\) and \((\mathbf o,w)\in E_{\mathrm{test}}\), then
\[
\tau(P_\xi)-w
\ge
\left(\frac12-\frac{h_{\mathrm{sep}}}{4}\right)
-\left(\frac12-\frac{h_{\mathrm{sep}}}{2}\right)
=\frac{h_{\mathrm{sep}}}{4}.
\]
Consequently, for every \(\xi\),
\[
\frac{h_{\mathrm{sep}}^2}{16}
(Q_\xi\ltimes\mathsf T)(E_{\mathrm{test}})
\le
\mathcal R_{\mathsf T}(P_\xi)
+\frac{h_{\mathrm{sep}}^2}{16}
\mathbf1_{F_{\mathrm{bad}}}(\xi).
\]
For exceptional vectors this follows from the probability bound by one and nonnegativity of risk; for the other vectors it follows by integrating the displayed pointwise loss bound. Since the same kernel is used for every \(\xi\), finite averaging gives
\[
(Q_{\mathrm{mix}}\ltimes\mathsf T)(E_{\mathrm{test}})
=
\int
(Q_\xi\ltimes\mathsf T)(E_{\mathrm{test}})
\,\nu_{\mathrm{cell}}(d\xi).
\]
Averaging the member-wise risk inequality therefore yields
\[
\frac{h_{\mathrm{sep}}^2}{16}
(Q_{\mathrm{mix}}\ltimes\mathsf T)(E_{\mathrm{test}})
\le
\mathcal R_{\mathrm{prior}}(\mathsf T)
+\frac{h_{\mathrm{sep}}^2}{16}\nu_{\mathrm{cell}}(F_{\mathrm{bad}}).
\]

Under \(P_\star\), an action in \(E_{\mathrm{test}}^c\) satisfies
\[
w-\tau(P_\star)
>
\left(\frac12-\frac{h_{\mathrm{sep}}}{2}\right)
-\left(\frac12-h_{\mathrm{sep}}\right)
=\frac{h_{\mathrm{sep}}}{2}.
\]
In particular, its squared loss is at least \(h_{\mathrm{sep}}^2/16\), so
\[
\frac{h_{\mathrm{sep}}^2}{16}
(Q_\star\ltimes\mathsf T)(E_{\mathrm{test}}^c)
\le\mathcal R_\star(\mathsf T).
\]
% lean: CausalSmith.Stat.MarRareqLogfrontier.largeAlphabetLaw_ate
% lean: CausalSmith.Stat.MarRareqLogfrontier.largeAlphabet_prior_test_risk
% lean: CausalSmith.Stat.MarRareqLogfrontier.largeAlphabet_comparison_test_risk

\item We next compare the two sample laws. For a single treated visit to any baseline cell \(x\), averaging the fair coordinate \(\xi_x\) gives the following observed-mark probabilities:
\[
\begin{array}{c|c|c}
(R,RY)
&
\displaystyle\int P_\xi((R,RY)\mid X=x,A=1)\,
\nu_{\mathrm{cell}}(d\xi)
&
P_\star((R,RY)\mid X=x,A=1)
\\ \hline
(1,1)&q/2&
\displaystyle\frac{1+q}{2}\frac q{1+q}=q/2
\\
(1,0)&1/2&
\displaystyle\frac{1+q}{2}\left(1-\frac q{1+q}\right)=1/2
\\
(0,0)&\delta/2&
\displaystyle1-\frac{1+q}{2}=\delta/2
\\
(0,1)&0&0
\end{array}
\]
In the prior column, the first probability comes from \(\xi_x=1\) and arrival, the second from \(\xi_x=0\), and the third from \(\xi_x=1\) and nonarrival. A control visit has mark \((R,RY)=(1,0)\) with probability one under every law. All surrogates equal zero.

Define the event of distinct treated baseline cells by
\[
E_{\mathrm{distinct}}
:=
\bigcap_{1\le i<j\le n}
\{A_i\ne1\ \text{or}\ A_j\ne1\ \text{or}\ X_i\ne X_j\}.
\]
For any sample atom \(\mathbf o=(o_1,\ldots,o_n)\) in this event, each treated-record probability depends on a different prior coordinate, while each control-record probability is constant in \(\xi\). Independence of the prior coordinates and the one-visit calculation therefore give
\[
\begin{aligned}
Q_{\mathrm{mix}}(\{\mathbf o\})
&=
\int
\prod_{i=1}^n
\mathcal L_{P_\xi}(O_1)(\{o_i\})\,
\nu_{\mathrm{cell}}(d\xi)
\\
&=
\prod_{i=1}^n
\int
\mathcal L_{P_\xi}(O_1)(\{o_i\})\,
\nu_{\mathrm{cell}}(d\xi)
\\
&=
\prod_{i=1}^n
\mathcal L_{P_\star}(O_1)(\{o_i\})
=
Q_\star(\{\mathbf o\}).
\end{aligned}
\]
The equality also covers atoms with impossible surrogate or arrived-outcome marks, whose probabilities are zero under both constructions. % lean: CausalSmith.Stat.MarRareqLogfrontier.largeAlphabetSampleMixture_matching

\item Under each \(P_\xi\) and under \(P_\star\), the probability of any treated baseline cell is
\[
P_\xi(A=1,X=x)=P_\star(A=1,X=x)=\frac1{2d},
\qquad x\in\{0,\ldots,d-1\}.
\]
Thus independence of distinct observations gives, for \(1\le i<j\le n\),
\[
\begin{aligned}
Q_\xi\{A_i=A_j=1,\ X_i=X_j\}
&\le
\sum_{x=0}^{d-1}\left(\frac1{2d}\right)^2
=\frac1{4d},
\\
Q_\star\{A_i=A_j=1,\ X_i=X_j\}
&\le\frac1{4d}.
\end{aligned}
\]
Taking a union over the \(\binom n2\) pairs yields
\[
Q_\xi(E_{\mathrm{distinct}}^c),
\ Q_\star(E_{\mathrm{distinct}}^c)
\le
\binom n2\frac1{4d}
\le\frac{n^2}{8d}
\le\frac1{8192}.
\]
When \(n=1\), the pair set is empty and these collision probabilities are zero. Averaging the first bound over the prior gives
\[
Q_{\mathrm{mix}}(E_{\mathrm{distinct}}^c)
=
\int Q_\xi(E_{\mathrm{distinct}}^c)\,
\nu_{\mathrm{cell}}(d\xi)
\le\frac1{8192}.
\]

On the finite sample space, agreement of the atom masses on
\(E_{\mathrm{distinct}}\) implies
\[
\begin{aligned}
\operatorname{TV}(Q_{\mathrm{mix}},Q_\star)
&\le
\frac12\sum_{\mathbf o}
\left|Q_{\mathrm{mix}}(\{\mathbf o\})
-Q_\star(\{\mathbf o\})\right|
\\
&\le
\frac12\left\{
Q_{\mathrm{mix}}(E_{\mathrm{distinct}}^c)
+Q_\star(E_{\mathrm{distinct}}^c)
\right\}
\le\frac1{8192}.
\end{aligned}
\]
Indeed, the summands vanish on \(E_{\mathrm{distinct}}\), and elsewhere each absolute difference is bounded by the sum of the two atom masses. Applying
\cref{obj:lem:randomized-kernel-tv-comparison} to the common kernel and the measurable event \(E_{\mathrm{test}}\) now gives
\[
1-\frac1{8192}
\le
1-\operatorname{TV}(Q_{\mathrm{mix}},Q_\star)
\le
(Q_{\mathrm{mix}}\ltimes\mathsf T)(E_{\mathrm{test}})
+
(Q_\star\ltimes\mathsf T)(E_{\mathrm{test}}^c).
\]
% lean: CausalSmith.Stat.MarRareqLogfrontier.largeAlphabet_actual_collision_testing

\item Under \(\nu_{\mathrm{cell}}\), the independent fair coordinates have
\[
E_{\nu_{\mathrm{cell}}}[\xi_x]=\frac12,
\qquad
\operatorname{Var}_{\nu_{\mathrm{cell}}}(\xi_x)=\frac14.
\]
Consequently,
\[
E_{\nu_{\mathrm{cell}}}\!\left[\sum_{x=0}^{d-1}\xi_x\right]
=\frac d2,
\qquad
\operatorname{Var}_{\nu_{\mathrm{cell}}}
\!\left(\sum_{x=0}^{d-1}\xi_x\right)
=\frac d4.
\]
The regime \(n^{-1}\le\delta^2\), together with \(\delta\ge0\), implies
\[
\delta>0,\qquad h_{\mathrm{sep}}\ge\frac{\delta}{4}>0,
\qquad
1\le n\delta^2,\qquad
\delta^2\le16h_{\mathrm{sep}}^2.
\]
These inequalities and the alphabet-size condition give the concentration budget
\[
d h_{\mathrm{sep}}^2
\ge1024n^2h_{\mathrm{sep}}^2
\ge64n^2\delta^2
\ge64n
\ge64.
\]
Since \(\tau(P_\xi)=d^{-1}\sum_x\xi_x\), Chebyshev's inequality applied to the sum of the prior coordinates yields
\[
\begin{aligned}
\nu_{\mathrm{cell}}(F_{\mathrm{bad}})
&=
\nu_{\mathrm{cell}}
\left\{
\left|\sum_{x=0}^{d-1}\xi_x-\frac d2\right|
>\frac{d h_{\mathrm{sep}}}{4}
\right\}
\\
&\le
\frac{d/4}{(d h_{\mathrm{sep}}/4)^2}
=
\frac4{d h_{\mathrm{sep}}^2}
\le\frac1{16}.
\end{aligned}
\]
% lean: CausalSmith.Stat.MarRareqLogfrontier.largeAlphabet_center_prior_concentration

\item Multiplying the testing and concentration bounds by the nonnegative factor \(h_{\mathrm{sep}}^2/16\) gives
\[
\frac{h_{\mathrm{sep}}^2}{16}
\left(1-\frac1{8192}\right)
\le
\frac{h_{\mathrm{sep}}^2}{16}
\left\{
(Q_{\mathrm{mix}}\ltimes\mathsf T)(E_{\mathrm{test}})
+
(Q_\star\ltimes\mathsf T)(E_{\mathrm{test}}^c)
\right\},
\]
\[
\frac{h_{\mathrm{sep}}^2}{16}
\nu_{\mathrm{cell}}(F_{\mathrm{bad}})
\le\frac{h_{\mathrm{sep}}^2}{256}.
\]
Combining these with the two test-risk inequalities established above, we obtain
\[
\begin{aligned}
\mathcal R_{\mathrm{prior}}(\mathsf T)
+\mathcal R_\star(\mathsf T)
&\ge
\frac{h_{\mathrm{sep}}^2}{16}
\left\{
1-\frac1{8192}-\frac1{16}
\right\}
\\
&\ge\frac{h_{\mathrm{sep}}^2}{32},
\end{aligned}
\]
where the last inequality uses
\(1-1/8192-1/16\ge1/2\). % lean: CausalSmith.Stat.MarRareqLogfrontier.largeAlphabet_deficit_minimax_lower

\item To pass to worst-case risk, define
\[
\mathcal R_{\mathrm{sup}}(\mathsf T)
:=
\sup_{P\in\mathcal M^{+}_{n,d,q}}
\mathcal R_{\mathsf T}(P).
\]
Binary potential outcomes give \(\tau(P)\in[-1,1]\). Since the kernel takes values in \([-1,1]\), its loss and risk satisfy
\[
0\le(w-\tau(P))^2\le4
\quad(w\in[-1,1]),
\qquad
0\le\mathcal R_{\mathsf T}(P)\le4.
\]
Thus the member-law risks are bounded above.

Every constructed law belongs to the model class in
\cref{obj:def:unrestricted-arrival-model-class}. The independent fair treatment draw verifies
\cref{obj:ass:randomized-independence,obj:ass:balanced-randomization}, and the defining identities for the realized variables verify
\cref{obj:ass:surrogate-consistency,obj:ass:outcome-consistency}.
Within an occupied cell, the outcome under \(P_\xi\) is deterministic, while under \(P_\star\) its arrival draw is independent of the outcome; these properties verify
\cref{obj:ass:arrival-mar}. For every \(a\in\{0,1\}\) and \(x\in\{0,\ldots,d-1\}\), the cell probabilities of
\cref{obj:synth_16} are
\[
p_{ax0}(P_\xi)=p_{ax0}(P_\star)=\frac1{2d},
\qquad
p_{ax1}(P_\xi)=p_{ax1}(P_\star)=0.
\]
On the occupied cells, their arrival probabilities are
\[
\rho_{ax0}(P_\xi)=1-a\delta\xi_x\ge q,
\qquad
\rho_{ax0}(P_\star)
=
\begin{cases}
1,&a=0,\\
(1+q)/2,&a=1,
\end{cases}
\ge q.
\]
This verifies \cref{obj:ass:occupied-cell-arrival}, and hence
\[
P_\xi,P_\star\in\mathcal M^{+}_{n,d,q}.
\]

Each member-law risk is therefore at most
\(\mathcal R_{\mathrm{sup}}(\mathsf T)\). Since
\(\nu_{\mathrm{cell}}\) is a probability law,
\[
\mathcal R_{\mathrm{prior}}(\mathsf T)
\le
\int\mathcal R_{\mathrm{sup}}(\mathsf T)\,
\nu_{\mathrm{cell}}(d\xi)
=
\mathcal R_{\mathrm{sup}}(\mathsf T),
\qquad
\mathcal R_\star(\mathsf T)
\le\mathcal R_{\mathrm{sup}}(\mathsf T).
\]
It follows that
\[
2\mathcal R_{\mathrm{sup}}(\mathsf T)
\ge
\mathcal R_{\mathrm{prior}}(\mathsf T)+\mathcal R_\star(\mathsf T)
\ge\frac{h_{\mathrm{sep}}^2}{32}.
\]
Using \(h_{\mathrm{sep}}\ge\delta/4\ge0\), we conclude
\[
\mathcal R_{\mathrm{sup}}(\mathsf T)
\ge\frac{h_{\mathrm{sep}}^2}{64}
\ge\frac{\delta^2}{1024}.
\]
This holds for every kernel in the stated decision class. Taking its infimum, as in
\cref{obj:def:unrestricted-minimax-risk}, proves
\[
\frac{\delta^2}{1024}\le\mathfrak R^{+}_{n,d,q}
\]
in the remaining regime. % lean: CausalSmith.Stat.MarRareqLogfrontier.largeAlphabet_deficit_minimax_lower

\item The two cases establish the deficit lower bound throughout the stated parameter range. Combining it with the initial bounds gives
\[
\max\left\{\frac1{256n},\frac{\delta^2}{1024}\right\}
\le\mathfrak R^{+}_{n,d,q}
\le\frac4n+\delta^2.
\]
Adding the two lower inequalities and dividing by two yields
\[
\mathfrak R^{+}_{n,d,q}
\ge
\frac1{512n}+\frac{\delta^2}{2048}
\ge
\frac1{2048}\left\{\frac1n+\delta^2\right\}.
\]
Finally, \(\delta^2\ge0\) gives
\[
\mathfrak R^{+}_{n,d,q}
\le\frac4n+\delta^2
\le4\left\{\frac1n+\delta^2\right\}.
\]
Substituting \(\delta=1-q\) proves the asserted envelope. % lean: CausalSmith.Stat.MarRareqLogfrontier.unrestricted_large_alphabet_deficit_lower
\end{enumerate}
\end{proof}

\begin{proof}[Proof of \cref{obj:prop:unrestricted-uniform-near-complete-elbow}]
\begin{enumerate}
\item We first record two bounds. By \cref{obj:thm:unrestricted-near-complete-arrival-envelope}, for every \(n,d\ge1\),
\[
\mathfrak R^{+}_{n,d,q_n}\le \frac4n+(1-q_n)^2.
\]
For \(n\ge1\), the alphabet size \(1024n^2\) satisfies
\[
1\le1024n^2.
\]
Thus \cref{obj:thm:unrestricted-large-alphabet-deficit-lower}, applied at this alphabet size, supplies
\[
\frac1{2048}\left\{\frac1n+(1-q_n)^2\right\}
\le \mathfrak R^{+}_{n,1024n^2,q_n}.
\]
Multiplication by \(n>0\), using
\[
n\cdot\frac1{2048}\left\{\frac1n+(1-q_n)^2\right\}
=\frac1{2048}\left\{1+n(1-q_n)^2\right\},
\]
gives
\[
\frac1{2048}\left\{1+n(1-q_n)^2\right\}
\le n\mathfrak R^{+}_{n,1024n^2,q_n}.
\]
% lean: CausalSmith.Stat.MarRareqLogfrontier.unrestricted_uniform_near_complete_elbow

\item Multiplying the upper bound from the preceding step by \(n>0\) yields, for every \(d\ge1\),
\[
n\mathfrak R^{+}_{n,d,q_n}
\le n\left\{\frac4n+(1-q_n)^2\right\}
=4+n(1-q_n)^2.
\]
For any real \(b\) satisfying the stated strict inequality, the same step gives
\[
b<
\frac1{2048}\left\{1+n(1-q_n)^2\right\}
\le n\mathfrak R^{+}_{n,1024n^2,q_n}.
\]
Hence \(d=1024n^2\) is the required integer witness.
% lean: CausalSmith.Stat.MarRareqLogfrontier.unrestricted_uniform_near_complete_elbow

\item Suppose first that the arrival-rate condition holds with \(M\ge0\). For all sufficiently large \(n\ge1\),
\[
1-q_n\le\frac{M}{\sqrt n}.
\]
Since \(q_n\le1\) and \(\sqrt n>0\), this implies
\[
0\le(1-q_n)\sqrt n\le M.
\]
Squaring and using \((\sqrt n)^2=n\), we obtain
\[
n(1-q_n)^2
=\bigl((1-q_n)\sqrt n\bigr)^2
\le M^2,
\qquad
(1-q_n)^2\le\frac{M^2}{n}.
\]
Consequently, for every \(d\ge1\) and all such \(n\),
\[
\mathfrak R^{+}_{n,d,q_n}
\le\frac4n+(1-q_n)^2
\le\frac4n+\frac{M^2}{n}
=\frac{4+M^2}{n}.
\]
This proves the uniform risk bound with \(D=4+M^2\).

Conversely, suppose the uniform risk bound holds with a real constant \(D\). Define the nonnegative real constant
\[
M_{\mathrm{risk}}:=\sqrt{\max\{2048D,0\}}.
\]
For all sufficiently large \(n\ge1\), applying that bound at \(d=1024n^2\) and multiplying by \(n\) gives
\[
n\mathfrak R^{+}_{n,1024n^2,q_n}\le D.
\]
Combining this with the lower bound from the first step yields
\[
\frac1{2048}\left\{1+n(1-q_n)^2\right\}\le D.
\]
In particular,
\[
\bigl(\sqrt n(1-q_n)\bigr)^2
=n(1-q_n)^2
\le2048D
\le\max\{2048D,0\}
=M_{\mathrm{risk}}^2.
\]
Both \(\sqrt n(1-q_n)\) and \(M_{\mathrm{risk}}\) are nonnegative, so
\[
\sqrt n(1-q_n)\le M_{\mathrm{risk}},
\qquad
1-q_n\le\frac{M_{\mathrm{risk}}}{\sqrt n}.
\]
This is the arrival-rate condition.
% lean: CausalSmith.Stat.MarRareqLogfrontier.unrestricted_uniform_near_complete_elbow

\item The lower bound in \cref{obj:thm:unrestricted-near-complete-arrival-envelope} directly supplies, for every \(n,d\ge1\),
\[
\frac1{256n}\le\mathfrak R^{+}_{n,d,q_n}.
\]
% lean: CausalSmith.Stat.MarRareqLogfrontier.unrestricted_near_complete_arrival_envelope

\item Finally, assume that \(\bigl(\sqrt n(1-q_n)\bigr)_{n\ge0}\) is unbounded above, and fix \(C\in\mathbb R\). Suppose, for contradiction, that the claimed strict risk inequality holds for only finitely many integers \(n\). Choose an integer \(n_{\mathrm{tail}}\ge1\) such that
\[
n\mathfrak R^{+}_{n,1024n^2,q_n}\le C
\qquad\text{for every }n\ge n_{\mathrm{tail}}.
\]
Define the nonnegative real constant
\[
M_{\mathrm{tail}}:=\sqrt{\max\{2048C,0\}}.
\]
For every \(n\ge n_{\mathrm{tail}}\), the first step implies
\[
\frac1{2048}\left\{1+n(1-q_n)^2\right\}\le C,
\]
and hence
\[
\bigl(\sqrt n(1-q_n)\bigr)^2
=n(1-q_n)^2
\le2048C
\le\max\{2048C,0\}
=M_{\mathrm{tail}}^2.
\]
Nonnegativity therefore gives
\[
\sqrt n(1-q_n)\le M_{\mathrm{tail}}
\qquad(n\ge n_{\mathrm{tail}}).
\]
The remaining terms form a finite set. Explicitly, the finite maximum below exists and bounds every term:
\[
\sqrt n(1-q_n)
\le
\max\left(
\{M_{\mathrm{tail}}\}
\cup
\left\{\sqrt{n'}(1-q_{n'}):
n'\in\mathbb N,\ n'<n_{\mathrm{tail}}\right\}
\right)
\qquad(n\in\mathbb N).
\]
This contradicts unboundedness above. Thus there are infinitely many integers \(n\) for which
\[
C<n\mathfrak R^{+}_{n,1024n^2,q_n}.
\]
% lean: CausalSmith.Stat.MarRareqLogfrontier.unrestricted_uniform_near_complete_elbow
\end{enumerate}
\end{proof}

\begin{proof}[Proof of \cref{obj:prop:unrestricted-near-complete-phase-boundary}]
\begin{enumerate}
\item For all sufficiently large \(n\), we have
\[
n\ge1,
\qquad
1-q_n\le\frac{M}{\sqrt n}.
\]
Fix such an \(n\). Since \(d_n\ge1\) and \(0<q_n\le1\), \cref{obj:thm:unrestricted-near-complete-arrival-envelope} applies. Retaining the last term in its upper-bound minimum gives
\[
\frac{1}{256n}
\le \mathfrak R^{+}_{n,d_n,q_n}
\le \frac4n+(1-q_n)^2.
\]
% lean: CausalSmith.Stat.MarRareqLogfrontier.unrestricted_near_complete_arrival_envelope

\item Since \(n>0\), we have \(\sqrt n>0\), so multiplying the deficit bound by \(\sqrt n\) yields
\[
(1-q_n)\sqrt n\le M.
\]
Moreover, \(q_n\le1\) implies
\[
0\le(1-q_n)\sqrt n.
\]
Thus, using \(M\ge0\), we may square the inequality to obtain
\[
\bigl((1-q_n)\sqrt n\bigr)^2\le M^2.
\]
The identity \((\sqrt n)^2=n\) then gives
\[
n(1-q_n)^2
=
\bigl((1-q_n)\sqrt n\bigr)^2
\le M^2.
\]
Dividing by \(n>0\), we conclude that
\[
(1-q_n)^2\le\frac{M^2}{n}.
\]
% lean: CausalSmith.Stat.MarRareqLogfrontier.unrestricted_near_complete_phase_boundary

\item Substituting this bound into the upper estimate gives
\[
\mathfrak R^{+}_{n,d_n,q_n}
\le \frac4n+(1-q_n)^2
\le \frac4n+\frac{M^2}{n}
=\frac{4+M^2}{n}.
\]
Together with the lower estimate, this proves both asserted inequalities for all sufficiently large \(n\).
% lean: CausalSmith.Stat.MarRareqLogfrontier.unrestricted_near_complete_phase_boundary
\end{enumerate}
\end{proof}

\begin{proof}[Proof of \cref{obj:thm:unrestricted-complete-arrival-frontier}]
Fix integers \(n,d\ge1\), with sampling as in \cref{obj:ass:iid-sampling}.

\begin{enumerate}
\item Since \(0<1\le1\), specialize \cref{obj:thm:unrestricted-near-complete-arrival-envelope} to \(q=1\). Its lower bound and observed-outcome contrast bound give
\[
\frac{1}{256n}\le \mathfrak R^{+}_{n,d,1},
\qquad
\sup_{P\in\mathcal M^{+}_{n,d,1}}
E_P\!\left[
\left\{\widehat\tau^{\mathrm{HT}}_n(\mathbf O_n)-\tau(P)\right\}^{2}
\right]
\le \frac4n+(1-1)^2=\frac4n.
\]
% lean: CausalSmith.Stat.MarRareqLogfrontier.unrestricted_near_complete_arrival_envelope

\item The estimator in \cref{obj:def:complete-arrival-ht-estimator} is measurable, takes values in \([-1,1]\), and depends on the observed sample and the fixed parameters \(n,d\). Define its deterministic probability kernel by
\[
\mathsf T_{\mathrm{HT}}(\mathbf o,E_{\mathrm{out}})
:=
\mathbf 1\!\left\{
\widehat\tau^{\mathrm{HT}}_n(\mathbf o)\in E_{\mathrm{out}}
\right\},
\quad
\mathbf o\in
\bigl(\{0,\ldots,d-1\}\times\{0,1\}^{4}\bigr)^n,
\quad
E_{\mathrm{out}}\subseteq[-1,1]\ \text{Borel}.
\]
This parameter-independent kernel belongs to the decision class \(\mathcal T_n\) of \cref{obj:def:unrestricted-minimax-risk}. For every \(\mathsf T\in\mathcal T_n\) and \(P\in\mathcal M^{+}_{n,d,1}\), squared loss gives
\[
0\le
\int\!\int_{[-1,1]}
(h-\tau(P))^2\,
\mathsf T(\mathbf o,dh)\,
\mathcal L_P(\mathbf O_n)(d\mathbf o).
\]
Evaluating the defining minimax infimum at \(\mathsf T_{\mathrm{HT}}\) therefore yields
\[
\mathfrak R^{+}_{n,d,1}
\le
\sup_{P\in\mathcal M^{+}_{n,d,1}}
\int\!\int_{[-1,1]}
(h-\tau(P))^2\,
\mathsf T_{\mathrm{HT}}(\mathbf o,dh)\,
\mathcal L_P(\mathbf O_n)(d\mathbf o).
\]
By the definition of this kernel, for every such \(P\),
\[
\int\!\int_{[-1,1]}
(h-\tau(P))^2\,
\mathsf T_{\mathrm{HT}}(\mathbf o,dh)\,
\mathcal L_P(\mathbf O_n)(d\mathbf o)
=
E_P\!\left[
\left\{\widehat\tau^{\mathrm{HT}}_n(\mathbf O_n)-\tau(P)\right\}^{2}
\right].
\]
Consequently,
\[
\mathfrak R^{+}_{n,d,1}
\le
\sup_{P\in\mathcal M^{+}_{n,d,1}}
E_P\!\left[
\left\{\widehat\tau^{\mathrm{HT}}_n(\mathbf O_n)-\tau(P)\right\}^{2}
\right].
\]
% lean: CausalSmith.Stat.MarRareqLogfrontier.complete_arrival_frontier_direct

\item Combining these inequalities gives
\[
\frac{1}{256n}
\le \mathfrak R^{+}_{n,d,1}
\le
\sup_{P\in\mathcal M^{+}_{n,d,1}}
E_P\!\left[
\left\{\widehat\tau^{\mathrm{HT}}_n(\mathbf O_n)-\tau(P)\right\}^{2}
\right]
\le \frac4n.
\]
The constants \(1/256\) and \(4\) are independent of \(n\) and \(d\), establishing order \(n^{-1}\) uniformly over \(d\ge1\).
% lean: CausalSmith.Stat.MarRareqLogfrontier.unrestricted_complete_arrival_frontier
\end{enumerate}
\end{proof}

\begin{proof}[Proof of \cref{obj:prop:unrestricted-complete-arrival-phase-boundary}]
\begin{enumerate}
\item Choose the real comparison constants
\[
c:=\frac{1}{256},
\qquad C:=4.
\]
Then \(0<c\le C\). Fix any sequence of natural alphabet sizes \((d_n)_{n\ge0}\) satisfying \(d_n\ge1\) for every \(n\). Define the arrival-floor sequence and deficit constant by
\[
q_n:=1\quad(n\ge0),
\qquad M:=0.
\]
Thus \(M\ge0\) and \(q_n\in(0,1]\) for every \(n\). For every \(n\ge1\),
\[
1-q_n=0=\frac{M}{\sqrt n},
\]
so the arrival-deficit condition holds for all sufficiently large \(n\).
% lean: CausalSmith.Stat.MarRareqLogfrontier.unrestricted_complete_arrival_phase_boundary

\item Applying \cref{obj:prop:unrestricted-near-complete-phase-boundary} with these choices gives, for all sufficiently large \(n\),
\[
\frac{1}{256n}
\le \mathfrak R^{+}_{n,d_n,1}
\le \frac{4}{n},
\]
since \(q_n=1\) and \(4+M^2=4\). These are precisely the required bounds with the chosen \(c\) and \(C\). Both constants were fixed before choosing the alphabet-size sequence, so they are independent of that sequence.
% lean: CausalSmith.Stat.MarRareqLogfrontier.unrestricted_near_complete_phase_boundary
% lean: CausalSmith.Stat.MarRareqLogfrontier.unrestricted_complete_arrival_phase_boundary
\end{enumerate}
\end{proof}

\begin{proof}[Proof of \cref{obj:thm:matched-minimax-frontier}]
\begin{enumerate}
\item Choose universal constants \(\underline c\) and \(C_0\) supplied by
\cref{obj:thm:full-experiment-lower,obj:thm:mixed-count-upper}, respectively, and set
\[
0<\underline c,\qquad 0<C_0,\qquad
\overline C:=\max\{\underline c,C_0\}.
\]
Thus
\[
0<\underline c\le\overline C<\infty,
\qquad C_0\le\overline C.
\]
Fix \(n,d,q\) satisfying the hypotheses of the theorem.
By \cref{obj:thm:full-experiment-lower},
\[
\underline c\,r_{n,d,q}\le\mathfrak R_{n,d,q}.
\]
% lean: CausalSmith.Stat.MarRareqLogfrontier.matched_minimax_frontier

\item We next construct an admissible decision kernel from the deterministic estimator in
\cref{obj:def:mixed-count-estimator}. For each observed sample \(o\), its defining finite average is
\[
\widehat\tau^{\mathrm{mix}}_{n,d,q}(o)
=
\sum_{\nu=0}^{n}
\Pr_\omega(K_0=\nu)\,
E_\omega\!\left[
\widetilde\tau(o;\omega)\mid K_0=\nu
\right].
\]
The omitted Poisson tail contributes zero by the estimator's definition.
Each conditional expectation is the uniform average over the \(3^\nu\) stream assignments.
Clipping, together with the zero-output branches, gives
\[
-1\le \widetilde\tau(o;\omega)\le1,
\qquad
-1\le
E_\omega\!\left[
\widetilde\tau(o;\omega)\mid K_0=\nu
\right]
\le1
\quad(0\le\nu\le n).
\]
The prefix weights satisfy
\[
\Pr_\omega(K_0=\nu)\ge0,\qquad
\sum_{\nu=0}^{n}\Pr_\omega(K_0=\nu)
=\Pr_\omega(K_0\le n)\le1.
\]
Consequently, at every observed sample,
\[
-1
\le-\Pr_\omega(K_0\le n)
\le\widehat\tau^{\mathrm{mix}}_{n,d,q}(o)
\le\Pr_\omega(K_0\le n)
\le1.
\]
The observed-sample space is finite with measurable singletons, so this deterministic map is measurable. Its Dirac kernel,
\[
\mathsf T_{\mathrm{mix}}(o,dh)
:=
\delta_{\widehat\tau^{\mathrm{mix}}_{n,d,q}(o)}(dh),
\]
therefore belongs to the decision class \(\mathcal T_n\) of
\cref{obj:def:minimax-risk}.
% lean: CausalSmith.Stat.MarRareqLogfrontier.mixed_count_estimator_regular

\item For every \(\mathsf T\in\mathcal T_n\) and \(P\in\mathcal M_{n,d,q}\), the joint squared-error risk is nonnegative:
\[
E_P\!\left[(\mathsf T-\tau(P))^2\right]
=
\int\!\int_{[-1,1]}
(h-\tau(P))^2\,\mathsf T(o,dh)\,
\mathcal L_P(\mathbf O_n)(do)
\ge0,
\]
because both integrals have nonnegative integrands. For the Dirac kernel just constructed,
\[
E_P\!\left[(\mathsf T_{\mathrm{mix}}-\tau(P))^2\right]
=
E_P\!\left[
\left(
\widehat\tau^{\mathrm{mix}}_{n,d,q}(\mathbf O_n)-\tau(P)
\right)^2
\right].
\]
Moreover,
\[
nq>0,\qquad
r_{n,d,q}
=
\min\left\{
1,\,
\frac1{nq}
+
\left[\frac{d}{nq\log(e+nq)}\right]^2
\right\}
\ge0.
\]
For each \(P\in\mathcal M_{n,d,q}\),
\cref{obj:thm:mixed-count-upper} now yields the uniform bound
\[
E_P\!\left[(\mathsf T_{\mathrm{mix}}-\tau(P))^2\right]
\le \mathfrak U_{n,d,q}
\le C_0\,r_{n,d,q}
\le\overline C\,r_{n,d,q}.
\]
The last inequality uses \(C_0\le\overline C\) and \(r_{n,d,q}\ge0\).
% lean: CausalSmith.Stat.MarRareqLogfrontier.mixed_count_upper

\item Nonnegativity of all the risks permits comparison of the minimax value with the worst-case risk of any admissible estimator. In particular,
\[
\mathfrak R_{n,d,q}
\le
\sup_{P\in\mathcal M_{n,d,q}}
E_P\!\left[(\mathsf T_{\mathrm{mix}}-\tau(P))^2\right].
\]
If \(\mathcal M_{n,d,q}\) is nonempty, the uniform bound from the preceding step gives
\[
\sup_{P\in\mathcal M_{n,d,q}}
E_P\!\left[(\mathsf T_{\mathrm{mix}}-\tau(P))^2\right]
\le\overline C\,r_{n,d,q}.
\]
If \(\mathcal M_{n,d,q}\) is empty, the real-supremum convention assigns its worst-case risk the value zero. Since \(\overline C\ge C_0>0\) and \(r_{n,d,q}\ge0\), in this case
\[
\sup_{P\in\mathcal M_{n,d,q}}
E_P\!\left[(\mathsf T_{\mathrm{mix}}-\tau(P))^2\right]
=0\le\overline C\,r_{n,d,q}.
\]
Thus, in either case,
\[
\underline c\,r_{n,d,q}
\le\mathfrak R_{n,d,q}
\le\overline C\,r_{n,d,q}.
\]
The constants were chosen independently of \(n,d,q\), as required.
% lean: CausalSmith.Stat.MarRareqLogfrontier.matched_minimax_frontier
\end{enumerate}
\end{proof}

\begin{proof}[Proof of \cref{obj:prop:fixed-d-effective-sample-reduction}]
Write
\[
N(t)=n(t)q(t)>0,
\qquad
\ell(t)=\log(e+N(t)).
\]
Define the restricted model class and its minimax risk by
\[
\mathcal M^{(0)}_{n,d,q}
:=
\left\{
P\in\mathcal M_{n,d,q}:
P(S(0)=0,S(1)=0)=1
\right\},
\qquad
\mathfrak R^{(0)}_{n,d,q}
:=
\inf_{\mathsf T\in\mathcal T_n}
\sup_{P\in\mathcal M^{(0)}_{n,d,q}}
E_P[(\mathsf T-\tau(P))^2].
\]
Here the estimator class and joint sample-and-kernel expectation are those of
\cref{obj:def:minimax-risk}.

\begin{enumerate}
\item Since \(N(t)\to\infty\), eventually
\[
N(t)\ge \max\{2,d^2\}.
\]
Also \(\ell(t)\ge1\), because \(e+N(t)\ge e\). Consequently, on this eventual range,
\[
\left[\frac{d}{N(t)\ell(t)}\right]^2
=
\frac{d^2}{N(t)^2\ell(t)^2}
\le
\frac{1}{N(t)\ell(t)^2}
\le
N(t)^{-1}.
\]
The expression in \cref{obj:def:frontier-rate} is nonnegative for every \(t\).
Eventually \(N(t)^{-1}\le1\), and both entries of its minimum are at least
\(N(t)^{-1}\). Thus
\[
N(t)^{-1}
\le
r_{n(t),d,q(t)}
=
\min\left\{
1,\,
N(t)^{-1}+
\left[\frac{d}{N(t)\ell(t)}\right]^2
\right\}
\le
2N(t)^{-1}.
\]
These inequalities prove the first order comparison.
% lean: CausalSmith.Stat.MarRareqLogfrontier.fixed_d_frontier_rate_theta

\item Let
\[
0<\underline c_{\mathrm{match}}
\le \overline C_{\mathrm{match}}<\infty
\]
be the constants supplied by \cref{obj:thm:matched-minimax-frontier}.
Its hypotheses hold at every \(t\): the sample and alphabet sizes are positive,
\(0<q(t)\le1\), and the rare-arrival restriction is assumed. Hence
\[
\underline c_{\mathrm{match}}\,r_{n(t),d,q(t)}
\le
\mathfrak R_{n(t),d,q(t)}
\le
\overline C_{\mathrm{match}}\,r_{n(t),d,q(t)}
\qquad(t\in\mathbb N).
\]
In particular, both quantities are nonnegative, and the lower inequality also gives
\[
r_{n(t),d,q(t)}
\le
\underline c_{\mathrm{match}}^{-1}
\mathfrak R_{n(t),d,q(t)}.
\]
Combining these bounds with the eventual frontier bounds from the preceding step yields
\[
\underline c_{\mathrm{match}}N(t)^{-1}
\le
\mathfrak R_{n(t),d,q(t)}
\le
2\overline C_{\mathrm{match}}N(t)^{-1}.
\]
% lean: CausalSmith.Stat.MarRareqLogfrontier.matched_minimax_frontier

\item We next obtain a pointwise upper bound for the restricted risk.
For each \(t\), use the estimator of \cref{obj:def:mixed-count-estimator}
with parameters \(n(t),d,q(t)\). Its auxiliary output is zero or a value clipped
to \([-1,1]\), so its conditional average satisfies
\[
-1
=
E_\omega[-1\mid\mathbf O_{n(t)}]
\le
\widehat\tau^{\mathrm{mix}}_{n(t),d,q(t)}(\mathbf O_{n(t)})
\le
E_\omega[1\mid\mathbf O_{n(t)}]
=
1.
\]
The average is measurable because it is a finite sum over the admissible prefix
lengths and assignments. Therefore the deterministic kernel
\[
\mathsf T^{\mathrm{mix}}_t(o,dh)
:=
\operatorname{Dirac}_{\widehat\tau^{\mathrm{mix}}_{n(t),d,q(t)}(o)}(dh)
\]
belongs to \(\mathcal T_{n(t)}\). Its joint risk equals its deterministic risk:
\[
E_P[(\mathsf T^{\mathrm{mix}}_t-\tau(P))^2]
=
E_P\!\left[
\left(
\widehat\tau^{\mathrm{mix}}_{n(t),d,q(t)}(\mathbf O_{n(t)})
-\tau(P)
\right)^2
\right].
\]
Let
\[
0<\overline C_{\mathrm{mix}}<\infty
\]
be supplied by \cref{obj:thm:mixed-count-upper}. For every
\(P\in\mathcal M^{(0)}_{n(t),d,q(t)}\), membership in the full model class
allows us to apply its risk-envelope bound:
\[
E_P[(\mathsf T^{\mathrm{mix}}_t-\tau(P))^2]
\le
\mathfrak U_{n(t),d,q(t)}
\le
\overline C_{\mathrm{mix}}\,r_{n(t),d,q(t)}.
\]
All squared risks are nonnegative, so the minimax value is bounded above by
the worst-case risk of this admissible kernel. If the restricted class is
nonempty, taking its supremum gives
\begin{equation}\label{eq:fixed-d-restricted-risk-upper}
\mathfrak R^{(0)}_{n(t),d,q(t)}
\le
\sup_{P\in\mathcal M^{(0)}_{n(t),d,q(t)}}
E_P[(\mathsf T^{\mathrm{mix}}_t-\tau(P))^2]
\le
\overline C_{\mathrm{mix}}\,r_{n(t),d,q(t)}.
\end{equation}
If the restricted class is empty, its real-valued worst-case risk is zero
under the empty-supremum convention. The same inequality follows from
\[
0\le \overline C_{\mathrm{mix}}\,r_{n(t),d,q(t)}.
\]
Thus the upper bound holds for every \(t\).
% lean: CausalSmith.Stat.MarRareqLogfrontier.mixed_count_upper
% lean: CausalSmith.Stat.MarRareqLogfrontier.fixed_d_effective_sample_reduction

\item To derive the restricted lower bound, fix \(t\). For
\(\mu_{\mathrm{test}}\in[1/4,3/4]\), define \(P_{\mu_{\mathrm{test}}}^{(t)}\) by
\[
X=0,\qquad S(0)=S(1)=Y(0)=0,
\qquad
A\sim\operatorname{Bernoulli}(1/2),\quad
Y(1)\sim\operatorname{Bernoulli}(\mu_{\mathrm{test}}),\quad
R\sim\operatorname{Bernoulli}(q(t)),
\]
where \(A,Y(1),R\) are mutually independent, and set
\(S=S(A)\), \(Y=Y(A)\). The category \(0\) exists because \(d\ge1\).
Take the observed sample independently from the induced observed-record law.

This construction satisfies randomization and consistency, and independence
of \(R\) from \((X,A,S,Y)\) implies the required conditional arrival independence.
Its occupied cells satisfy
\[
p_{a,0,0}(P_{\mu_{\mathrm{test}}}^{(t)})=\frac12,
\qquad
\rho_{a,0,0}(P_{\mu_{\mathrm{test}}}^{(t)})=q(t)
\qquad(a\in\{0,1\}),
\]
while every other cell has zero membership probability. Together with the
assumed parameter restrictions and rare-arrival slice, these facts establish
\[
P_{\mu_{\mathrm{test}}}^{(t)}\in\mathcal M^{(0)}_{n(t),d,q(t)},
\qquad
\tau(P_{\mu_{\mathrm{test}}}^{(t)})
=
E_{P_{\mu_{\mathrm{test}}}^{(t)}}[Y(1)-Y(0)]
=
\mu_{\mathrm{test}}.
\]

Define the separation and the two outcome means by
\[
\delta(t)
:=
\min\left\{\frac14,\frac{1}{4\sqrt{N(t)}}\right\},
\qquad
u_0(t):=\frac12,
\qquad
u_1(t):=\frac12+\delta(t).
\]
Since \(N(t)>0\),
\[
0<\delta(t)\le\frac14,
\qquad
u_0(t),u_1(t)\in[1/4,3/4].
\]
For \(i\in\{0,1\}\), define the single-record and sample laws by
\[
Q_i^{\mathrm{obs}}(t)
:=
\mathcal L_{P_{u_i(t)}^{(t)}}(O_1),
\qquad
Q_i^{\mathrm{sam}}(t)
:=
\mathcal L_{P_{u_i(t)}^{(t)}}(\mathbf O_{n(t)})
=
\bigl(Q_i^{\mathrm{obs}}(t)\bigr)^{\otimes n(t)}.
\]

For finite probability laws \(Q_1\ll Q_0\), use the divergence and
total-variation conventions
\[
\chi^2(Q_1\Vert Q_0)
:=
\sum_{o:\,Q_0(\{o\})>0}
\frac{\bigl(Q_1(\{o\})-Q_0(\{o\})\bigr)^2}{Q_0(\{o\})},
\qquad
\operatorname{TV}(Q_1,Q_0)
:=
\sup_{\mathcal A}|Q_1(\mathcal A)-Q_0(\mathcal A)|
=
\frac12\sum_o|Q_1(\{o\})-Q_0(\{o\})|.
\]
Under \(P_{\mu_{\mathrm{test}}}^{(t)}\), every observed record has \(X=S=0\). The remaining
coordinates, with arrived outcome \(RY=R\,Y\), have probabilities
\[
\begin{array}{c|ccccc}
(A,R,RY)
&(0,0,0)&(1,0,0)&(0,1,0)&(1,1,0)&(1,1,1)\\ \hline
\text{Probability}
&\dfrac{1-q(t)}2
&\dfrac{1-q(t)}2
&\dfrac{q(t)}2
&\dfrac{q(t)(1-\mu_{\mathrm{test}})}2
&\dfrac{q(t)\mu_{\mathrm{test}}}2
\end{array}
\]
and all other records have probability zero. At \(q(t)=1\), the first two
atoms have zero probability under both laws and are omitted from the
divergence sum. Thus \(Q_1^{\mathrm{obs}}(t)\ll Q_0^{\mathrm{obs}}(t)\), and summing
the atom probabilities gives
\[
\begin{aligned}
1+\chi^2\!\left(Q_1^{\mathrm{obs}}(t)\Vert Q_0^{\mathrm{obs}}(t)\right)
&=
1-q(t)+\frac{q(t)}2
+q(t)\bigl((1-u_1(t))^2+u_1(t)^2\bigr)\\
&=
1+2q(t)\delta(t)^2.
\end{aligned}
\]
The definition of \(\delta(t)\) gives
\[
16N(t)\delta(t)^2\le1,
\qquad
\chi^2\!\left(Q_1^{\mathrm{obs}}(t)\Vert Q_0^{\mathrm{obs}}(t)\right)
=
2q(t)\delta(t)^2
\le
\frac{1}{8n(t)}.
\]
Factoring the finite product sums and using \(1+x\le e^x\) now yields
\[
\begin{aligned}
1+\chi^2\!\left(Q_1^{\mathrm{sam}}(t)\Vert Q_0^{\mathrm{sam}}(t)\right)
&=
\left(
1+\chi^2\!\left(Q_1^{\mathrm{obs}}(t)\Vert Q_0^{\mathrm{obs}}(t)\right)
\right)^{n(t)}\\
&\le
\exp\!\left(
n(t)\chi^2\!\left(Q_1^{\mathrm{obs}}(t)\Vert Q_0^{\mathrm{obs}}(t)\right)
\right)
\le e^{1/8}.
\end{aligned}
\]
Since
\[
e^{1/2}e^{1/2}=e<3,
\qquad
e^{1/8}\le e^{1/2}<\sqrt3<2,
\]
the sample divergence is less than \(1\). Cauchy--Schwarz applied to the
finite total-variation sum therefore gives
\[
\operatorname{TV}\!\left(Q_1^{\mathrm{sam}}(t),Q_0^{\mathrm{sam}}(t)\right)
\le
\frac12
\sqrt{\chi^2\!\left(Q_1^{\mathrm{sam}}(t)\Vert Q_0^{\mathrm{sam}}(t)\right)}
\le\frac12.
\]
Consequently, every measurable sample event \(\mathcal A\) satisfies
\[
Q_1^{\mathrm{sam}}(t)(\mathcal A^c)
+
Q_0^{\mathrm{sam}}(t)(\mathcal A)
\ge
1-\operatorname{TV}\!\left(Q_1^{\mathrm{sam}}(t),Q_0^{\mathrm{sam}}(t)\right)
\ge\frac12.
\]
% lean: CausalSmith.Stat.MarRareqLogfrontier.oneCell_product_testing_half

\item Let \(\mathsf T\in\mathcal T_{n(t)}\) be arbitrary, and define its
barycenter and the relevant sample events by
\[
b_{\mathsf T}(o):=\int_{[-1,1]}h\,\mathsf T(o,dh),
\]
\[
\mathcal E_{\mathsf T,i}
:=
\left\{o:\frac{\delta(t)}2\le
|b_{\mathsf T}(o)-u_i(t)|\right\}
\quad(i\in\{0,1\}),
\qquad
\mathcal A_{\mathsf T}
:=
\left\{o:
|b_{\mathsf T}(o)-u_1(t)|
\le
|b_{\mathsf T}(o)-u_0(t)|
\right\}.
\]
These events are measurable. The triangle inequality gives, at every sample,
\[
\delta(t)
=
|u_1(t)-u_0(t)|
\le
|b_{\mathsf T}(o)-u_0(t)|
+
|b_{\mathsf T}(o)-u_1(t)|.
\]
It follows that
\[
\mathcal A_{\mathsf T}\subseteq\mathcal E_{\mathsf T,0},
\qquad
\mathcal A_{\mathsf T}^{c}\subseteq\mathcal E_{\mathsf T,1}.
\]
Applying the testing bound from the preceding step and then these inclusions,
\[
\frac12
\le
Q_1^{\mathrm{sam}}(t)(\mathcal A_{\mathsf T}^{c})
+
Q_0^{\mathrm{sam}}(t)(\mathcal A_{\mathsf T})
\le
Q_1^{\mathrm{sam}}(t)(\mathcal E_{\mathsf T,1})
+
Q_0^{\mathrm{sam}}(t)(\mathcal E_{\mathsf T,0}).
\]
Thus
\[
\frac14
\le
\max_{i\in\{0,1\}}
Q_i^{\mathrm{sam}}(t)(\mathcal E_{\mathsf T,i}).
\]
On \(\mathcal E_{\mathsf T,i}\), the barycenter's squared error is at least
\((\delta(t)/2)^2\). The barycenter inequality of
\cref{obj:lem:randomized-kernel-barycenter} therefore supplies
\[
\left(\frac{\delta(t)}2\right)^2
Q_i^{\mathrm{sam}}(t)(\mathcal E_{\mathsf T,i})
\le
\int
(b_{\mathsf T}(o)-u_i(t))^2\,Q_i^{\mathrm{sam}}(t)(do)
\le
E_{P_{u_i(t)}^{(t)}}[(\mathsf T-u_i(t))^2].
\]
Whether the probability for \(i=0\) is at most the probability for \(i=1\),
or conversely, multiplication by the nonnegative factor
\((\delta(t)/2)^2\) preserves their order. Hence
\[
\begin{aligned}
\frac{\delta(t)^2}{16}
&\le
\left(\frac{\delta(t)}2\right)^2
\max_{i\in\{0,1\}}
Q_i^{\mathrm{sam}}(t)(\mathcal E_{\mathsf T,i})\\
&=
\max_{i\in\{0,1\}}
\left\{
\left(\frac{\delta(t)}2\right)^2
Q_i^{\mathrm{sam}}(t)(\mathcal E_{\mathsf T,i})
\right\}\\
&\le
\max_{i\in\{0,1\}}
E_{P_{u_i(t)}^{(t)}}[(\mathsf T-u_i(t))^2].
\end{aligned}
\]
% lean: CausalSmith.Stat.MarRareqLogfrontier.complete_arrival_two_point_decision_reduction

\item The separation also has the required lower calibration.
If \(1/4\le1/(4\sqrt{N(t)})\), then
\[
\delta(t)^2=\frac1{16}
\ge
\frac1{16}\min\{1,N(t)^{-1}\}.
\]
In the other case,
\[
\delta(t)^2
=
\frac{1}{16N(t)}
\ge
\frac1{16}\min\{1,N(t)^{-1}\}.
\]
Combining either inequality with the preceding risk bound gives, for every
\(\mathsf T\in\mathcal T_{n(t)}\),
\[
\frac1{256}\min\{1,N(t)^{-1}\}
\le
\frac{\delta(t)^2}{16}
\le
\max_{i\in\{0,1\}}
E_{P_{u_i(t)}^{(t)}}[(\mathsf T-u_i(t))^2].
\]

The estimator class is nonempty: the constant kernel
\[
\mathsf T^{0}(o,dh):=\operatorname{Dirac}_{0}(dh)
\]
is admissible. Moreover, binary potential outcomes imply
\(\tau(P)\in[-1,1]\), and the kernel outputs lie in \([-1,1]\), so
\[
0\le E_P[(\mathsf T-\tau(P))^2]\le4
\]
for every full-data law \(P\) and admissible \(\mathsf T\).
Both testing laws belong to the full model and to its restricted subclass.
Their maximum risk therefore bounds each corresponding worst-case risk
from below; the uniform bound by \(4\) ensures these suprema are finite.
Taking the infimum over estimators gives
\[
\frac1{256}\min\{1,N(t)^{-1}\}
\le
\mathfrak R_{n(t),d,q(t)},
\qquad
\frac1{256}\min\{1,N(t)^{-1}\}
\le
\mathfrak R^{(0)}_{n(t),d,q(t)}.
\]
% lean: CausalSmith.Stat.MarRareqLogfrontier.oneCell_minimax_risk_floors

\item The last lower bound establishes nonnegativity of the restricted risk
for every \(t\). Its upper bound in \cref{eq:fixed-d-restricted-risk-upper}, combined with the eventual
frontier bound, gives
\[
0\le
\mathfrak R^{(0)}_{n(t),d,q(t)}
\le
2\overline C_{\mathrm{mix}}N(t)^{-1}
\qquad\text{eventually}.
\]
Also \(N(t)\to\infty\) implies \(N(t)\ge1\) eventually, so
\[
\min\{1,N(t)^{-1}\}=N(t)^{-1}.
\]
The lower bound consequently becomes
\[
\frac1{256}N(t)^{-1}
\le
\mathfrak R^{(0)}_{n(t),d,q(t)},
\qquad\text{equivalently}\qquad
N(t)^{-1}\le256\,\mathfrak R^{(0)}_{n(t),d,q(t)}.
\]
Together these inequalities prove the restricted-risk order comparison,
and hence all three assertions.
% lean: CausalSmith.Stat.MarRareqLogfrontier.fixed_d_effective_sample_reduction
\end{enumerate}
\end{proof}

\begin{proof}[Proof of \cref{obj:prop:phase-boundaries}]
\begin{enumerate}
\item\label{proof:phase-boundaries:comparison} For every \(t\in\mathbb N\), write
\[
\ell_t:=\log(e+N_t),\qquad
r_t:=r_{n_t,d_t,q_t}
=\min\left\{1,N_t^{-1}
+\left(\frac{d_t}{N_t\ell_t}\right)^2\right\},
\qquad
\mathfrak R_t:=\mathfrak R_{n_t,d_t,q_t}.
\]
The expression for \(r_t\) is supplied by \cref{obj:def:frontier-rate}. The size and arrival assumptions give
\[
N_t=n_tq_t>0,\qquad
N_t^{-1}+\left(\frac{d_t}{N_t\ell_t}\right)^2\ge0,
\qquad r_t\ge0.
\]
All hypotheses of \cref{obj:thm:matched-minimax-frontier} hold at every \(t\). Consequently, there are constants
\(0<\underline c\le\overline C<\infty\), independent of \(t\), such that
\[
\underline c\,r_t\le\mathfrak R_t\le\overline C\,r_t.
\]
In particular,
\[
0\le r_t\le\underline c^{-1}\mathfrak R_t,
\qquad
0\le\mathfrak R_t\le\overline C\,r_t.
\]
The first inequality and the squeeze theorem show that
\(\mathfrak R_t\to0\) implies \(r_t\to0\); the second gives the converse. Thus
\[
\mathfrak R_t\longrightarrow0
\quad\Longleftrightarrow\quad
r_t\longrightarrow0.
\]
% lean: CausalSmith.Stat.MarRareqLogfrontier.phase_boundaries

\item\label{proof:phase-boundaries:convergence} Define, for every \(t\in\mathbb N\),
\[
\nu_t:=N_t^{-1},
\qquad
\xi_t:=\frac{d_t}{N_t\ell_t}.
\]
Since \(e>1\) and \(N_t>0\),
\[
e+N_t>1,\qquad
\ell_t>0,\qquad
N_t\ell_t>0,\qquad
\nu_t\ge0,\qquad \xi_t\ge0.
\]
Effective sample growth gives
\[
\nu_t\longrightarrow0.
\]
Continuity of the square root, together with
\[
\sqrt{\xi_t^2}=|\xi_t|=\xi_t,
\]
shows that \(\xi_t^2\to0\) implies \(\xi_t\to0\). Continuity of squaring gives the reverse implication. Hence
\[
\xi_t^2\longrightarrow0
\quad\Longleftrightarrow\quad
\xi_t\longrightarrow0.
\]
Because \(N_t\ell_t>0\) for every \(t\), the quotient characterization of little \(o\) gives
\[
d_t=o(N_t\ell_t)
\quad\Longleftrightarrow\quad
\xi_t\longrightarrow0.
\]

We next establish the truncation equivalence
\[
r_t=\min\{1,\nu_t+\xi_t^2\}\longrightarrow0
\quad\Longleftrightarrow\quad
\xi_t^2\longrightarrow0.
\]
For the forward direction, \(r_t\to0\) implies \(r_t<1\) eventually. At each such index, \(\nu_t+\xi_t^2\ge1\) would give \(r_t=1\). Therefore, eventually,
\[
\nu_t+\xi_t^2<1,\qquad
r_t=\nu_t+\xi_t^2.
\]
The sum consequently tends to zero, and
\[
0\le\xi_t^2\le\nu_t+\xi_t^2
\]
implies \(\xi_t^2\to0\). Conversely, if \(\xi_t^2\to0\), then
\[
\nu_t+\xi_t^2\longrightarrow0,
\qquad
\min\{1,\nu_t+\xi_t^2\}\longrightarrow\min\{1,0\}=0,
\]
by continuity of addition and minimum. Combining these equivalences with \cref{proof:phase-boundaries:comparison} yields
\[
\mathfrak R_t\longrightarrow0
\quad\Longleftrightarrow\quad
d_t=o(N_t\ell_t)
=o\!\left(N_t\log(e+N_t)\right).
\]
% lean: CausalSmith.Stat.MarRareqLogfrontier.phase_min_rate_tendsto
% lean: CausalSmith.Stat.MarRareqLogfrontier.phase_frontier_rate_tendsto

\item\label{proof:phase-boundaries:theta-comparison} The comparison in \cref{proof:phase-boundaries:comparison} also gives a two-sided asymptotic comparison. Indeed, nonnegativity allows us to write, for every \(t\),
\[
|\mathfrak R_t|
\le\overline C\,|r_t|,
\qquad
|r_t|
\le\underline c^{-1}|\mathfrak R_t|.
\]
These are the defining bounds for
\[
\mathfrak R_t=O(r_t),
\qquad
r_t=O(\mathfrak R_t),
\qquad\text{and hence}\qquad
\mathfrak R_t=\Theta(r_t).
\]
% lean: CausalSmith.Stat.MarRareqLogfrontier.phase_boundaries

\item\label{proof:phase-boundaries:baseline} Since \(N_t\to\infty\), eventually \(N_t\ge1\). At those indices,
\[
0\le\nu_t\le1,
\qquad
\nu_t\le\nu_t+\xi_t^2.
\]
Thus both entries in the minimum defining \(r_t\) are at least \(\nu_t\), and
\[
0\le\nu_t\le\min\{1,\nu_t+\xi_t^2\}=r_t.
\]
This proves the baseline bound
\[
\nu_t=O(r_t).
\]
% lean: CausalSmith.Stat.MarRareqLogfrontier.phase_rate_lower_bigO

\item\label{proof:phase-boundaries:dimension-necessity} Define the dimension scale, for every \(t\in\mathbb N\), by
\[
B_{\mathrm{phase},t}:=\sqrt{N_t}\,\ell_t.
\]
The positivity established in \cref{proof:phase-boundaries:convergence} gives
\[
B_{\mathrm{phase},t}>0,
\qquad
N_t\ell_t>0.
\]
Moreover, using \((\sqrt{N_t})^2=N_t\),
\[
\nu_t(N_t\ell_t)^2
=N_t^{-1}N_t^2\ell_t^2
=N_t\ell_t^2
=B_{\mathrm{phase},t}^2.
\]

Suppose first that \(r_t=O(\nu_t)\). Choose a real constant
\(\beta_{\mathrm{rate}}\) such that, eventually,
\[
0\le r_t\le\beta_{\mathrm{rate}}\nu_t.
\]
Since \(\beta_{\mathrm{rate}}\nu_t\to0\), the squeeze theorem gives \(r_t\to0\), and hence \(r_t<1\) eventually. As in \cref{proof:phase-boundaries:convergence}, this forces
\[
\nu_t+\xi_t^2<1,
\qquad
r_t=\nu_t+\xi_t^2
\]
eventually. Intersecting these eventual conditions with the assumed bound gives
\[
\nu_t+\xi_t^2\le\beta_{\mathrm{rate}}\nu_t,
\qquad
\xi_t^2\le\beta_{\mathrm{rate}}\nu_t,
\]
where the latter inequality uses \(\nu_t\ge0\). Multiplication by the positive quantity \((N_t\ell_t)^2\) then yields
\[
d_t^2
=\xi_t^2(N_t\ell_t)^2
\le\beta_{\mathrm{rate}}\nu_t(N_t\ell_t)^2
=\beta_{\mathrm{rate}}B_{\mathrm{phase},t}^2.
\]
Set
\[
\beta_{\mathrm{bound}}:=\max\{1,\beta_{\mathrm{rate}}\}\in\mathbb R.
\]
Then
\[
\beta_{\mathrm{bound}}\ge1,\qquad
\beta_{\mathrm{rate}}\le\beta_{\mathrm{bound}}
\le\beta_{\mathrm{bound}}^2.
\]
Consequently, eventually,
\[
d_t^2
\le\beta_{\mathrm{bound}}^2B_{\mathrm{phase},t}^2.
\]
Both \(d_t\) and \(\beta_{\mathrm{bound}}B_{\mathrm{phase},t}\) are nonnegative, so this squared inequality gives
\[
d_t\le\beta_{\mathrm{bound}}B_{\mathrm{phase},t}.
\]
Therefore
\[
r_t=O(\nu_t)
\quad\Longrightarrow\quad
d_t=O(B_{\mathrm{phase},t}).
\]
% lean: CausalSmith.Stat.MarRareqLogfrontier.phase_dimension_rate_bigO

\item\label{proof:phase-boundaries:dimension-sufficiency} Conversely, suppose that \(d_t=O(B_{\mathrm{phase},t})\). Choose a real constant
\(\beta_{\mathrm{dim}}\) such that, eventually,
\[
0\le d_t\le\beta_{\mathrm{dim}}B_{\mathrm{phase},t}.
\]
The right-hand side is consequently nonnegative, so squaring gives
\[
d_t^2\le\beta_{\mathrm{dim}}^2B_{\mathrm{phase},t}^2.
\]
Dividing by \((N_t\ell_t)^2>0\) and using the identity from \cref{proof:phase-boundaries:dimension-necessity} yields
\[
\xi_t^2
=\frac{d_t^2}{(N_t\ell_t)^2}
\le
\beta_{\mathrm{dim}}^2
\frac{B_{\mathrm{phase},t}^2}{(N_t\ell_t)^2}
=\beta_{\mathrm{dim}}^2\nu_t.
\]
It follows that, eventually,
\[
0\le r_t
=\min\{1,\nu_t+\xi_t^2\}
\le\nu_t+\xi_t^2
\le(1+\beta_{\mathrm{dim}}^2)\nu_t.
\]
Thus \(r_t=O(\nu_t)\). Together with \cref{proof:phase-boundaries:dimension-necessity}, this establishes
\[
r_t=O(\nu_t)
\quad\Longleftrightarrow\quad
d_t=O(B_{\mathrm{phase},t})
=O\!\left(\sqrt{N_t}\log(e+N_t)\right).
\]
% lean: CausalSmith.Stat.MarRareqLogfrontier.phase_dimension_rate_bigO

\item If \(\mathfrak R_t=\Theta(\nu_t)\), then
\[
r_t=O(\mathfrak R_t),
\qquad
\mathfrak R_t=O(\nu_t)
\quad\Longrightarrow\quad
r_t=O(\nu_t),
\]
by transitivity of big \(O\). The equivalence in \cref{proof:phase-boundaries:dimension-sufficiency} therefore gives
\(d_t=O(B_{\mathrm{phase},t})\).

Conversely, if \(d_t=O(B_{\mathrm{phase},t})\), then \cref{proof:phase-boundaries:dimension-sufficiency} gives \(r_t=O(\nu_t)\). Combining this with \cref{proof:phase-boundaries:theta-comparison,proof:phase-boundaries:baseline} gives
\[
\mathfrak R_t=O(r_t),\quad r_t=O(\nu_t)
\quad\Longrightarrow\quad
\mathfrak R_t=O(\nu_t),
\]
and
\[
\nu_t=O(r_t),\quad r_t=O(\mathfrak R_t)
\quad\Longrightarrow\quad
\nu_t=O(\mathfrak R_t).
\]
These two bounds prove
\[
\mathfrak R_t=\Theta(N_t^{-1})
\quad\Longleftrightarrow\quad
d_t=O\!\left(\sqrt{N_t}\log(e+N_t)\right),
\]
which is the second claimed equivalence.
% lean: CausalSmith.Stat.MarRareqLogfrontier.phase_boundaries
\end{enumerate}
\end{proof}

\begin{proof}[Proof of \cref{obj:thm:unrestricted-fixed-q-minimax-frontier}]
Fix \(0<q<1/2\). For \(n,d\ge1\), put
\[
\mathfrak h_{n,d}^{\circ}
:=
\frac1n+\left(\frac{d}{n\log(e+n)}\right)^2,
\qquad
\mathfrak h_{n,d}:=\min\{1,\mathfrak h_{n,d}^{\circ}\}.
\]
For a bounded estimator kernel \(\mathsf T\) and a full-data law \(P\), write
\[
\mathcal R_P(\mathsf T)
:=
\int\!\int_{[-1,1]}
(h-\tau(P))^2\,\mathsf T(o,dh)\,
\mathcal L_P(\mathbf O_n)(do).
\]
Thus \(\mathcal R_P(\mathsf T)\) integrates both the sample and the estimator draw.

\begin{enumerate}
\item
The deterministic estimator in \cref{obj:def:mixed-count-estimator} defines an admissible decision kernel
\[
\mathsf T_{\mathrm{mix}}(o,dh)
:=
\delta_{\widehat\tau^{\mathrm{mix}}_{n,d,q}(o)}(dh).
\]
By \cref{obj:thm:unrestricted-mixed-count-upper}, there is a universal constant \(\overline C_{\mathrm{mix}}>0\) such that
\[
0\le \mathcal R_P(\mathsf T_{\mathrm{mix}})
\le \overline C_{\mathrm{mix}}\,r_{n,d,q}
\qquad
\text{for every }P\in\mathcal M^+_{n,d,q}.
\]
Taking the supremum over the model and then using this kernel in the defining infimum gives
\[
\mathfrak R^+_{n,d,q}
\le
\sup_{P\in\mathcal M^+_{n,d,q}}
\mathcal R_P(\mathsf T_{\mathrm{mix}})
\le
\overline C_{\mathrm{mix}}\,r_{n,d,q}.
\]
If the model class is empty, its real-valued worst-case risk is taken to be zero, and the same inequality follows from
\[
r_{n,d,q}
=
\min\left\{1,\frac1{nq}
+\left(\frac{d}{nq\log(e+nq)}\right)^2\right\}
\ge0.
\]
% lean: CausalSmith.Stat.MarRareqLogfrontier.unrestricted_fixed_q_frontier_upper

\item
We next convert this upper bound to the fixed-\(q\) rate. Define
\[
\beta_q:=\frac{1-\log q}{q},
\qquad
N:=nq,
\qquad
\ell:=\log(e+N),
\qquad
\ell_{n,1}:=\log(e+n).
\]
Throughout the present domain,
\[
N>0,\qquad \ell\ge1,\qquad \ell_{n,1}>0,
\qquad \log q\le0,\qquad \beta_q\ge1.
\]
Since
\[
e+n\le\frac{e+N}{q},
\]
monotonicity of the logarithm yields
\[
\ell_{n,1}\le\ell-\log q
\le(1-\log q)\ell.
\]
The second inequality follows from
\[
(-\log q)(\ell-1)\ge0.
\]
Consequently,
\[
n\ell_{n,1}
\le n(1-\log q)\ell
=\beta_q N\ell.
\]
Moreover,
\[
q\beta_q=1-\log q\ge1,
\qquad
q\beta_q^2=(1-\log q)\beta_q\ge1,
\]
so
\[
\frac1N\le\frac{\beta_q^2}{n}.
\]
The positive denominator comparison also gives
\[
0\le\frac{d}{N\ell}
\le\beta_q\frac{d}{n\ell_{n,1}},
\qquad
\left(\frac{d}{N\ell}\right)^2
\le\beta_q^2
\left(\frac{d}{n\ell_{n,1}}\right)^2.
\]
Adding these bounds,
\[
\frac1N+\left(\frac{d}{N\ell}\right)^2
\le\beta_q^2\mathfrak h_{n,d}^{\circ}.
\]
Truncation therefore gives
\[
r_{n,d,q}
\le\min\{1,\beta_q^2\mathfrak h_{n,d}^{\circ}\}
\le\beta_q^2\mathfrak h_{n,d}.
\]
Indeed, if \(\mathfrak h_{n,d}^{\circ}\le1\), the last inequality is
\[
\min\{1,\beta_q^2\mathfrak h_{n,d}^{\circ}\}
\le\beta_q^2\mathfrak h_{n,d}^{\circ}
=\beta_q^2\mathfrak h_{n,d};
\]
if \(\mathfrak h_{n,d}^{\circ}\ge1\), it is
\[
\min\{1,\beta_q^2\mathfrak h_{n,d}^{\circ}\}
\le1\le\beta_q^2
=\beta_q^2\mathfrak h_{n,d}.
\]
Define the positive fixed-\(q\) constant
\[
\overline C_{\mathrm{poly}}(q)
:=\overline C_{\mathrm{mix}}\beta_q^2.
\]
The upper bound from the preceding step now implies
\[
\mathfrak R^+_{n,d,q}
\le\overline C_{\mathrm{mix}}r_{n,d,q}
\le\overline C_{\mathrm{poly}}(q)\mathfrak h_{n,d}.
\]
% lean: CausalSmith.Stat.MarRareqLogfrontier.frontierRate_le_fixedQRate_mul; CausalSmith.Stat.MarRareqLogfrontier.unrestricted_fixed_q_polynomial_upper

\item
For the lower bound, invoke the published fixed-positivity lower bound on the control-zero observational subclass
\citep[Theorem~2, Section~3.2, p.~9, with its proof in Section~C.5, pp.~39--44]{ZengBalakrishnanHanKennedy2026}.
Its truncated, all-\(n,d\) comparison is supplied by
\cref{obj:thm:zeng-fixed-positivity-polynomial-frontier}: at \(\epsilon=q\), there is a constant \(\underline c_Z(q)>0\) such that
\[
\underline c_Z(q)\mathfrak h_{n,d}
\le\mathfrak R^Z_{n,d,q}
\qquad(n,d\ge1).
\]

We transfer this lower bound to the unrestricted experiment. For an observational estimator kernel \(\widehat\theta\), define
\[
\mathcal R^Z_{P_Z}(\widehat\theta)
:=
\int\!\int_{[-1,1]}
(h-\theta_Z(P_Z))^2\,
\widehat\theta(\mathbf z,dh)\,
\mathcal L_{P_Z}(\mathbf Z_n)(d\mathbf z).
\]
Given any admissible unrestricted kernel \(\mathsf T\), define its observational pullback by
\[
\widehat\theta_{\mathsf T}(\mathbf z,\mathcal V)
:=
2^{-n}
\sum_{\mathbf u\in\{0,1\}^n}
\mathsf T\!\left(
\bigl(o(x_i,b_i,y_i;u_i)\bigr)_{i=1}^n,\mathcal V
\right),
\]
where
\[
\mathbf z=((x_i,b_i,y_i))_{i=1}^n,
\qquad
\mathbf u=(u_i)_{i=1}^n,
\qquad
\mathcal V\in\mathscr B([-1,1]),
\]
and \(o(x,b,y;u)\) is the synthetic observation map of
\cref{obj:lem:synthetic-randomization-kernel}. This finite mixture composes the parameter-independent fair-coin transformation with \(\mathsf T\), and hence is an admissible observational kernel.

For every \(P_Z\in\mathcal D_{d,q}\),
\cref{obj:lem:synthetic-randomization-kernel} supplies a completion
\(P_Z^*\in\mathcal M^+_{n,d,q}\) whose observed product law is exactly the transformed sample law and whose target satisfies
\[
\tau(P_Z^*)=\theta_Z(P_Z).
\]
Thus
\[
\mathcal R^Z_{P_Z}(\widehat\theta_{\mathsf T})
=
\mathcal R_{P_Z^*}(\mathsf T).
\]
Both targets and all kernel outputs lie in \([-1,1]\), giving
\[
0\le\mathcal R^Z_{P_Z}(\widehat\theta_{\mathsf T})\le4,
\qquad
0\le\mathcal R_P(\mathsf T)\le4.
\]
The supremum comparison and the observational minimax definition therefore give, for every \(\mathsf T\),
\[
\mathfrak R^Z_{n,d,q}
\le
\sup_{P_Z\in\mathcal D_{d,q}}
\mathcal R^Z_{P_Z}(\widehat\theta_{\mathsf T})
\le
\sup_{P\in\mathcal M^+_{n,d,q}}
\mathcal R_P(\mathsf T).
\]
Taking the infimum over \(\mathsf T\) yields
\[
\mathfrak R^Z_{n,d,q}\le\mathfrak R^+_{n,d,q}.
\]

Now define
\[
\underline c(q):=\frac{\underline c_Z(q)}{\beta_q^2}>0.
\]
Using the rate comparison established in the preceding step,
\[
\underline c(q)r_{n,d,q}
\le
\frac{\underline c_Z(q)}{\beta_q^2}
\,\beta_q^2\mathfrak h_{n,d}
=
\underline c_Z(q)\mathfrak h_{n,d}
\le
\mathfrak R^Z_{n,d,q}
\le
\mathfrak R^+_{n,d,q}.
\]
% lean: CausalSmith.Stat.MarRareqLogfrontier.zeng_fixed_positivity_polynomial_frontier; CausalSmith.Stat.MarRareqLogfrontier.zengMinimaxRisk_le_unrestrictedMinimaxRisk; CausalSmith.Stat.MarRareqLogfrontier.unrestricted_fixed_q_minimax_frontier

\item
Choose the common upper constant
\[
\overline C(q)
:=
\max\{\overline C_{\mathrm{mix}},
       \overline C_{\mathrm{poly}}(q)\}
+\underline c(q).
\]
These are finite real constants depending only on \(q\), and
\[
0<\underline c(q)\le\overline C(q).
\]
The rates are nonnegative:
\[
\mathfrak h_{n,d}\ge0,
\qquad r_{n,d,q}\ge0,
\]
and the chosen constant satisfies
\[
\overline C_{\mathrm{mix}}\le\overline C(q),
\qquad
\overline C_{\mathrm{poly}}(q)\le\overline C(q).
\]

To obtain the fixed-rate lower bound with the same constant, observe that
\[
0<N\le n,
\qquad
0<\ell\le\ell_{n,1},
\qquad
0<N\ell\le n\ell_{n,1}.
\]
It follows that
\[
\frac1n\le\frac1N,
\qquad
0\le\frac{d}{n\ell_{n,1}}\le\frac{d}{N\ell},
\]
and hence
\[
\left(\frac{d}{n\ell_{n,1}}\right)^2
\le
\left(\frac{d}{N\ell}\right)^2.
\]
Adding and truncating at one proves
\[
\mathfrak h_{n,d}
=
\min\left\{1,\frac1n+
\left(\frac{d}{n\ell_{n,1}}\right)^2\right\}
\le
\min\left\{1,\frac1N+
\left(\frac{d}{N\ell}\right)^2\right\}
=r_{n,d,q}.
\]
Combining this comparison with the preceding bounds gives
\[
\begin{aligned}
\underline c(q)\mathfrak h_{n,d}
&\le \underline c(q)r_{n,d,q}
 \le\mathfrak R^+_{n,d,q},\\
\mathfrak R^+_{n,d,q}
&\le\overline C_{\mathrm{poly}}(q)\mathfrak h_{n,d}
 \le\overline C(q)\mathfrak h_{n,d},\\
\underline c(q)r_{n,d,q}
&\le\mathfrak R^+_{n,d,q},\\
\mathfrak R^+_{n,d,q}
&\le\overline C_{\mathrm{mix}}r_{n,d,q}
 \le\overline C(q)r_{n,d,q}.
\end{aligned}
\]
These are the two asserted comparisons.
% lean: CausalSmith.Stat.MarRareqLogfrontier.fixedQRate_le_frontierRate; CausalSmith.Stat.MarRareqLogfrontier.unrestricted_fixed_q_minimax_frontier
\end{enumerate}
\end{proof}

\begin{proof}[Proof of \cref{obj:prop:unrestricted-fixed-q-phase-boundaries}]
Fix \(q\), the arrival floor, with \(0<q<1/2\). Throughout the proof, \(n\) ranges over positive integers, and all limits are taken as \(n\to\infty\).

\begin{enumerate}
\item For a sequence of integers \(d_n\ge1\), define
\[
\ell_n^{\mathrm{fix}}:=\log(e+n),
\qquad
r_n^{\mathrm{fix}}
:=
\min\left\{1,\frac1n+
\left(\frac{d_n}{n\ell_n^{\mathrm{fix}}}\right)^2\right\}.
\]
Apply \cref{obj:thm:unrestricted-fixed-q-minimax-frontier}, whose lower comparison uses the published zero-control-regression lower bound
\citep[Theorem~2, Section~3.2, p.~9, and its proof in Section~C.5, pp.~39--44]{ZengBalakrishnanHanKennedy2026}.
It supplies constants \(0<\underline c(q)\le\overline C(q)\), depending only on \(q\), such that
\[
\underline c(q)\,r_n^{\mathrm{fix}}
\le \mathfrak R^{+}_{n,d_n,q}
\le \overline C(q)\,r_n^{\mathrm{fix}}
\qquad(n\ge1).
\]
These constants are common to every dimension sequence. Moreover,
\[
0\le r_n^{\mathrm{fix}},
\qquad
0\le\overline C(q).
\]
% lean: CausalSmith.Stat.MarRareqLogfrontier.unrestricted_fixed_q_minimax_frontier

\item We first establish the consistency equivalence. The comparison above gives
\[
0\le r_n^{\mathrm{fix}}
\le \underline c(q)^{-1}\mathfrak R^{+}_{n,d_n,q},
\qquad
0\le\mathfrak R^{+}_{n,d_n,q}
\le\overline C(q)\,r_n^{\mathrm{fix}}.
\]
Applying the squeeze theorem in each direction yields
\[
\mathfrak R^{+}_{n,d_n,q}\longrightarrow0
\quad\Longleftrightarrow\quad
r_n^{\mathrm{fix}}\longrightarrow0.
\]

To characterize the latter limit, define
\[
\nu_n:=\frac1n,
\qquad
\xi_n:=\frac{d_n}{n\ell_n^{\mathrm{fix}}}.
\]
Since \(e+n>1\), we have
\[
\ell_n^{\mathrm{fix}}>0,
\qquad
n\ell_n^{\mathrm{fix}}>0,
\qquad
\xi_n\ge0,
\qquad
\nu_n\ge0,
\qquad
\nu_n\longrightarrow0.
\]
The last limit follows from \(n\to\infty\) and continuity of inversion at infinity. Continuity of the square root and nonnegativity of \(\xi_n\) give
\[
\xi_n^2\longrightarrow0
\quad\Longrightarrow\quad
\xi_n=\sqrt{\xi_n^2}\longrightarrow0.
\]
Conversely, continuity of squaring gives
\[
\xi_n\longrightarrow0
\quad\Longrightarrow\quad
\xi_n^2\longrightarrow0.
\]
The positive denominator also gives the usual quotient characterization of little-\(o\):
\[
d_n=o\!\left(n\ell_n^{\mathrm{fix}}\right)
\quad\Longleftrightarrow\quad
\xi_n\longrightarrow0.
\]

It remains to handle the truncation. If
\(r_n^{\mathrm{fix}}\to0\), then eventually
\[
r_n^{\mathrm{fix}}=\min\{1,\nu_n+\xi_n^2\}<1.
\]
At every such index,
\[
\nu_n+\xi_n^2<1,
\qquad
r_n^{\mathrm{fix}}=\nu_n+\xi_n^2:
\]
indeed, \(\nu_n+\xi_n^2\ge1\) would make the minimum equal to \(1\).
Thus \(\nu_n+\xi_n^2\to0\), and
\[
0\le\xi_n^2\le\nu_n+\xi_n^2
\]
implies \(\xi_n^2\to0\). Conversely, if \(\xi_n^2\to0\), then
\[
\nu_n+\xi_n^2\longrightarrow0,
\qquad
\min\{1,\nu_n+\xi_n^2\}\longrightarrow\min\{1,0\}=0,
\]
by continuity of addition and the minimum.
% lean: CausalSmith.Stat.MarRareqLogfrontier.phase_min_rate_tendsto

Consequently,
\[
r_n^{\mathrm{fix}}\longrightarrow0
\quad\Longleftrightarrow\quad
\xi_n^2\longrightarrow0
\quad\Longleftrightarrow\quad
\xi_n\longrightarrow0
\quad\Longleftrightarrow\quad
d_n=o\!\left(n\log(e+n)\right).
\]
Together with the risk comparison, this proves the first assertion.
% lean: CausalSmith.Stat.MarRareqLogfrontier.fixedQRate_tendsto_iff_dimension_littleO

\item For the sufficient condition for the parametric rate, choose
\[
c:=\underline c(q)>0.
\]
Given \(M>0\), define
\[
C:=\overline C(q)(1+M^2).
\]
These constants satisfy
\[
c=\underline c(q)
\le\overline C(q)
\le\overline C(q)(1+M^2)=C.
\]
Suppose that eventually
\[
d_n\le M\sqrt n\,\ell_n^{\mathrm{fix}}.
\]
Both sides are nonnegative, so squaring this inequality and using
\((\sqrt n)^2=n\) gives
\[
d_n^2
\le\left(M\sqrt n\,\ell_n^{\mathrm{fix}}\right)^2
=M^2n\left(\ell_n^{\mathrm{fix}}\right)^2.
\]
Dividing by the positive squared denominator yields
\[
\xi_n^2
=\frac{d_n^2}{(n\ell_n^{\mathrm{fix}})^2}
\le\frac{M^2}{n}.
\]
Here the cancellation uses
\[
\frac1n(n\ell_n^{\mathrm{fix}})^2
=n\left(\ell_n^{\mathrm{fix}}\right)^2.
\]
Since \(1/n\le1\) for \(n\ge1\), both entries of the minimum defining the rate are at least \(1/n\). Hence
\[
\frac1n
\le r_n^{\mathrm{fix}}
=\min\left\{1,\frac1n+\xi_n^2\right\}
\le\frac1n+\xi_n^2
\le\frac{1+M^2}{n}.
\]
Multiplying the lower and upper bounds by the corresponding frontier constants gives, eventually,
\[
\frac{c}{n}
\le\underline c(q)\,r_n^{\mathrm{fix}}
\le\mathfrak R^{+}_{n,d_n,q}
\le\overline C(q)\,r_n^{\mathrm{fix}}
\le\frac{C}{n}.
\]
The choice of \(c\) depends only on \(q\), and the choice of \(C\) depends only on \(q\) and \(M\). The eventual dimension bound determines the threshold along each sequence.
% lean: CausalSmith.Stat.MarRareqLogfrontier.unrestricted_fixed_q_phase_boundaries

\item For the converse, suppose there exist constants \(0<c_*\le C_*\) such that eventually
\[
\frac{c_*}{n}
\le\mathfrak R^{+}_{n,d_n,q}
\le\frac{C_*}{n}.
\]
In particular, \(C_*\ge0\). Combining the upper risk bound with the lower frontier comparison gives
\[
\underline c(q)\,r_n^{\mathrm{fix}}
\le\mathfrak R^{+}_{n,d_n,q}
\le\frac{C_*}{n},
\]
and therefore
\[
0\le r_n^{\mathrm{fix}}
\le\frac{C_*}{\underline c(q)}\,\frac1n
\longrightarrow0.
\]
Thus eventually \(r_n^{\mathrm{fix}}<1\).

Define the dimension-bound constant
\[
K_{\mathrm{dim}}
:=\max\left\{1,\frac{C_*}{\underline c(q)}\right\}.
\]
Its calibration gives
\[
0\le K_{\mathrm{dim}},
\qquad
1\le K_{\mathrm{dim}},
\qquad
\frac{C_*}{\underline c(q)}
\le K_{\mathrm{dim}}
\le K_{\mathrm{dim}}^2.
\]
At all sufficiently large indices, \(r_n^{\mathrm{fix}}<1\) implies
\[
\frac1n+\xi_n^2<1,
\qquad
r_n^{\mathrm{fix}}=\frac1n+\xi_n^2,
\]
because a sum at least \(1\) would make the truncated rate equal to \(1\).
The preceding rate bound and \(1/n\ge0\) now yield
\[
\xi_n^2
\le\frac1n+\xi_n^2
=r_n^{\mathrm{fix}}
\le\frac{C_*}{\underline c(q)}\,\frac1n.
\]
Multiplying by \((n\ell_n^{\mathrm{fix}})^2>0\) gives
\[
d_n^2
\le\frac{C_*}{\underline c(q)}
\,n\left(\ell_n^{\mathrm{fix}}\right)^2
\le K_{\mathrm{dim}}^2
\,n\left(\ell_n^{\mathrm{fix}}\right)^2.
\]
The second inequality uses
\(n(\ell_n^{\mathrm{fix}})^2\ge0\). Also,
\[
K_{\mathrm{dim}}\sqrt n\,\ell_n^{\mathrm{fix}}\ge0,
\qquad
\left(K_{\mathrm{dim}}\sqrt n\,\ell_n^{\mathrm{fix}}\right)^2
=K_{\mathrm{dim}}^2n\left(\ell_n^{\mathrm{fix}}\right)^2.
\]
Since \(d_n\ge0\), comparison of these squares implies
\[
d_n\le K_{\mathrm{dim}}\sqrt n\,\ell_n^{\mathrm{fix}}
\qquad\text{eventually}.
\]
Both the dimension and the comparison scale are nonnegative, so this is precisely the required bound
\[
d_n=O\!\left(\sqrt n\,\log(e+n)\right).
\]
% lean: CausalSmith.Stat.MarRareqLogfrontier.unrestricted_fixed_q_phase_boundaries
\end{enumerate}
\end{proof}

\begin{proof}[Proof of \cref{obj:thm:unrestricted-mixed-count-upper}]
\begin{enumerate}
\item By \cref{obj:thm:unrestricted-uniform-mixed-count-certificate}, choose a universal constant \(0<\overline C<\infty\). For every \(n,d,q,P\) satisfying the stated domain, model, and sampling conditions, this constant gives
\[
E_P\!\left[
 \bigl(\widehat\tau^{\mathrm{mix}}_{n,d,q}(\mathbf O_n)-\Phi(P)\bigr)^2
\right]
\le \mathfrak U_{n,d,q}
\le \overline C\,r_{n,d,q},
\]
where \(\Phi(P)\) is the signed cell functional of
\cref{obj:def:cell-functional}. The envelope here has precisely the branch formulas in the statement.
% lean: CausalSmith.Stat.MarRareqLogfrontier.unrestricted_uniform_mixed_count_certificate

\item The model hypothesis permits application of
\cref{obj:lem:unrestricted-cell-identification}, which supplies
\[
\tau(P)=\Phi(P).
\]
Substituting this identity into the preceding risk bound yields
\[
E_P\!\left[
 \bigl(\widehat\tau^{\mathrm{mix}}_{n,d,q}(\mathbf O_n)-\tau(P)\bigr)^2
\right]
\le \mathfrak U_{n,d,q}
\le \overline C\,r_{n,d,q}.
\]
% lean: CausalSmith.Stat.MarRareqLogfrontier.unrestricted_cell_identification

\item To establish the auxiliary-risk comparison, fix an observed sample
\[
\mathbf o=(o_i)_{i=1}^n
\in
\bigl(\{0,\ldots,d-1\}\times\{0,1\}^{4}\bigr)^n,
\qquad
o_i=(X_i,A_i,S_i,R_i,R_iY_i).
\]
Define the prefix probabilities and remaining probability by
\[
w_{k_0}:=P_\omega(K_0=k_0)
=e^{-n/2}\frac{(n/2)^{k_0}}{k_0!}
\quad(k_0\in\mathbb Z_{\ge0}),
\qquad
\delta_{\mathrm{tail}}
:=1-\sum_{k_0=0}^{n}w_{k_0}
=P_\omega(K_0>n).
\]
The disjoint prefix events imply
\[
w_{k_0}\ge0,\qquad
\sum_{k_0=0}^{n}w_{k_0}\le1,\qquad
\delta_{\mathrm{tail}}\ge0.
\]
For each \(0\le k_0\le n\), define the assignment set and its output average by
\[
\Sigma_{k_0}
:=\bigl\{\sigma:\{1,\ldots,k_0\}\to\{1,2,3\}\bigr\},
\qquad
\overline\tau_{k_0}(\mathbf o)
:=3^{-k_0}\sum_{\sigma\in\Sigma_{k_0}}
 \widetilde\tau(\mathbf o;(k_0,\sigma)),
\]
where the three assignment values denote the membership, pilot, and estimation streams. The notation \((k_0,\sigma)\) specifies the auxiliary prefix and assignment in \cref{obj:def:mixed-count-estimator}.

The weights \(w_0,\ldots,w_n,\delta_{\mathrm{tail}}\) sum to one, and
\[
\sum_{k_0=0}^{n}w_{k_0}\overline\tau_{k_0}(\mathbf o)-\tau(P)
=
\delta_{\mathrm{tail}}\bigl(-\tau(P)\bigr)
+\sum_{k_0=0}^{n}w_{k_0}
 \bigl(\overline\tau_{k_0}(\mathbf o)-\tau(P)\bigr).
\]
Convexity of the square therefore gives
\[
\left(
 \sum_{k_0=0}^{n}w_{k_0}\overline\tau_{k_0}(\mathbf o)-\tau(P)
\right)^2
\le
\delta_{\mathrm{tail}}\tau(P)^2
+\sum_{k_0=0}^{n}w_{k_0}
 \bigl(\overline\tau_{k_0}(\mathbf o)-\tau(P)\bigr)^2.
\]
% lean: CausalSmith.Stat.MarRareqLogfrontier.finite_weighted_square_with_remainder

\item Since \(|\Sigma_{k_0}|=3^{k_0}\), centering the assignment average gives
\[
\overline\tau_{k_0}(\mathbf o)-\tau(P)
=
3^{-k_0}\sum_{\sigma\in\Sigma_{k_0}}
 \bigl(\widetilde\tau(\mathbf o;(k_0,\sigma))-\tau(P)\bigr).
\]
The finite mean-square inequality, obtained by Cauchy--Schwarz, consequently yields
\[
\bigl(\overline\tau_{k_0}(\mathbf o)-\tau(P)\bigr)^2
\le
3^{-k_0}\sum_{\sigma\in\Sigma_{k_0}}
 \bigl(\widetilde\tau(\mathbf o;(k_0,\sigma))-\tau(P)\bigr)^2.
\]
This also covers \(k_0=0\): the empty prefix has one assignment, so the averaging denominator equals one.
% lean: CausalSmith.Stat.MarRareqLogfrontier.finite_uniform_square

\item Multiplying these inequalities by the nonnegative prefix probabilities and summing gives
\[
\begin{aligned}
&\sum_{k_0=0}^{n}w_{k_0}
 \bigl(\overline\tau_{k_0}(\mathbf o)-\tau(P)\bigr)^2\\
&\qquad\le
\sum_{k_0=0}^{n}w_{k_0}3^{-k_0}
 \sum_{\sigma\in\Sigma_{k_0}}
 \bigl(\widetilde\tau(\mathbf o;(k_0,\sigma))-\tau(P)\bigr)^2.
\end{aligned}
\]
The estimator's finite averaging formula in
\cref{obj:def:mixed-count-estimator} is
\[
\widehat\tau^{\mathrm{mix}}_{n,d,q}(\mathbf o)
=
\sum_{k_0=0}^{n}w_{k_0}\overline\tau_{k_0}(\mathbf o).
\]
Combining this identity with the preceding two bounds establishes, for every \(\mathbf o\),
\[
\begin{aligned}
&\bigl(\widehat\tau^{\mathrm{mix}}_{n,d,q}(\mathbf o)-\tau(P)\bigr)^2\\
&\quad\le
\sum_{k_0=0}^{n}w_{k_0}3^{-k_0}
 \sum_{\sigma\in\Sigma_{k_0}}
 \bigl(\widetilde\tau(\mathbf o;(k_0,\sigma))-\tau(P)\bigr)^2
+\delta_{\mathrm{tail}}\tau(P)^2\\
&\quad=
E_\omega\!\left[
 \bigl(\widetilde\tau(\mathbf O_n;\omega)-\tau(P)\bigr)^2
 \,\middle|\,\mathbf O_n=\mathbf o
\right].
\end{aligned}
\]
The last term in the sum is the squared loss of the zero output on \(K_0>n\). All functions in this inequality are measurable and integrable on the finite observed-sample space. Integrating under the canonical product law yields
\[
E_P\!\left[
 \bigl(\widehat\tau^{\mathrm{mix}}_{n,d,q}(\mathbf O_n)-\tau(P)\bigr)^2
\right]
\le
E_P\!\left[
 E_\omega\!\left[
  \bigl(\widetilde\tau(\mathbf O_n;\omega)-\tau(P)\bigr)^2
  \,\middle|\,\mathbf O_n
 \right]
\right],
\]
as required.
% lean: CausalSmith.Stat.MarRareqLogfrontier.mixed_count_derandomization
\end{enumerate}
\end{proof}

\begin{proof}[Proof of \cref{obj:thm:connected-interval-frontier}]
Fix \(\alpha\in(0,1)\). Throughout the proof, write
\[
\beta:=1-\alpha>0,\qquad
N:=nq>0,\qquad
\ell:=\log(e+N)\ge1,
\]
where \(n,d\ge1\), \(0<q\le1\), and \(q\ell^2\le1/64\).
Define the two length scales by
\[
s_{\mathrm{par}}:=\min\{1,N^{-1/2}\},
\qquad
s_{\mathrm{many}}:=\min\left\{1,\frac{d}{N\ell}\right\}.
\]

\begin{enumerate}
\item \label{proof:connected-interval-upper}\textbf{Upper bound for the risk-envelope interval.}
By \cref{obj:thm:mixed-count-upper}, there is a universal
\(\overline C>0\) such that, for every \(P\in\mathcal M_{n,d,q}\),
\[
E_P\!\left[
  \bigl(\widehat\tau^{\mathrm{mix}}_{n,d,q}-\tau(P)\bigr)^2
\right]
\le \mathfrak U_{n,d,q}
\le \overline C\,r_{n,d,q}.
\]
For every observed sample, define the envelope radius by
\[
\rho_{\mathrm{env}}
:=\sqrt{\frac{\mathfrak U_{n,d,q}}{\alpha}}.
\]
The endpoints in \cref{obj:def:risk-envelope-interval} satisfy
\[
\operatorname{length}(I^{\mathrm{mix}}_\alpha)
=
\min\{1,\widehat\tau^{\mathrm{mix}}_{n,d,q}+\rho_{\mathrm{env}}\}
-
\max\{-1,\widehat\tau^{\mathrm{mix}}_{n,d,q}-\rho_{\mathrm{env}}\}
\le 2\rho_{\mathrm{env}}.
\]
Consequently, with
\[
C_\alpha:=2\sqrt{\frac{\overline C}{\alpha}}>0,
\]
we have
\[
E_P[\operatorname{length}(I^{\mathrm{mix}}_\alpha)]
\le
2\sqrt{\frac{\mathfrak U_{n,d,q}}{\alpha}}
\le C_\alpha\sqrt{r_{n,d,q}}.
\]
Taking the supremum gives
\[
\sup_{P\in\mathcal M_{n,d,q}}
E_P[\operatorname{length}(I^{\mathrm{mix}}_\alpha)]
\le C_\alpha\sqrt{r_{n,d,q}}.
\]
For an empty model class, the supremum is zero by the convention in
\cref{obj:def:interval-length-risk}, and the same inequality follows from
\(r_{n,d,q}\ge0\).
% lean: CausalSmith.Stat.MarRareqLogfrontier.riskEnvelope_interval_worst_upper

\item \label{proof:connected-interval-calibration}\textbf{Calibration of the activated regime.}
Define constants depending only on \(\alpha\) by
\[
\eta:=e^{-32},\qquad
M_{\mathrm{act},\alpha}:=\left\lceil\frac8\beta\right\rceil+1,
\qquad
D_{\mathrm{act},\alpha}
:=
\left\lceil
\frac{18432M_{\mathrm{act},\alpha}^2}{\beta}
\right\rceil+1,
\]
\[
T_{\mathrm{act},\alpha}
:=
\max\left\{
1,e^2,4\eta,2\eta D_{\mathrm{act},\alpha},
\exp\left(\frac{32-\log(\beta/16)}{28}\right)
\right\},
\qquad
c_{\mathrm{act},\alpha}
:=\frac{21\beta}{32}\frac{\eta}{24}>0.
\]
The ceiling inequalities give
\[
M_{\mathrm{act},\alpha}\ge1,\qquad
\frac1{M_{\mathrm{act},\alpha}}\le\frac\beta8,\qquad
D_{\mathrm{act},\alpha}\ge1,\qquad
\frac{18432M_{\mathrm{act},\alpha}^2}{\beta}
\le D_{\mathrm{act},\alpha}.
\]
Use the quantities from \cref{obj:def:activation-lower-handle}:
\[
K:=\lceil\ell\rceil,\qquad H:=K^2,\qquad
b_0:=\frac{\eta}{N\ell},\qquad
J:=\min\left\{d-1,\left\lfloor\frac1{2b_0}\right\rfloor\right\}.
\]
They are defined for all the parameters under consideration, including
\(d=1\), when \(J=0\).

Suppose first that \(N\ge T_{\mathrm{act},\alpha}\). Since \(\ell\ge1\),
\[
D_{\mathrm{act},\alpha}
\le \frac{N}{2\eta}
\le \frac{N\ell}{2\eta}
=\frac1{2b_0}.
\]
Also,
\[
b_0\le\frac{\eta}{N}\le\frac14,\qquad
\ell\ge2,\qquad K\ge2.
\]
The ceiling bound \(K\le\ell+1\le2\ell\), together with the slice restriction,
gives
\[
qH=qK^2\le4q\ell^2\le\frac1{16}.
\]
By \cref{obj:lem:moment-matched-reciprocal-priors}, choose finitely supported
probability laws \(\Pi_0,\Pi_1\) on \([1,H]\) with matching moments through
degree \(K\). Interchange their names if necessary and define
\[
\mu_{\mathrm{rec},i}:=\int z^{-1}\,d\Pi_i(z)
\quad(i=0,1).
\]
Then
\[
\int z^v\,d\Pi_0(z)=\int z^v\,d\Pi_1(z)
\quad(0\le v\le K),
\qquad
\mu_{\mathrm{rec},1}-\mu_{\mathrm{rec},0}\ge\frac1{12}.
\]

We will show that whenever \(J\ge D_{\mathrm{act},\alpha}\),
\[
c_{\mathrm{act},\alpha}s_{\mathrm{many}}
\le\mathfrak L_{n,d,q,\alpha}.
\]
% lean: CausalSmith.Stat.MarRareqLogfrontier.interval_activated_large_regime

\item \textbf{The activated grid and its likelihood comparison.}
Continue under
\(N\ge T_{\mathrm{act},\alpha}\) and \(J\ge D_{\mathrm{act},\alpha}\).
For each integer \(0\le\iota\le M_{\mathrm{act},\alpha}\), define
\[
\Pi^{(\iota)}
:=
\left(1-\frac{\iota}{M_{\mathrm{act},\alpha}}\right)\Pi_0
+\frac{\iota}{M_{\mathrm{act},\alpha}}\Pi_1.
\]
These are finitely supported probability laws on \([1,H]\), with the same
moments through degree \(K\).

For completeness, the full-data laws underlying the activated experiment
can be specified as follows. For a fixed
\(\mathbf z=(z_0,\ldots,z_{d-1})\in[1,H]^d\), let
\(P^{\mathrm{act}}_{\mathbf z}\) have baseline probabilities
\[
P^{\mathrm{act}}_{\mathbf z}(X=x)
=
\begin{cases}
b_0,&0\le x<J,\\
1-Jb_0,&x=J,\\
0,&J<x<d.
\end{cases}
\]
Let treatment be an independent fair coin, and set
\[
S(0)=S(1)=Y(0)=0,\qquad S=0,\qquad Y=AY(1).
\]
Conditional on \(X=x\), generate the treated potential outcome and arrival
independently, with
\[
P^{\mathrm{act}}_{\mathbf z}(Y(1)=1\mid X=x)
=
\begin{cases}
z_x^{-1},&x<J,\\
0,&x\ge J,
\end{cases}
\qquad
P^{\mathrm{act}}_{\mathbf z}(R=1\mid X=x,A)
=
\begin{cases}
qz_x,&x<J,\\
qH,&x\ge J.
\end{cases}
\]
The inequalities \(Jb_0\le1/2\) and \(qH\le1/16\) ensure valid
probabilities. Independence of treatment, consistency, conditional
independence of arrival and outcome, and the arrival floor follow from
this construction. Thus
\[
P^{\mathrm{act}}_{\mathbf z}\in\mathcal M_{n,d,q},
\qquad
\tau(P^{\mathrm{act}}_{\mathbf z})
=b_0\sum_{x=0}^{J-1}z_x^{-1}.
\]

Let \(\widetilde{\mathbb Q}_{\mathrm{act},\iota}\) be the augmented
sample law obtained from the rule in
\cref{obj:def:activation-lower-handle}, using
\((\Pi^{(\iota)})^{\otimes d}\) for the latent vector. Let
\(\mathbb Q_{\mathrm{act},\iota}\) be its projection onto the observed
sample. The conditional probabilities in that construction give
\[
\mathbb Q_{\mathrm{act},\iota}
=
\int
\mathcal L_{P^{\mathrm{act}}_{\mathbf z}}(\mathbf O_n)\,
d(\Pi^{(\iota)})^{\otimes d}(\mathbf z).
\]
Define the activation count in each rare category and the common cutoff
event by
\[
V_x:=\sum_{i=1}^n
\mathbf1\{X_i=x,\ \mathsf B_{\mathrm{act},i}=1\}
\quad(0\le x<J),
\qquad
G_{\mathrm{act}}:=\{V_x\le K\text{ for every }x<J\}.
\]
Conditional on the membership, treatment, and activation pattern, every
activated observation in category \(x\) contributes an affine factor in
\(z_x\), and every inactive observation contributes a constant factor.
The conditional likelihood therefore has degree at most \(V_x\) in each
\(z_x\). On \(G_{\mathrm{act}}\), integration against the product priors
and moment matching imply equality of all augmented sample atoms:
\[
\widetilde{\mathbb Q}_{\mathrm{act},\iota}(\{\mathbf o_{\mathrm{aug}}\})
=
\widetilde{\mathbb Q}_{\mathrm{act},0}(\{\mathbf o_{\mathrm{aug}}\})
\qquad(\mathbf o_{\mathrm{aug}}\in G_{\mathrm{act}}).
\]
Hence the two laws also assign the same mass to \(G_{\mathrm{act}}\)
and its complement. Splitting any event into its intersections with
these two sets, and then projecting, yields
\[
\operatorname{TV}
\bigl(\mathbb Q_{\mathrm{act},0},
      \mathbb Q_{\mathrm{act},\iota}\bigr)
\le
\widetilde{\mathbb Q}_{\mathrm{act},0}(G_{\mathrm{act}}^c).
\]

To bound this last probability, define
\[
K_{\mathrm{tail}}:=K+1,\qquad
\lambda_{\mathrm{act}}:=nb_0qH=\frac{\eta K^2}{\ell}.
\]
The membership and activation pattern has a common law under all the
priors, and each \(V_x\) has distribution
\(\operatorname{Binomial}(n,b_0qH)\). Counting subsets of
\(K_{\mathrm{tail}}\) successful observations gives
\[
E\binom{V_x}{K_{\mathrm{tail}}}
=
\binom{n}{K_{\mathrm{tail}}}(b_0qH)^{K_{\mathrm{tail}}},
\]
with both sides zero when \(K_{\mathrm{tail}}>n\). Since
\(\mathbf1\{V_x>K\}\le\binom{V_x}{K_{\mathrm{tail}}}\),
\[
\Pr(V_x>K)
\le\frac{\lambda_{\mathrm{act}}^{K_{\mathrm{tail}}}}
         {K_{\mathrm{tail}}!}.
\]
The ceiling inequalities imply
\[
\lambda_{\mathrm{act}}
=\frac{\eta K^2}{\ell}
\le2\eta K_{\mathrm{tail}}
\le e^{-31}K_{\mathrm{tail}},
\qquad
K_{\mathrm{tail}}\ge\ell.
\]
Using the \(K_{\mathrm{tail}}\)-th term of the exponential series,
\[
\frac{\lambda_{\mathrm{act}}^{K_{\mathrm{tail}}}}
     {K_{\mathrm{tail}}!}
\le
e^{-31K_{\mathrm{tail}}}
\frac{K_{\mathrm{tail}}^{K_{\mathrm{tail}}}}
     {K_{\mathrm{tail}}!}
\le e^{-30K_{\mathrm{tail}}}
\le e^{-30\ell}.
\]
A union bound, followed by the mass cap and
\(N\le e^\ell\), \(\ell\le e^\ell\), now gives
\[
\widetilde{\mathbb Q}_{\mathrm{act},0}(G_{\mathrm{act}}^c)
\le Je^{-30\ell}
\le \frac{N\ell}{2e^{-32}}e^{-30\ell}
\le e^{32-28\ell}.
\]
Finally, the definition of \(T_{\mathrm{act},\alpha}\) implies
\[
\ell\ge\frac{32-\log(\beta/16)}{28},
\qquad
e^{32-28\ell}\le\frac\beta{16}.
\]
Thus, for every grid index,
\[
\operatorname{TV}
\bigl(\mathbb Q_{\mathrm{act},0},
      \mathbb Q_{\mathrm{act},\iota}\bigr)
\le\frac\beta{16}.
\]
% lean: CausalSmith.Stat.MarRareqLogfrontier.activatedAugmentedMixture_interval_minimax_calibrated

\item \textbf{Coverage of the prior neighborhoods and their length cost.}
Define the target separation, centers, and neighborhood radius by
\[
\Delta_{\mathrm{act}}
:=Jb_0(\mu_{\mathrm{rec},1}-\mu_{\mathrm{rec},0})
\ge\frac{Jb_0}{12}>0,
\]
\[
\theta_{\mathrm{act},\iota}
:=
Jb_0\left[
\left(1-\frac{\iota}{M_{\mathrm{act},\alpha}}\right)
\mu_{\mathrm{rec},0}
+\frac{\iota}{M_{\mathrm{act},\alpha}}\mu_{\mathrm{rec},1}
\right]
=
\theta_{\mathrm{act},0}
+\frac{\iota\Delta_{\mathrm{act}}}{M_{\mathrm{act},\alpha}},
\qquad
w_{\mathrm{act}}
:=\frac{\Delta_{\mathrm{act}}}{8M_{\mathrm{act},\alpha}}.
\]
Under the \(\iota\)-th product prior, the actual target has mean
\(\theta_{\mathrm{act},\iota}\). Independence of the latent coordinates
and \(z_x^{-1}\in[0,1]\) give
\[
\operatorname{Var}\!\left(
b_0\sum_{x=0}^{J-1}z_x^{-1}
\right)\le\frac{Jb_0^2}{4}.
\]
Indeed, a variable in \([0,1]\) has variance at most \(1/4\), and the
variances of these independent summands add. Chebyshev's inequality
therefore gives
\[
\Pr_{\Pi^{(\iota)\otimes d}}\!\left(
\left|\tau(P^{\mathrm{act}}_{\mathbf z})
-\theta_{\mathrm{act},\iota}\right|>w_{\mathrm{act}}
\right)
\le
\frac{Jb_0^2/4}{w_{\mathrm{act}}^2}
\le
\frac{2304M_{\mathrm{act},\alpha}^2}{J}
\le\frac\beta8.
\]

Fix any \(\mathsf I\in\mathcal C_{n,d,q,\alpha}\), and let
\(\nu_{\mathrm{act},\iota}\) denote its returned-interval law when the
sample has law \(\mathbb Q_{\mathrm{act},\iota}\):
\[
\nu_{\mathrm{act},\iota}
:=\mathbb Q_{\mathrm{act},\iota}\mathsf I.
\]
For each grid index, define the interval event
\[
F_{\mathrm{act},\iota}
:=
\left\{
[l,u]:
[l,u]\cap
[\theta_{\mathrm{act},\iota}-w_{\mathrm{act}},
 \theta_{\mathrm{act},\iota}+w_{\mathrm{act}}]
\ne\varnothing
\right\}.
\]
Honesty holds for every fixed completion. Averaging that coverage over
the finite prior and subtracting the target-concentration failure
probability yields
\[
\nu_{\mathrm{act},\iota}(F_{\mathrm{act},\iota})
\ge\beta-\frac\beta8=\frac{7\beta}{8}.
\]
By \cref{obj:lem:randomized-kernel-tv-comparison} and the preceding
likelihood comparison,
\[
\operatorname{TV}
(\nu_{\mathrm{act},0},\nu_{\mathrm{act},\iota})
\le\frac\beta{16},
\qquad
\nu_{\mathrm{act},0}(F_{\mathrm{act},\iota})
\ge\frac{13\beta}{16}.
\]

Define the number of neighborhoods met by a returned interval by
\[
\mathcal N_{\mathrm{hit}}([l,u])
:=
\sum_{\iota=0}^{M_{\mathrm{act},\alpha}}
\mathbf1_{F_{\mathrm{act},\iota}}([l,u]).
\]
Meeting the \(\iota\)-th neighborhood is equivalent to
\(\theta_{\mathrm{act},\iota}\in[l-w_{\mathrm{act}},u+w_{\mathrm{act}}]\).
The centers are spaced by
\(\Delta_{\mathrm{act}}/M_{\mathrm{act},\alpha}\). Thus the distance
between the extreme centers contained in this expanded interval gives
\[
(u-l)+2w_{\mathrm{act}}
\ge
\frac{\Delta_{\mathrm{act}}}{M_{\mathrm{act},\alpha}}
\bigl(\mathcal N_{\mathrm{hit}}([l,u])-1\bigr).
\]
For zero or one contained center, the right side is nonpositive or
zero, so the inequality covers those cases as well. Taking expectations,
\[
\begin{aligned}
E_{\nu_{\mathrm{act},0}}[u-l]
&\ge
\frac{\Delta_{\mathrm{act}}}{M_{\mathrm{act},\alpha}}
\left[
(M_{\mathrm{act},\alpha}+1)\frac{13\beta}{16}-1
\right]
-2w_{\mathrm{act}}\\
&=
\Delta_{\mathrm{act}}
\left[
\frac{13\beta}{16}
+\frac{13\beta}{16M_{\mathrm{act},\alpha}}
-\frac{5}{4M_{\mathrm{act},\alpha}}
\right]\\
&\ge\frac{21\beta}{32}\Delta_{\mathrm{act}},
\end{aligned}
\]
where the last step uses \(1/M_{\mathrm{act},\alpha}\le\beta/8\).
The left side is a prior average of the expected lengths of this same
procedure over laws in \(\mathcal M_{n,d,q}\). Therefore
\[
\frac{21\beta}{32}\Delta_{\mathrm{act}}
\le
\sup_{P\in\mathcal M_{n,d,q}}E_P[\operatorname{length}(\mathsf I)].
\]
The constant procedure returning \([-1,1]\) is honest, so the procedure
class is nonempty. Taking its defining infimum gives
\[
\frac{21\beta}{32}\Delta_{\mathrm{act}}
\le\mathfrak L_{n,d,q,\alpha}.
\]

To express the separation in terms of \(d\), split according to the
minimum defining \(J\). If
\(d-1\le\lfloor(2b_0)^{-1}\rfloor\), then \(J=d-1\), and
\(J\ge D_{\mathrm{act},\alpha}\ge1\) implies \(d\ge2\) and
\(J\ge d/2\). Otherwise,
\(J=\lfloor(2b_0)^{-1}\rfloor\), and
\[
Jb_0\ge\frac12-b_0\ge\frac14.
\]
In both cases,
\[
Jb_0\ge\min\left\{\frac{db_0}{2},\frac14\right\}.
\]
Since \(\eta\le1/2\), the two inequalities
\[
\frac{\eta}{24}s_{\mathrm{many}}
\le\frac{db_0}{24},
\qquad
\frac{\eta}{24}s_{\mathrm{many}}
\le\frac1{48}
\]
give
\[
\frac{\eta}{24}s_{\mathrm{many}}
\le
\frac1{12}\min\left\{\frac{db_0}{2},\frac14\right\}
\le\frac{Jb_0}{12}
\le\Delta_{\mathrm{act}}.
\]
Consequently,
\[
c_{\mathrm{act},\alpha}s_{\mathrm{many}}
\le\mathfrak L_{n,d,q,\alpha}
\qquad
(N\ge T_{\mathrm{act},\alpha},\ J\ge D_{\mathrm{act},\alpha}).
\]
% lean: CausalSmith.Stat.MarRareqLogfrontier.interval_activated_large_regime

\item \textbf{An explicit one-cell floor.}
Use the family \(P_u\) supplied by
\cref{obj:lem:one-cell-testing-family}. We record the calibration needed
here:
\[
c_{\mathrm{par},\alpha}:=\frac{3\beta^2}{512}>0,
\qquad
c_{\mathrm{par},\alpha}s_{\mathrm{par}}
\le\mathfrak L_{n,d,q,\alpha}.
\]
To obtain this explicit constant, define
\[
\eta_{\mathrm{cell}}:=\frac\beta8,\qquad
w_{\mathrm{cell}}
:=\min\left\{\frac14,\frac{\eta_{\mathrm{cell}}}{4\sqrt N}\right\}.
\]
Choose an integer \(M_{\mathrm{cell},\alpha}\ge8/\beta\), and set
\[
\delta_{\mathrm{cell}}
:=\frac{w_{\mathrm{cell}}}{M_{\mathrm{cell},\alpha}}>0,
\qquad
u_{\mathrm{cell},\iota}
:=\frac12+\iota\delta_{\mathrm{cell}}
\quad(0\le\iota\le M_{\mathrm{cell},\alpha}).
\]
These grid points are distinct and belong to \([1/4,3/4]\), and
\[
2Nw_{\mathrm{cell}}^2\le\frac{\eta_{\mathrm{cell}}^2}{8},
\qquad
w_{\mathrm{cell}}\ge
\frac{\eta_{\mathrm{cell}}}{4}s_{\mathrm{par}}.
\]
Indeed, \(N>0\) and \(0<\eta_{\mathrm{cell}}=\beta/8<1\), so
\(w_{\mathrm{cell}}>0\). The defining minimum gives
\[
w_{\mathrm{cell}}\le\frac{\eta_{\mathrm{cell}}}{4\sqrt N},
\qquad
2Nw_{\mathrm{cell}}^2
\le 2N\left(\frac{\eta_{\mathrm{cell}}}{4\sqrt N}\right)^2
=\frac{\eta_{\mathrm{cell}}^2}{8},
\]
where the equality uses \((\sqrt N)^2=N\). For the lower bound, if
\(1/4\le\eta_{\mathrm{cell}}/(4\sqrt N)\), then
\(w_{\mathrm{cell}}=1/4\), and \(s_{\mathrm{par}}\le1\) gives
\[
\frac{\eta_{\mathrm{cell}}}{4}s_{\mathrm{par}}
\le\frac{\eta_{\mathrm{cell}}}{4}
\le\frac14=w_{\mathrm{cell}}.
\]
Otherwise, \(w_{\mathrm{cell}}=\eta_{\mathrm{cell}}/(4\sqrt N)\),
and \(s_{\mathrm{par}}\le N^{-1/2}\) gives
\[
\frac{\eta_{\mathrm{cell}}}{4}s_{\mathrm{par}}
\le\frac{\eta_{\mathrm{cell}}}{4\sqrt N}
=w_{\mathrm{cell}}.
\]

Write the one-record and sample laws of this family as
\[
\mathbb Q_{\mathrm{cell},u}^{(1)}:=\mathcal L_{P_u}(O_1),
\qquad
\mathbb Q_{\mathrm{cell},u}^{(n)}
:=(\mathbb Q_{\mathrm{cell},u}^{(1)})^{\otimes n}.
\]
In this family, the arrival floor and \(P_u(R=1)=q\), together with
balanced assignment, force arrival probability \(q\) in each treatment
arm. Conditional independence of arrival and outcome then gives the
two treated arrived-outcome atom probabilities \(qu/2\) and
\(q(1-u)/2\); all other observed atom probabilities are constant in \(u\).
Consequently, defining the one-record chi-squared divergence by its
finite sum,
\[
\mathcal X_{\mathrm{cell},\iota}^{(1)}
:=
\sum_{\substack{o:\\
\mathbb Q_{\mathrm{cell},1/2}^{(1)}(\{o\})>0}}
\frac{
\left[
\mathbb Q_{\mathrm{cell},u_{\mathrm{cell},\iota}}^{(1)}(\{o\})
-\mathbb Q_{\mathrm{cell},1/2}^{(1)}(\{o\})
\right]^2}
{\mathbb Q_{\mathrm{cell},1/2}^{(1)}(\{o\})}
=
2q(\iota\delta_{\mathrm{cell}})^2.
\]
This computation also covers \(q=1\), since the vanished nonarrival
atoms are omitted from the sum. Define
\(\mathcal X_{\mathrm{cell},\iota}^{(n)}\) by the analogous sum for the
sample laws, and put
\[
t_{\mathrm{cell},\iota}
:=2N(\iota\delta_{\mathrm{cell}})^2
\le\frac{\eta_{\mathrm{cell}}^2}{8}.
\]
Independence of the observations gives
\[
1+\mathcal X_{\mathrm{cell},\iota}^{(n)}
=
\bigl(1+\mathcal X_{\mathrm{cell},\iota}^{(1)}\bigr)^n
\le e^{t_{\mathrm{cell},\iota}}.
\]
For \(0\le t_{\mathrm{cell},\iota}<1\), comparison of the exponential
and geometric series gives
\[
e^{t_{\mathrm{cell},\iota}}
\le\frac1{1-t_{\mathrm{cell},\iota}},
\qquad
\frac{t_{\mathrm{cell},\iota}}{1-t_{\mathrm{cell},\iota}}
\le\eta_{\mathrm{cell}}^2,
\]
using \(t_{\mathrm{cell},\iota}\le\eta_{\mathrm{cell}}^2/8\) and
\(\eta_{\mathrm{cell}}\le1\). Thus Cauchy--Schwarz applied to the
finite likelihood-ratio sum yields
\[
\operatorname{TV}
\bigl(\mathbb Q_{\mathrm{cell},1/2}^{(n)},
      \mathbb Q_{\mathrm{cell},u_{\mathrm{cell},\iota}}^{(n)}\bigr)
\le\frac12\sqrt{\mathcal X_{\mathrm{cell},\iota}^{(n)}}
\le\eta_{\mathrm{cell}}.
\]

For any honest procedure \(\mathsf I\), coverage of each grid point is
at least \(\beta\) under its corresponding sample law. By
\cref{obj:lem:randomized-kernel-tv-comparison}, under the reference law
\(P_{1/2}\) each point is contained in the returned interval with
probability at least
\(\beta-\eta_{\mathrm{cell}}=7\beta/8\). A connected interval
containing a given number of grid points has length at least their
spacing times that number minus one. Summing the inclusion
probabilities therefore gives
\[
\begin{aligned}
E_{P_{1/2}}[\operatorname{length}(\mathsf I)]
&\ge
\delta_{\mathrm{cell}}
\left[(M_{\mathrm{cell},\alpha}+1)\frac{7\beta}{8}-1\right]\\
&\ge\frac{3\beta}{4}w_{\mathrm{cell}},
\end{aligned}
\]
because \(1/M_{\mathrm{cell},\alpha}\le\beta/8\).
All interval risks lie in \([0,2]\), and the constant interval
\([-1,1]\) is honest. Hence taking the supremum over models and the
infimum over honest procedures preserves this lower bound. Finally,
\[
c_{\mathrm{par},\alpha}s_{\mathrm{par}}
\le
\frac{3\beta}{4}\frac{\eta_{\mathrm{cell}}}{4}s_{\mathrm{par}}
\le
\frac{3\beta}{4}w_{\mathrm{cell}}
\le\mathfrak L_{n,d,q,\alpha},
\]
which proves the asserted explicit floor.
% lean: CausalSmith.Stat.MarRareqLogfrontier.oneCell_interval_minimax_floor

\item \textbf{Extension of the many-cell scale to all regimes.}
Define
\[
c_{\mathrm{many},\alpha}
:=
\min\left\{
c_{\mathrm{act},\alpha},
\ c_{\mathrm{par},\alpha}
       \min\{1,T_{\mathrm{act},\alpha}^{-1/2}\},
\ \frac{c_{\mathrm{par},\alpha}}{D_{\mathrm{act},\alpha}}
\right\}>0.
\]
We claim
\[
c_{\mathrm{many},\alpha}s_{\mathrm{many}}
\le\mathfrak L_{n,d,q,\alpha}
\]
throughout the parameter range.

First suppose \(N\ge T_{\mathrm{act},\alpha}\). If
\(J\ge D_{\mathrm{act},\alpha}\), the activated bound gives
\[
c_{\mathrm{many},\alpha}s_{\mathrm{many}}
\le c_{\mathrm{act},\alpha}s_{\mathrm{many}}
\le\mathfrak L_{n,d,q,\alpha}.
\]
If \(J<D_{\mathrm{act},\alpha}\), the capacity inequality from \cref{proof:connected-interval-calibration}
implies
\[
\left\lfloor\frac1{2b_0}\right\rfloor
\ge D_{\mathrm{act},\alpha}.
\]
Thus the minimum defining \(J\) forces
\(d-1<D_{\mathrm{act},\alpha}\), and hence
\(d\le D_{\mathrm{act},\alpha}\). Since \(N\ge1\) and \(\ell\ge1\),
\[
s_{\mathrm{many}}
\le\frac{d}{N\ell}
\le\frac{D_{\mathrm{act},\alpha}}{\sqrt N}
=
D_{\mathrm{act},\alpha}s_{\mathrm{par}}.
\]
It follows that
\[
c_{\mathrm{many},\alpha}s_{\mathrm{many}}
\le
c_{\mathrm{many},\alpha}D_{\mathrm{act},\alpha}s_{\mathrm{par}}
\le c_{\mathrm{par},\alpha}s_{\mathrm{par}}
\le\mathfrak L_{n,d,q,\alpha}.
\]
This case includes \(d=1\) and \(J=0\).

Finally, if \(N<T_{\mathrm{act},\alpha}\), monotonicity of the square
root and reciprocal gives
\[
\min\{1,T_{\mathrm{act},\alpha}^{-1/2}\}
\le s_{\mathrm{par}}.
\]
Using \(s_{\mathrm{many}}\le1\), we obtain
\[
c_{\mathrm{many},\alpha}s_{\mathrm{many}}
\le c_{\mathrm{many},\alpha}
\le
c_{\mathrm{par},\alpha}
\min\{1,T_{\mathrm{act},\alpha}^{-1/2}\}
\le\mathfrak L_{n,d,q,\alpha}.
\]
% lean: CausalSmith.Stat.MarRareqLogfrontier.interval_many_scale_of_large_regime

\item \label{proof:connected-interval-lower-scales}\textbf{Combination of the lower scales.}
The identities
\[
\sqrt{\min\{1,N^{-1}\}}=s_{\mathrm{par}},
\qquad
\sqrt{\min\left\{1,\left(\frac{d}{N\ell}\right)^2\right\}}
=s_{\mathrm{many}}
\]
follow by splitting at \(N^{-1}=1\) and \(d/(N\ell)=1\), respectively.
Define
\[
c_{\mathrm{both},\alpha}
:=\min\{c_{\mathrm{par},\alpha},c_{\mathrm{many},\alpha}\}>0.
\]
The preceding bounds give
\[
c_{\mathrm{both},\alpha}
\max\{s_{\mathrm{par}},s_{\mathrm{many}}\}
\le\mathfrak L_{n,d,q,\alpha}.
\]
Moreover,
\[
r_{n,d,q}
\le2\max\{s_{\mathrm{par}}^2,s_{\mathrm{many}}^2\}.
\]
Indeed, when both untruncated terms \(N^{-1}\) and
\((d/(N\ell))^2\) are at most one, their sum is at most twice their
maximum. If either exceeds one, the maximum of the truncated terms
equals one, while \(r_{n,d,q}\le1\). Since the maximum is nonnegative,
\[
\sqrt{r_{n,d,q}}
\le2\max\{s_{\mathrm{par}},s_{\mathrm{many}}\}.
\]
Therefore, defining
\[
c_{\mathrm{low},\alpha}:=\frac{c_{\mathrm{both},\alpha}}2>0,
\]
we obtain
\[
c_{\mathrm{low},\alpha}\sqrt{r_{n,d,q}}
\le\mathfrak L_{n,d,q,\alpha}.
\]
% lean: CausalSmith.Stat.MarRareqLogfrontier.frontier_interval_lower_of_two_scales

\item \textbf{Honesty and the minimax comparison.}
Set
\[
\underline c_\alpha
:=\min\{c_{\mathrm{low},\alpha},C_\alpha\},
\qquad
\overline C_\alpha:=C_\alpha.
\]
Then \(0<\underline c_\alpha\le\overline C_\alpha<\infty\), and
\cref{proof:connected-interval-lower-scales} supplies the required lower inequality.

Fix \(P\in\mathcal M_{n,d,q}\). Define the squared error by
\[
V_{\mathrm{env}}
:=
\bigl(\widehat\tau^{\mathrm{mix}}_{n,d,q}(\mathbf O_n)-\tau(P)\bigr)^2.
\]
The risk bound implies
\[
0\le E_P[V_{\mathrm{env}}]\le\mathfrak U_{n,d,q}.
\]
If \(\mathfrak U_{n,d,q}=0\), nonnegativity gives
\(V_{\mathrm{env}}=0\) almost surely, so the interval contains
\(\tau(P)\) almost surely. If \(\mathfrak U_{n,d,q}>0\), define
\[
t_{\mathrm{env}}:=\frac{\mathfrak U_{n,d,q}}{\alpha}>0.
\]
Integrating the pointwise inequality
\(V_{\mathrm{env}}\ge
t_{\mathrm{env}}\mathbf1\{V_{\mathrm{env}}\ge t_{\mathrm{env}}\}\)
gives
\[
t_{\mathrm{env}}\Pr_P(V_{\mathrm{env}}\ge t_{\mathrm{env}})
\le E_P[V_{\mathrm{env}}]\le\mathfrak U_{n,d,q},
\qquad
\Pr_P(V_{\mathrm{env}}\ge t_{\mathrm{env}})\le\alpha.
\]
Because \(\tau(P)\in[-1,1]\), squared error at most
\(t_{\mathrm{env}}\) implies inclusion in the clipped interval.
Consequently,
\[
\{\tau(P)\notin I^{\mathrm{mix}}_\alpha\}
\subseteq\{V_{\mathrm{env}}\ge t_{\mathrm{env}}\},
\qquad
\Pr_P\{\tau(P)\in I^{\mathrm{mix}}_\alpha\}\ge1-\alpha.
\]

The mixed-count estimator takes values in \([-1,1]\) by
\cref{obj:def:mixed-count-estimator}. Its envelope interval therefore
has ordered endpoints in \([-1,1]\). The finite observed-sample space
makes its endpoint map measurable, and its Dirac kernel is an honest
member of \(\mathcal C_{n,d,q,\alpha}\). For this kernel, joint
sample-and-kernel expected length equals
\(E_P[\operatorname{length}(I^{\mathrm{mix}}_\alpha)]\).
The definition in \cref{obj:def:interval-length-risk} hence gives
\[
\mathfrak L_{n,d,q,\alpha}
\le
\sup_{P\in\mathcal M_{n,d,q}}
E_P[\operatorname{length}(I^{\mathrm{mix}}_\alpha)].
\]
Combining this comparison with \cref{proof:connected-interval-upper,proof:connected-interval-lower-scales} proves
\[
\underline c_\alpha\sqrt{r_{n,d,q}}
\le\mathfrak L_{n,d,q,\alpha}
\le
\sup_{P\in\mathcal M_{n,d,q}}
E_P[\operatorname{length}(I^{\mathrm{mix}}_\alpha)]
\le\overline C_\alpha\sqrt{r_{n,d,q}},
\]
together with the asserted coverage.
% lean: CausalSmith.Stat.MarRareqLogfrontier.connected_interval_frontier
\end{enumerate}
\end{proof}

\begin{proof}[Proof of \cref{obj:thm:zeng-fixed-positivity-polynomial-frontier}]
\begin{enumerate}
\item \textit{Uniform risk bound for the deterministic estimator.}
Let \(\overline C>0\) be the universal constant supplied by
\cref{obj:thm:unrestricted-mixed-count-upper}. Fix \(n,d\ge1\),
\(0<\epsilon<1/2\), and \(P_Z\in\mathcal D_{d,\epsilon}\).
Define the auxiliary assignment vector and synthetic sample by
\[
\mathbf u\sim\operatorname{Unif}(\{0,1\}^n),
\qquad
\mathbf O_n^{\mathrm{syn}}
:=(o_i(u_i))_{i=1}^n,
\]
with \(\mathbf u\) independent of \(\mathbf Z_n\).
By \cref{obj:lem:synthetic-randomization-kernel}, there exists
\(P_Z^*\in\mathcal M^+_{n,d,\epsilon}\) such that
\[
\mathcal L_{P_Z,\mathbf u}(\mathbf O_n^{\mathrm{syn}})
=
\bigl(\mathcal L_{P_Z^*}(X,A,S,R,RY)\bigr)^{\otimes n},
\qquad
\tau(P_Z^*)=\theta_Z(P_Z).
\]
Thus the synthetic sample satisfies the sampling and model hypotheses of
\cref{obj:thm:unrestricted-mixed-count-upper} at \(q=\epsilon\).

For each fixed sample realization
\(\mathbf z=((x_i,b_i,y_i))_{i=1}^n\), Cauchy--Schwarz on the
\(2^n\) assignment vectors gives
\[
\begin{split}
&\left[
\sum_{\mathbf u\in\{0,1\}^n}
\left\{
\widehat\tau^{\mathrm{mix}}_{n,d,\epsilon}
\bigl(o_1(u_1),\ldots,o_n(u_n)\bigr)
-\theta_Z(P_Z)
\right\}
\right]^2\\
&\qquad\le
2^n\sum_{\mathbf u\in\{0,1\}^n}
\left\{
\widehat\tau^{\mathrm{mix}}_{n,d,\epsilon}
\bigl(o_1(u_1),\ldots,o_n(u_n)\bigr)
-\theta_Z(P_Z)
\right\}^2.
\end{split}
\]
Dividing by \(2^{2n}\) and using the defining finite average yields
\[
\begin{split}
&\left(
\widehat\theta^{\mathrm{poly}}_{n,d,\epsilon}(\mathbf z)
-\theta_Z(P_Z)
\right)^2\\
&\qquad\le
\frac1{2^n}
\sum_{\mathbf u\in\{0,1\}^n}
\left\{
\widehat\tau^{\mathrm{mix}}_{n,d,\epsilon}
\bigl(o_1(u_1),\ldots,o_n(u_n)\bigr)
-\theta_Z(P_Z)
\right\}^2.
\end{split}
\]
Integrating this inequality, identifying the uniform coin average with
the synthetic sample law, and then applying
\cref{obj:thm:unrestricted-mixed-count-upper}, we obtain
\[
\begin{split}
&E_{P_Z}\!\left[
\left(
\widehat\theta^{\mathrm{poly}}_{n,d,\epsilon}(\mathbf Z_n)
-\theta_Z(P_Z)
\right)^2
\right]\\
&\quad\le
E_{P_Z,\mathbf u}\!\left[
\left(
\widehat\tau^{\mathrm{mix}}_{n,d,\epsilon}
(\mathbf O_n^{\mathrm{syn}})
-\theta_Z(P_Z)
\right)^2
\right]\\
&\quad=
E_{P_Z^*}\!\left[
\left(
\widehat\tau^{\mathrm{mix}}_{n,d,\epsilon}(\mathbf O_n)
-\tau(P_Z^*)
\right)^2
\right]
\le \mathfrak U_{n,d,\epsilon}
\le \overline C\,r_{n,d,\epsilon}.
\end{split}
\]
Taking the supremum over a nonempty observational class proves the
first assertion. If that class is empty, its real-valued supremum is
zero, and the same bound follows from
\[
0\le r_{n,d,\epsilon}
=
\min\left\{
1,\frac1{n\epsilon}
+\left[\frac d{n\epsilon\log(e+n\epsilon)}\right]^2
\right\},
\qquad
0<\overline C.
\]
% lean: CausalSmith.Stat.MarRareqLogfrontier.thetaPoly_uniform_upper

\item \textit{Upper bound at fixed positivity.}
Fix \(0<\epsilon<1/2\). For \(n,d\ge1\), define
\[
r^{\mathrm{raw}}_{n,d}
:=
\frac1n+\left[\frac d{n\log(e+n)}\right]^2,
\qquad
r^{\mathrm{fix}}_{n,d}:=\min\{1,r^{\mathrm{raw}}_{n,d}\},
\]
and set
\[
\ell_{\epsilon,n}:=\log(e+n\epsilon),
\qquad
\ell_{1,n}:=\log(e+n),
\qquad
\gamma_{\epsilon,\mathrm{cmp}}
:=\frac{1-\log\epsilon}{\epsilon}.
\]
All denominators below are positive. In particular,
\[
\ell_{\epsilon,n}\ge1,
\qquad
\ell_{1,n}>0,
\qquad
\log\epsilon\le0,
\qquad
\gamma_{\epsilon,\mathrm{cmp}}\ge1.
\]
Since \(\epsilon\le1\),
\[
e+n\le\frac{e+n\epsilon}{\epsilon}.
\]
Monotonicity of the logarithm therefore gives
\[
\ell_{1,n}
\le \ell_{\epsilon,n}-\log\epsilon
\le (1-\log\epsilon)\ell_{\epsilon,n},
\]
where the second inequality follows from
\[
(-\log\epsilon)(\ell_{\epsilon,n}-1)\ge0.
\]
Consequently,
\[
n\ell_{1,n}
\le
\gamma_{\epsilon,\mathrm{cmp}}\,
n\epsilon\ell_{\epsilon,n}.
\]
Also,
\[
\epsilon\gamma_{\epsilon,\mathrm{cmp}}^2
\ge
\epsilon\gamma_{\epsilon,\mathrm{cmp}}
=
1-\log\epsilon
\ge1.
\]
These two comparisons imply, respectively,
\[
\left[\frac d{n\epsilon\ell_{\epsilon,n}}\right]^2
\le
\gamma_{\epsilon,\mathrm{cmp}}^2
\left[\frac d{n\ell_{1,n}}\right]^2,
\qquad
\frac1{n\epsilon}
\le
\gamma_{\epsilon,\mathrm{cmp}}^2\frac1n.
\]
Adding gives
\[
\frac1{n\epsilon}
+\left[\frac d{n\epsilon\ell_{\epsilon,n}}\right]^2
\le
\gamma_{\epsilon,\mathrm{cmp}}^2 r^{\mathrm{raw}}_{n,d}.
\]
Truncation preserves the needed comparison. Indeed, if
\(r^{\mathrm{raw}}_{n,d}\le1\), then
\[
\min\{1,\gamma_{\epsilon,\mathrm{cmp}}^2r^{\mathrm{raw}}_{n,d}\}
\le
\gamma_{\epsilon,\mathrm{cmp}}^2r^{\mathrm{raw}}_{n,d}
=
\gamma_{\epsilon,\mathrm{cmp}}^2r^{\mathrm{fix}}_{n,d};
\]
if \(r^{\mathrm{raw}}_{n,d}\ge1\), then
\[
\min\{1,\gamma_{\epsilon,\mathrm{cmp}}^2r^{\mathrm{raw}}_{n,d}\}
\le1
\le
\gamma_{\epsilon,\mathrm{cmp}}^2
=
\gamma_{\epsilon,\mathrm{cmp}}^2r^{\mathrm{fix}}_{n,d}.
\]
Thus
\[
r_{n,d,\epsilon}
\le
\gamma_{\epsilon,\mathrm{cmp}}^2r^{\mathrm{fix}}_{n,d}.
\]

The mixed-count estimator takes values in \([-1,1]\), by its clipped
output and averaging construction in \cref{obj:def:mixed-count-estimator}.
Hence, for every sample,
\[
-2^n
\le
\sum_{\mathbf u\in\{0,1\}^n}
\widehat\tau^{\mathrm{mix}}_{n,d,\epsilon}
\bigl(o_1(u_1),\ldots,o_n(u_n)\bigr)
\le2^n,
\]
so
\[
-1\le
\widehat\theta^{\mathrm{poly}}_{n,d,\epsilon}(\mathbf z)
\le1.
\]
This finite-sample map is measurable and is an admissible deterministic
procedure in \cref{obj:def:zeng-minimax-risk}. The defining infimum and
the preceding bounds therefore give
\[
\begin{split}
\mathfrak R^Z_{n,d,\epsilon}
&\le
\sup_{P_Z\in\mathcal D_{d,\epsilon}}
E_{P_Z}\!\left[
\left(
\widehat\theta^{\mathrm{poly}}_{n,d,\epsilon}(\mathbf Z_n)
-\theta_Z(P_Z)
\right)^2
\right]\\
&\le \overline C\,r_{n,d,\epsilon}
\le C_{\epsilon,\mathrm{up}}\,r^{\mathrm{fix}}_{n,d},
\qquad
C_{\epsilon,\mathrm{up}}
:=
\overline C\,\gamma_{\epsilon,\mathrm{cmp}}^2>0.
\end{split}
\]
% lean: CausalSmith.Stat.MarRareqLogfrontier.zeng_minimax_upper_fixedQ

\item \textit{Published lower bound on the zero-control subclass.}
Define
\[
\mathcal D^{(0)}_{d,\epsilon}
:=
\left\{
P_Z\in\mathcal D_{d,\epsilon}:
\mu_{0x}=0\ \text{for every }x\text{ with }p_x>0
\right\},
\]
and its minimax risk by
\[
\mathfrak R^{Z,0}_{n,d,\epsilon}
:=
\inf_{\widehat\theta}
\sup_{P_Z\in\mathcal D^{(0)}_{d,\epsilon}}
E_{P_Z}\!\left[
\bigl(\widehat\theta-\theta_Z(P_Z)\bigr)^2
\right],
\]
using the same procedure class as in
\cref{obj:def:zeng-minimax-risk}. The subclass inclusion gives
\[
\mathfrak R^{Z,0}_{n,d,\epsilon}
\le
\mathfrak R^Z_{n,d,\epsilon}.
\]
The published lower bound
\citep[Theorem~2, Section~3.2, p.~9, and its zero-control construction
in Section~C.5, pp.~39--44]{ZengBalakrishnanHanKennedy2026}
supplies constants
\[
c_{\epsilon,\mathrm{lit}}>0,
\qquad
c_{\epsilon,\mathrm{dim}}>0,
\qquad
n_{\epsilon,\mathrm{lit}}\in\mathbb N
\]
such that
\[
\mathfrak R^{Z,0}_{n,d,\epsilon}
\ge
c_{\epsilon,\mathrm{lit}}
\left\{
\frac1n+\left[\frac d{n\log n}\right]^2
\right\}
\]
whenever
\[
n\ge n_{\epsilon,\mathrm{lit}},
\qquad
1\le d\le c_{\epsilon,\mathrm{dim}}\,n\log n.
\]
The target belongs to \([-1,1]\), since it is a probability-weighted
average of differences of two means in \([0,1]\). By
\cref{obj:lem:randomized-kernel-barycenter}, clipping and then taking
conditional barycenters preserve or reduce squared-error risk at
every law. Thus the deterministic and randomized minimax values agree,
and the published bound applies to the procedure class above.
% lean: CausalSmith.Stat.MarRareqLogfrontier.ZengTheoremTwoLower

\item \textit{A uniform parametric lower bound.}
We next establish a universal constant \(c_{\mathrm{par}}>0\) such that
\[
\frac{c_{\mathrm{par}}}{n}
\le
\mathfrak R^{Z,0}_{n,d,\epsilon}
\le
\mathfrak R^Z_{n,d,\epsilon}
\qquad
(n,d\ge1,\ 0<\epsilon<1/2).
\]
For \(n\ge1\), define
\[
\delta_n:=\frac1{4\sqrt n},
\qquad
\vartheta^{\mathrm{test}}_0:=\frac12,
\qquad
\vartheta^{\mathrm{test}}_1:=\frac12+\delta_n.
\]
Then \(0<\delta_n\le1/4\), so both target values lie in \([0,1]\).
For \(\iota\in\{0,1\}\), define the one-record law
\(P^{\mathrm{test}}_{Z,\iota}\) on
\(\{1,\ldots,d\}\times\{0,1\}^2\) by
\[
P^{\mathrm{test}}_{Z,\iota}\{(x,b,y)\}
=
\begin{cases}
1/2,&(x,b,y)=(1,0,0),\\
(1-\vartheta^{\mathrm{test}}_\iota)/2,
 &(x,b,y)=(1,1,0),\\
\vartheta^{\mathrm{test}}_\iota/2,
 &(x,b,y)=(1,1,1),\\
0,&\text{otherwise}.
\end{cases}
\]
Its sole occupied baseline category has treatment probability \(1/2\),
control mean zero, and treated mean \(\vartheta^{\mathrm{test}}_\iota\).
Thus
\[
P^{\mathrm{test}}_{Z,\iota}\in\mathcal D^{(0)}_{d,\epsilon},
\qquad
\theta_Z(P^{\mathrm{test}}_{Z,\iota})
=\vartheta^{\mathrm{test}}_\iota.
\]

The two laws have common support. Define their chi-square divergence by
\[
\chi^2\!\left(
P^{\mathrm{test}}_{Z,1}\Vert P^{\mathrm{test}}_{Z,0}
\right)
:=
\sum_{\substack{z\in\{1,\ldots,d\}\times\{0,1\}^2:\\
P^{\mathrm{test}}_{Z,0}\{z\}>0}}
\frac{
\bigl(P^{\mathrm{test}}_{Z,1}\{z\}
-P^{\mathrm{test}}_{Z,0}\{z\}\bigr)^2
}{
P^{\mathrm{test}}_{Z,0}\{z\}
}.
\]
Only the two treated-outcome probabilities change. Each changes in
absolute value by \(\delta_n/2\), and each has reference probability
\(1/4\). Hence
\[
\chi^2\!\left(
P^{\mathrm{test}}_{Z,1}\Vert P^{\mathrm{test}}_{Z,0}
\right)
=2\delta_n^2
=\frac1{8n}.
\]
Write the sample laws and their likelihood ratio as
\[
Q^{\mathrm{test}}_\iota
:=
\bigl(P^{\mathrm{test}}_{Z,\iota}\bigr)^{\otimes n},
\qquad \iota\in\{0,1\},
\]
\[
L_{\mathrm{test}}(\mathbf z)
:=
\begin{cases}
Q^{\mathrm{test}}_1\{\mathbf z\}/
Q^{\mathrm{test}}_0\{\mathbf z\},
 &Q^{\mathrm{test}}_0\{\mathbf z\}>0,\\
0,&Q^{\mathrm{test}}_0\{\mathbf z\}=0.
\end{cases}
\]
Independence gives the second-moment identity and bound
\[
E_{Q^{\mathrm{test}}_0}[L_{\mathrm{test}}^2]
=
\left[
1+\chi^2\!\left(
P^{\mathrm{test}}_{Z,1}\Vert P^{\mathrm{test}}_{Z,0}
\right)
\right]^n
=
\left(1+\frac1{8n}\right)^n
\le e^{1/8}.
\]

Define the universal testing constant by
\[
c_{\mathrm{test}}:=\frac1{4e^{1/8}}>0.
\]
For any event \(\mathcal A\) in the finite sample space, set
\[
\eta_{\mathcal A}
:=
Q^{\mathrm{test}}_1(\mathcal A^c)
+
Q^{\mathrm{test}}_0(\mathcal A).
\]
Cauchy--Schwarz yields
\[
1-\eta_{\mathcal A}
\le Q^{\mathrm{test}}_1(\mathcal A)
=
E_{Q^{\mathrm{test}}_0}
[L_{\mathrm{test}}\mathbf1_{\mathcal A}]
\le
\sqrt{e^{1/8}Q^{\mathrm{test}}_0(\mathcal A)}.
\]
If \(\eta_{\mathcal A}<c_{\mathrm{test}}\), then
\(Q^{\mathrm{test}}_0(\mathcal A)\le\eta_{\mathcal A}\) implies
\[
\sqrt{e^{1/8}Q^{\mathrm{test}}_0(\mathcal A)}
<\frac12,
\qquad
1-\eta_{\mathcal A}>\frac34,
\]
because \(c_{\mathrm{test}}\le1/4\). This contradicts the preceding
inequality. Therefore every event satisfies
\[
Q^{\mathrm{test}}_1(\mathcal A^c)
+
Q^{\mathrm{test}}_0(\mathcal A)
\ge c_{\mathrm{test}}.
\]

Consider an arbitrary estimator. By
\cref{obj:lem:randomized-kernel-barycenter}, it suffices to consider its
clipped conditional action kernel \(\kappa\), supported on \([-1,1]\).
Define its barycenter by
\[
b_\kappa(\mathbf z)
:=
\int_{[-1,1]}h\,\kappa(\mathbf z,dh).
\]
The same cited result gives
\[
b_\kappa(\mathbf z)\in[-1,1],
\qquad
E_{Q^{\mathrm{test}}_\iota}
\left[
(b_\kappa-\vartheta^{\mathrm{test}}_\iota)^2
\right]
\le
E_{P^{\mathrm{test}}_{Z,\iota}}
\left[
(\widehat\theta-\vartheta^{\mathrm{test}}_\iota)^2
\right],
\qquad \iota\in\{0,1\},
\]
where the right-hand expectation includes the estimator's
randomization.

Define the nearest-target test and its large-error events by
\[
\mathcal A_{\mathrm{test}}
:=
\left\{
\mathbf z:
|b_\kappa(\mathbf z)-\vartheta^{\mathrm{test}}_1|
\le
|b_\kappa(\mathbf z)-\vartheta^{\mathrm{test}}_0|
\right\},
\]
\[
\mathcal A_{\mathrm{bad},\iota}
:=
\left\{
\mathbf z:
|b_\kappa(\mathbf z)-\vartheta^{\mathrm{test}}_\iota|
\ge\frac{\delta_n}{2}
\right\},
\qquad \iota\in\{0,1\}.
\]
The triangle inequality gives
\[
\delta_n
=
|\vartheta^{\mathrm{test}}_1-\vartheta^{\mathrm{test}}_0|
\le
|b_\kappa-\vartheta^{\mathrm{test}}_0|
+
|b_\kappa-\vartheta^{\mathrm{test}}_1|.
\]
Consequently,
\[
\mathcal A_{\mathrm{test}}\subseteq\mathcal A_{\mathrm{bad},0},
\qquad
\mathcal A_{\mathrm{test}}^c\subseteq\mathcal A_{\mathrm{bad},1}.
\]
Applying the testing bound and then these inclusions yields
\[
c_{\mathrm{test}}
\le
Q^{\mathrm{test}}_1(\mathcal A_{\mathrm{bad},1})
+
Q^{\mathrm{test}}_0(\mathcal A_{\mathrm{bad},0}),
\]
and hence
\[
\frac{c_{\mathrm{test}}}{2}
\le
\max_{\iota\in\{0,1\}}
Q^{\mathrm{test}}_\iota(\mathcal A_{\mathrm{bad},\iota}).
\]
On each large-error event, squared loss is at least
\((\delta_n/2)^2\). Therefore
\[
\begin{split}
\frac{c_{\mathrm{test}}\delta_n^2}{8}
&\le
\left(\frac{\delta_n}{2}\right)^2
\max_{\iota\in\{0,1\}}
Q^{\mathrm{test}}_\iota(\mathcal A_{\mathrm{bad},\iota})\\
&\le
\max_{\iota\in\{0,1\}}
E_{Q^{\mathrm{test}}_\iota}
\left[
(b_\kappa-\vartheta^{\mathrm{test}}_\iota)^2
\right]\\
&\le
\max_{\iota\in\{0,1\}}
E_{P^{\mathrm{test}}_{Z,\iota}}
\left[
(\widehat\theta-\vartheta^{\mathrm{test}}_\iota)^2
\right].
\end{split}
\]
Both testing laws belong to the zero-control subclass, so taking the
minimax infimum gives
\[
\frac{c_{\mathrm{par}}}{n}
=
\frac{c_{\mathrm{test}}\delta_n^2}{8}
\le
\mathfrak R^{Z,0}_{n,d,\epsilon}
\le
\mathfrak R^Z_{n,d,\epsilon},
\qquad
c_{\mathrm{par}}:=\frac{c_{\mathrm{test}}}{128}>0.
\]
% lean: CausalSmith.Stat.MarRareqLogfrontier.zeng_parametric_minimax_lower

\item \textit{Calibration and the large-sample bound within the published range.}
Choose an integer \(n_{\epsilon,1}\) satisfying
\[
n_{\epsilon,1}\ge\frac2{c_{\epsilon,\mathrm{dim}}},
\]
and define
\[
n_{\epsilon,*}
:=
\max\{n_{\epsilon,\mathrm{lit}},3,n_{\epsilon,1}\},
\qquad
c_{\epsilon,\mathrm{sat}}
:=
c_{\epsilon,\mathrm{lit}}
\left(\frac{c_{\epsilon,\mathrm{dim}}}{2}\right)^2,
\]
\[
c_\epsilon
:=
\min\left\{
\frac{c_{\mathrm{par}}}{n_{\epsilon,*}},
\min\{c_{\epsilon,\mathrm{lit}},c_{\epsilon,\mathrm{sat}}\}
\right\}>0.
\]
For all \(n,d\ge1\),
\[
0\le r^{\mathrm{fix}}_{n,d}\le1.
\]

First suppose \(n\ge n_{\epsilon,*}\). Then
\[
n\ge n_{\epsilon,\mathrm{lit}},
\qquad
n\ge3,
\qquad
\log n>1,
\qquad
n\log n>0.
\]
Furthermore, the choice of \(n_{\epsilon,1}\) gives
\[
2\le c_{\epsilon,\mathrm{dim}}\,n
\le c_{\epsilon,\mathrm{dim}}\,n\log n.
\]
If
\[
d\le c_{\epsilon,\mathrm{dim}}\,n\log n,
\]
the published bound and subclass inclusion imply
\[
c_{\epsilon,\mathrm{lit}}
\left\{
\frac1n+\left[\frac d{n\log n}\right]^2
\right\}
\le
\mathfrak R^{Z,0}_{n,d,\epsilon}
\le
\mathfrak R^Z_{n,d,\epsilon}.
\]
Since
\[
0<\log n\le\log(e+n),
\qquad
0\le\frac d{n\log(e+n)}
\le\frac d{n\log n},
\]
we have
\[
r^{\mathrm{fix}}_{n,d}
\le
\frac1n+\left[\frac d{n\log(e+n)}\right]^2
\le
\frac1n+\left[\frac d{n\log n}\right]^2.
\]
Using \(c_\epsilon\le c_{\epsilon,\mathrm{lit}}\) gives
\[
\begin{split}
c_\epsilon r^{\mathrm{fix}}_{n,d}
&\le
c_\epsilon
\left\{
\frac1n+\left[\frac d{n\log n}\right]^2
\right\}\\
&\le
c_{\epsilon,\mathrm{lit}}
\left\{
\frac1n+\left[\frac d{n\log n}\right]^2
\right\}
\le
\mathfrak R^Z_{n,d,\epsilon}.
\end{split}
\]
% lean: CausalSmith.Stat.MarRareqLogfrontier.zeng_minimax_lower_fixedQ

\item \textit{The large-alphabet case.}
Continue to assume \(n\ge n_{\epsilon,*}\), and now suppose
\[
d>c_{\epsilon,\mathrm{dim}}\,n\log n.
\]
Define the smaller alphabet size by
\[
d_{\epsilon,n}^{\mathrm{pad}}
:=
\left\lfloor c_{\epsilon,\mathrm{dim}}\,n\log n\right\rfloor.
\]
The preceding lower bound of two on the quantity inside the floor gives
\[
1\le d_{\epsilon,n}^{\mathrm{pad}}\le d,
\qquad
d_{\epsilon,n}^{\mathrm{pad}}
\le c_{\epsilon,\mathrm{dim}}\,n\log n.
\]
Also, the elementary floor bound gives
\[
\begin{split}
\frac{c_{\epsilon,\mathrm{dim}}}{2}\,n\log n
&\le c_{\epsilon,\mathrm{dim}}\,n\log n-1\\
&<d_{\epsilon,n}^{\mathrm{pad}}.
\end{split}
\]
The published lower bound therefore applies at this smaller alphabet:
\[
c_{\epsilon,\mathrm{lit}}
\left\{
\frac1n+
\left[
\frac{d_{\epsilon,n}^{\mathrm{pad}}}{n\log n}
\right]^2
\right\}
\le
\mathfrak R^{Z,0}_{n,d_{\epsilon,n}^{\mathrm{pad}},\epsilon}
\le
\mathfrak R^Z_{n,d_{\epsilon,n}^{\mathrm{pad}},\epsilon}.
\]

To compare the two alphabet sizes, extend each law
\(P_Z\in\mathcal D_{d_{\epsilon,n}^{\mathrm{pad}},\epsilon}\)
by zero masses. Explicitly, define the \(d\)-category law
\(\operatorname{pad}_{d_{\epsilon,n}^{\mathrm{pad}},d}P_Z\) by
\[
\bigl(\operatorname{pad}_{d_{\epsilon,n}^{\mathrm{pad}},d}P_Z\bigr)
\{(x,b,y)\}
:=
\begin{cases}
P_Z\{(x,b,y)\},&1\le x\le d_{\epsilon,n}^{\mathrm{pad}},\\
0,&d_{\epsilon,n}^{\mathrm{pad}}<x\le d,
\end{cases}
\qquad b,y\in\{0,1\}.
\]
On the original categories, all baseline, treatment, and outcome
probabilities are preserved; the added categories have zero mass.
Thus
\[
\operatorname{pad}_{d_{\epsilon,n}^{\mathrm{pad}},d}P_Z
\in\mathcal D_{d,\epsilon},
\qquad
\theta_Z\!\left(
\operatorname{pad}_{d_{\epsilon,n}^{\mathrm{pad}},d}P_Z
\right)
=\theta_Z(P_Z).
\]
For a \(d\)-category estimator \(\widehat\theta\), define
\(\widehat\theta^{\mathrm{res}}\) on the smaller sample space by
\[
\widehat\theta^{\mathrm{res}}
\bigl((x_i,b_i,y_i)_{i=1}^n\bigr)
:=
\widehat\theta
\bigl((x_i,b_i,y_i)_{i=1}^n\bigr),
\qquad
x_i\in\{1,\ldots,d_{\epsilon,n}^{\mathrm{pad}}\},
\]
preserving its conditional action distribution if it is randomized.
The product sample law is preserved under this inclusion, so
\[
\begin{split}
&E_{P_Z}\!\left[
\left(\widehat\theta^{\mathrm{res}}-\theta_Z(P_Z)\right)^2
\right]\\
&\qquad=
E_{\operatorname{pad}_{d_{\epsilon,n}^{\mathrm{pad}},d}P_Z}
\!\left[
\left(
\widehat\theta-
\theta_Z\!\left(
\operatorname{pad}_{d_{\epsilon,n}^{\mathrm{pad}},d}P_Z
\right)
\right)^2
\right].
\end{split}
\]
Taking suprema and then infima proves
\[
\mathfrak R^Z_{n,d_{\epsilon,n}^{\mathrm{pad}},\epsilon}
\le
\mathfrak R^Z_{n,d,\epsilon}.
\]

Since \(n\log n>0\), the floor bound also gives
\[
\frac{c_{\epsilon,\mathrm{dim}}}{2}
\le
\frac{d_{\epsilon,n}^{\mathrm{pad}}}{n\log n}.
\]
Squaring and using \(1/n\ge0\), we obtain
\[
c_{\epsilon,\mathrm{sat}}
\le
c_{\epsilon,\mathrm{lit}}
\left\{
\frac1n+
\left[
\frac{d_{\epsilon,n}^{\mathrm{pad}}}{n\log n}
\right]^2
\right\}.
\]
Finally, \(r^{\mathrm{fix}}_{n,d}\le1\) and
\(c_\epsilon\le c_{\epsilon,\mathrm{sat}}\) give
\[
\begin{split}
c_\epsilon r^{\mathrm{fix}}_{n,d}
&\le c_\epsilon
\le c_{\epsilon,\mathrm{sat}}\\
&\le
c_{\epsilon,\mathrm{lit}}
\left\{
\frac1n+
\left[
\frac{d_{\epsilon,n}^{\mathrm{pad}}}{n\log n}
\right]^2
\right\}\\
&\le
\mathfrak R^{Z,0}_{n,d_{\epsilon,n}^{\mathrm{pad}},\epsilon}
\le
\mathfrak R^Z_{n,d_{\epsilon,n}^{\mathrm{pad}},\epsilon}
\le
\mathfrak R^Z_{n,d,\epsilon}.
\end{split}
\]
% lean: CausalSmith.Stat.MarRareqLogfrontier.zeng_minimax_lower_fixedQ

\item \textit{Small samples and the final constants.}
If \(1\le n<n_{\epsilon,*}\), the parametric bound and the calibration
of \(c_\epsilon\) imply
\[
c_\epsilon r^{\mathrm{fix}}_{n,d}
\le c_\epsilon
\le\frac{c_{\mathrm{par}}}{n_{\epsilon,*}}
\le\frac{c_{\mathrm{par}}}{n}
\le\mathfrak R^Z_{n,d,\epsilon}.
\]
Together with the two large-sample cases, this establishes the lower
bound for every \(n,d\ge1\). Define
\[
C_\epsilon:=\max\{C_{\epsilon,\mathrm{up}},c_\epsilon\}.
\]
Then \(0<c_\epsilon\le C_\epsilon<\infty\). Since
\(r^{\mathrm{fix}}_{n,d}\ge0\), the upper bound already proved gives
\[
\mathfrak R^Z_{n,d,\epsilon}
\le
C_{\epsilon,\mathrm{up}}r^{\mathrm{fix}}_{n,d}
\le
C_\epsilon r^{\mathrm{fix}}_{n,d}.
\]
Thus, for every \(n,d\ge1\),
\[
c_\epsilon
\min\left\{
1,\frac1n+\left[\frac d{n\log(e+n)}\right]^2
\right\}
\le
\mathfrak R^Z_{n,d,\epsilon}
\le
C_\epsilon
\min\left\{
1,\frac1n+\left[\frac d{n\log(e+n)}\right]^2
\right\},
\]
as claimed.
% lean: CausalSmith.Stat.MarRareqLogfrontier.zeng_fixed_positivity_polynomial_frontier
\end{enumerate}
\end{proof}

\begin{proof}[Proof of \cref{obj:thm:unrestricted-uniform-mixed-count-certificate}]
We take \(\overline C=10^{30}\). We first compare the deterministic envelope with the rate, and then establish the risk bound.

\begin{enumerate}
\item Define the untruncated rate expression by
\[
r^\circ_{n,d,q}
:=N^{-1}+\left(\frac{d}{N\ell}\right)^2,
\qquad
N=nq,\quad m=\frac n6,\quad N_0=\frac N6,
\quad \ell=\log(e+N).
\]
The domain assumptions give
\[
0<N\le n,\qquad m>0,\qquad N_0>0,
\qquad r_{n,d,q}=\min\{1,r^\circ_{n,d,q}\}.
\]
If \(N\le1\), then \(N^{-1}\ge1\), so
\[
\mathfrak U_{n,d,q}=r_{n,d,q}=1
\le10^{30}r_{n,d,q}.
\]
For the envelope comparison, henceforth suppose \(N>1\).
% lean: CausalSmith.Stat.MarRareqLogfrontier.riskEnvelope_le_frontierRate

\item The logarithmic scale and stream sizes satisfy
\[
1<\ell,\qquad e^\ell=e+N,\qquad e^{-\ell}\le N^{-1},
\qquad
\frac1{N_0}=6N^{-1},\qquad
\frac1m=\frac6n\le6N^{-1}.
\]
Indeed, \(e+N>e\), and inversion of \(e^\ell=e+N>N\) gives the exponential bound.

The elementary numerical bound for \(\log2\) gives
\[
\gamma_0=\log2-\frac12\ge\frac18.
\]
For every real \(t_{\mathrm{dec}}>0\), the inequality
\(t_{\mathrm{dec}}\le e^{t_{\mathrm{dec}}}\) implies
\(e^{-t_{\mathrm{dec}}}\le t_{\mathrm{dec}}^{-1}\). Applying this with
\(t_{\mathrm{dec}}=n\gamma_0\) yields
\[
e^{-n\gamma_0}
\le\frac1{n\gamma_0}
\le\frac8n
\le8N^{-1}.
\]
% lean: CausalSmith.Stat.MarRareqLogfrontier.poisson_tail_rate

\item Suppose first that \(\ell\ge128\) and \(d\ge\sqrt N\). On this branch,
\[
k=\lfloor\ell/64\rfloor,\qquad B=256\ell,
\qquad
\frac{\ell}{128}\le k\le\frac{\ell}{64}.
\]
For completeness, for every real \(t_{\mathrm{floor}}\ge1\),
\[
\lfloor t_{\mathrm{floor}}\rfloor\ge\frac{t_{\mathrm{floor}}}{2}.
\]
If \(t_{\mathrm{floor}}<2\), its floor is at least one; if
\(t_{\mathrm{floor}}\ge2\), use
\(\lfloor t_{\mathrm{floor}}\rfloor>t_{\mathrm{floor}}-1
\ge t_{\mathrm{floor}}/2\). This proves the lower bound on \(k\).

The inequality \(2t_{\mathrm{exp}}\le e^{t_{\mathrm{exp}}}\), for real
\(t_{\mathrm{exp}}\), gives
\[
\ell\le16e^{\ell/32},
\qquad
\ell^4\le65536e^{\ell/8}.
\]
Since \(e<3\) and \(N>1\),
\[
e^\ell=e+N\le4N\le4d^2,
\qquad e^{\ell/2}\le2d.
\]
Consequently,
\[
\begin{aligned}
(e^{\ell/8}+1)\ell^4
&\le2e^{\ell/8}\,65536e^{\ell/8}\\
&=131072e^{\ell/4}
\le131072e^{\ell/2}
\le262144d,
\end{aligned}
\qquad
\ell\le e^{\ell/2}\le2d.
\]
% lean: CausalSmith.Stat.MarRareqLogfrontier.needle_growth

\item On the same branch, the preceding bounds imply
\[
e^{-16\ell}\le N^{-1},
\qquad e^{-32\ell}\le N^{-1},
\qquad N^{-1}\le\frac{2d}{N\ell}.
\]
Moreover,
\[
\begin{aligned}
\frac{4dB}{N_0k^2}
&=\frac{6144d\ell}{Nk^2}\\
&\le\frac{6144d\ell}{N(\ell/128)^2}
=100663296\,\frac{d}{N\ell}.
\end{aligned}
\]
Thus
\[
\frac{4dB}{N_0k^2}+3e^{-16\ell}
\le100663302\,\frac{d}{N\ell},
\]
and, because both sides are nonnegative,
\[
b_{n,d,q}^2
\le40532401478172816
\left(\frac{d}{N\ell}\right)^2.
\]
% lean: CausalSmith.Stat.MarRareqLogfrontier.riskEnvelope_le_frontierRate

\item Since \(N_0\le m\) and \(B\ge1\),
\[
\frac{B}{mN_0}
\le\frac{B}{N_0^2}
\le\frac{B^2}{N_0^2}.
\]
Therefore
\[
\frac{B^2}{N_0^2}+\frac{B}{mN_0}
\le\frac{2B^2}{N_0^2}
=4718592\,\frac{\ell^2}{N^2}.
\]
The growth estimate above also gives
\[
\frac{d(e^{\ell/8}+1)\ell^2}{N^2}
=
\frac{d(e^{\ell/8}+1)\ell^4}{N^2\ell^2}
\le262144\left(\frac{d}{N\ell}\right)^2.
\]
Combining these inequalities,
\[
64d(e^{\ell/8}+1)
\left(\frac{B^2}{N_0^2}+\frac{B}{mN_0}\right)
\le79164837199872
\left(\frac{d}{N\ell}\right)^2.
\]
The remaining variance terms satisfy
\[
8\left(\frac4{N_0}+\frac1m\right)\le240N^{-1},
\]
and, using \(N^{-1}\le1\) and \(1/m\le6\),
\[
32e^{-32\ell}\left(1+\frac1m\right)\le224N^{-1}.
\]
Together with the prefix-tail bound, this proves
\[
b_{n,d,q}^2+v_{n,d,q}+4e^{-n\gamma_0}
\le40611566315372688
\left(\frac{d}{N\ell}\right)^2+496N^{-1}
\le10^{30}r^\circ_{n,d,q}.
\]
If \(r^\circ_{n,d,q}\le1\), take the minimum with four on the left.
If \(r^\circ_{n,d,q}>1\), use
\(\mathfrak U_{n,d,q}\le4\le10^{30}\). In either case,
\[
\mathfrak U_{n,d,q}\le10^{30}r_{n,d,q}.
\]
% lean: CausalSmith.Stat.MarRareqLogfrontier.riskEnvelope_le_frontierRate

\item On the remaining branch, \(\ell<128\) or \(d<\sqrt N\). Since \(e\ge1\),
\[
0\le\frac{8d}{eN_0}\le\frac{48d}{N},
\qquad
\left(\frac{8d}{eN_0}\right)^2
\le2304\left(\frac dN\right)^2.
\]
If \(\ell<128\), then
\[
\left(\frac dN\right)^2
=\ell^2\left(\frac{d}{N\ell}\right)^2
\le16384\left(\frac{d}{N\ell}\right)^2.
\]
If \(d<\sqrt N\), then \(d^2\le N\), so
\[
\left(\frac dN\right)^2\le N^{-1}.
\]
Hence throughout this branch,
\[
\left(\frac dN\right)^2
\le N^{-1}+16384\left(\frac{d}{N\ell}\right)^2.
\]
Also,
\[
4\left(\frac4{N_0}+\frac1m\right)\le120N^{-1}.
\]
It follows that
\[
\begin{aligned}
\left(\frac{8d}{eN_0}\right)^2
+4\left(\frac4{N_0}+\frac1m\right)+4e^{-n\gamma_0}
&\le2456N^{-1}
+37748736\left(\frac{d}{N\ell}\right)^2\\
&\le10^{30}r^\circ_{n,d,q}.
\end{aligned}
\]
The same split at \(r^\circ_{n,d,q}=1\) proves
\[
\mathfrak U_{n,d,q}\le10^{30}r_{n,d,q}.
\]
This completes the deterministic comparison over the stated domain.
% lean: CausalSmith.Stat.MarRareqLogfrontier.riskEnvelope_le_frontierRate

\item Fix now \(P\in\mathcal M^+_{n,d,q}\). By
\cref{obj:lem:unrestricted-cell-identification,obj:def:ate-functional},
\[
\Phi(P)=\tau(P)
=P(Y(1)=1)-P(Y(0)=1)\in[-1,1].
\]
The auxiliary output in \cref{obj:def:mixed-count-estimator} lies in
\([-1,1]\). Its conditional average lies in that interval as well: for each
prefix length \(\nu\le n\), there are \(3^\nu\) equally weighted stream
assignments, and
\[
\Pr(K_0=\nu)\ge0,
\qquad
\sum_{\nu=0}^{n}\Pr(K_0=\nu)\le1,
\]
with the remaining probability assigned output zero. Thus
\[
\widehat\tau^{\mathrm{mix}}_{n,d,q}(\mathbf O_n)\in[-1,1],
\qquad
E_P\!\left[
(\widehat\tau^{\mathrm{mix}}_{n,d,q}-\Phi(P))^2
\right]\le4.
\]
If \(N\le1\), the auxiliary output and its conditional average are zero, so
\[
E_P\!\left[
(\widehat\tau^{\mathrm{mix}}_{n,d,q}-\Phi(P))^2
\right]=\Phi(P)^2\le1=\mathfrak U_{n,d,q}.
\]
If \(N>1\) and \(\mathfrak U_{n,d,q}=4\), the displayed bound by four proves
the required inequality. We henceforth consider
\(N>1\) and \(\mathfrak U_{n,d,q}\ne4\).
% lean: CausalSmith.Stat.MarRareqLogfrontier.unrestricted_auxiliary_risk_envelope

\item We first establish the transfer to independent Poisson streams.
Write
\[
\mathsf Q_{\mathrm{str}}
:=
\bigotimes_{h_{\mathrm{str}}\in\{0,1,2\}}
\left[
\sum_{\nu=0}^{\infty}
e^{-m}\frac{m^\nu}{\nu!}
\bigl(\mathcal L_P(O_1)\bigr)^{\otimes\nu}
\right],
\]
where each sum is a probability law on the disjoint union of finite
observed-record sequences. The stream labels \(0,1,2\) denote membership,
pilot, and estimation, respectively.

An uncapped prefix of length \(K_0\sim\operatorname{Poisson}(n/2)\), marked
independently with these three labels, has law \(\mathsf Q_{\mathrm{str}}\)
after separation into streams. Indeed, for any prescribed nonnegative
integer stream sizes \(\nu_0,\nu_1,\nu_2\), put
\(\nu_{\mathrm{tot}}:=\nu_0+\nu_1+\nu_2\). Their joint probability is
\[
e^{-n/2}\frac{(n/2)^{\nu_{\mathrm{tot}}}}{\nu_{\mathrm{tot}}!}
\frac{\nu_{\mathrm{tot}}!}{\nu_0!\nu_1!\nu_2!}
3^{-\nu_{\mathrm{tot}}}
=
\prod_{h_{\mathrm{str}}=0}^{2}
e^{-m}\frac{m^{\nu_{h_{\mathrm{str}}}}}
{\nu_{h_{\mathrm{str}}}!}.
\]
Conditional on these sizes, the records in the three streams have their
respective product observed-data laws. Equality on all count fibres proves
the asserted stream law.

Let
\[
\widetilde\tau_{\mathrm{Pois}}
:=
\operatorname{clip}_{[-1,1]}
\left(2\sum_{j=(a,x,s)\in\mathcal J_d}g(a)W_jG_j\right)
\]
be the uncapped output, using the statistics of
\cref{obj:synth_9,obj:synth_10,obj:synth_11,obj:synth_12,obj:synth_14,obj:synth_15,obj:synth_17}
on these finite streams. On the ratio branch, \(G_j=D_j\).
Here and below \(E_{\mathrm{Pois}}\) and
\(\Pr_{\mathrm{Pois}}\) denote expectation and probability under
\(\mathsf Q_{\mathrm{str}}\).

For the capped auxiliary output, conditional averaging gives
\[
\begin{aligned}
&E_\omega\!\left[
(\widetilde\tau(\mathbf O_n;\omega)-\Phi(P))^2
\mid\mathbf O_n\right]\\
&\quad=
(\widehat\tau^{\mathrm{mix}}_{n,d,q}(\mathbf O_n)-\Phi(P))^2
+
E_\omega\!\left[
(\widetilde\tau(\mathbf O_n;\omega)
-\widehat\tau^{\mathrm{mix}}_{n,d,q}(\mathbf O_n))^2
\mid\mathbf O_n\right].
\end{aligned}
\]
On each prefix-count fibre \(\nu\le n\), its stream law agrees with the
corresponding uncapped law. Summing these fibres, discarding the
nonnegative uncapped loss on the remaining fibres, and bounding the
overflow loss by four yields
\[
E_P\!\left[
(\widehat\tau^{\mathrm{mix}}_{n,d,q}-\Phi(P))^2
\right]
\le
E_{\mathrm{Pois}}\!\left[
(\widetilde\tau_{\mathrm{Pois}}-\Phi(P))^2
\right]+4\Pr(K_0>n).
\]
The Poisson probability series evaluates the exponential moment as
\[
E[2^{K_0}]
=e^{-n/2}\sum_{\nu=0}^{\infty}\frac{n^\nu}{\nu!}
=e^{n/2}.
\]
Since \(\{K_0>n\}\subseteq\{2^{K_0}\ge2^n\}\), exponential Markov gives
\[
\Pr(K_0>n)
\le\Pr(2^{K_0}\ge2^n)
\le2^{-n}E[2^{K_0}]
=e^{-n\log2+n/2}.
\]
Multiplying by \(e^{-n/2}\) cancels the intensity factor, so
\[
\Pr(K_0>n)e^{-n/2}
\le e^{-n\log2},
\qquad
\Pr(K_0>n)\le e^{-n\gamma_0}.
\]
Consequently,
\[
E_P\!\left[
(\widehat\tau^{\mathrm{mix}}_{n,d,q}-\Phi(P))^2
\right]
\le
E_{\mathrm{Pois}}\!\left[
(\widetilde\tau_{\mathrm{Pois}}-\Phi(P))^2
\right]+4e^{-n\gamma_0}.
\]
All these expectations are well defined; the capped sample space is finite
and the clipped outputs are bounded.
% lean: CausalSmith.Stat.MarRareqLogfrontier.deterministicRisk_mixedCountEstimator_le_streamRisk_add_tail

\item For every \(j=(a,x,s)\in\mathcal J_d\), define
\[
p_j:=p_{axs}(P),\qquad
\mu_j:=\mu_{axs}(P),\qquad
t_j:=mP(A=a,X=x,S=s,R=1).
\]
These definitions include null cells, with the mean convention of
\cref{obj:def:cell-functional}. Probability monotonicity and the
occupied-cell arrival condition give
\[
0\le\mu_j\le1,\qquad
t_j\ge N_0p_j,\qquad
\sum_jp_j=1,\qquad |\mathcal J_d|=4d.
\]
For \(p_j=0\), the membership, arrival, and arrived-one probabilities are
zero, and all the associated counts vanish almost surely.

Applying the same Poisson partition calculation to stream and cell
categories gives independent counts across distinct cells and streams,
with
\[
M_j\sim\operatorname{Poisson}(mp_j),
\qquad
C'_j,C_j\sim\operatorname{Poisson}(t_j).
\]
The estimation arrived-one and arrived-zero intensities are, respectively,
\(t_j\mu_j\) and \(t_j(1-\mu_j)\). Poisson moments therefore yield
\[
E_{\mathrm{Pois}}W_j=p_j,\qquad
E_{\mathrm{Pois}}W_j^2=p_j^2+\frac{p_j}{m}.
\]
Also, the membership statistic is independent of every statistic formed
from the pilot and estimation streams.
% lean: CausalSmith.Stat.MarRareqLogfrontier.integral_poissonMemberWeight_mul_selectedCellEstimate

\item Consider the polynomial branch
\(\ell\ge128\), \(d\ge\sqrt N\). Define the pilot-selection event by
\[
\mathsf L_j:=\{C'_j\le B/4\}.
\]
The coefficient and cell-statistic definitions in
\cref{obj:synth_13,obj:synth_14,obj:synth_15} give
\[
1-Q_k(t)=\sum_{v=1}^{k-1}a_vt^v,
\qquad
G_j=\mathbf1_{\mathsf L_j}H_j+
\mathbf1_{\mathsf L_j^c}D_j
\quad(t\ge0).
\]
To justify the expansion at zero and its degree, the numerator
\(1-T_k(1-2t/B)\) vanishes at zero and has degree at most \(k\).
The recurrence for \(T_k\), differentiated at one, gives
\[
T_0'(1)=0,\qquad T_1'(1)=1,\qquad
T_{h+1}'(1)=2+2T_h'(1)-T_{h-1}'(1)
\quad(h\ge1).
\]
Induction yields \(T_h'(1)=h^2\). Division of the numerator by \(t\)
therefore gives
\[
Q_k(0)=1,\qquad \deg(1-Q_k)\le k-1,
\]
which proves the displayed finite expansion.

For integers \(v\ge1\), write
\[
(C_j)_v:=\prod_{h=0}^{v-1}\max\{C_j-h,0\},
\qquad
\mathsf F_{j,v}:=
U_j\prod_{h=1}^{v-1}\max\{C_j-h,0\}.
\]
Empty products equal one. Counting ordered distinct tuples with the first
record an arrived one and the remaining records arrived gives
\[
E_{\mathrm{Pois}}\mathsf F_{j,v}
=t_j^v\mu_j.
\]
More explicitly, condition on the marked prefix size, count its ordered
distinct \(v\)-tuples, and then use the Poisson factorial moment
\(E(K_0)_v=(n/2)^v\). The one-record probabilities are
\(P(A=a,X=x,S=s,R=1,Y=1)/3\) for the distinguished record and
\(P(A=a,X=x,S=s,R=1)/3\) for the others. Their product gives the stated
identity. When the arrived-cell probability is zero, both sides vanish.
Thus
\[
E_{\mathrm{Pois}}H_j
=\mu_j\{1-Q_k(t_j)\}.
\]

For the ratio statistic, conditioning on the pattern of estimation
arrivals gives mean \(\mu_j\) when \(C_j>0\); its zero-count value is zero.
Consequently,
\[
E_{\mathrm{Pois}}D_j=\mu_j(1-e^{-t_j}),
\qquad 0\le D_j\le1.
\]
Pilot independence now gives the exact selected-statistic bias:
\[
E_{\mathrm{Pois}}G_j-\mu_j
=-\mu_j\left\{
\Pr_{\mathrm{Pois}}(\mathsf L_j)Q_k(t_j)
+\Pr_{\mathrm{Pois}}(\mathsf L_j^c)e^{-t_j}
\right\}.
\]
% lean: CausalSmith.Stat.MarRareqLogfrontier.markedPoissonSelectedCellEstimate_bias

\item We next derive the coefficient bound needed for the polynomial
moments. For a real polynomial \(F_{\mathrm{pol}}\), define
\[
\|F_{\mathrm{pol}}\|_{\mathrm{coef}}
:=\sum_{h\ge0}|[t_{\mathrm{pol}}^h]F_{\mathrm{pol}}|,
\]
where \(t_{\mathrm{pol}}\) is the polynomial indeterminate and the sum is
finite. For a second real polynomial and coefficient indices, take
\[
F_{\mathrm{in}}\in\mathbb R[t_{\mathrm{pol}}],
\qquad
h_{\mathrm{coef}},i_{\mathrm{coef}},j_{\mathrm{coef}}\in\mathbb N_0.
\]
All coefficient sums below are finite, with coefficients beyond the
polynomial degree equal to zero. The coefficientwise triangle inequality
 gives
\[
\begin{aligned}
\|F_{\mathrm{pol}}+F_{\mathrm{in}}\|_{\mathrm{coef}}
&=\sum_{h_{\mathrm{coef}}\ge0}
\left|[t_{\mathrm{pol}}^{h_{\mathrm{coef}}}]F_{\mathrm{pol}}
+[t_{\mathrm{pol}}^{h_{\mathrm{coef}}}]F_{\mathrm{in}}\right|\\
&\le\sum_{h_{\mathrm{coef}}\ge0}
\left|[t_{\mathrm{pol}}^{h_{\mathrm{coef}}}]F_{\mathrm{pol}}\right|
+\sum_{h_{\mathrm{coef}}\ge0}
\left|[t_{\mathrm{pol}}^{h_{\mathrm{coef}}}]F_{\mathrm{in}}\right|\\
&=\|F_{\mathrm{pol}}\|_{\mathrm{coef}}
+\|F_{\mathrm{in}}\|_{\mathrm{coef}}.
\end{aligned}
\]
Applying this inequality successively to the finite monomial expansion
of the product yields
\[
\begin{aligned}
\|F_{\mathrm{pol}}F_{\mathrm{in}}\|_{\mathrm{coef}}
&=\left\|
\sum_{i_{\mathrm{coef}}\ge0}\sum_{j_{\mathrm{coef}}\ge0}
\bigl([t_{\mathrm{pol}}^{i_{\mathrm{coef}}}]F_{\mathrm{pol}}\bigr)
\bigl([t_{\mathrm{pol}}^{j_{\mathrm{coef}}}]F_{\mathrm{in}}\bigr)
 t_{\mathrm{pol}}^{i_{\mathrm{coef}}+j_{\mathrm{coef}}}
\right\|_{\mathrm{coef}}\\
&\le\sum_{i_{\mathrm{coef}}\ge0}\sum_{j_{\mathrm{coef}}\ge0}
\left\|
\bigl([t_{\mathrm{pol}}^{i_{\mathrm{coef}}}]F_{\mathrm{pol}}\bigr)
\bigl([t_{\mathrm{pol}}^{j_{\mathrm{coef}}}]F_{\mathrm{in}}\bigr)
 t_{\mathrm{pol}}^{i_{\mathrm{coef}}+j_{\mathrm{coef}}}
\right\|_{\mathrm{coef}}\\
&=\sum_{i_{\mathrm{coef}}\ge0}\sum_{j_{\mathrm{coef}}\ge0}
\left|[t_{\mathrm{pol}}^{i_{\mathrm{coef}}}]F_{\mathrm{pol}}\right|
\left|[t_{\mathrm{pol}}^{j_{\mathrm{coef}}}]F_{\mathrm{in}}\right|\\
&=\left(\sum_{i_{\mathrm{coef}}\ge0}
\left|[t_{\mathrm{pol}}^{i_{\mathrm{coef}}}]F_{\mathrm{pol}}\right|\right)
\left(\sum_{j_{\mathrm{coef}}\ge0}
\left|[t_{\mathrm{pol}}^{j_{\mathrm{coef}}}]F_{\mathrm{in}}\right|\right)\\
&=\|F_{\mathrm{pol}}\|_{\mathrm{coef}}
\|F_{\mathrm{in}}\|_{\mathrm{coef}}.
\end{aligned}
\]
Here the norm of each monomial is the absolute value of its coefficient,
by the defining coefficient sum. For every real polynomial \(F_{\mathrm{in}}\), submultiplicativity gives by induction
\[
\|F_{\mathrm{in}}^h\|_{\mathrm{coef}}
\le\|F_{\mathrm{in}}\|_{\mathrm{coef}}^h
\qquad(h\in\mathbb N_0).
\]
For \(F_{\mathrm{pol}}\ne0\), expand the composition as the finite sum
\[
F_{\mathrm{pol}}\circ F_{\mathrm{in}}
=
\sum_{h=0}^{\deg F_{\mathrm{pol}}}
\bigl([t_{\mathrm{pol}}^h]F_{\mathrm{pol}}\bigr)F_{\mathrm{in}}^h.
\]
For the composition bound, assume
\(\|F_{\mathrm{in}}\|_{\mathrm{coef}}\ge1\).
Subadditivity and the power bound then give
\[
\begin{aligned}
\|F_{\mathrm{pol}}\circ F_{\mathrm{in}}\|_{\mathrm{coef}}
&\le
\sum_{h=0}^{\deg F_{\mathrm{pol}}}
|[t_{\mathrm{pol}}^h]F_{\mathrm{pol}}|
\|F_{\mathrm{in}}^h\|_{\mathrm{coef}}\\
&\le
\sum_{h=0}^{\deg F_{\mathrm{pol}}}
|[t_{\mathrm{pol}}^h]F_{\mathrm{pol}}|
\|F_{\mathrm{in}}\|_{\mathrm{coef}}^h\\
&\le
\sum_{h=0}^{\deg F_{\mathrm{pol}}}
|[t_{\mathrm{pol}}^h]F_{\mathrm{pol}}|
\|F_{\mathrm{in}}\|_{\mathrm{coef}}^{\deg F_{\mathrm{pol}}}\\
&=
\|F_{\mathrm{pol}}\|_{\mathrm{coef}}
\|F_{\mathrm{in}}\|_{\mathrm{coef}}^{\deg F_{\mathrm{pol}}}.
\end{aligned}
\]
When \(F_{\mathrm{pol}}=0\), the composition has coefficient norm zero.
The Chebyshev recurrence implies
\[
\|T_{h+2}\|_{\mathrm{coef}}
\le2\|T_{h+1}\|_{\mathrm{coef}}+\|T_h\|_{\mathrm{coef}},
\qquad
\|T_0\|_{\mathrm{coef}}=\|T_1\|_{\mathrm{coef}}=1.
\]
Since \((1+\sqrt2)^2=2(1+\sqrt2)+1\), induction gives
\[
\|T_h\|_{\mathrm{coef}}\le(1+\sqrt2)^h
\quad(h\in\mathbb N).
\]
Define
\[
\mathcal P_k(t_{\mathrm{pol}})
:=1-T_k(1-2t_{\mathrm{pol}}).
\]
Then
\[
\|\mathcal P_k\|_{\mathrm{coef}}
\le1+(1+\sqrt2)^k3^k\le1+8^k,
\]
using \(\|1-2t_{\mathrm{pol}}\|_{\mathrm{coef}}=3\) and
\(\sqrt2\le3/2\). The scaled quotient satisfies
\[
Q_k(Bt_{\mathrm{pol}})
=\frac1{2k^2}\,
\frac{\mathcal P_k(t_{\mathrm{pol}})}{t_{\mathrm{pol}}},
\]
where division is exact because \(\mathcal P_k(0)=0\).
Removing the zero constant coefficient cannot increase the coefficient
norm. Hence, using \(k\ge2\) and \(8^k\ge64\),
\[
\begin{aligned}
\|1-Q_k(Bt_{\mathrm{pol}})\|_{\mathrm{coef}}
&\le1+\frac{1+8^k}{2k^2}\\
&\le1+\frac{1+8^k}{8}
\le8^k.
\end{aligned}
\]
In particular,
\[
\sum_{v=1}^{k-1}|a_vB^v|\le8^k.
\]
% lean: CausalSmith.Stat.MarRareqLogfrontier.needle_scaled_coefficient_oneNorm_le

\item For \(t\ge0\) and integers \(v,w\ge1\), define the factorial overlap
\[
\mathcal A_{v,w}(t)
:=
\sum_{h=0}^{\min\{v,w\}}
\binom vh\binom wh h!\,t^{v+w-h}.
\]
Counting the \(h\) shared entries of two ordered distinct tuples gives
\[
(C_j)_v(C_j)_w
=
\sum_{h=0}^{\min\{v,w\}}
\binom vh\binom wh h!(C_j)_{v+w-h}.
\]
Taking Poisson factorial moments yields
\[
E_{\mathrm{Pois}}[(C_j)_v(C_j)_w]
=\mathcal A_{v,w}(t_j).
\]
Since \(0\le\mathsf F_{j,v}\le(C_j)_v\), expansion of \(H_j^2\) gives
\[
E_{\mathrm{Pois}}H_j^2
\le
\sum_{v=1}^{k-1}\sum_{w=1}^{k-1}
|a_va_w|\mathcal A_{v,w}(t_j).
\]

For \(0\le t\le B\) and \(v,w\le k\), the bound
\[
\binom vh\binom wh h!\le\frac{(vw)^h}{h!}
\]
implies
\[
\begin{aligned}
\mathcal A_{v,w}(t)
&\le B^{v+w}
\sum_{h=0}^{\min\{v,w\}}
\frac{(k^2/B)^h}{h!}\\
&\le B^{v+w}e^{k^2/B}.
\end{aligned}
\]
Therefore, when \(t_j\le B\),
\[
E_{\mathrm{Pois}}H_j^2
\le e^{k^2/B}
\left(\sum_{v=1}^{k-1}|a_vB^v|\right)^2
\le e^{k^2/B}64^k.
\]
The tuning bounds give
\[
\frac{k^2}{B}\le\frac{\ell}{64},
\qquad 64^k=2^{6k}\le e^{6k},
\]
so
\[
E_{\mathrm{Pois}}H_j^2
\le e^{k^2/B+6k}\le e^{\ell/8}
\quad(t_j\le B).
\]
% lean: CausalSmith.Stat.MarRareqLogfrontier.integral_sq_markedPoissonFactorialBranch_le_exp_of_light

\item A second overlap estimate controls arbitrarily large intensities.
Using \(\binom whh!\le w^h\) and the binomial theorem,
\[
\begin{aligned}
\mathcal A_{v,w}(t)
&\le
\sum_{h=0}^{v}\binom vh w^h t^{v-h}t^w\\
&=(t+w)^vt^w
\le(t+v+w)^{v+w}.
\end{aligned}
\]
For \(v,w\le k\), split according as \((t+2k)/B\le1\) or exceeds one.
In both cases,
\[
\left(\frac{t+2k}{B}\right)^{v+w}
\le
\max\left\{1,\left(\frac{t+2k}{B}\right)^{2k}\right\}.
\]
The coefficient bound thus gives
\[
E_{\mathrm{Pois}}H_j^2
\le64^k
\max\left\{1,
\left(\frac{t_j+2k}{B}\right)^{2k}\right\}.
\]
% lean: CausalSmith.Stat.MarRareqLogfrontier.integral_sq_markedPoissonFactorialBranch_le_full_growth

\item Suppose \(t_j>B\). Define the integer pilot cutoff by
\[
\nu_{\mathrm{cut}}:=\lfloor B/4\rfloor.
\]
Then \(\nu_{\mathrm{cut}}<t_j/4\). The Poisson lower-tail estimate used here
follows from the centered exponential moment. To see its relevant form,
put
\[
a_{\mathrm{tail}}:=\frac1{\sqrt2}.
\]
The elementary exponential inequality
\[
e^{-a_{\mathrm{tail}}}
\le1-a_{\mathrm{tail}}+\frac{a_{\mathrm{tail}}^2}{2}
\]
follows by multiplying
\(e^{a_{\mathrm{tail}}}\ge
1+a_{\mathrm{tail}}+a_{\mathrm{tail}}^2/2\) by the nonnegative quadratic
on the right and using
\[
\left(1+a_{\mathrm{tail}}+\frac{a_{\mathrm{tail}}^2}{2}\right)
\left(1-a_{\mathrm{tail}}+\frac{a_{\mathrm{tail}}^2}{2}\right)
=1+\frac{a_{\mathrm{tail}}^4}{4}\ge1.
\]
For \(C'_j\sim\operatorname{Poisson}(t_j)\), exponential Markov therefore
gives
\[
\begin{aligned}
\Pr_{\mathrm{Pois}}\left(t_j-C'_j>\frac{t_j}{\sqrt2}\right)
&\le
\exp\left\{-\frac{t_j}{2}
+t_j(e^{-a_{\mathrm{tail}}}-1+a_{\mathrm{tail}})\right\}\\
&\le e^{-t_j/4}.
\end{aligned}
\]
On \(\{C'_j\le\nu_{\mathrm{cut}}\}\),
\(t_j-C'_j>3t_j/4>t_j/\sqrt2\). Consequently,
\[
\Pr_{\mathrm{Pois}}(\mathsf L_j)\le e^{-t_j/4}.
\]

Since \(t_j>B\), the maximum in the preceding moment bound selects its
second argument. The inequalities \(64^k\le e^{6k}\) and
\(t_{\mathrm{grow}}\le e^{t_{\mathrm{grow}}}\) for
\(t_{\mathrm{grow}}:=(t_j+2k)/B\) give
\[
\Pr_{\mathrm{Pois}}(\mathsf L_j)E_{\mathrm{Pois}}H_j^2
\le
\exp\left\{-\frac{t_j}{4}
+6k+\frac{2k(t_j+2k)}{B}\right\}.
\]
Here
\[
k\le\frac{t_j}{16384},
\qquad \frac{2k}{B}\le\frac1{8192},
\]
and hence
\[
6k+\frac{2k(t_j+2k)}B
\le
t_j\left(\frac6{16384}+\frac1{8192}
+\frac1{67108864}\right)
\le\frac{t_j}{8}.
\]
Thus
\[
\Pr_{\mathrm{Pois}}(\mathsf L_j)E_{\mathrm{Pois}}H_j^2
\le e^{-t_j/8}\le e^{-32\ell}.
\]
Also \(\Pr_{\mathrm{Pois}}(\mathsf L_j)\le e^{-32\ell}\).
% lean: CausalSmith.Stat.MarRareqLogfrontier.markedPoisson_pilot_light_mul_factorial_secondMoment_le_intensity

\item We now bound the two terms in the selected-statistic bias.
For \(0\le t\le B\), the polynomial quotient in
\cref{obj:synth_18} satisfies
\[
tQ_k(t)=\frac{B}{2k^2}\{1-T_k(1-2t/B)\}.
\]
Because \(|1-2t/B|\le1\), the cosine characterization of \(T_k\) gives
\(|T_k(1-2t/B)|\le1\), and therefore
\[
t|Q_k(t)|\le\frac B{k^2}.
\]
For \(t_j\le B\), use \(N_0p_j\mu_j\le t_j\) to obtain
\[
p_j\mu_j
\Pr_{\mathrm{Pois}}(\mathsf L_j)|Q_k(t_j)|
\le\frac{B}{N_0k^2}.
\]
This includes \(t_j=0\), because then \(p_j=0\).

For \(t_j>B\), Cauchy--Schwarz and the heavy moment estimate give
\[
\begin{aligned}
\left\{
\Pr_{\mathrm{Pois}}(\mathsf L_j)
|E_{\mathrm{Pois}}H_j|
\right\}^2
&\le
\Pr_{\mathrm{Pois}}(\mathsf L_j)
E_{\mathrm{Pois}}H_j^2\\
&\le e^{-t_j/8}.
\end{aligned}
\]
Using
\(\mu_jQ_k(t_j)=\mu_j-E_{\mathrm{Pois}}H_j\), it follows that
\[
\begin{aligned}
\mu_j\Pr_{\mathrm{Pois}}(\mathsf L_j)|Q_k(t_j)|
&\le
\Pr_{\mathrm{Pois}}(\mathsf L_j)
+\Pr_{\mathrm{Pois}}(\mathsf L_j)|E_{\mathrm{Pois}}H_j|\\
&\le2e^{-t_j/16}\le2e^{-16\ell}.
\end{aligned}
\]
Splitting the cell sum at \(t_j=B\), using \(|g(a)|=1\) from
\cref{obj:synth_19}, and then using \(4d\) cells and \(\sum_jp_j=1\), yields
\[
\left|
2\sum_jg(a)p_j\mu_j
\Pr_{\mathrm{Pois}}(\mathsf L_j)Q_k(t_j)
\right|
\le\frac{8dB}{N_0k^2}+4e^{-16\ell}.
\]
% lean: CausalSmith.Stat.MarRareqLogfrontier.abs_sum_pilot_needle_bias_le

\item The complementary-pilot term uses an upper-tail bound with the same
integer cutoff. The Poisson probability series gives
\[
E_{\mathrm{Pois}}[2^{C'_j}]
=e^{-t_j}\sum_{\nu=0}^{\infty}\frac{(2t_j)^\nu}{\nu!}
=e^{t_j}.
\]
The threshold inclusion and exponential Markov inequality imply
\[
\begin{aligned}
\Pr_{\mathrm{Pois}}(C'_j>\nu_{\mathrm{cut}})
&\le\Pr_{\mathrm{Pois}}(2^{C'_j}\ge2^{\nu_{\mathrm{cut}}})\\
&\le2^{-\nu_{\mathrm{cut}}}E_{\mathrm{Pois}}[2^{C'_j}]
=e^{-\nu_{\mathrm{cut}}\log2+t_j}.
\end{aligned}
\]
Multiplying by \(e^{-t_j}\) cancels the intensity factor and yields
\[
\Pr_{\mathrm{Pois}}(C'_j>\nu_{\mathrm{cut}})e^{-t_j}
\le e^{-\nu_{\mathrm{cut}}\log2}.
\]
Since
\[
64\ell<\nu_{\mathrm{cut}}+1,\qquad
\ell\ge128,\qquad \log2>\frac58,
\]
we have \(\nu_{\mathrm{cut}}\log2\ge16\ell\). Thus
\[
\Pr_{\mathrm{Pois}}(\mathsf L_j^c)e^{-t_j}\le e^{-16\ell},
\]
including \(t_j=0\), when the upper-tail probability is zero.
Consequently,
\[
\left|
2\sum_jg(a)p_j\mu_j
\Pr_{\mathrm{Pois}}(\mathsf L_j^c)e^{-t_j}
\right|
\le2e^{-16\ell}.
\]

Define the ideal raw contrast by
\[
\tau_{\mathrm{raw}}
:=2\sum_jg(a)W_jG_j.
\]
Membership independence gives
\(E_{\mathrm{Pois}}(W_jG_j)=p_jE_{\mathrm{Pois}}G_j\).
Combining the exact bias identity with the last two aggregate bounds,
\[
|E_{\mathrm{Pois}}\tau_{\mathrm{raw}}-\Phi(P)|
\le
\frac{8dB}{N_0k^2}+6e^{-16\ell}
=b_{n,d,q}.
\]
% lean: CausalSmith.Stat.MarRareqLogfrontier.abs_integral_raw_poissonSelected_estimator_bias_le

\item We next establish the ratio variance contribution.
For a fixed marked prefix length \(\nu\), define the estimation-arrival
index set by
\[
\mathsf I_j^{(\nu)}
:=
\{i\in\{1,\ldots,\nu\}:
i\text{ has estimation label and is an arrival in cell }j\}.
\]
On a prescribed pattern \(\mathsf I_j^{(\nu)}=\mathsf I_{\mathrm{arr}}\) with
\(\nu_{\mathrm{arr}}:=|\mathsf I_{\mathrm{arr}}|>0\), the restricted product law makes the
arrived-outcome indicators independent with mean \(\mu_j\). Their centered
first moments vanish, and their centered squares are at most one.
Expanding the square therefore gives, for patterns of positive
probability,
\[
E\!\left[
(U_j-\nu_{\mathrm{arr}}\mu_j)^2
\mid\mathsf I_j^{(\nu)}=\mathsf I_{\mathrm{arr}}
\right]\le\nu_{\mathrm{arr}},
\]
and hence
\[
E\!\left[
(D_j-\mu_j)^2
\mid\mathsf I_j^{(\nu)}=\mathsf I_{\mathrm{arr}}
\right]\le\nu_{\mathrm{arr}}^{-1}.
\]
Zero-probability patterns contribute zero. On the zero-count pattern,
\((D_j-\mu_j)^2=\mu_j^2\le1\). Summing first over arrival patterns and then
over the Poisson prefix size gives
\[
E_{\mathrm{Pois}}(D_j-\mu_j)^2
\le
E_{\mathrm{Pois}}\!\left[
\mathbf1_{\{C_j=0\}}
+\frac{\mathbf1_{\{C_j>0\}}}{C_j}
\right],
\]
where the quotient is zero at \(C_j=0\).

If \(t_j\le1\), the integrand on the right is at most one, so its expectation
is at most \(1\le4/(1+t_j)\). If \(t_j>1\), its value is at most
\(2/(C_j+1)\). For every nonnegative integer \(h_{\mathrm{count}}\), the Poisson probabilities satisfy
\[
\Pr_{\mathrm{Pois}}(C_j=h_{\mathrm{count}})
=e^{-t_j}\frac{t_j^{h_{\mathrm{count}}}}{h_{\mathrm{count}}!},
\qquad
\frac{t_j}{h_{\mathrm{count}}+1}
\Pr_{\mathrm{Pois}}(C_j=h_{\mathrm{count}})
=
\Pr_{\mathrm{Pois}}(C_j=h_{\mathrm{count}}+1).
\]
The probability series sums to one, and the reciprocal series is bounded termwise by it. Summing the recurrence and shifting the index therefore gives
\[
\begin{aligned}
t_jE_{\mathrm{Pois}}\frac1{C_j+1}
&=
\sum_{h_{\mathrm{count}}=0}^{\infty}
\frac{t_j}{h_{\mathrm{count}}+1}
\Pr_{\mathrm{Pois}}(C_j=h_{\mathrm{count}})\\
&=
\sum_{h_{\mathrm{count}}=0}^{\infty}
\Pr_{\mathrm{Pois}}(C_j=h_{\mathrm{count}}+1)\\
&=\Pr_{\mathrm{Pois}}(C_j>0)
=1-e^{-t_j}.
\end{aligned}
\]
Dividing by \(t_j>1\) yields
\[
E_{\mathrm{Pois}}\frac1{C_j+1}
=\frac{1-e^{-t_j}}{t_j}.
\]
Therefore
\[
E_{\mathrm{Pois}}(D_j-\mu_j)^2
\le\frac2{t_j}\le\frac4{1+t_j}
\quad(t_j>1).
\]
When the arrived-cell probability is zero, \(C_j=U_j=\mu_j=0\)
almost surely, so the centered second moment is zero. In all cases,
\[
\operatorname{Var}_{\mathrm{Pois}}(D_j)
=
\operatorname{Var}_{\mathrm{Pois}}(D_j-\mu_j)
\le E_{\mathrm{Pois}}(D_j-\mu_j)^2
\le\frac4{1+t_j}.
\]

Expanding the variance of the independent product \(W_jD_j\) gives
\[
\operatorname{Var}_{\mathrm{Pois}}(W_jD_j)
=p_j^2\operatorname{Var}_{\mathrm{Pois}}(D_j)
+\frac{p_j}{m}E_{\mathrm{Pois}}D_j^2.
\]
Using \(D_j^2\le1\) and \(N_0p_j\le t_j\),
\[
\frac{p_j^2}{1+t_j}\le\frac{p_j}{N_0},
\]
including \(p_j=0\). Summing over cells proves
\[
\sum_j\operatorname{Var}_{\mathrm{Pois}}(W_jD_j)
\le\frac4{N_0}+\frac1m.
\]
% lean: CausalSmith.Stat.MarRareqLogfrontier.sum_variance_poissonMemberWeight_mul_ratioBranch_le

\item On the polynomial branch, define the weighted correction by
\[
\Delta_j:=W_j(G_j-D_j)
=W_j\mathbf1_{\mathsf L_j}(H_j-D_j).
\]
Pilot and membership independence give
\[
E_{\mathrm{Pois}}\Delta_j^2
=
\left(p_j^2+\frac{p_j}{m}\right)
\Pr_{\mathrm{Pois}}(\mathsf L_j)
E_{\mathrm{Pois}}(H_j-D_j)^2.
\]
Since \((H_j-D_j)^2\le2(H_j^2+D_j^2)\), for \(t_j\le B\),
\[
E_{\mathrm{Pois}}\Delta_j^2
\le2\left(p_j^2+\frac{p_j}{m}\right)(e^{\ell/8}+1).
\]
Here \(p_j\le B/N_0\), so
\[
p_j^2+\frac{p_j}{m}
\le\frac{B^2}{N_0^2}+\frac{B}{mN_0}.
\]
Summing over at most \(4d\) cells,
\[
8\sum_{j:t_j\le B}E_{\mathrm{Pois}}\Delta_j^2
\le
64d(e^{\ell/8}+1)
\left(\frac{B^2}{N_0^2}+\frac{B}{mN_0}\right).
\]
For \(t_j>B\), the two heavy-pilot estimates instead yield
\[
\Pr_{\mathrm{Pois}}(\mathsf L_j)
E_{\mathrm{Pois}}(H_j-D_j)^2
\le4e^{-32\ell},
\]
and therefore
\[
8\sum_{j:t_j>B}E_{\mathrm{Pois}}\Delta_j^2
\le32e^{-32\ell}\left(1+\frac1m\right).
\]
In the last step we used
\[
\sum_j\left(p_j^2+\frac{p_j}{m}\right)
\le1+\frac1m,
\]
because \(0\le p_j\le1\) and \(\sum_jp_j=1\).
% lean: CausalSmith.Stat.MarRareqLogfrontier.sum_integral_sq_memberWeight_mul_selectedCorrection_le

\item Distinct cells are independent in the uncapped marked experiment.
Since \(g(a)^2=1\),
\[
\operatorname{Var}_{\mathrm{Pois}}(\tau_{\mathrm{raw}})
=4\sum_j\operatorname{Var}_{\mathrm{Pois}}(W_jG_j).
\]
For each cell, \(W_jG_j=W_jD_j+\Delta_j\). Nonnegativity of
\(\operatorname{Var}_{\mathrm{Pois}}(W_jD_j-\Delta_j)\) implies
\[
\operatorname{Var}_{\mathrm{Pois}}(W_jG_j)
\le
2\operatorname{Var}_{\mathrm{Pois}}(W_jD_j)
+2\operatorname{Var}_{\mathrm{Pois}}(\Delta_j)
\le
2\operatorname{Var}_{\mathrm{Pois}}(W_jD_j)
+2E_{\mathrm{Pois}}\Delta_j^2.
\]
The preceding aggregate estimates thus give
\[
\operatorname{Var}_{\mathrm{Pois}}(\tau_{\mathrm{raw}})
\le v_{n,d,q}.
\]
The raw contrast has finite second moment by the finite polynomial moment
bounds already established, and its squared error decomposes as
\[
E_{\mathrm{Pois}}(\tau_{\mathrm{raw}}-\Phi(P))^2
=
\operatorname{Var}_{\mathrm{Pois}}(\tau_{\mathrm{raw}})
+(E_{\mathrm{Pois}}\tau_{\mathrm{raw}}-\Phi(P))^2
\le v_{n,d,q}+b_{n,d,q}^2.
\]
Clipping decreases squared distance to every point of \([-1,1]\): below
\(-1\) it replaces the value by \(-1\), above \(1\) by \(1\), and inside the
interval it preserves the value. Since \(\Phi(P)\in[-1,1]\),
\[
E_{\mathrm{Pois}}(\widetilde\tau_{\mathrm{Pois}}-\Phi(P))^2
\le b_{n,d,q}^2+v_{n,d,q}.
\]
Combining this with the prefix transfer and the bound by four proves
\[
E_P\!\left[
(\widehat\tau^{\mathrm{mix}}_{n,d,q}-\Phi(P))^2
\right]
\le
\min\{4,b_{n,d,q}^2+v_{n,d,q}+4e^{-n\gamma_0}\}
=\mathfrak U_{n,d,q}.
\]
% lean: CausalSmith.Stat.MarRareqLogfrontier.unrestricted_auxiliary_risk_envelope_active

\item Finally, consider the ratio branch. Its ideal raw contrast is
\[
\tau_{\mathrm{raw}}^{\mathrm{ratio}}
:=2\sum_jg(a)W_jD_j.
\]
The ratio mean and membership independence give
\[
E_{\mathrm{Pois}}\tau_{\mathrm{raw}}^{\mathrm{ratio}}-\Phi(P)
=-2\sum_jg(a)p_j\mu_je^{-t_j}.
\]
Because \(t_j\ge N_0p_j\), \(0\le\mu_j\le1\), and
\(t_{\mathrm{bias}}e^{-t_{\mathrm{bias}}}\le e^{-1}\) for every real
\(t_{\mathrm{bias}}\ge0\),
\[
p_j\mu_je^{-t_j}
\le p_je^{-N_0p_j}
\le\frac1{eN_0}.
\]
The exponential inequality follows from
\(e^{t_{\mathrm{bias}}-1}\ge t_{\mathrm{bias}}\).
Thus
\[
|E_{\mathrm{Pois}}\tau_{\mathrm{raw}}^{\mathrm{ratio}}-\Phi(P)|
\le\frac{8d}{eN_0}.
\]
Independence across cells and the ratio variance aggregate give
\[
\operatorname{Var}_{\mathrm{Pois}}
(\tau_{\mathrm{raw}}^{\mathrm{ratio}})
=
4\sum_j\operatorname{Var}_{\mathrm{Pois}}(W_jD_j)
\le4\left(\frac4{N_0}+\frac1m\right).
\]
The squared-error decomposition, clipping, prefix transfer, and bound by
four therefore imply
\[
\begin{aligned}
&E_P\!\left[
(\widehat\tau^{\mathrm{mix}}_{n,d,q}-\Phi(P))^2
\right]\\
&\quad\le
\min\left\{4,
\left(\frac{8d}{eN_0}\right)^2
+4\left(\frac4{N_0}+\frac1m\right)
+4e^{-n\gamma_0}\right\}
=\mathfrak U_{n,d,q}.
\end{aligned}
\]
Together with the zero and capped-envelope cases, this establishes the
risk bound for every admissible \(P\). Combining it with the deterministic
comparison proved above gives
\[
E_P\!\left[
\left(\widehat\tau^{\mathrm{mix}}_{n,d,q}(\mathbf O_n)-\Phi(P)\right)^2
\right]
\le\mathfrak U_{n,d,q}
\le10^{30}r_{n,d,q}.
\]
The expectation is under the canonical product observed-sample law
specified in the statement, and \(10^{30}\) is positive, finite, and
independent of \(P,n,d,q\).
% lean: CausalSmith.Stat.MarRareqLogfrontier.unrestricted_uniform_mixed_count_certificate
\end{enumerate}
\end{proof}

\begin{proof}[Proof of \cref{obj:thm:uniform-mixed-count-certificate}]
\begin{enumerate}
\item Choose the universal constant supplied by
\cref{obj:thm:unrestricted-uniform-mixed-count-certificate}, so that
\[
0<\overline C<\infty.
\]
Fix \(n,d,q\) and \(P\in\mathcal M_{n,d,q}\) satisfying the hypotheses.
% lean: CausalSmith.Stat.MarRareqLogfrontier.unrestricted_uniform_mixed_count_certificate

\item The parameter, randomization, consistency, and arrival conditions in
\cref{obj:def:model-class} give
\[
P\in\mathcal M^{+}_{n,d,q},
\]
by \cref{obj:def:unrestricted-arrival-model-class}. Moreover,
\cref{obj:def:model-class} supplies
\cref{obj:ass:iid-sampling} for the observed sample. Applying
\cref{obj:thm:unrestricted-uniform-mixed-count-certificate} therefore yields
\[
E_P\!\left[
\left(\widehat\tau^{\mathrm{mix}}_{n,d,q}(\mathbf O_n)-\Phi(P)\right)^2
\right]
\le \mathfrak U_{n,d,q},
\qquad
\mathfrak U_{n,d,q}\le \overline C\,r_{n,d,q},
\]
where \(\Phi(P)\) is the signed cell functional of
\cref{obj:def:cell-functional}. The estimator, rate, and deterministic
envelope in that certificate are precisely those in the present statement,
with the same formulas on each branch.
% lean: CausalSmith.Stat.MarRareqLogfrontier.uniform_mixed_count_certificate

\item Since \(P\in\mathcal M^{+}_{n,d,q}\),
\cref{obj:lem:unrestricted-cell-identification} gives
\[
\tau(P)=\Phi(P).
\]
Substituting this identity into the risk bound proves
\[
E_P\!\left[
\left(\widehat\tau^{\mathrm{mix}}_{n,d,q}(\mathbf O_n)-\tau(P)\right)^2
\right]
=
E_P\!\left[
\left(\widehat\tau^{\mathrm{mix}}_{n,d,q}(\mathbf O_n)-\Phi(P)\right)^2
\right]
\le \mathfrak U_{n,d,q}
\le \overline C\,r_{n,d,q}.
\]
The chosen constant is universal, so these inequalities hold throughout
the stated model class.
% lean: CausalSmith.Stat.MarRareqLogfrontier.unrestricted_cell_identification
\end{enumerate}
\end{proof}

\begin{proof}[Proof of \cref{obj:thm:mixed-count-upper}]
\begin{enumerate}
\item Choose the universal constant \(\overline C>0\) supplied by
\cref{obj:thm:uniform-mixed-count-certificate}. Fix \(n,d,q\) and a full-data law \(P\) satisfying the stated conditions. Membership in \(\mathcal M_{n,d,q}\) supplies all hypotheses of that certificate.
By \cref{obj:lem:cell-identification-background}, the causal target equals the total observed-cell functional on this model class, which is the identification step used by the certificate. The certificate therefore gives
\[
E_P\!\left[
\left(
\widehat\tau^{\mathrm{mix}}_{n,d,q}(\mathbf O_n)-\tau(P)
\right)^2
\right]
\le \mathfrak U_{n,d,q}
\le \overline C\,r_{n,d,q}.
\]
The envelope and branch rules in that certificate are exactly those in the present statement. It remains to establish the comparison with the auxiliary randomized risk.
% lean: CausalSmith.Stat.MarRareqLogfrontier.uniform_mixed_count_certificate

\item Write the prefix probabilities and their remaining mass as
\[
w_{k_0}:=\Pr(K_0=k_0)
=e^{-n/2}\frac{(n/2)^{k_0}}{k_0!},
\qquad k_0\in\mathbb N,
\qquad
w_{\mathrm{tail}}:=1-\sum_{k_0=0}^{n}w_{k_0}.
\]
The singleton prefix events are disjoint, so
\[
w_{k_0}\ge0,
\qquad
\sum_{k_0=0}^{n}w_{k_0}\le1,
\qquad
w_{\mathrm{tail}}=\Pr(K_0>n)\ge0.
\]
For each integer \(0\le k_0\le n\), define the assignment set by
\[
\mathcal A_{\mathrm{str},k_0}
:=
\left\{
\sigma:\{1,\ldots,k_0\}
\longrightarrow
\{\mathrm{membership},\mathrm{pilot},\mathrm{estimation}\}
\right\},
\qquad
|\mathcal A_{\mathrm{str},k_0}|=3^{k_0}.
\]
Let the output at a prescribed prefix and assignment, and its assignment average, be
\[
\widetilde\tau_{k_0,\sigma}(\mathbf O_n)
:=
\widetilde\tau(\mathbf O_n;\omega)
\quad\text{when }\omega\text{ prescribes }K_0=k_0
\text{ and assignment }\sigma,
\]
\[
\overline\tau_{\mathrm{str},k_0}(\mathbf O_n)
:=
\frac{1}{3^{k_0}}
\sum_{\sigma\in\mathcal A_{\mathrm{str},k_0}}
\widetilde\tau_{k_0,\sigma}(\mathbf O_n).
\]
These outputs use the branch rules of
\cref{obj:def:mixed-count-estimator}. For \(k_0=0\), the assignment set contains the unique empty assignment.

The zero output on the prefix tail and the conditional-average definition give
\[
\widehat\tau^{\mathrm{mix}}_{n,d,q}(\mathbf O_n)
=
\sum_{k_0=0}^{n}
w_{k_0}\overline\tau_{\mathrm{str},k_0}(\mathbf O_n).
\]
Pointwise in the observed sample, its centered value therefore satisfies
\[
\widehat\tau^{\mathrm{mix}}_{n,d,q}(\mathbf O_n)-\tau(P)
=
w_{\mathrm{tail}}\bigl(-\tau(P)\bigr)
+
\sum_{k_0=0}^{n}
w_{k_0}
\bigl(\overline\tau_{\mathrm{str},k_0}(\mathbf O_n)-\tau(P)\bigr).
\]
The weights on the right are nonnegative and sum to one. Convexity of the square yields
\[
\left(
\widehat\tau^{\mathrm{mix}}_{n,d,q}(\mathbf O_n)-\tau(P)
\right)^2
\le
w_{\mathrm{tail}}\tau(P)^2
+
\sum_{k_0=0}^{n}
w_{k_0}
\left(
\overline\tau_{\mathrm{str},k_0}(\mathbf O_n)-\tau(P)
\right)^2.
\]
% lean: CausalSmith.Stat.MarRareqLogfrontier.finite_weighted_square_with_remainder

\item For every \(0\le k_0\le n\), centering the assignment average gives
\[
\overline\tau_{\mathrm{str},k_0}(\mathbf O_n)-\tau(P)
=
\frac{1}{3^{k_0}}
\sum_{\sigma\in\mathcal A_{\mathrm{str},k_0}}
\left(
\widetilde\tau_{k_0,\sigma}(\mathbf O_n)-\tau(P)
\right).
\]
The finite Cauchy--Schwarz inequality consequently gives
\[
\left(
\overline\tau_{\mathrm{str},k_0}(\mathbf O_n)-\tau(P)
\right)^2
\le
\frac{1}{3^{k_0}}
\sum_{\sigma\in\mathcal A_{\mathrm{str},k_0}}
\left(
\widetilde\tau_{k_0,\sigma}(\mathbf O_n)-\tau(P)
\right)^2.
\]
Multiplying these inequalities by the nonnegative \(w_{k_0}\), summing, and inserting the result into the preceding bound yields
\[
\begin{aligned}
&\left(
\widehat\tau^{\mathrm{mix}}_{n,d,q}(\mathbf O_n)-\tau(P)
\right)^2\\
&\quad\le
w_{\mathrm{tail}}\tau(P)^2
+
\sum_{k_0=0}^{n}\frac{w_{k_0}}{3^{k_0}}
\sum_{\sigma\in\mathcal A_{\mathrm{str},k_0}}
\left(
\widetilde\tau_{k_0,\sigma}(\mathbf O_n)-\tau(P)
\right)^2.
\end{aligned}
\]
% lean: CausalSmith.Stat.MarRareqLogfrontier.finite_uniform_square

\item Conditional on the observed sample, the prefix has probabilities \(w_{k_0}\), each assignment has probability \(3^{-k_0}\), and the tail output is zero. Thus
\[
\begin{aligned}
&E_\omega\!\left[
\left(\widetilde\tau(\mathbf O_n;\omega)-\tau(P)\right)^2
\mid\mathbf O_n
\right]\\
&\quad=
w_{\mathrm{tail}}\tau(P)^2
+
\sum_{k_0=0}^{n}\frac{w_{k_0}}{3^{k_0}}
\sum_{\sigma\in\mathcal A_{\mathrm{str},k_0}}
\left(
\widetilde\tau_{k_0,\sigma}(\mathbf O_n)-\tau(P)
\right)^2.
\end{aligned}
\]
The observed-sample space is finite, so the displayed functions are measurable and integrable. Integrating the pointwise inequality establishes
\[
E_P\!\left[
\left(
\widehat\tau^{\mathrm{mix}}_{n,d,q}(\mathbf O_n)-\tau(P)
\right)^2
\right]
\le
E_{P,\omega}\!\left[
\left(
\widetilde\tau(\mathbf O_n;\omega)-\tau(P)
\right)^2
\right].
\]
Together with the certificate bound, this proves all the asserted inequalities with the same universal constant \(\overline C\).
% lean: CausalSmith.Stat.MarRareqLogfrontier.mixed_count_derandomization
% lean: CausalSmith.Stat.MarRareqLogfrontier.mixed_count_upper
\end{enumerate}
\end{proof}

\begin{proof}[Proof of \cref{obj:thm:full-experiment-lower}]
Let \(\underline c_{\mathrm{one}}>0\) be the universal constant supplied by
\cref{obj:lem:one-cell-testing-family}. Define the universal constants
\[
\eta=e^{-32},\qquad
N_{\mathrm{act}}
=\max\left\{1,e^2,4\eta,4096\eta,
e^{(32+\log 8)/28}\right\},
\]
\[
\underline c_{\mathrm{act}}
=\min\left\{
\frac{1}{256N_{\mathrm{act}}},
\frac{1}{256\cdot2048^2},
\frac{\eta^2}{16384}
\right\},
\qquad
\underline c_*=\min\{\underline c_{\mathrm{one}},
\underline c_{\mathrm{act}}\},
\qquad
\underline c=\frac{\underline c_*}{4}.
\]
These constants are positive and independent of \(n,d,q\).

Fix parameters satisfying the hypotheses, and define
\[
N=nq,\qquad \ell=\log(e+N),\qquad
u_{\mathrm{act}}=\frac{d}{N\ell},\qquad
a_{\mathrm{eff}}=N^{-1},\qquad
b_{\mathrm{dim}}=u_{\mathrm{act}}^2.
\]
Then
\[
N>0,\qquad
\ell\ge\log e=1,\qquad
a_{\mathrm{eff}}\ge0,\qquad b_{\mathrm{dim}}\ge0.
\]
By \cref{obj:def:frontier-rate},
\[
r_{n,d,q}=\min\{1,a_{\mathrm{eff}}+b_{\mathrm{dim}}\}.
\]

\begin{enumerate}
\item \emph{The one-cell contribution.}
By \cref{obj:lem:one-cell-testing-family},
\[
\underline c_{\mathrm{one}}\min\{1,a_{\mathrm{eff}}\}
\le \mathfrak R_{n,d,q}.
\]
We next establish the complementary bound
\[
\sqrt N\le d
\quad\Longrightarrow\quad
\underline c_{\mathrm{act}}\min\{1,b_{\mathrm{dim}}\}
\le \mathfrak R_{n,d,q}.
\]
% lean: CausalSmith.Stat.MarRareqLogfrontier.one_cell_testing_family

\item \emph{A numerical one-cell estimate for constant calibration.}
We record explicitly the numerical estimate used below:
\[
\frac1{256}\min\{1,N^{-1}\}\le\mathfrak R_{n,d,q}.
\]
For \(\mu_{\mathrm{one}}\in[1/4,3/4]\), define \(P_{\mathrm{one},\mu_{\mathrm{one}}}\) by
\[
X=0,\qquad S(0)=S(1)=Y(0)=0,\qquad
(A,Y(1),R)\sim
\operatorname{Bernoulli}(1/2)\otimes
\operatorname{Bernoulli}(\mu_{\mathrm{one}})\otimes
\operatorname{Bernoulli}(q),
\]
\[
S=0,\qquad Y=AY(1).
\]
Treatment is independent of the potential variables, and arrival is independent
of the outcome conditional on the observed cell. Every occupied cell has arrival
probability \(q\). Consequently,
\[
P_{\mathrm{one},\mu_{\mathrm{one}}}\in\mathcal M_{n,d,q},
\qquad
\tau(P_{\mathrm{one},\mu_{\mathrm{one}}})=\mu_{\mathrm{one}}.
\]
Write
\[
Q_{\mathrm{one},\mu_{\mathrm{one}}}=\mathcal L_{P_{\mathrm{one},\mu_{\mathrm{one}}}}(O_1),
\qquad
\delta_{\mathrm{one}}
=\min\left\{\frac14,\frac{1}{4\sqrt N}\right\}.
\]
Thus \(0<\delta_{\mathrm{one}}\le1/4\), and the two target values
\(1/2\) and \(1/2+\delta_{\mathrm{one}}\) belong to \([1/4,3/4]\).
Considering the two branches of this minimum gives
\[
\delta_{\mathrm{one}}^2
\ge \frac1{16}\min\{1,N^{-1}\},
\qquad
16N\delta_{\mathrm{one}}^2\le1.
\]

For probability laws \(Q'\ll Q\), use the convention
\[
\chi^2(Q'\Vert Q)
=\int\left(\frac{dQ'}{dQ}-1\right)^2\,dQ.
\]
The two one-record laws differ precisely at the treated, arrived records with
outcome zero or one. Each of these atoms has probability \(q/4\) under
\(Q_{\mathrm{one},1/2}\), and its probability changes in absolute value by
\(q\delta_{\mathrm{one}}/2\). Summing their contributions gives
\[
\chi^2\!\left(
Q_{\mathrm{one},1/2+\delta_{\mathrm{one}}}
\Vert Q_{\mathrm{one},1/2}\right)
=2q\delta_{\mathrm{one}}^2
\le\frac{1}{8n}.
\]
Independence of the records factors the second moment of the likelihood ratio,
so
\[
1+\chi^2\!\left(
Q_{\mathrm{one},1/2+\delta_{\mathrm{one}}}^{\otimes n}
\Vert Q_{\mathrm{one},1/2}^{\otimes n}\right)
=
\left[
1+\chi^2\!\left(
Q_{\mathrm{one},1/2+\delta_{\mathrm{one}}}
\Vert Q_{\mathrm{one},1/2}\right)
\right]^n
\le e^{1/8}.
\]
Since \(e^{1/8}\le e^{1/2}<\sqrt3<2\), Cauchy--Schwarz applied to the
likelihood ratio yields
\[
\operatorname{TV}\!\left(
Q_{\mathrm{one},1/2+\delta_{\mathrm{one}}}^{\otimes n},
Q_{\mathrm{one},1/2}^{\otimes n}\right)
\le
\frac12\sqrt{
\chi^2\!\left(
Q_{\mathrm{one},1/2+\delta_{\mathrm{one}}}^{\otimes n}
\Vert Q_{\mathrm{one},1/2}^{\otimes n}\right)}
\le\frac12.
\]

For every \(\mathsf T\in\mathcal T_n\), define its barycenter on the
observed-sample space by
\[
b_{\mathsf T}(o)=\int_{[-1,1]}h\,\mathsf T(o,dh).
\]
By \cref{obj:lem:randomized-kernel-barycenter},
\[
b_{\mathsf T}(o)\in[-1,1],
\qquad
E_P[(b_{\mathsf T}(\mathbf O_n)-\tau(P))^2]
\le E_P[(\mathsf T-\tau(P))^2].
\]
Define the nearer-target test
\[
\psi_{\mathrm{one}}(o)
=
\mathbf1\left\{
\left|b_{\mathsf T}(o)-\left(\frac12+\delta_{\mathrm{one}}\right)\right|
\le
\left|b_{\mathsf T}(o)-\frac12\right|
\right\}.
\]
An error under either target entails absolute estimation error at least
\(\delta_{\mathrm{one}}/2\). The testing inequality in
\cref{obj:lem:randomized-kernel-tv-comparison}, followed by the barycenter
risk comparison, therefore gives
\[
\begin{aligned}
&\max\left\{
E_{P_{\mathrm{one},1/2}}\!\left[(\mathsf T-\tfrac12)^2\right],
E_{P_{\mathrm{one},1/2+\delta_{\mathrm{one}}}}\!
\left[(\mathsf T-\tfrac12-\delta_{\mathrm{one}})^2\right]
\right\}\\
&\qquad\ge
\frac{\delta_{\mathrm{one}}^2}{8}
\left[
1-\operatorname{TV}\!\left(
Q_{\mathrm{one},1/2+\delta_{\mathrm{one}}}^{\otimes n},
Q_{\mathrm{one},1/2}^{\otimes n}\right)
\right]
\ge\frac{\delta_{\mathrm{one}}^2}{16}
\ge\frac1{256}\min\{1,N^{-1}\}.
\end{aligned}
\]
Both laws belong to the model class. Taking the supremum over that class and
then the infimum over \(\mathsf T\) proves the displayed numerical estimate.
% lean: CausalSmith.Stat.MarRareqLogfrontier.oneCell_minimax_risk_floors

\item \emph{Large effective size and many rare cells.}
To prove the complementary bound, suppose \(\sqrt N\le d\).
Use the construction parameters of
\cref{obj:def:activation-lower-handle}:
\[
K=\lceil\ell\rceil,\qquad H=K^2,\qquad
b_0=\frac{\eta}{N\ell},\qquad
J=\min\left\{d-1,\left\lfloor\frac1{2b_0}\right\rfloor\right\}.
\]
Here \(b_0>0\), \(0\le J\le d-1\), and
\[
Jb_0\le\frac12.
\]
The ceiling bound and \(\ell\ge1\) give
\[
\ell\le K<\ell+1\le2\ell,\qquad H\ge1.
\]
The rare-arrival restriction then implies
\[
0<qH\le4q\ell^2\le\frac1{16}.
\]

First suppose \(N_{\mathrm{act}}\le N\). In particular,
\[
N\ge1,\qquad N\ge e^2,\qquad
N\ge4\eta,\qquad N\ge4096\eta,\qquad
N\ge e^{(32+\log8)/28}.
\]
The capacity bound is
\[
2048(2b_0)=\frac{4096\eta}{N\ell}\le1,
\qquad
2048\le\frac1{2b_0}.
\]
Within this branch, first suppose \(J\ge2048\). Since \(N\ge e^2\),
\[
\ell=\log(e+N)\ge2,\qquad K\ge2.
\]
By \cref{obj:lem:moment-matched-reciprocal-priors}, there are finitely
supported probability laws \(\Pi_0,\Pi_1\) on \([1,H]\) satisfying
\[
\int z^v\,d\Pi_0(z)=\int z^v\,d\Pi_1(z)
\qquad(0\le v\le K,\ v\in\mathbb N).
\]
If the reciprocal mean under \(\Pi_0\) exceeds that under \(\Pi_1\), interchange
the priors; otherwise retain their order. Their separation then reads
\[
\int z^{-1}\,d\Pi_1(z)-\int z^{-1}\,d\Pi_0(z)
\ge\frac1{12}.
\]

For each \(\mathbf z=(z_x)_{x=0}^{d-1}\in[1,H]^d\), define a full-data law
\(P_{\mathrm{act},\mathbf z}\) as follows. Its baseline probabilities are
\[
p_{\mathrm{act},x}=
\begin{cases}
b_0,&0\le x<J,\\
1-Jb_0,&x=J,\\
0,&J<x<d.
\end{cases}
\]
Set \(S(0)=S(1)=Y(0)=0\), draw \(A\sim\operatorname{Bernoulli}(1/2)\)
independently of the baseline and potential variables, and impose \(S=0\) and
\(Y=AY(1)\). Conditional on \(X=x\), draw \(Y(1)\) and \(R\) independently with
\[
Y(1)\mid X=x\sim
\begin{cases}
\operatorname{Bernoulli}(z_x^{-1}),&x<J,\\
\operatorname{Bernoulli}(0),&x\ge J,
\end{cases}
\qquad
R\mid X=x\sim
\begin{cases}
\operatorname{Bernoulli}(qz_x),&x<J,\\
\operatorname{Bernoulli}(qH),&x\ge J.
\end{cases}
\]
The baseline masses sum to one and are nonnegative. Since \(1\le z_x\le H\),
the occupied-cell arrival probabilities belong to \([q,1]\). The specified
independences give randomized assignment and conditional independence of
arrival and outcome, and the definitions of \(S,Y\) give consistency.
With independent sampling from this fixed law, all model conditions hold:
\[
P_{\mathrm{act},\mathbf z}\in\mathcal M_{n,d,q},
\qquad
\tau_{\mathrm{act}}(\mathbf z)
:=\tau(P_{\mathrm{act},\mathbf z})
=b_0\sum_{x=0}^{J-1}z_x^{-1}.
\]

For \(\iota\in\{0,1\}\), draw
\[
\mathbf Z=(Z_x)_{x=0}^{d-1}\sim\Pi_\iota^{\otimes d}
\]
once for the entire sample. Define the observed mixture law
\[
Q_{\mathrm{act},\iota}
=\int_{[1,H]^d}
\mathcal L_{P_{\mathrm{act},\mathbf z}}(\mathbf O_n)\,
d\Pi_\iota^{\otimes d}(\mathbf z),
\]
and let \(\widetilde Q_{\mathrm{act},\iota}\) be the corresponding augmented
law constructed in \cref{obj:def:activation-lower-handle}. Thus
\[
\widetilde Q_{\mathrm{act},\iota}
=
\int_{[1,H]^d}
\mathcal L\!\left(
(\mathsf B_{\mathrm{act},t},O_t)_{t=1}^{n}
\,\middle|\,\mathbf Z=\mathbf z\right)
d\Pi_\iota^{\otimes d}(\mathbf z).
\]
The activation flags have probability \(qH\), independently of the latent
vector and the baseline and treatment pattern.

At the record level, forgetting the flag gives the following probabilities
for the three marks \((R,RY)=(1,1),(1,0),(0,0)\):
\[
\begin{array}{c|ccc}
& (1,1)&(1,0)&(0,0)\\ \hline
A=1,\ X=x<J&q&q(z_x-1)&1-qz_x\\
A=0,\ X=x<J&0&qz_x&1-qz_x\\
X\ge J&0&qH&1-qH
\end{array}
\]
These are exactly the probabilities induced by \(P_{\mathrm{act},\mathbf z}\).
Define the forgetting map by
\[
\operatorname{pr}_{\mathrm{obs}}
\bigl((\mathsf B_{\mathrm{act},t},O_t)_{t=1}^{n}\bigr)
=(O_t)_{t=1}^{n}.
\]
Conditional independence of the records and integration over the common
latent vector therefore give
\[
Q_{\mathrm{act},\iota}
=\widetilde Q_{\mathrm{act},\iota}
\circ\operatorname{pr}_{\mathrm{obs}}^{-1}.
\]
% lean: CausalSmith.Stat.MarRareqLogfrontier.interval_ordered_reciprocal_priors

\item \emph{Equality on the moment-matching event.}
Continue in the branch \(N\ge N_{\mathrm{act}}\), \(J\ge2048\).
For an augmented sample, define the activated count in cell \(x\) and the
cutoff event by
\[
V_x=\sum_{t=1}^{n}
\mathbf1\{X_t=x,\ \mathsf B_{\mathrm{act},t}=1\},
\qquad
\mathcal G=\{V_x\le K\text{ for every }0\le x<J\}.
\]

Fix a possible augmented sample \(\mathbf o_{\mathrm{aug}}\), and denote its
coordinates by the corresponding record symbols. Its baseline, treatment,
surrogate, and flag pattern has the common probability
\[
\varpi_{\mathbf o_{\mathrm{aug}}}
=
\mathbf1\{S_t=0\text{ for every }t\}
\prod_{t=1}^{n}
\frac{p_{\mathrm{act},X_t}}2
(qH)^{\mathsf B_{\mathrm{act},t}}
(1-qH)^{1-\mathsf B_{\mathrm{act},t}}.
\]
For each cell \(x\), define \(\mathcal P_{\mathbf o_{\mathrm{aug}},x}(z)\)
as the product, over records with \(X_t=x\), of their conditional mark
probabilities from \cref{obj:def:activation-lower-handle}:
\[
\mathcal P_{\mathbf o_{\mathrm{aug}},x}(z)
=
\prod_{t:X_t=x}
\Pr\!\left(
(R,RY)=(R_t,R_tY_t)
\,\middle|\,
X=x,\ A=A_t,\ \mathsf B_{\mathrm{act}}=\mathsf B_{\mathrm{act},t},
\ Z_x=z
\right).
\]
Empty products equal one. The row entries extend these expressions to
polynomials on \(\mathbb R\); their probability interpretation applies on
\([1,H]\). Impossible marks have weight zero. For a rare cell, an inactive
factor is constant and an activated factor is affine in \(z\). For every
other cell the factors are constant. Hence
\[
\deg\mathcal P_{\mathbf o_{\mathrm{aug}},x}\le V_x
\quad(x<J),\qquad
\deg\mathcal P_{\mathbf o_{\mathrm{aug}},x}\le0
\quad(x\ge J),
\]
where a zero polynomial satisfies each stated degree bound.

On \(\mathcal G\), define the coefficients
\(\beta_{\mathbf o_{\mathrm{aug}},x,v}\) by the polynomial expansion
\[
\mathcal P_{\mathbf o_{\mathrm{aug}},x}(z)
=\sum_{v=0}^{K}
\beta_{\mathbf o_{\mathrm{aug}},x,v}z^v
\qquad(z\in\mathbb R),
\]
padding with zero coefficients when necessary. Product-prior independence
then gives
\[
\widetilde Q_{\mathrm{act},\iota}\{\mathbf o_{\mathrm{aug}}\}
=
\varpi_{\mathbf o_{\mathrm{aug}}}
\prod_{x=0}^{d-1}
\left[
\sum_{v=0}^{K}
\beta_{\mathbf o_{\mathrm{aug}},x,v}
\int z^v\,d\Pi_\iota(z)
\right]
\qquad(\mathbf o_{\mathrm{aug}}\in\mathcal G).
\]
The matched moments make the right-hand side identical for
\(\iota=0,1\). Thus
\[
\widetilde Q_{\mathrm{act},0}\{\mathbf o_{\mathrm{aug}}\}
=
\widetilde Q_{\mathrm{act},1}\{\mathbf o_{\mathrm{aug}}\}
\qquad(\mathbf o_{\mathrm{aug}}\in\mathcal G).
\]
% lean: CausalSmith.Stat.MarRareqLogfrontier.activatedAugmentedSample_prior_singleton_eq

\item \emph{The cutoff probability and total variation.}
The retained baseline, treatment, and flag pattern has the same distribution
under both priors. In particular, for \(x<J\),
\[
V_x\sim\operatorname{Binomial}(n,b_0qH).
\]
Define
\[
k_{\mathrm{tail}}=K+1,\qquad
\nu_{\mathrm{act}}=nb_0qH=\frac{\eta K^2}{\ell}.
\]
For an integer \(v\ge0\), let its falling factorial be
\[
(v)_{k_{\mathrm{tail}}}
=
\begin{cases}
\displaystyle\prod_{t=0}^{k_{\mathrm{tail}}-1}(v-t),
&v\ge k_{\mathrm{tail}},\\
0,&v<k_{\mathrm{tail}}.
\end{cases}
\]
Expanding over ordered distinct record indices gives
\[
E[(V_x)_{k_{\mathrm{tail}}}]
=(n)_{k_{\mathrm{tail}}}(b_0qH)^{k_{\mathrm{tail}}}
\le\nu_{\mathrm{act}}^{k_{\mathrm{tail}}}.
\]
When \(k_{\mathrm{tail}}>n\), both falling-factorial moments are zero.
On \(V_x>K\), the falling factorial is at least \(k_{\mathrm{tail}}!\).
Therefore
\[
\widetilde Q_{\mathrm{act},\iota}(V_x>K)
\le
\frac{\nu_{\mathrm{act}}^{k_{\mathrm{tail}}}}
{k_{\mathrm{tail}}!}.
\]

The bounds \(K\le2\ell\) and \(K\le k_{\mathrm{tail}}\) imply
\[
K^2\le2\ell k_{\mathrm{tail}},
\qquad
\nu_{\mathrm{act}}\le2\eta k_{\mathrm{tail}}
\le e^{-31}k_{\mathrm{tail}},
\qquad
k_{\mathrm{tail}}\ge\ell,
\]
where \(2\le e\) gives the second inequality involving the exponential.
The exponential series gives
\(k_{\mathrm{tail}}^{k_{\mathrm{tail}}}/k_{\mathrm{tail}}!
\le e^{k_{\mathrm{tail}}}\), whence
\[
\frac{\nu_{\mathrm{act}}^{k_{\mathrm{tail}}}}
{k_{\mathrm{tail}}!}
\le
e^{-31k_{\mathrm{tail}}}
\frac{k_{\mathrm{tail}}^{k_{\mathrm{tail}}}}
{k_{\mathrm{tail}}!}
\le e^{-30k_{\mathrm{tail}}}
\le e^{-30\ell}.
\]
A union bound over the rare cells yields
\[
\widetilde Q_{\mathrm{act},\iota}(\mathcal G^c)
\le Je^{-30\ell}.
\]
The mass cap and elementary exponential bounds give
\[
J\le\frac{N\ell}{2\eta},\qquad
N\le e^\ell,\qquad \ell\le e^\ell,\qquad
N\ell\le e^{2\ell}.
\]
Consequently,
\[
\widetilde Q_{\mathrm{act},\iota}(\mathcal G^c)
\le
\frac{N\ell}{2\eta}e^{-30\ell}
\le
\frac{e^{2\ell}}{\eta}e^{-30\ell}
=e^{32-28\ell}.
\]
The threshold \(N\ge e^{(32+\log8)/28}\) implies
\[
\ell\ge\frac{32+\log8}{28},
\qquad
\widetilde Q_{\mathrm{act},\iota}(\mathcal G^c)\le\frac18.
\]

Equality of the atoms on \(\mathcal G\) implies equality of the probabilities
of \(\mathcal G\), and hence of its complement. For every augmented-sample
event \(\mathcal B\),
\[
\begin{aligned}
&\left|
\widetilde Q_{\mathrm{act},0}(\mathcal B)
-\widetilde Q_{\mathrm{act},1}(\mathcal B)
\right|\\
&\quad=
\left|
\widetilde Q_{\mathrm{act},0}(\mathcal B\cap\mathcal G^c)
-\widetilde Q_{\mathrm{act},1}(\mathcal B\cap\mathcal G^c)
\right|
\le\widetilde Q_{\mathrm{act},0}(\mathcal G^c).
\end{aligned}
\]
Applying this inequality to preimages under the forgetting map gives
\[
\operatorname{TV}(Q_{\mathrm{act},0},Q_{\mathrm{act},1})
\le
\operatorname{TV}(\widetilde Q_{\mathrm{act},0},
\widetilde Q_{\mathrm{act},1})
\le\frac18.
\]
% lean: CausalSmith.Stat.MarRareqLogfrontier.activatedAugmentedMixture_cutoff_budget

\item \emph{From center separation to squared-error risk.}
For \(\iota\in\{0,1\}\), define the prior center and its separation by
\[
\theta_{\mathrm{act},\iota}
=E_{\Pi_\iota^{\otimes d}}[\tau_{\mathrm{act}}(\mathbf Z)]
=Jb_0\int z^{-1}\,d\Pi_\iota(z),
\qquad
\Delta_{\mathrm{act}}
=|\theta_{\mathrm{act},1}-\theta_{\mathrm{act},0}|.
\]
The reciprocal-mean separation implies
\[
\Delta_{\mathrm{act}}\ge\frac{Jb_0}{12}.
\]
Since \(Z_x^{-1}\in[0,1]\), its variance satisfies
\[
\operatorname{Var}_{\Pi_\iota}(Z_x^{-1})
=
E_{\Pi_\iota}\!\left[(Z_x^{-1}-\tfrac12)^2\right]
-\left(E_{\Pi_\iota}[Z_x^{-1}]-\tfrac12\right)^2
\le\frac14.
\]
Independence of the latent coordinates therefore yields
\[
E_{\Pi_\iota^{\otimes d}}\!\left[
(\tau_{\mathrm{act}}(\mathbf Z)-\theta_{\mathrm{act},\iota})^2
\right]
=
b_0^2\sum_{x=0}^{J-1}
\operatorname{Var}_{\Pi_\iota}(Z_x^{-1})
\le\frac{Jb_0^2}{4}.
\]

Fix \(\mathsf T\in\mathcal T_n\), and define its maximal model risk and its
centered mixture risks by
\[
\mathcal W(\mathsf T)
=
\sup_{P\in\mathcal M_{n,d,q}}
E_P[(\mathsf T-\tau(P))^2],
\qquad
C_{\mathrm{act},\iota}(\mathsf T)
=
\int
(b_{\mathsf T}(o)-\theta_{\mathrm{act},\iota})^2\,
dQ_{\mathrm{act},\iota}(o).
\]
The bounded action and target make these risks finite. Class membership of
every prior support law gives
\[
\int
E_{P_{\mathrm{act},\mathbf z}}\!\left[
(\mathsf T-\tau_{\mathrm{act}}(\mathbf z))^2
\right]
d\Pi_\iota^{\otimes d}(\mathbf z)
\le\mathcal W(\mathsf T).
\]
Using the barycenter comparison from
\cref{obj:lem:randomized-kernel-barycenter} and
\[
(b_{\mathsf T}(o)-\theta_{\mathrm{act},\iota})^2
\le
2(b_{\mathsf T}(o)-\tau_{\mathrm{act}}(\mathbf z))^2
+
2(\tau_{\mathrm{act}}(\mathbf z)-\theta_{\mathrm{act},\iota})^2
\]
inside the mixture integral gives
\[
C_{\mathrm{act},\iota}(\mathsf T)
\le2\mathcal W(\mathsf T)+\frac{Jb_0^2}{2}.
\]

Define the nearer-center test
\[
\psi_{\mathrm{act}}(o)
=
\mathbf1\left\{
|b_{\mathsf T}(o)-\theta_{\mathrm{act},1}|
\le
|b_{\mathsf T}(o)-\theta_{\mathrm{act},0}|
\right\}.
\]
The triangle inequality shows that a testing error under center \(\iota\)
entails
\[
|b_{\mathsf T}(o)-\theta_{\mathrm{act},\iota}|
\ge\frac{\Delta_{\mathrm{act}}}{2}.
\]
By \cref{obj:lem:randomized-kernel-tv-comparison},
\[
Q_{\mathrm{act},0}(\psi_{\mathrm{act}}=1)
+
Q_{\mathrm{act},1}(\psi_{\mathrm{act}}=0)
\ge
1-\operatorname{TV}(Q_{\mathrm{act},0},Q_{\mathrm{act},1}).
\]
Bounding each error probability by its centered squared risk consequently
gives
\[
\frac{\Delta_{\mathrm{act}}^2}{8}
\left[
1-\operatorname{TV}(Q_{\mathrm{act},0},Q_{\mathrm{act},1})
\right]
\le
\max_{\iota\in\{0,1\}}C_{\mathrm{act},\iota}(\mathsf T).
\]
Together with the total variation and centered-risk bounds, this implies
\[
\frac{7\Delta_{\mathrm{act}}^2}{128}
-\frac{Jb_0^2}{4}
\le\mathcal W(\mathsf T).
\]
Now \(J\ge2048\) and \(\Delta_{\mathrm{act}}\ge Jb_0/12\), so
\[
\mathcal W(\mathsf T)
\ge
(Jb_0)^2
\left(\frac7{18432}-\frac1{4J}\right)
\ge\frac{(Jb_0)^2}{4096}.
\]
Taking the infimum over \(\mathsf T\) gives the calibrated fuzzy bound
\[
\frac{(Jb_0)^2}{4096}\le\mathfrak R_{n,d,q}.
\]
% lean: CausalSmith.Stat.MarRareqLogfrontier.activatedAugmentedMixture_minimaxRisk_calibrated

\item \emph{Converting rare-cell mass to the dimensional contribution.}
Still suppose \(N\ge N_{\mathrm{act}}\) and \(J\ge2048\). We have
\[
b_0=\frac{\eta}{N\ell}\le\frac14,
\qquad
0<\eta=e^{-32}\le\frac12,
\qquad d\ge2.
\]
The first inequality uses \(N\ge4\eta\) and \(\ell\ge1\); the second follows
from \(e^{32}\ge1+32\).

If \(d-1\le\lfloor(2b_0)^{-1}\rfloor\), then
\[
J=d-1\ge\frac d2,\qquad Jb_0\ge\frac{db_0}{2}.
\]
Otherwise the floor determines \(J\), and
\[
Jb_0
=\left\lfloor\frac1{2b_0}\right\rfloor b_0
\ge\frac12-b_0\ge\frac14.
\]
Thus both branches give
\[
\min\left\{\frac{db_0}{2},\frac14\right\}\le Jb_0.
\]
Moreover,
\[
\frac{\eta}{24}\min\{1,u_{\mathrm{act}}\}
\le\frac{\eta u_{\mathrm{act}}}{24}
=\frac{db_0}{24},
\qquad
\frac{\eta}{24}\min\{1,u_{\mathrm{act}}\}
\le\frac{\eta}{24}\le\frac1{48}.
\]
It follows that
\[
\frac{\eta}{24}\min\{1,u_{\mathrm{act}}\}
\le
\frac1{12}\min\left\{\frac{db_0}{2},\frac14\right\}
\le\frac{Jb_0}{12},
\]
and hence
\[
\frac{\eta}{2}\min\{1,u_{\mathrm{act}}\}\le Jb_0.
\]
For \(u_{\mathrm{act}}\ge0\), considering \(u_{\mathrm{act}}\le1\) and
\(u_{\mathrm{act}}\ge1\) gives
\[
\bigl(\min\{1,u_{\mathrm{act}}\}\bigr)^2
=\min\{1,u_{\mathrm{act}}^2\}.
\]
Squaring the mass bound therefore yields
\[
\frac{\eta^2}{4}\min\{1,u_{\mathrm{act}}^2\}
\le(Jb_0)^2.
\]
The choice of \(\underline c_{\mathrm{act}}\) and the calibrated fuzzy bound
now imply
\[
\underline c_{\mathrm{act}}\min\{1,b_{\mathrm{dim}}\}
\le
\frac{\eta^2}{16384}\min\{1,u_{\mathrm{act}}^2\}
\le\frac{(Jb_0)^2}{4096}
\le\mathfrak R_{n,d,q}.
\]
% lean: CausalSmith.Stat.MarRareqLogfrontier.interval_prior_segment_width_lower

\item \emph{Large effective size with fewer rare cells.}
Suppose instead that \(N\ge N_{\mathrm{act}}\) and \(J<2048\).
The capacity bound established above gives
\[
\left\lfloor\frac1{2b_0}\right\rfloor\ge2048.
\]
The definition of \(J\) therefore forces
\[
d-1<2048,\qquad d\le2048.
\]
This includes \(J=0\), for which the same argument gives \(d=1\).

Since \(N\ge1\) and \(\ell\ge1\),
\[
\sqrt N\le N\le N\ell.
\]
Consequently,
\[
\min\{1,u_{\mathrm{act}}\}
\le\frac{d}{N\ell}
\le\frac{2048}{\sqrt N},
\qquad
\min\{1,u_{\mathrm{act}}^2\}
\le\frac{2048^2}{N}.
\]
Also,
\[
\min\{1,N^{-1}\}=N^{-1},
\qquad
\underline c_{\mathrm{act}}\,2048^2\le\frac1{256}.
\]
The numerical one-cell estimate thus gives
\[
\underline c_{\mathrm{act}}\min\{1,b_{\mathrm{dim}}\}
\le
\frac{\underline c_{\mathrm{act}}\,2048^2}{N}
\le\frac1{256N}
\le\mathfrak R_{n,d,q}.
\]
% lean: CausalSmith.Stat.MarRareqLogfrontier.activated_fuzzy_lower

\item \emph{Small effective size.}
In the remaining effective-size branch, \(N<N_{\mathrm{act}}\).
Because \(N>0\) and \(N_{\mathrm{act}}\ge1\),
\[
N_{\mathrm{act}}^{-1}\le1,\qquad
N_{\mathrm{act}}^{-1}\le N^{-1},
\qquad
N_{\mathrm{act}}^{-1}\le\min\{1,N^{-1}\}.
\]
The definition of \(\underline c_{\mathrm{act}}\) and the numerical
one-cell estimate give
\[
\underline c_{\mathrm{act}}\min\{1,b_{\mathrm{dim}}\}
\le\underline c_{\mathrm{act}}
\le\frac1{256N_{\mathrm{act}}}
\le\frac1{256}\min\{1,N^{-1}\}
\le\mathfrak R_{n,d,q}.
\]
The three effective-size and rare-count branches establish the complementary
bound whenever \(\sqrt N\le d\).
% lean: CausalSmith.Stat.MarRareqLogfrontier.activated_fuzzy_lower

\item \emph{Combining the bounds when \(d\ge\sqrt N\).}
First suppose \(\sqrt N\le d\). Both component bounds apply, and their
nonnegative capped factors give
\[
\underline c_*\min\{1,a_{\mathrm{eff}}\}
\le
\underline c_{\mathrm{one}}\min\{1,a_{\mathrm{eff}}\}
\le\mathfrak R_{n,d,q},
\]
\[
\underline c_*\min\{1,b_{\mathrm{dim}}\}
\le
\underline c_{\mathrm{act}}\min\{1,b_{\mathrm{dim}}\}
\le\mathfrak R_{n,d,q}.
\]

For the nonnegative quantities \(a_{\mathrm{eff}},b_{\mathrm{dim}}\),
\[
\min\{1,a_{\mathrm{eff}}+b_{\mathrm{dim}}\}
\le
\min\{1,a_{\mathrm{eff}}\}+\min\{1,b_{\mathrm{dim}}\}.
\]
Indeed, if both quantities are at most one, the right-hand side is their sum;
if either is at least one, the right-hand side is at least one.
Each capped quantity is at most their maximum, so
\[
\min\{1,a_{\mathrm{eff}}+b_{\mathrm{dim}}\}
\le
2\max\{\min\{1,a_{\mathrm{eff}}\},\min\{1,b_{\mathrm{dim}}\}\},
\]
while the component risk bounds imply
\[
\underline c_*
\max\{\min\{1,a_{\mathrm{eff}}\},\min\{1,b_{\mathrm{dim}}\}\}
\le\mathfrak R_{n,d,q}.
\]
Since \(\mathfrak R_{n,d,q}\ge0\), we conclude that
\[
\underline c\,r_{n,d,q}
\le
\frac{\underline c_*}{2}
\max\{\min\{1,a_{\mathrm{eff}}\},\min\{1,b_{\mathrm{dim}}\}\}
\le\frac12\mathfrak R_{n,d,q}
\le\mathfrak R_{n,d,q}.
\]
% lean: CausalSmith.Stat.MarRareqLogfrontier.full_experiment_lower

\item \emph{Combining the bounds when \(d<\sqrt N\).}
Suppose \(\sqrt N>d\). Since \(d\ge0\),
\[
d^2\le(\sqrt N)^2=N.
\]
Also \(\ell\ge1\), so
\[
\ell^2\ge1,\qquad
(N\ell)^2>0,\qquad
N^{-1}(N\ell)^2=N\ell^2.
\]
Therefore
\[
b_{\mathrm{dim}}
=\frac{d^2}{(N\ell)^2}
\le\frac{N\ell^2}{(N\ell)^2}
=N^{-1}
=a_{\mathrm{eff}},
\]
and monotonicity of the minimum gives
\[
\min\{1,b_{\mathrm{dim}}\}
\le\min\{1,a_{\mathrm{eff}}\}.
\]
Applying the capped-sum inequality established above,
\[
r_{n,d,q}
\le
\min\{1,a_{\mathrm{eff}}\}+\min\{1,b_{\mathrm{dim}}\}
\le2\min\{1,a_{\mathrm{eff}}\}.
\]
Finally, the one-cell contribution and
\(\underline c_*\le\underline c_{\mathrm{one}}\) yield
\[
\underline c\,r_{n,d,q}
\le
\frac{\underline c_*}{2}\min\{1,a_{\mathrm{eff}}\}
\le
\underline c_{\mathrm{one}}\min\{1,a_{\mathrm{eff}}\}
\le\mathfrak R_{n,d,q}.
\]
The two dimension branches cover all admissible parameters.
% lean: CausalSmith.Stat.MarRareqLogfrontier.full_experiment_lower
\end{enumerate}
\end{proof}

\bibliographystyle{plainnat}

\bibliography{references}

\end{document}
